% 本文件由 LaTeX 排版 Agent 根据有效 translation.json 自动生成。
% author=Augustine of Hippo; work_id=1408; title=论圣洗，驳多纳徒派

\chapter{\WorkTitle{论圣洗，驳多纳徒派}}

% structure: p0001 source=h1 target=omitted reason=page-title-already-rendered-by-parent

\Emph{这部论著写于公元400年左右。关于它，\Person{奥古斯丁}在\newline{}\WorkTitle{订正录}卷二第十八章中说：\Quote{我写了七卷书论\newline{}圣洗，驳多纳徒派，他们试图以最有福的主教兼殉道者\Person{西彼廉}的权威来为自己辩护；在\newline{}这些书卷中，我指出，对于驳倒\newline{}多纳徒派，并直接堵住他们维护\newline{}自己分裂天主教会之口，没有什么比\newline{}作为\Person{西彼廉}的书信和行传\newline{}更有效的了。}}\par

卷一\newline{}卷二\newline{}卷三\newline{}卷四\newline{}卷五\newline{}卷六\newline{}卷七\par

\SourceRecord{Translated by J.R. King and revised by Chester D. Hartranft.；Nicene and Post-Nicene Fathers, First Series；Vol. 4.；Edited by Philip Schaff.；Buffalo, NY: Christian Literature Publishing Co.,；1887.；<http://www.newadvent.org/fathers/1408.htm>.}

% structure: p0001 source=h1 target=section reason=page-title

\section{\WorkTitle{论圣洗圣事，驳多纳徒派（卷一）}}

\EditorialNote{本书写于约公元400年。关于它，奥古斯丁在\WorkTitle{修正录}卷二第十八章中说：\Quote{我写了七卷论圣洗驳多纳徒派的书，他们试图以最可敬的主教和殉道者\Person{西彼廉}的权威为自己辩护；我指出，没有什么比\Person{西彼廉}的书信和行传更有效地驳斥多纳徒派，直接堵住他们以分裂对抗天主教会的口。}}\par

\Emph{他证明，圣洗圣事可以在天主教会共融之外由异端者或分裂者施行，但不应该从他们手中领受；并且，任何人在异端或分裂状态中，圣洗对他们都毫无益处。}\par

% structure: p0004 source=h2 target=subsection reason=relative-heading-rank; base=h2

\subsection{第一章}

1. 在先前我们为驳斥\Person{巴马尼安}致\Person{提科纽}的公开信而写的论著中，我们曾允诺日后将更详尽地探讨圣洗问题；的确，即使我们未曾做出这一允诺，我们也不会忘记，这是我们对弟兄们的祈祷所应尽的本分。因此，在本论著中，我们藉天主的助佑，不仅致力于驳斥多纳徒派在此事上一贯反对我们的异议，而且也要推进天主可能使我们就真福殉道者\Person{西彼廉}的权威所说的话，他们企图利用这一权威作为支柱，以阻止他们的乖戾在真理的攻击之前倒塌。我们提出这一点，是为了让所有不为党派偏见蒙蔽判断的人明白，\Person{西彼廉}的权威非但不利于他们，反而直接导致他们的驳斥和挫败。\par

2. 在上述论著中，我们已经说过，圣洗的恩宠可以在天主教会共融之外被赋予，正如它也可以在那里被持守一样。但多纳徒派中无人否认，即使背教者也持守圣洗的恩宠；因为当他们重返教会的怀抱，通过悔改而皈依时，圣洗从未再授予他们，由此可知它从未失去。同样，那些在分裂的亵渎中离开教会共融的人，当然也持守他们在离开前所领受的圣洗的恩宠，因为如果他们返回，圣洗不会再次授予他们——这证明他们在教会合一内所领受的，在分离后并未失去。但如果在外面可以持守，为什么不能在那里授予呢？如果你说：\Quote{在教会之外，圣洗不可正当地施行}；我们回答说：\Quote{正如它不可正当地持守，却仍在某种意义下被持守，同样，它确实不可正当地施行，但仍旧被施行。}然而，正如通过和好归于一，那些在分离中无益地持守的，开始有益地拥有；同样，通过这同一和好，那在分离外无益地施行的，开始变得有益。但不可允许说那曾经施行的没有被施行，也不可有人指责某人没有施行此事，同时却承认那人已经施行了他自己所领受的。因为圣洗圣事是受洗者所拥有的；而授予圣洗的圣事是受圣秩者所拥有的。正如受洗者若离开教会的合一，并不因此失去圣洗圣事，同样，受圣秩者若离开教会的合一，也不因此失去授予圣洗的圣事。因为两种圣事都不可被损害。如果圣事在恶人身上必然失效，那么两者都必须失效；如果它在恶人身上保持有效，那么对两者都必须如此。因此，如果承认那在分离后不能失去的圣洗，那么也必须承认那由在分离中未失去授予圣洗的圣事的人所施行的圣洗。因为正如那些回到教会的人，如果他们在分离前已受洗，不会重新受洗；同样，那些回到教会的人，如果在分离前已受圣秩，当然不会重新接受圣秩；而是要么若教会需要，他们重新行使他们从前的职务，要么若不行使，他们至少保留圣秩的圣事；因此，当为了标记他们的和好而给他们覆手时，他们不置于平信徒之列。因为当\Person{斐理西安}与\Person{玛西米安}一同分离时，连多纳徒派本身也不认为他失去了圣洗圣事或授予圣洗的圣事。因为现在他是他们自己团体中公认的成员，连同那些他在\Person{玛西米安}分裂期间与他们分离时所施洗的人。因此，其他人可以从他们那里领受，即使他们尚未加入我们的团体，而他们自己因从我们的团体分离而并未失去的东西。由此可见，那些试图在天主教合一内重新施洗的人犯了不虔之罪；而我们做得正确，即使是在分裂中施行的天主的圣事，我们也不敢拒绝。因为在他们与我们意见一致的各个方面，他们也与我们共融；只有在他们与我们意见分歧的方面，他们才与我们分离。因为物质或精神亲缘的联系与分离不能用不同的法则来衡量。正如物体的合一源于位置的连续性，同样，在意志的一致中，灵魂之间也有某种接触。因此，如果一个人从合一中分离出来，想要做与他在合一状态中所领受的不同之事，在这一点上他确实分离了，不再是有机整体的一部分；但无论在何处，他想要按照在合一状态中习惯的方式行事——即他自己在合一状态中学到并接受的那些教训——在这一点上他仍然是一个肢体，与整个机体联合。\par

% structure: p0007 source=h2 target=subsection reason=relative-heading-rank; base=h2

\subsection{第二章}

3. 因此，\OriginalTerm{多纳徒派}{Donatists}在一些事上与我们同在，在一些事上已离开了我们。因此，凡他们与我们一致的事，我们不禁止他们去做；但凡他们与我们不同的事，我们恳切地鼓励他们来从我们这里接受，或者视情况返回并重新获得；我们以一切可能的方式，满怀爱心地忙碌，使他们摆脱错误、得到纠正，选择这条道路。因此我们不对他们说：\Quote{不要施洗}，而是说：\Quote{不要处于分裂中施洗。}我们也不对那些我们看见即将受洗的人说：\Quote{不要领受圣洗}，而是说：\Quote{不要在分裂中领受它。}因为如果有人因紧急需要所迫，找不到一位天主教徒从他领受圣洗，于是心中保持天主教的和睦，从一位在天主教合一范围之外的人那里领受了那本来打算在合一之内领受的圣事，此人若立即离世，我们视他无非就是一位天主教徒。但如果他得以脱离肉身的死亡，亲自以身体回到那他心中从未离开过的天主教团体，我们非但不责备他的行为，反而要极其真实和确信地赞扬它；因为他相信天主临在于他的心，当他努力保持合一时，他不愿在没有圣洗圣事的情况下离开此生，他知道这圣事属于天主，而非属于人，无论他在何处找到它。但如果有人有能力在天主教会内领受圣洗，却出于某种心灵乖戾，宁愿在分裂中受洗，即使后来他想到要来天主教会，因为他确信在那里那圣事将对他有益——这圣事确实可以领受，但在别处无法有益——那么毫无疑问，他是乖戾的，是有罪的，而且其罪过因蓄意为之而更加昭彰。因为他毫不怀疑圣事在教会内领受才是正确的，这一点由他的信念所证明：即使是从别处领受的，他也必须到教会里寻求益处。\par

% structure: p0009 source=h2 target=subsection reason=relative-heading-rank; base=h2

\subsection{第三章}

4. 我们肯定两项主张——在天主教会内有圣洗圣事，且惟独在其中方可正当地领受——这两项唐纳派都否认。同样，我们肯定另外两项主张——在唐纳派中有圣洗圣事，但在他们那里不能正当地领受——对于这两项，他们竭力肯定前一项，即在他们中有圣洗圣事；但拒绝承认后一项，即在他们的教会中不能正当地领受。在这四项主张中，三项为我们所特有，一项与我们双方一致。因为说在天主教会内有圣洗圣事，在那里可正当地领受，以及在唐纳派中不能正当地领受，这些只是我们的断言；但说在唐纳派中也有圣洗圣事，则是他们的断言，也是我们所同意的。因此，若有人愿意受洗，且已确信他应选择我们的教会作为得救的途径，并确信基督的圣洗即使在其他地方领受，也只在我们内有益，却仍希望在唐纳派的裂教中受洗，因为不仅他们，不仅我们，而且双方都承认在他们那里有圣洗圣事，那么他应当暂停并察看其他三点。因为如果他决定在他们是所否认的事上跟随我们，尽管他更相信我们所共同承认的而非我们所独断的，这于我们的目的已足：他更相信他们不肯定而惟独我们肯定的，胜于他们独断的。我们在天主教会内有圣洗圣事，这是我们的断言，他们的否认。在天主教会内可正当地领受圣洗圣事，这是我们的断言，他们的否认。在唐纳派的裂教中不能正当地领受圣洗圣事，这是我们的断言，他们的否认。因此，正如他更愿意相信我们独断应信之事，同样让他更愿意去做我们独断应行之事。但让他更坚定地相信双方都断言应信之事，胜于我们独断之事——倘若他如此倾向。因为他更倾向于相信在唐纳派的裂教中有基督的圣洗，因为这为双方所承认，而非在天主教会内有，后者仅为公教徒所断言。但同样，他更愿意相信基督的圣洗也在我们内，如我们所独断，而非不在我们内，如他们所独断。因为他已决定并完全确信，我们之间的分歧之处，我们的权威应优先于他们的。因此，他更愿意相信我们所独断的，即圣洗在我们内可正当地领受，而非不能正当地领受，因后者仅基于他们的断言。按同一规则，他更愿意相信我们所独断的，即圣洗在他们内不能正当地领受，而非如他们所独断的，即能正当地领受。因此，他发现在唐纳派中受洗的信心多少无益，因为虽然我们双方都承认在他们中有圣洗圣事，但我们并未双方都宣告应从他们那里领受。但他已决意在分歧之处更依附我们。因此，让他放心在我们的共融中领受圣洗，在那里他确知它既存在又可正当地领受；不要在一个共融中领受，在那里他所决定追随其意见的人虽承认它存在，但说不能正当地领受。甚至，即使他视之为疑问，即在唐纳派中是否不可能正当地领受那在天主教会中确知可正当地领受的，那么单凭他偏爱不确定者而不确定者，就已在关乎其灵魂得救之事上犯了大罪。无论如何，他必定确信一个人在天主教会中可正当地受洗，因为他已决定即使在其他地方受洗也要归附我们。但至少让他承认，在唐纳派中是否不得当地受洗是不确定的，因为他发现这是他所确信其意见应优先于他们的人所断言的；因此，以确定者优先于不确定者，让他在这里受洗，在这里他有充分理由确信这合法，因为当他曾考虑在其他地方受洗时，他仍决定之后应归附此方。\par

% structure: p0011 source=h2 target=subsection reason=relative-heading-rank; base=h2

\subsection{第四章}

5. 此外，若有人不明白我们为何断言多纳徒派所行的圣事并非合法，却又承认圣事确实存在于他们中间，那么就让他留意：我们也否认圣事在他们那里是合法存在的，正如他们否认圣事在那些脱离他们共融的人那里是合法存在一样。也让他思考关于军人印记的类比：这个印记既可以被逃兵保留，也可以被不在军队中的人接受，但在军队之外，它既不应被接受，也不应被保留；同时，当一个人入伍或重返服役时，这印记并不会被改变或更新。然而，我们必须区分两种情况：一种是那些无意中加入这些异端者行列的人，他们以为自己是进入基督的真教会；另一种是那些知道除那按照应许传遍全世界、甚至延伸到地极的天主教会以外，别无其他天主教会的人；这教会生于莠草之中，渴望将来从罪恶的疲倦中得到安息，在圣咏集里说：\BibleQuote{我从地极向你呼号，我的心困倦；你把我高举在磐石上。}但这磐石就是基督，宗徒说我们在祂里面已经被提升，并与祂一同坐在天上，虽然不是实际上，而只是在希望中。\BibleRef{弗 2:6} 因此圣咏继续说：\BibleQuote{你作了我的引导，因为你成了我的希望，一座坚固的堡垒，以抵御仇敌的面目。}藉着祂的应许，这些应许如同储存在坚固堡垒中的矛和标枪，仇敌不仅被防御，而且被推翻，因为他披着羊皮，\BibleRef{玛 7:15} 从而他们可以说：\BibleQuote{看，基督在这里，或在那里；}\BibleRef{玛 24:23} 并且他们可能将许多人从那建在山上的天主之城中分离出来，诱入他们自己的网罗的孤独中，从而彻底毁灭他们。这些人明知如此，却选择在基督身体的共融统一之外接受基督的圣洗，尽管他们打算以后带着在别处领受的圣事进入这同一共融。因为他们提议在对抗基督教会的情况下接受基督的圣洗，并且在他们接受圣洗的那一天也清楚地知道这一点。如果这是罪，那么谁会允许说：\Quote{饶恕我犯罪一天}呢？因为如果他打算转投天主教会，我倒想问为什么。除了回答\Quote{属于多纳徒派而不属于天主教会的一体性是不好的}之外，他还能给出什么答案呢？那么，你作恶多少天，你就犯了多少天的罪。可以说，天数越多，罪越大；天数越少，罪越小；但绝不能说什么罪也没犯。然而，有什么必要允许这受诅咒的恶行持续一天或一个小时呢？因为那希望获得这种许可的人，还不如向教会或向天主本人请求允许他背教一天。因为如果他不在乎成为一时片刻的分裂者或异端者，就没有理由害怕成为一天背教者。\par

% structure: p0013 source=h2 target=subsection reason=relative-heading-rank; base=h2

\subsection{第五章}

6. 他说：我宁愿在双方都承认圣洗存在的地方领受基督的圣洗。但那些你打算加入的人说，在那里不能合法地领受；而那些说可以合法地领受的人，正是你打算离开的一方。因此，那些你本人认为权威较低的人所说的话，与那些你本人更偏爱的人所说的话相对立，如果并非虚假，至少用较温和的话说，是不确定的。所以我恳求你，宁可选真弃假，或选定弃疑。因为不仅那些你将要加入的人，而且你自己（你将要加入他们）都承认，你想要的圣洗可以在那个团体中合法地领受，而你却要在别处领受后再加入它。因为如果你对那里是否可以合法领受有任何怀疑，你也会对是否应该改变产生怀疑。因此，如果从多纳徒派领受圣洗是否有罪尚且存疑，谁还能怀疑不优先选择在确定无罪的地方领受圣洗是确定的罪呢？而那些因无知在那里受洗、以为那是基督真教会的人，比起这些人，罪过较轻；尽管连他们也因分裂的不虔敬而受伤，也不能免于严重的伤害，因为别人所受的甚至更重。因为当某些人说：\BibleQuote{在审判的日子，索多玛地所受的惩罚也要比你们容易忍受}\BibleRef{玛 11:24}时，并不是说索多玛人将逃脱刑罚，而只是说那些人将受到更重的刑罚。\par

7. 然而，这一点从前也许曾隐晦不明。那为留心并接受改正之人带来健康的事，对那些虽不再被容忍无知、却仍执迷不悟以致自取灭亡的人，只是加重了他们的罪过。因为谴责\Person{马克西米安}派、并在此后恢复他们（连同那些他们——用他们自己会议的话说——在分裂中亵渎地施洗的人）的地位，了结了这整个争议，消除了所有争论。在我们自己与那些与\Person{普里米安}共融的多纳徒派之间，没有任何争议点可能引起疑问，以为基督的圣洗不仅可以为那些与教会分离的人所保留，甚至可以被他们所授予。因为正如他们自己不得不承认，那些\Person{斐利西安}在分裂中施洗的人确实领受了真正的圣洗——因为他们现今承认他们为自身团体的成员，所领受的圣洗除了在分裂中领受的外别无其他；同样，我们说，即使在天主教共融之外，那些与其共融断绝者所授予的圣洗，仍是基督的圣洗，因为他们断绝时并未失去这圣洗。当他们接纳那些\Person{斐利西安}在分裂中施洗的人与自己和好时，他们自认为所赐予这些人的（即不应使这些人领受他们尚未拥有的，而应使他们在分裂中已领受却无益的、并且已经拥有的圣洗对他们有益），这正是天主通过天主教共融，真正赐予和恩许给那些从任何异端或分裂（在其中领受了基督的圣洗）而来的人的：即他们不应开始领受圣洗圣事，仿佛未曾拥有过它，而应使他们已经拥有的圣洗现在开始对他们有益。\par

% structure: p0016 source=h2 target=subsection reason=relative-heading-rank; base=h2

\subsection{第六章}

8. 因此，在我们与那些可称为真正多纳徒派（其主教是\Place{迦太基}的\Person{普里米安}）之间，如今在这一问题上已无争议。因为天主愿意通过\Person{马克西米安}的追随者了结此事，使他们因他的先例而被强迫承认那些他们在基督徒爱德的劝导下原不肯承认的事。但这使我们接下来思考：那些拒绝与\Person{普里米安}派共融的人，声称在他们团体中，按着人数的稀少比例，保留着更大的多纳徒主义纯正——这些人是否似乎有些可为自己辩护的理由呢？即使他们只是\Person{马克西米安}派，我们也不应轻视他们的救恩。何况，当我们发现这一多纳徒派已分裂成许多极微小的碎片，所有这些小部分都指责那个以\Person{普里米安}为首的、大得多的部分接纳了\Person{马克西米安}追随者的圣洗；而每一小部分都竭力主张自己是唯一真圣洗的保管者，这圣洗不存在于别处——既不存在于天主教会所延伸的整个世界，也不存在于多纳徒派那个较大的主体，甚至不存在于其他微小部分，而只存在于自身。然而，如果所有这些碎片都愿听从的不是人的声音，而是最不容置疑的真理彰显，并愿意约束自身乖僻的烈性，他们就会从自身的荒芜返回——不是返回到多纳徒的主体（他们不过是这主体碎片的更小碎片），而是返回到天主教会之根的永不枯竭的丰饶。因为他们中凡不反对我们的，就是支持我们的；但当他们不与我们一同聚集时，便是分散。\par

% structure: p0018 source=h2 target=subsection reason=relative-heading-rank; base=h2

\subsection{第七章}

9. 首先，为了不使我似乎仅仅依靠人为的论证——因为这个问题中有如此多的晦暗，以致在教会早期，即在多纳图派分裂之前，就已经使一些极为重要的人物，甚至我们的教父们，即那些心中充满爱德的主教们，在彼此之间争论和怀疑，但始终维护教会的和平，以至于他们各自地区的会议的法令长期彼此不同，直到最终通过一次全球性的全体大公会议，确立了对所有疑虑的最有益的观点：——因此，我从福音中带来明确的证据，借着天主的助佑，我打算证明，在天主眼中，这决定是多么正确和真实，即对于每一个分裂者和异端者，那导致他分离的创伤应通过教会的医药来医治；但在他身上仍然健康的部分，应得到赞许的承认，而不是通过谴责来伤害它。的确，主在福音中说：\BibleQuote{不与我相合的，就是反对我的；不同我收集的，就是分散的。}\BibleRef{玛 12:30}然而，当门徒们告诉主，他们看见一个人因主的名驱逐魔鬼，就禁止了他，因为他不跟从他们，主却说：\BibleQuote{不要禁止他：因为不反对我们的，就是支持我们的。因为没有人因我的名行奇迹，会轻易毁谤我。}确实，如果这个人身上没有任何需要纠正的地方，那么任何人，若置身于教会的共融之外，与所有基督徒兄弟情谊分离，却因基督的名收集，就安全了；这样，\BibleQuote{不与我相合的，就是反对我的；不同我收集的，就是分散的}这话就不真实了。但如果在门徒们因无知而急于阻止他的那一点上，他需要纠正，那么为什么我们的主通过说\BibleQuote{不要禁止他}来阻止这种阻止呢？祂随后所说的\BibleQuote{不反对你们的，就是支持你们的}又如何能真实呢？因为在这点上，他不反对他们，而是支持他们，因为他因基督的名行治愈的奇迹。因此，为了使两者都真实——正如两者都真实一样——既宣言\BibleQuote{不与我相合的，就是反对我的；不同我收集的，就是分散的}，又命令\BibleQuote{不要禁止他；因为不反对你们的，就是支持你们的}——我们必须理解什么呢？除非是：那人在他对那大能名字的尊崇中应被坚定，因为在这方面他不反对教会，而是支持它；然而，他应因与教会分离而受到责备，因为他的收集变成了分散；并且，如果碰巧他寻求与教会合一，他不应领受他已经拥有的，而应纠正他走错的地方。\par

% structure: p0020 source=h2 target=subsection reason=relative-heading-rank; base=h2

\subsection{第八章}

10. 的确，外邦人科尔乃略的祈祷并未被忽视，他的施舍也未被拒绝；他反而有资格被派一位天使到他那里，并见到那使者，通过他他无疑可以学到一切必要的事，而无需任何人到他那里。但是，因为他祈祷和施舍中的一切善行，除非他借着基督徒兄弟情谊与和平的纽带加入教会，否则不能使他受益，所以他奉命派人去请伯多禄，并借着他认识了基督；并且，按照他的命令领受圣洗后，他通过共融的纽带与基督徒的团体相连，此前他只是因善工的相似而与之相连。\EditorialNote{宗十章}的确，轻视他尚未拥有的，并夸耀他已拥有的，是极为致命的。同样，那些因分离而破坏爱德，从而打破合一纽带的人，如果他们不遵守他们在那个团体中所领受的任何东西，那么他们在一切上都是分离的；因此，任何被他们纳入他们团体的人，若后来希望来到教会，应当领受他尚未领受的一切。但如果他们遵守一些相同的东西，那么在这些方面他们并未分离；就此而言，他们仍然是教会的框架的一部分，而在所有其他方面他们被切断了。因此，任何被他们与自己联合的人，在他们未与教会分离的所有方面，都与教会联合。因此，如果他希望来到教会，那么在他不健康、走错的地方，他得到医治；但在他与教会联合的健康之处，他得到的不是医治，而是承认——免得我们因试图医治健康之处反而造成创伤。因此，那些他们施洗的人，他们治好了拜偶像或不信的创伤；但他们用自己的分裂给这些人造成了更严重的创伤。因为主的人民中的拜偶像者被刀剑击杀；\EditorialNote{出三十二章}但分裂者被地开口吞没。\EditorialNote{户十六章}而宗徒说：\BibleQuote{我若有全备的信德，甚至能移山，却没有爱德，我什么也不算。}\BibleRef{格前 13:2}\par

11. 若有人被带到外科医生那里，身体某要害部位受了重伤，医生说不治好这伤必定致死，那送他来的人，我想，不会愚蠢到向医生数算他所有健康的肢体，并指着它们回答说：\Quote{难道所有这些健康的肢体都不能救他的命，而只有一处受伤的肢体就足以使他死亡吗？}他们当然不会这样说，而是把他交给医生医治。同样，他们也并不因为这样交托，就请医生也医治那些健康的肢体；而是希望医生尽心地给那一个部位敷药，因为从那部位，死亡正威胁着其他健康的部分，而且若不治好伤口，死亡是必然的。那么，一个人若拥有健全的信德，或只是在信德的圣事上健全，而他的爱德健全却因分裂的致命创伤而毁坏，以致随着爱德的倾覆，其他本身健全的点也被拖入死亡的感染中，这对他有什么益处呢？为防止这事，天主的仁慈，藉着祂圣教会的合一，不断地努力使他们通过和平的联系，藉着修和的药来获得医治。他们不要以为我们承认他们身上有某些健全之处，他们就完全健全了；也不要以为我们需要医治那健全的部分，因为我们指出在某些部分有创伤。因此，就圣事的健全而言，因为他们不反对我们，他们就是支持我们的；但就分裂的创伤而言，因为他们不与基督一起收集，他们就分散。不要因他们所拥有的而自高。为什么他们只以骄傲的目光看那些健全的部分？也要谦卑地看他们的创伤，不仅注意到他们所拥有的，也要注意到他们所缺乏的。\par

% structure: p0023 source=h2 target=subsection reason=relative-heading-rank; base=h2

\subsection{第九章}

12. 让他们看看，有多少重要的事情，若缺少某一件，便毫无益处，也让他们看看那件是什么。在此，请他们听，不是我的话，而是宗徒的话：\BibleQuote{我若能说人间的语言，和天使的语言；但我若没有爱德，我就成了个发声的锣，或发响的钹。我若有先知之恩，又明白一切奥秘和各种知识；我若有全备的信德，甚至能移山；但我若没有爱德，我什么也不算。}\BibleRef{格前 13:1} 因此，如果他们在神圣奥秘中拥有天使般的言语，又有先知之恩，如同\Person{盖法}（见：\BibleRef{若 11:51}）和\Person{撒乌耳}（见：\BibleRef{撒上 18:10}）那样，圣经明证他们是该受谴责的，但仍能预言，这对他们有什么益处呢？如果他们不仅知道，甚至也拥有圣事，如同\Person{西满}术士那样（见：\BibleRef{宗 8:13}）；如果他们拥有信德，如同魔鬼也承认基督（因为我们不能以为他们不信，当他们说：\BibleQuote{祢与我们有什么相干？祢是天主子，我们知道祢是谁。}\BibleRef{谷 1:24}）；如果他们散尽自己的财产周济穷人，如同许多人在天主教会内，也在不同异端团体中所做的；如果在任何逼迫的压力下，他们与我们一同为所承认的信德而舍身被焚：却因为他们所做这一切都离开了教会，没有\BibleQuote{以爱德彼此担待}，也没有\BibleQuote{尽力以和平的联系，保持心神的合一}（见：\BibleRef{弗 4:2}），既然他们没有爱德，即使有所有那些对他们无益的好处，他们也不能获得永远的救恩。\par

% structure: p0025 source=h2 target=subsection reason=relative-heading-rank; base=h2

\subsection{第十章}

13. 但他们自认为在以下问题上表现出极大的精细：他们问，多纳徒派中的基督的圣洗是否使人成为子女？这样，如果我们承认它使人成为子女，他们就可以宣称他们的教会是母亲，能在基督的圣洗中生养子女；而且既然教会是唯一的，他们就可以说我们的教会不是教会。但如果我们说它不能使人成为子女，他们就问：\Quote{那么，为什么你们不使那些从我们这里到你们那里去的人，在我们这里受洗之后，再领受圣洗而重生呢？如果他们在我们这里还没有因此重生的话？}\par

14. 正如他们的派别并非因使其分裂的因素而获得生产能力，而是因使其与教会结合的因素。因为它与和平与爱德的纽带分离，却在同一圣洗中结合。因此，只有一个教会，唯独被称为公教；每当它在这些与自身分离的不同团体的共融中拥有任何属于它自己的东西时，那生产能力必然属于它，而非它们，因着它在每一团体中属于自己的东西。因为生产并非来自他们的分离，而是来自他们所保存的教会本质；如果他们继续放弃这本质，就会失去生产能力。因此，每一次生产都来自教会，教会的圣事得以保存，任何这样的出生只能从这圣事而来——尽管并非所有接受其出生的人都属于它的合一，这合一将拯救那些坚持到底的人。不仅是那些公开犯有分裂之明显亵渎的人不属于它，而且那些在外表上与其合一相连，但被罪恶生活分开的人也不属于它。因为教会本身曾通过圣洗的圣事生了\Person{西满术士}；然而却向他宣告，他在基督的继承中无分。他在圣洗、福音、圣事方面缺少什么吗？但因他缺少爱德，他白生了；或许他从未出生对他更好。那些宗徒对他们说\BibleQuote{我曾用奶哺养你们，没有用饭，如同在基督里的婴孩}的人，他们的出生缺少什么？然而，宗徒又因他们是属血肉的，将他们从他们正冲入的分裂之亵渎中唤回：\BibleQuote{我曾用奶哺养你们，没有用饭：因为那时你们还不能承受，如今你们还是不能。因为你们仍是属血肉的：因为在你们中间有嫉妒和纷争，你们岂不是属血肉的，照着人的样子行事吗？因为有人说，我是属保禄的；另有人说，我是属阿颇罗的；你们岂不是人吗？}\BibleRef{格前 3:1} 关于这些人，他上文说：\BibleQuote{弟兄们，我藉我们主耶稣基督的名，劝你们都说一样的话，你们中间不可有分裂，只要一心一意，彼此完全合而为一。因为\Person{克罗厄}家里的人曾对我提起你们，我的弟兄们，说你们中间有争竞。我的意思就是，你们各人说，我是属保禄的，我是属阿颇罗的，我是属刻法的，我是属基督的。基督分开了吗？保禄为你们钉死在十字架上吗？你们是受洗归入保禄的名下吗？}\BibleRef{格前 1:10} 因此，这些人如果继续停留在同样乖张的顽固中，他们无疑是被生了，却不会因着和平与合一的纽带而属于那使他们出生的教会本身。因此，她本身藉着同一圣事，如同藉着丈夫的种子，在自己怀中和在婢女的怀中孕育他们。因为宗徒说这一切事都是预像，并非没有意义。但那些过于骄傲，不与合法母亲相连的人，就像\Person{依市玛耳}，关于他经上说：\BibleQuote{把这婢女和她的儿子赶出去，因为婢女的儿子不可与我儿子\Person{依撒格}一同继承产业。}\BibleRef{创 21:10} 而那些和平地爱戴父亲合法的妻子、因合法血统是她儿子的人，就像\Person{雅各伯}的儿子们，确实由婢女所生，却同样承受继承。\BibleRef{创 30:3} 但在家庭内、由母亲本身子宫所生，却忽略了所领受恩宠的人，就像\Person{依撒格}的儿子\Person{厄撒乌}，他被弃绝了，天主亲自为此作证，说：\Quote{我爱了\Person{雅各伯}，恨了\Person{厄撒乌}}；尽管他们是双胞胎兄弟，同一子宫所生。\par

% structure: p0028 source=h2 target=subsection reason=relative-heading-rank; base=h2

\subsection{第十一章}

15. 他们也问：\Quote{罪在\Person{多纳图}派的团体中是否在圣洗中获得赦免？}因此，如果我们说在那里罪得赦免，他们就会回答说，那么圣神就在那里；因为主在嘘气时将圣神赐给门徒时，接着说道：\BibleQuote{你们要去使万民成为门徒，因父及子及圣神之名给他们授洗。}\BibleRef{玛 28:19} \BibleQuote{你们赦免谁的罪，就给谁赦免；你们保留谁的罪，就给谁保留。}\BibleRef{若 20:23} 如果这样，他们说，那么我们的共融就是基督的教会；因为圣神除非在教会内，不行赦罪之事。如果我们的共融是基督的教会，那么你们的共融就不是基督的教会。因为无论在哪里，那独一的教会，经上论它说：\BibleQuote{我的鸽子，我的完全人，只有这一个；是她母亲独生的。}\BibleRef{歌 6:9} 而且教会不会像分裂那样多。但如果我们说罪在那里不得赦免，那么，他们说，那里就没有真正的圣洗；因此你们应该给那些从我们这里接受的人施行圣洗。既然你们没有这样做，你们就承认你们不在基督的教会内。\par

16. 我们依照圣经回答他们，请他们回答他们自己所问我们的问题。因为我恳求他们告诉我们：没有爱德的地方，能否有罪的赦免？因为罪是灵魂的黑暗。我们读到圣若望说：\BibleQuote{恨自己弟兄的，仍然在黑暗中} \BibleRef{若一 2:11} 但若不是被对弟兄的仇恨所蒙蔽，没有人会制造分裂。因此，如果我们说在那里罪不得赦免，那么在他们当中受洗的人怎能重生呢？圣洗中的重生，不就是从旧人的败坏中更新吗？一个人的过去之罪没有得赦免，又怎能如此更新呢？但若他没有重生，他也就没有穿上基督；由此似乎得出他应当重新受洗。因为宗徒说：\BibleQuote{凡你们受洗归于基督的，就是穿上了基督} \BibleRef{迦 3:27} 如果他还没有这样穿上基督，那么他也不应被视为已在基督内受洗。此外，既然我们说他在基督内受了洗，我们就承认他穿上了基督；若我们承认这一点，我们就承认他重生了。若果真如此，圣若望说\BibleQuote{恨自己弟兄的，仍然留在黑暗中}，而他的罪却已经赦免了，这怎么可能呢？难道分裂不包含仇恨自己的弟兄吗？谁会坚持这一点呢？因为分裂的起源和坚持，除了仇恨弟兄之外，别无他物。\par

17. 他们认为自己解答了这个问题，说：\Quote{这样，分裂中就没有罪的赦免，因此也没有藉重生的新人创造，所以也没有基督的洗礼。} 但既然我们承认基督的洗礼存在于分裂中，我们就向他们提出这个问题请他们解答：\Person{西满术士}是否领受了基督的真洗礼？他们必回答：是，受圣经权威的迫使。我问他们是否承认他得到了罪的赦免。他们必定承认。于是我质问，为何\Person{伯多禄}对他说，他在圣徒的份上没有份。因为他们说，他后来犯了罪，想用钱购买天主的恩赐，他认为宗徒们能出卖这恩赐。\par

% structure: p0032 source=h2 target=subsection reason=relative-heading-rank; base=h2

\subsection{第十二章。}

18. 假如他是在欺骗中接近圣洗本身呢？他的罪是赦免了，还是没有赦免？让他们选择吧。无论他们选择哪一项，都将符合我们的目的。如果他们说是赦免了，那么\BibleQuote{惩戒的圣神怎能逃避欺骗} \BibleRef{智 1:5} 呢，如果在他这个满心欺诈的人身上，圣神竟能产生罪的赦免？如果他们说是没有赦免，那我就要问，如果他后来怀着痛悔之心和真正的忧伤承认了自己的罪，是否应该被认为需要重新受洗？如果主张这一点简直是疯狂，那么让他们承认：一个人可以领受基督的真圣洗，但他的心若坚持恶意或亵渎，可能不允许罪的赦免赐予；这样，他们就能明白：在脱离教会的团体中，人可以受洗，在那里基督的圣洗在所述圣事的举行中被给予和领受，但只有当领受者与教会的合一和好，并从欺骗的亵渎中被净化时，这圣洗才会对赦免罪过有效，因为欺骗使罪被保留，阻止赦免。因为，如同那个在欺骗中接近圣事的人，没有第二次圣洗，而是通过信德的纪律和真诚的告解被净化——这没有圣洗他做不到——使得先前给予的圣洗，当先前的欺骗被真诚的告解除去时，才变得有能力促成他的得救；同样，对于那个与基督的和平与爱为敌，在任何异端或分裂中领受了基督圣洗的人——那分裂者并未失去这圣洗——尽管因他的亵渎其罪未得赦免，但当他纠正错误，来到教会的共融与合一时，他不应再受洗：因为正是通过他与教会和平的修和，他得到了这一恩惠，即这圣事开始在合一中有效于赦免他的罪，而这在分裂中领受时无法如此有效。\par

19. 但若他们说，在欺骗中接近圣事的人，在他领受的那一刻，他的罪确实因如此伟大圣事的神圣德能被除去了，但因他的欺骗而立即返回；这样，圣神既在圣洗中临在于他以除去他的罪，又因他坚持欺骗而逃逸，致使罪回来：这样两句经文就都成为真实——\BibleQuote{凡你们受洗归于基督的，就是穿上了基督}，以及\BibleQuote{惩戒的圣神将逃离欺骗}——这就是说，一方面圣洗的圣德给他穿上基督，另一方面欺骗的罪恶给他脱去基督；好比一个人从黑暗穿过光明又回到黑暗，他的眼睛始终朝向黑暗，尽管当他穿过时光明不得不照射到他——如果他们这样说，那就让他们明白：对于那些在教会之外受洗，但领受了教会圣洗的人（这圣洗在自身内是神圣的，无论在哪里），也是如此；这圣洗不属于那些分离者，而属于那被分离的身体；然而它甚至在分离者中也有功效，使他们穿过它的光明而回到自己的黑暗，他们的罪在那瞬间曾因圣洗的圣德而被驱散，却立刻返回他们身上，仿佛黑夜返回，而光明在他们穿过时曾驱散它。\par

20. 主以那欠他一万塔冷通的仆人——他应其请求而全部豁免——最清楚地教导我们：在没有兄弟之爱的地方，已经赦免的罪会重返于人。但当这仆人拒绝怜悯那欠他一百银钱的同仆时，主便命令他偿还所豁免的债务。因此，藉圣洗领受赦免的时刻，犹如结算账目之时，所有查明的债务皆可豁免。但后来那仆人并没有借钱给他的同仆，而只是在他无力偿还时无情地追讨；其实同仆早已欠他这笔债，他自己在向主人结算时却获得如此巨额债务的豁免。他并非先豁免了同仆，然后才领受主人的宽恕。同仆的话证明了这一点：\Quote{请你容忍我，我必还清你这一切。}否则他会说：\Quote{你先前已豁免了我，为何又追讨？}主自己的话更清楚地表明这一点，因为祂说：\Quote{但那仆人出去，遇见一个同仆，欠他一百银钱。}\BibleRef{玛 18:23}祂并没有说：\Quote{那人已豁免他一百银钱的债。}既然祂说\Quote{欠他}，显然仆人并未豁免同仆的债。确实，对于一个即将结算如此巨债、并期待主人宽恕的人来说，更妥当、更合宜的做法是：他应先豁免同仆的债务，然后再在亟需祈求主人怜悯的情况下去结算。然而，他尚未豁免同仆这一事实，并不妨碍主在收账之日豁免他所有的债务。但这对他有何益处呢？由于他顽固地缺乏爱德，所有这些债务立刻加倍地回到他头上。同样，即使接受圣洗者心中仍对弟兄怀有仇恨，圣洗的恩宠也并不会阻止所有罪的赦免。因为昨天的罪以及之前的一切，甚至受洗前的那一刻和受洗时的罪都得以赦免。但随后他立即又开始承担责任，不仅是为接下来的日子、时辰、片刻，也为过去的罪——所有已赦免之罪的罪咎重返于他，这在教会中屡见不鲜。\par

% structure: p0036 source=h2 target=subsection reason=relative-heading-rank; base=h2

\subsection{第十三章}

21. 因为常有这样的事：一个人极其不义地仇恨一个敌人，尽管我们被命令去爱甚至不义的敌人，并为他们祈祷。但在某种突发死亡危险中，他开始不安，渴望领受圣洗，匆忙领受，以致情况紧急，几乎无法进行必要的简短考查，更不用说长时间交谈来将这种仇恨从心中驱除，即使施洗者知道此事。这种情形确实不仅在我们中间，也在他们中间常有发生。那么，我们该怎么说呢？这个人的罪被赦免了吗？让他们选择他们认为正确的一种。如果它们被赦免，它们会立刻返回：这是福音的教导，真理的权威宣告。因此，无论它们被赦免与否，随后都需要治疗；然而，如果这人活了下来，并知道他的过失需要改正，且改正了，他既不在他们中间，也不在我们中间重新受洗。所以，在分裂者和异端者与我们持相同见解或遵行相同做法的事上，当他们加入我们时，我们不纠正他们，反而称赞我们所见到的。因为凡与我不相异的，就不是与我是分离的。但由于这些事在他们仍是分裂者或异端者时对他们毫无益处，因为他们在其他方面与我们不同，更不用说分离本身所犯的最严重的罪，所以，无论他们的罪是留存在他们身上，还是在赦免后立即重返，我们在两种情况下都劝勉他们来获得平安和基督徒爱德的健全，不仅使他们获得先前没有的东西，也使已有的开始对他们有益。\par

% structure: p0038 source=h2 target=subsection reason=relative-heading-rank; base=h2

\subsection{第十四章}

22. 因此，他们对我们说：\Quote{如果你们承认我们的圣洗，我们缺少什么，以致你们认为我们应该认真考虑加入你们的团体？}这是徒劳的。因为我们回答说：我们不承认任何你们的圣洗；因为那不是分裂者或异端者的圣洗，而是天主和教会的圣洗，无论它在哪里被发现，无论它被转移到哪里。但它绝不属于你们，只是因为你们持错误见解、行亵渎之事，并邪恶地使自己与教会分离。因为即使你们在实践和见解中一切都正确，却仍然坚持这种分离，违反兄弟平安的纽带，违反所有弟兄的合一——这些弟兄按照应许显现在全世界，他们的历史细节和内心的秘密，你们根本不可能全部知道或考虑，以致有权定他们的罪；而且，他们不可能因为服从教会的法官而非争议的一方而受责备——在这样的事上，至少在这一方面，你们有所欠缺，就像缺乏爱德的人有所欠缺一样。我们何必再重述我们的论证呢？你们在宗徒身上看看你们自己吧，你们所欠缺的是那么多。因为对于一个缺乏爱德的人，无论他是立即被某种诱惑的狂风吹到教会之外，还是留在主的麦场里，直到最后的簸扬才被分离，这有什么分别呢？然而，这样的人一旦在圣洗中重生，就不必再重生。\par

% structure: p0040 source=h2 target=subsection reason=relative-heading-rank; base=h2

\subsection{第十五章}

23. 因为正是教会生了所有的人，或在她的圈内，由她自己的子宫；或在圈外，由她新郎的种子——或由她自己，或由她的婢女。但\Person{厄撒乌}虽然由合法的妻子所生，却因与兄弟争吵，而与天主的子民分离。\Person{阿协尔}虽然确实是由妻子的权力所生，却是出于婢女，并因兄弟之善意而被接纳进入\Place{应许之地}。由此，阻碍\Person{依市玛耳}、导致他与天主之民分离的，并不是生于婢女，而是他与兄弟的争吵。他并没有从妻子的权能中受益，他更像是妻子的儿子，因为正是凭着她婚姻的权利，他才能在婢女的子宫中被怀胎和出生。正如多纳徒派中，凡是出生的人都是借着存在于\OriginalTerm{圣洗圣事}{Baptism}中的教会之权利而获得出生；但如果他们与弟兄们和睦相处，借着和平的统一，他们来到\Place{应许之地}，不是为了再次被从他们真正母亲的怀中赶出，而是要在他们父亲的种子中被承认；但如果他们坚持不和睦，他们将属于\Person{依市玛耳}的血统。因为\Person{依市玛耳}是长子，以后才有\Person{依撒格}；\Person{厄撒乌}是年长的，\Person{雅各伯}是年幼的。这并不是说异端在教会之前生育，或者教会自己先生育属肉体或属生灵的人，然后才生育属神的人；而是因为，在我们实际有死的命运中，我们由亚当的种子而生，\BibleQuote{那属神的不是在先，而是属生灵的在先，以后才是属神的。}\BibleRef{格前 15:46}但仅由肉性的知觉，因为\BibleQuote{属血气的人不接受天主圣神的事}\BibleRef{格前 2:14}，就产生了所有的纷争和分裂。而宗徒说\EditorialNote{迦拉达书第四章}，所有坚持这种肉性知觉的人都属于旧约，即属于对地上许诺的渴望，这些确实是属神事物的预像；但\BibleQuote{属血气的人不接受天主圣神的事。}\BibleRef{格前 2:14}\par

24. 因此，不论何时，人此生开始具有这种本性，即虽然他们领受了为他们在其中生活的时期所指定的神圣圣事，却仍贪图属肉体的事物，并从天主那里盼望和渴求属肉体的事物，不论是今生还是来世，他们仍然是属肉体的。但教会，即是天主的子民，即使在此生的旅途中也是一个古老的体制，在某些人身上有属肉体的利益，在另一些人身上有属神的利益。属于属肉体的是旧约，属于属神的是新约。但在起初的日子，两者都是隐藏的，从亚当直到\Person{梅瑟}。借着\Person{梅瑟}，旧约被彰显，而新约隐藏在其中，因为它以一种隐秘的方式被预表。但主一降生成人，新约就被启示出来；然而，虽然旧约的圣事过去了，旧约特有的倾向却没有过去。因为它们仍然存在于那些人中，宗徒宣称他们虽已确实由新约的圣事所生，但作为属生灵的人，仍然有能力接受天主圣神的事。因为，如同在旧约的圣事中，有些人已经是属神的，秘密地属于当时隐藏的新约，同样现在，在这已经启示的新约圣事中，也有许多人是属生灵的。如果他们不前进去接受天主圣神的事，即宗徒的言论敦促他们去接受的，他们仍将属于旧约。但如果他们前进，甚至在接受它们之前，单凭他们的前进和亲近，他们就属于新约；如果在成为属神的人之前，他们被从此生夺去，借着圣事圣德之保护，他们仍被算在\Place{活人之地}，那里主是我们的希望和我们的福分。我也找不到对这圣经更真实的解释，即\BibleQuote{你的眼睛看见我尚未成形的体质}，考虑下文\BibleQuote{并且在你书卷中，一切都要被写下}。\par

% structure: p0043 source=h2 target=subsection reason=relative-heading-rank; base=h2

\subsection{第十六章}

25. 但是那生了\Person{亚伯尔}、\Person{哈诺客}、\Person{诺厄}和\Person{亚巴郎}的同一个母亲，也生了\Person{梅瑟}和继他之后的先知们，直到我们的主来临；那生了他们的母亲，也生了我们的宗徒和殉道者，以及所有好基督徒。因为所有已经出现的人，确实在不同的时代出生，但都被列入我们子民的团体；并且作为同一城邦的公民，他们经历了此世旅途的辛劳，有些人正在经历，另一些人将经历直到终结。同样，那生了\Person{加音}、\Person{含}、\Person{依市玛耳}和\Person{厄撒乌}的母亲，也生了\Person{达堂}和那同一子民中像他一样的人；那生了他们的母亲，也生了假宗徒\Person{犹达斯}和\Person{西满玛古斯}，以及所有其他假基督徒，他们直到此时仍顽固地坚持他们的肉体情欲，无论是混杂在教会的统一中，还是公开分离于分裂之中。但当这类人听到福音被宣讲，并从那些属神的人手中领受圣事时，就好像\Person{黎贝加}从她自己的子宫生了他们，如同她生了\Person{厄撒乌}一样；但当他们借着那些不真诚宣讲福音的人的手，在天主子民中被产生出来时，\Person{撒辣}确实是母亲，却是借着\Person{哈加尔}。因此，当好的、属神的门徒由那些属肉体的人的宣讲或圣洗而产出时，\Person{肋阿}或\Person{辣黑耳}确实以妻子的权利生了他们，却是从婢女的子宫。但当好的、忠信的门徒由那些在福音中属神的人所生，并且要么达到属神年龄的成熟，要么不停努力，或只是因能力不足而受阻，这些就像\Person{依撒格}出生于\Person{撒辣}的子宫，或\Person{雅各伯}出生于\Person{黎贝加}的子宫，在新生命和新约中出生。\par

% structure: p0045 source=h2 target=subsection reason=relative-heading-rank; base=h2

\subsection{第十七章}

26. 因此，不论他们似乎留在教会内，或是公开在外，凡属血肉的就是血肉，凡属糠秕的就是糠秕，无论他们坚持留在禾场上荒芜不结果，还是当诱惑临到时，像被风吹走一样被带走。甚至那以肉体的顽固与圣人的集会结合的人，也总是与那毫无玷污或皱纹的教会的合一分离。然而，我们不应绝望于任何人，无论他是表现为此类本性而在教会之内，还是更公开地从外反对教会。但属神的人，或那些以虔诚努力稳步迈向这一目标的人，不会迷途在外；因为即使由于人的某种乖谬或需要，他们似乎被驱逐出去，他们比留在教会内更受认可，因为他们绝不激动起来攻击教会，而是扎根于基督的爱德的最坚固根基，在合一的坚固磐石上。因为在此意义上，在\Person{亚巴郎}的祭献中所说的：\BibleQuote{但他没有把飞鸟分开。}\BibleRef{创 15:10}\par

% structure: p0047 source=h2 target=subsection reason=relative-heading-rank; base=h2

\subsection{第十八章}

27. 关于圣洗的争论，我认为已经论述得足够充分了；既然这个最明显的分裂被称为\OriginalTerm{多纳徒派}{Donatists}，那么在圣洗问题上，我们只剩下以虔诚的信德相信普世教会所维护的，远离分裂的亵渎。然而，如果在教会内部，不同的人在这点上仍持有不同意见，同时不破坏和平，那么直到某个大公会议通过明确而简单的法令之前，寻求合一的爱德应当遮盖人性软弱的错误，正如经上所记载的：\BibleQuote{因为爱德遮盖许多罪过。}\BibleRef{伯前 4:8}\par

28. 有关这一点，有来自\Person{真福殉道者西彼廉}的书信的极大证明——最后论及他，那些按肉身自夸拥有其权威的人，却因他的爱德而在精神上被推翻。因为在当时，在全教会尚未通过大公会议的谕令权威性地宣告就此事应采取何种做法之前，他，与大约八十位非洲教会的主教同仁一起，认为凡在\Emph{天主教会}共融之外受洗的人，在加入教会时，应重新受洗。我认为，主没有向如此卓越的人启示此事中的错误，其原因在于，他在守护教会的和平与健康时表现出的虔诚谦逊和爱德，得以彰显并被注意，从而可以说，不仅为当时的基督徒，也为后世的人，成为治愈能力的榜样。因为当一位如此重要教会的主教，本人功绩德行如此伟大，具有如此卓越的心灵和口才，对于洗礼持有一种与后来通过更勤勉探索真理所确认的观点不同的意见时；尽管他的许多同事所持的立场尚未由权威显明，但得到教会过去习惯的认可，后为整个\Emph{普世教会}所接纳；然而在这种情况下，他并没有因拒绝共融而将自己与那些持不同意见的人分开，反而不断劝勉他人，要他们\BibleQuote{以爱德互相宽容，尽力以和平的纽带保持圣神的合一}\BibleRef{弗 4:2}。因为如此，当身体的架构保持完整时，如果某些肢体出现了疾病，它更可能从整体的健康中恢复健康，而不是因为与身体分离而断绝了任何治愈照料的机会。如果他与他们分离，有多少人会跟随他！他会在人前为自己赢得何等名声！\Quote{西彼廉派}这个名字会比\Quote{多纳徒派}传播得远得多！但他不是丧亡之子，不是那些被描述为\Quote{你高举他们时，就将他们摔下}的人；他是教会和平之子，在他心灵清晰的光照中，只有一件事未能看见，只为了通过他能更卓越地看见另一件事。\Quote{然而，}宗徒说，\Quote{我还要指明一条更卓越的道路：我能说人间的语言和天使的语言，若没有爱德，我就成了鸣的锣、响的钹。}因此，他对圣事隐藏的奥秘洞察不全。但是，如果他通晓一切圣事的奥秘，却没有爱德，那也毫无价值。然而，由于他虽对奥秘洞察不全，却以一切勇毅、谦卑和信德谨慎保存爱德，他配得殉道的荣冠；这样，如果因人的软弱有任何阴影爬上了他理智的明晰，也可被他鲜血的荣耀光辉驱散。因为我们的主耶稣基督并非徒然宣告自己是葡萄树，门徒如枝条在葡萄树上，并命令那些不结果实的枝条应被砍掉，从葡萄树上移除，如同无用的枝条。\BibleRef{若 15:1}但真正的果实是什么呢？岂非那新的后裔，关于这后裔祂进一步说：\BibleQuote{我给你们一条新命令：你们要彼此相爱。}\BibleRef{若 13:34}这就是那爱德，没有它，其余的都无益处。宗徒也说：\BibleQuote{圣神的果实是爱、喜乐、平安、忍耐、慈祥、良善、忠诚、温和、节制；}\BibleRef{迦 5:22}这一切都以爱德开始，并与其余的结合形成一种奇妙团簇中的统一。主又并非徒然加上说：\BibleQuote{凡结果实的枝条，我父就修剪它，使它结更多的果实，}\BibleRef{若 15:2}这是因为那些在爱德果实上强壮的人，可能还有需要修剪的东西，是园丁不会放任不管的。因此，当这位圣人就洗礼持有一种与后来经过最勤勉考量后被彻底审查并确认的真观点不同的意见时，他的错误因他留在公教合一中以及他丰盈的爱德而得到补偿；最终，它被殉道的修枝刀清除了。\par

% structure: p0050 source=h2 target=subsection reason=relative-heading-rank; base=h2

\subsection{第十九章}

29. 但为了不让自己显得是在颂扬这位真福殉道者（其实这些颂扬不属于他，而属于那位因祂的恩宠使他成为他所是的天主），以逃避论证的责任，现在让我们从他的书信中提出那些最能堵住多纳徒派之口的证据。因为他们向无知者推崇他的权威，以显示他们那样做是好的，即当他们重新为前来投靠他们的信友施洗时。他们太可怜了——除非他们改正自己，甚至他们自己也被彻底定罪——他们在如此伟大的人所立的榜样中，选择模仿的就是那个缺点；那个缺点并没有伤害他，因为他一直稳步前行直到终点，而他们却偏离了那条有人没有认识和平的道路的道路。确实，基督的洗礼是神圣的；尽管它可能存在于异端者或分裂者之中，但它并不属于异端或分裂；因此，即使那些从那里来到天主教会的人，也不应重新受洗。然而，在这一点上犯错是一回事；那些偏离教会和平、一头栽进分裂深渊的人，却进而决定任何加入他们的人都应重新受洗，则是另一回事。因为前者是圣洁灵魂光辉上的一个斑点，丰盈的爱德本愿将其遮盖；后者则是他们卑下污秽中的污点，他们脸上对和平的憎恨炫耀地暴露出来。但关于我们进一步要考虑的、与\Person{真福西彼廉}的权威有关的主题，我们将重新开始。\par

\SourceRecord{Translated by J.R. King and revised by Chester D. Hartranft.；Nicene and Post-Nicene Fathers, First Series；Vol. 4.；Edited by Philip Schaff.；Buffalo, NY: Christian Literature Publishing Co.,；1887.；<http://www.newadvent.org/fathers/14081.htm>.}

% structure: p0001 source=h1 target=section reason=page-title

\section{\WorkTitle{论圣洗·驳多纳徒派（卷二）}}

\Emph{在这卷书中，奥斯定证明：多纳徒派引用\Person{居普良}主教与殉道者的权威是徒劳的，因为这权威实际上更反对他们而非支持公教徒。因为\Person{居普良}主张，关于在其前任\Person{阿格里皮努斯}治下，当异端者加入公教会共融时应为其施洗的看法，应仅在以下条件下接受：即应与持相反意见者保持和平，且教会合一绝不因任何裂教而破裂。}\par

% structure: p0003 source=h2 target=subsection reason=relative-heading-rank; base=h2

\subsection{第一章}

1. 在此后的书中，我打算在天主的助佑下，表明那些论据如何对我们——即对公教和平——有利，而多纳徒派声称要引述真福\Person{居普良}的权威来反对我们，以及这些论据如何反过来反对那些引述它们的人。因此，在论述过程中，如果我不得不重复我已在其他著作中说过的话（尽管我会尽量少重复），但这不应受到那些已读过并同意这些著作之人的反对；因为不仅有必要经常向迟钝之人灌输所需教导的内容，而且即使是对于那些禀赋较高之人，以多种不同方式处理并讨论同一要点，也大大有助于使他们的学习更容易、教导能力更熟练。因为我深知，当读者在手中之书中遇到任何棘手问题时，却被指望去参考另一本他恰好没有的书来寻求解答，这时他会多么沮丧。因此，若我因当前问题的紧迫而不得不重复我在其他书中已说过的话，我恳请那些已熟悉我其他著作之人的谅解，好使不熟悉者得以解除疑难；因为给予已有者，总比不去满足任何有需要之人要好。\par

2. 那么，当真理的力量堵住他们的口，而他们又不愿同意时，他们竟敢说什么呢？他们说：\Quote{\Person{西彼廉}，我们都知道他功勋卓著、学识渊博，曾在一次主教会议上与许多同僚主教各自发表意见，共同议定：所有异端者和分裂者，即所有与唯一教会的共融分离的人，都没有圣洗；因此，凡在受他们施洗后加入教会共融的人，必须在教会内受洗。} \Person{西彼廉}的权威并不使我害怕，因为他的谦逊使我安心。我们确实知道主教兼殉道者\Person{西彼廉}的伟大功绩；但它怎能比宗徒兼殉道者\Person{伯多禄}的更大呢？\Person{西彼廉}在致\Person{昆图斯}的书信中论及\Person{伯多禄}这样说：\Quote{因为主首先拣选了\Person{伯多禄}，并在他身上建立了祂的教会\BibleRef{玛 16:18}，后来\Person{保禄}就割损问题与他争辩时，\Person{伯多禄}并没有傲慢或狂妄地自以为是，声称自己拥有首席权，那些后来才来的人应当更加服从他。他也没有因为\Person{保禄}从前是教会的迫害者而轻视他，反而接纳了真理的劝告，欣然同意\Person{保禄}所持的合理理由；由此给我们留下了和谐与忍耐的榜样，叫我们不要顽固地喜爱自己的意见，而应将弟兄和同僚们为共同健康与益处所提出的任何真实而合理的建议视为己有。} 这里有一段话，\Person{西彼廉}记载了我们也在圣经中所学到的：宗徒之长\Person{伯多禄}，其宗徒首席地位闪耀着如此卓越的恩宠，却因在割损问题上采纳了一种有违真理要求的习俗，而被后来的宗徒\Person{保禄}所纠正。因此，如果\Person{伯多禄}在某个方面未能按福音真理正直行事，以致迫使外邦人犹太化，正如\Person{保禄}在那封呼求天主见证他并不说谎的书信中所写的那样；因为他（\Person{保禄}）说：\BibleQuote{我现在写给你们的，你看，我在天主面前，并不说谎；}\BibleRef{迦 1:20}在这样神圣而可畏地呼求天主见证后，他讲述了全部经过，其中说道：\BibleQuote{但我一看见他们行事不正直，不按福音的真理，就当众对\Person{伯多禄}说：\InnerQuote{你既然是犹太人，却按照外邦人的方式生活，而不是按照犹太人的方式，怎么还能强迫外邦人按照犹太人的方式生活呢？}}\BibleRef{迦 2:14}——我说，如果\Person{伯多禄}可以违反后来教会所持的真理准则，迫使外邦人按犹太人的方式生活，那么，为什么\Person{西彼廉}不能违反后来普世教会所持的信德准则，迫使异端者和分裂者重新受洗呢？我认为，将\Person{西彼廉}与\Person{伯多禄}就其殉道荣冠相比较，并没有贬低\Person{西彼廉}；倒是我应该担心这样是否对\Person{伯多禄}不敬。因为谁不知道，宗徒职位的首席地位应当胜过任何主教职呢？但是，即使他们宗座尊严有高下之分，他们在殉道中却有同样的荣耀。至于那些因爱德合一而宣认并死于真信德的人，其心灵是否在不同方面彼此超过，只有主自己知道——祂恩宠的隐藏而奇妙的安排，使十字架上的强盗在临终时一次宣认祂，并在同日被送往乐园\BibleRef{路 23:40}，而跟随主的\Person{伯多禄}却三次否认祂，他的荣冠被推迟了\BibleRef{玛 26:69}。我们若凭表面证据作判断，未免轻率。但若有人现在被发现强迫一个人按照犹太习俗接受割损，作为洗礼的必要预备，这会引起世人更普遍的反对，远比强迫一个人重新受洗更为人反对。因此，如果\Person{伯多禄}因做了这事而被后来的同僚\Person{保禄}纠正，却因和平与合一的纽带而被保全，直至被提升至殉道，那么，我们岂不应更加欣然且坚定地优先考虑由普世教会法令所确立的规范，而不是一位主教或一个省区会议的权威？因为这位\Person{西彼廉}，在主张自己观点时，仍然渴望与那些在这点上与他意见不同的人保持和平的合一，这从他本人被多纳徒派引用的那次主教会议的开场白中可以看出。其内容如下：\par

% structure: p0006 source=h2 target=subsection reason=relative-heading-rank; base=h2

\subsection{第二章}

3. \Quote{当九月初一，来自\Place{阿非利加}、\Place{努米底亚}和\Place{茅利塔尼亚}诸省的许多主教，连同他们的司铎和执事，在\Place{迦太基}聚集时，大部分平信徒也在场；当\Person{犹巴亚努斯}致\Person{西彼廉}的信，以及\Person{西彼廉}给\Person{犹巴亚努斯}关于为异端者施洗的答复，都被宣读之后，\Person{西彼廉}说：\InnerQuote{最亲爱的同僚们，你们已经听到我们的同僚主教\Person{犹巴亚努斯}写信给我，就异端者非法且亵圣的圣洗征询我微薄的能力，以及我所给的答复——我给出了一项我们一而再、再而三地经常给出的判断：来到教会的异端者应当受洗，并以教会的圣洗得到圣化。\Person{犹巴亚努斯}的另一封信也已读给你们听，在那封信中，他出于真诚而虔诚的热忱，答复我们的书信，不仅表示同意，而且还致感谢，承认他受到了教导。接下来需要我们各自就此问题发表意见，不判断任何人，也不剥夺任何人与我们意见不同者共融的权利。因为我们中间没有一人自封为主教中的主教，或以暴虐的恐怖强迫同僚必须服从，因为每位主教都自由运用其自主权，有权形成自己的判断，既不能受他人判断，也不能判断他人。但我们都要等待我们主\Person{耶稣基督}的审判，唯有祂有权将我们安置在祂教会的管理中，并判断我们在其中的行为。}}\par

% structure: p0008 source=h2 target=subsection reason=relative-heading-rank; base=h2

\subsection{第三章}

4. 现在让异端者骄傲膨胀的颈项昂起，如果他们敢，反对这篇讲话的神圣谦卑。你们疯狂的\OriginalTerm{多纳徒派}{Donatists}，我们热切渴望你们回归圣教会的和平与统一，以便在其中获得健康，你们对此有何话说？你们惯于拿\Person{西彼廉}的书信、他的意见、他的会议来反对我们；为什么你们为你们的裂教引用\Person{西彼廉}的权威，却拒绝他促进教会和平的榜样？但是谁不知道，新旧约的圣经正典被限制在其自身界限内，它绝对高踞于所有后来的主教书信之上，以至于我们对于其中所包含的是正确和真实这一点，不能有任何疑问或争议；但是所有圣经正典结束以后所写的或正在写的主教书信，如果其中包含任何偏离真理的内容，都可以被某位在此事上比他们更智慧的人的话语所驳斥，或者被其他主教更重要的权威和更有学识的经验所驳斥，被会议的权威所驳斥；而且，在各个地区和省份举行的会议本身，必须毫无疑问地服从为整个基督教世界举行的大公会议的权威；并且，即使是大公会议，较早的也常常被后来的所纠正，当通过实际经验，之前隐藏的事情被揭示出来，之前隐藏的变得已知时，这一切都没有亵渎的骄傲的旋风，没有傲慢的颈项膨胀，没有嫉妒仇恨的争斗，只有神圣的谦卑、公教的和平和基督徒的爱德。\par

% structure: p0010 source=h2 target=subsection reason=relative-heading-rank; base=h2

\subsection{第四章}

5. 因此，圣\Person{西彼廉}，他的尊荣因谦卑而增加，他如此热爱伯多禄的榜样，以至于使用了这些话：\Quote{由此赐给我们和谐与忍耐的榜样，叫我们不要固执地喜爱自己的意见，而应将我们弟兄和同僚的任何真实而正确的建议视为自己的，为共同的健康与益处。}——他，我说，充分表明他非常愿意纠正自己的意见，如果有人向他证明，正如脱离羊圈的人不能失去基督的洗礼一样，他们也能施予基督的洗礼，这一点是确定的；关于这个问题我们已经说了很多。而且我们自己也不敢断言任何此类事情，除非我们得到整个教会一致权威的支持，如果那时这个问题已经通过大公会议的调查和法令被确认为真理，他本人无疑也会服从。因为如果他引用伯多禄作为榜样，让他安静而和平地接受一位较年轻同僚的纠正，那么他自己与他的省会议会多么更容易服从整个世界的权威，当真理这样被揭示出来时？事实上，如此神圣而和平的灵魂会非常愿意听从任何能够向他证明真理的人的理由；也许他甚至这样做了，尽管我们不知道。因为当时在主教之间发生的所有事情既不可能都被记录下来，我们也不了解所有被记录的事情。因为一个被如此多的辩论迷雾所笼罩的问题，怎么可能得到大公会议的充分光照和权威决定，除非它首先在各个地区的许多讨论和比较主教观点的基础上经过相当长时间的讨论？这就是和平健全的一种效果：当有疑问的问题被长期研究，并且由于难以达到真理，在弟兄般的讨论中产生意见分歧，直到人们最终达到纯粹的真理时，团结的纽带仍然存在，以免在被切断的部分中发现致命错误的不可治愈的伤口。\par

% structure: p0012 source=h2 target=subsection reason=relative-heading-rank; base=h2

\subsection{第五章}

6. 常有这样的事：真理不完全启示给更有学识的人，好使他们那带来更大果实的忍耐和谦逊的爱德得到考验，无论是他们在对待比较隐晦的事情持有不同意见时如何保存合一，还是他们在得知真理被宣告为与自己所想相反时以何种态度接受真理。而这两方面中，我们在\Person{圣西彼廉}身上看到了其中一面，即他与那些意见不同的人保存合一的途径。因为他曾说：\Quote{不判断任何人，也不剥夺任何人与我们不同意见者的共融权。}另一面，即他在发现真理与自己所想不同时能以什么态度接受真理，虽然他的书信对此沉默，但他的功德却宣扬了这一点。如果没有现存的书信证明，他的殉道冠冕见证了它；如果主教会议没有宣布，众天使的军旅宣布了它。因为一个最和平的灵魂赢得殉道的冠冕，是在那他不愿分离、即使意见不同也不离开的合一中，这并非小小的证明。因为我们不过是人；因此在任何事情上持有与真理不同的观点，是人性所难免的诱惑。但因过于偏爱自己的意见，或因嫉妒比自己高明的人，甚至达到分裂教会共融并建立异端或分裂的亵圣地步，那是一种与魔鬼相称的狂妄。但在任何事情上永不离真理持有不同意见，这种完美只在天神身上找到。既然我们是人，但就希望而言，我们是天神，我们在复活时也将与他们一样\BibleRef{玛 22:30}，无论如何，只要我们尚缺天神的完美，至少让我们摆脱魔鬼的狂妄。因此宗徒说：\BibleQuote{你们所受的试探，无非是普通人所能受的。}\BibleRef{格前 10:13}因此，人有时犯错属于本性。因此他在另一处说：\BibleQuote{我们凡是成熟的人，就该有这种心情；若在什么事上你们有别的想法，天主也要启示给你们。}\BibleRef{斐 3:15}但祂按自己的旨意（无论是在今生还是在来世）将启示赐给谁呢？岂不赐给那些行在和平之道上、不偏离进入任何分裂的人吗？不是赐给那些不认识和平之道，或因其他原因破坏了合一纽带的人。因此，当宗徒说：\BibleQuote{若在什么事上你们有别的想法，天主也要启示给你们。}免得他们以为除了和平之道外，他们的错误观点也可能被启示给他们，便立刻补充说：\BibleQuote{但是，我们既已达到这一步，就该按这准则前进。}\BibleRef{斐 3:16}而\Person{圣西彼廉}，遵循这准则，凭借最坚韧的容忍，并非仅仅因为他流血殉道，而是因为他在合一流血（因为即使他舍身被焚，却没有爱德，对他毫无益处\BibleRef{格前 13:3}），通过殉道的宣认达到天神的光明中，并且即使之前没有，至少在那时，承认了在那一点上真理的启示：他在错误中，并未将坚持错误观点置于合一的纽带之上。\par

% structure: p0014 source=h2 target=subsection reason=relative-heading-rank; base=h2

\subsection{第六章}

7. 那么，你们这些多纳徒派，对此有何话说？如果关于圣洗我们的意见是真的，而所有在\Person{圣西彼廉}时代持有相反意见的人并未被从教会合一中剪除，直到天主将那点错误的真理启示给他们，那么你们为何通过亵圣的分裂打破了和平的纽带？但如果关于圣洗你们的意见是真的，\Person{圣西彼廉}和其他人——你们声称他与那些人一起召开了这样一次公会议——仍然与那些持有不同意见的人保持合一，那么你们为何打破了和平的纽带？任选其一，你们不得不作出反对你们分裂的判断。回答我，你们为何分裂自己？你们为何在与全世界对立的地方建起一座祭台？你们为何不与那些宗徒书信所寄送的教会共融？你们自己阅读并承认这些书信，并说你们按照它们的原则生活。回答我，你们为何分裂自己？我想是为了避免因与恶人共融而丧亡。那么，\Person{圣西彼廉}和他的许多同僚是如何没有丧亡的呢？因为他们虽然相信异端者和分裂者不拥有圣洗，但仍然选择与那些未受圣洗便被接纳入教会的人保持共融，尽管他们相信这些人昭彰的亵圣之罪还在他们头上，也不愿与教会的合一分离，正如\Person{圣西彼廉}的话：\Quote{不判断任何人，也不剥夺任何人与我们不同意见者的共融权。}\par

8. 因此，如果义人因与恶人共融而必致丧亡，那么教会早在\Person{西彼廉}时代就已经丧亡了。那么多纳徒的起源从何而来？他是在哪里受教、在哪里受洗、从哪里被祝圣的，既然教会早已因与恶人共融的传染而毁灭？但如果教会仍然存在，恶人就不能在与义人的共融中伤害义人。那么你们为何分离？看，我看见在合一中的\Person{西彼廉}及其同僚们，他们召开会议，决定那些在教会共融之外受洗的人没有真正的洗礼，因此当他们加入教会时必须给予洗礼。但再次，我看见在同一个合一中，有些人对这件事持有不同看法，他们承认从异端和裂教者那里来的人有基督的洗礼，不敢再为他们施洗。所有这些人，公教会用她慈母的胸怀拥抱，轮流彼此担当重担，竭力保持圣神所赐的合一，与和平的联系 \BibleRef{弗 4:3}，直到天主向其中一方启示他们在观点上的任何错误。如果一方持有真理，他们是否被另一方感染？或者，如果另一方持有真理，他们是否被第一方感染？由你选择。如果有污染，那么教会甚至在那时就已经不复存在了；所以告诉我，你们是从何处出来的？但如果教会仍然存在，善人在这种共融中决不会被恶人污染；因此告诉我，你们为何要破坏这纽带？\par

9. 或许可能这样：裂教者若未经洗礼而被接纳，并不带来污染，而是由那些交付圣书者带来污染？因为你们当中有交付圣书者，这由最清楚的历史证据所证明。如果你们当时对你们所控告的人提出了真实的证据，你们本应在全世界合一之前证明你们的理由，这样你们就会被保留而他们被排除。如果你们努力这样做却没有成功，世界不应受责备，因为它信任教会的审判者而不是诉讼中的败诉方；而如果你们没有提出诉讼，世界同样不应受责备，因为它不能未经听审就定人的罪。那么，你们为何与无辜者分离？你们无法为你们裂教的亵渎行为辩护。但我略过这点。不过我要说：如果交付圣书者能够污染你们（他们并未被你们定罪，反而你们被他们击败），那么裂教者和异端者的亵渎行为，按照你们的主张，未经洗礼而被接纳进入教会，岂不是更能污染\Person{西彼廉}？然而他并没有分离自己。既然教会继续存在，显然它不能被污染。那么，你们为何要分离？我并不是说从无辜者分离（事实证明了他们是无辜的），而是从从未证明是交付圣书者的那些人分离？交付圣书者的罪，如同我开始说的，比裂教者的罪更重吗？让我们不要引入欺骗的天平，我们可以随意放上我们想要的砝码，随意说\Quote{这是重的，这是轻的}，而是让我们从神圣的圣经中拿出神圣的天平，如同从主的宝库中拿出，用它来衡量，看哪个更重；或者更确切地说，让我们不要自己衡量，而是诵读主所宣告的重量。当主以近期的惩罚为例，表明需要防备古时的罪之时，偶像被制造和崇拜，先知书被嘲笑国王的愤怒焚烧，裂教被尝试，偶像崇拜被刀剑惩罚\EditorialNote{出谷纪第三十二章}，焚烧书卷被战争中的杀戮和异国俘虏惩罚\EditorialNote{耶肋米亚第三十六章}，裂教被大地裂开、活活吞下裂教领袖而其余人被天火烧尽惩罚\EditorialNote{户籍纪第十六章}。现在谁会怀疑那个受到更重惩罚的罪是更严重的？如果来自如此亵渎团体的人，按照你们的主张，未经洗礼，都不能污染\Person{西彼廉}，那么那些未被定罪而只是被怀疑交付圣书的人，又怎能污染你们？因为即使他们不仅交出书卷让其被焚烧，而且亲手烧毁它们，他们犯的罪也比犯裂教者轻；因为裂教受到更重的惩罚，而另一种受到较轻的惩罚，这不是出于人的判断，而是出于天主的审判。\par

% structure: p0018 source=h2 target=subsection reason=relative-heading-rank; base=h2

\subsection{第七章}

10. 那么，你们为什么分裂了？如果你们还有一点理智，必定会看到你们对这些论据找不出任何可能的回答。他们说：\Quote{我们并非完全没有办法，我们仍能回答说：\InnerQuote{这是我们的意愿。}\BibleQuote{你是谁，竟判断别人的仆人？他或站立或跌倒，自有他的主人在。}}\BibleRef{罗 14:4} 他们不明白这话是对那些愿意判断别人内心（而非公开事实）的人说的。因为宗徒本人怎能如此多地论及分裂和异端的罪？或者，圣咏中那节经文：\Quote{你们如果真爱正义，就该公平判断，人子啊！}又作何解释？但是，主亲自说：\BibleQuote{不要凭外表判断，但要公义判断}\BibleRef{若 7:24}，如果我们不能判断任何人？最后，对于那些他们不公判断的叛教者，他们自己怎敢判断别人的仆人？他们是站是倒，自有主。或者，在\Person{马克西米安}追随者的最近案例中，他们如何毫不迟疑地提出一个全体会议（据他们声称）以无误权威所作的判决，以至于将那些人比作最初被活吞的分裂者？然而，他们自己（他们无法否认）要么定罪了无辜者，要么接纳了有罪者。但当真理被提出而他们无法反驳时，他们咕哝着真正有益的喃喃：\Quote{这是我们的意愿：\BibleQuote{你是谁，竟判断别人的仆人？他或站立或跌倒，自有他的主人在。}} 但是当一只软弱的羊在旷野中被发现，而本应把它召回羊圈的牧人却不见踪影时，他们就咬牙切齿，折断那软弱的脖子：\Quote{你若不是叛教者，就是个好人。关心你灵魂的得救；成为基督徒。} 多么无理的疯狂！当对一个基督徒说\Quote{成为基督徒}时，除了否认他是基督徒之外，还教导了别的什么吗？那些基督的迫害者想要教导的不正是同一课吗？反抗他们才赢得了殉道的荣冠？我们是否应当认为，以刀剑威胁所犯的罪比以舌头诡诈所犯的罪更轻？\par

11. 回答我，你们这些凶残的豺狼，你们寻求披着羊皮（\BibleRef{玛 7:15}），以为真福\Person{西彼廉}的书信对你们有利。分裂者的亵渎玷污了\Person{西彼廉}吗？如果没有，那么别人犯的罪怎么可能玷污无辜者，如果分裂的亵渎都不能玷污他们？那么，你们为什么分裂了？你们为何躲避较轻的过错（那是你们自己的发明），却犯了最重的过错，即分裂的亵渎？你们现在会不会承认，那些在教会共融之外或在某种异端或分裂中领受了圣洗的人不再是分裂者或异端者，因为通过回归教会，弃绝他们先前的错误，他们已不再是他们从前所是的？那么，如何解释，虽然他们没有领受圣洗，他们的罪却没有留在他们头上？是因为圣洗是基督的，但如果没有教会共融，它就不能使他们获益；然而当他们回归，弃绝过去的错误，通过覆手被接纳进入教会共融时，他们现在在爱德中扎根并建立（没有爱德，一切皆无益处），才开始从那个圣事中领受罪过的赦免和生命的圣化，而他们当初在教会之外徒然拥有之？\par

12. 那么，请不要再提出\Person{西彼廉}的权威来支持重复圣洗，但要与我们一同持守\Person{西彼廉}维护合一的榜样。因为圣洗的问题当时尚未完全解决，但教会遵守了最有益的习惯：纠正错误的，但不重复已施行的，即使是对分裂者和异端者也是如此：她治愈受伤的部分，但不干涉完整的部分。这个习惯，我认为来自传承（如同许多其他事情，被认为是在实际认可下传下来的，因为它们在整个教会中得以保存，尽管既没有出现在他们的书信中，也没有出现在他们继任者的会议中）——我说，这个最有益的习惯，按照圣\Person{西彼廉}的说法，开始被他的前任\Person{阿格里皮努斯}所谓地修正了。但是，根据来自对真理更仔细研究的教导，在经历极大的怀疑和摇摆之后，最终由全体会议作出决定，我们应当相信，它被\Person{阿格里皮努斯}开始败坏，而非得到修正。因此，当如此重大的问题压在他身上，难以决定：罪过的赦免和人的精神重生能否发生在异端者或分裂者之中，而\Person{阿格里皮努斯}的权威（以及与他共同误解此事的少数人的权威）引导着他，宁愿尝试新事物而不愿维护他们不知道如何辩护的习惯；在这种情况下，或然性的考虑闯入他灵魂的眼目，阻碍了对真理的彻底探讨。\par

% structure: p0022 source=h2 target=subsection reason=relative-heading-rank; base=h2

\subsection{第八章}

13. 我也不认为有福的西彼廉在自由表达和较早说出他反对教会惯例的想法时，有任何其他动机，除了他应当感恩地接受任何能发现拥有更完整真理启示的人，并且他应当树立一个可效法的榜样，不仅在教导的勤勉上，也在学习的谦逊上；但是，如果找不到任何人提出任何论据来驳斥那些或然性的考虑，那么他应当坚持他的意见，完全意识到他既没有隐藏他认为是真理的东西，也没有破坏他所爱的合一。因为他如此理解宗徒的话：\BibleQuote{让先知们两个人或三个人说话，其余的人评判。若有什么启示给坐在旁边的另一个人，那先说话的就当安静。}\BibleRef{格前 14:29}\Quote{在这段话中，他教导并指明，许多事情为更好而被启示给个人，我们每个人不应固执于他一度吸收和坚持的东西，但如果有什么显得更好、更有用，他应当欣然接受。}无论如何，用这些话，他不仅劝告那些看不到更好途径的人同意他，也鼓励任何有能力的人提出论据，以便可以确认维持旧惯例的做法；如果这些论据的性质不允许被驳倒，他就可以亲自表明他所说的\Quote{我们每个人不应固执于他一度吸收和坚持的东西，但如果有什么显得更好、更有用，他应当欣然接受}是多么真诚。但是，既然除了单纯向他坚持惯例的人之外，没有人出现，而且他们提出支持惯例的论据也不足以说服像他这样的灵魂，这位伟大的辩论家并不满足于放弃他的意见——这些意见虽然不真实，但他自己无法看见，至少没有被驳倒——而去支持一个虽然有真理支持但尚未被确认的惯例。然而，如果不是他的前任阿格黎皮诺和在非洲的一些主教同伴们，甚至通过一次公会议的决定，首先引诱他离开这个惯例，他肯定不敢反对它。但是，在如此晦涩问题的困惑中，他看到周围到处是强大的普遍惯例，他宁愿通过祈祷和将心神伸向天主来约束自己，以便察觉到或教导那此后由全体大公会议决定为真理的东西。但是，当他在由阿格黎皮诺主持的前一次公会议的权威中找到了缓解疲惫的方法时，他宁愿坚持他那可以说是他前任发现的东西，而不愿再花力气去探究。因为，在他写给克文图斯的信末尾，他这样表明他如何在权威的休息处，如果可以这样说，为他的疲惫寻求安息。\par

% structure: p0024 source=h2 target=subsection reason=relative-heading-rank; base=h2

\subsection{第九章}

14. \Quote{而且，}他说，\Quote{阿格黎皮诺，一位记忆杰出的人，连同与他一起的主教们，他们当时在非洲和努米底亚省管理主的教会，经过共同公会议的调查研究加以权衡，确立并证实了这一点；他们的判决，既是宗教性的、合法的、有益的，符合公教信仰和教会，我们也予以遵循。}通过这个见证，他充分证明，如果举行过任何讨论此事的公会议，无论是涵盖整个教会，还是至少代表我们海外的弟兄们，他会多么乐意作证。但这样的公会议尚未举行，因为整个世界被惯例的强力纽带联系在一起；这被认为足以反对那些想要引入新事物的人，因为他们无法理解真理。然而后来，当问题在许多方面得到讨论和研究时，新做法不仅被发明，而且还被提交给一次全体大公会议的权威和权力——这是在西彼廉殉道之后，但发生在我们出生之前。但这确实是教会的惯例，后来由一次全体大公会议确认，在会议中真理被揭示，许多困难被清除，从有福的西彼廉本人写给犹巴安的同一封信中可以看出，这封信被引述为在公会议中宣读。因为他说：\Quote{但有人问，对于那些从异端出来进入教会，已经未经洗礼被接纳的人，该怎么办呢？}这里他显然充分说明了通常的做法，尽管他宁愿事情不是这样；而就在他引用阿格黎皮诺公会议的事实本身，他就清楚证明教会的惯例是不同的。事实上，他无需通过这次公会议来寻求建立这个做法，如果惯例已经批准；而且在公会议本身中，一些发言者在发表意见时明确声明，他们在决定他们认为正确的事情时，是反对教会的惯例的。因此，让多纳徒派考虑这一点，这无疑是任何人都能看到的事实：如果要遵循西彼廉的权威，那么更应该遵循它来维持合一，而不是改变教会的惯例；但如果尊重他的公会议，它必须让位于后来的普世教会的公会议，西彼廉很高兴成为其成员，经常警告他的同僚，他们应当效法他维护整个身体的连贯性。因为后来的公会议在后代中总是被优先于较早的公会议；而且整个整体，出于好的理由，总是被视为优于部分。\par

% structure: p0026 source=h2 target=subsection reason=relative-heading-rank; base=h2

\subsection{第十章}

15. 然而，当他们得知\OriginalTerm{圣西彼廉}{Cyprian}虽然自己并不接纳那些在异端或分裂中受洗的人为教会成员，却因为他明确声明\Quote{不判断任何人，也不剥夺任何人与我们意见不同者的共融权}而与那些接纳他们的人保持共融时，他们又采取什么态度呢？如果他与这样的人共融是被玷污的，那么他们在圣洗的问题上为何还听从他的权威？但若他与他们共融并未被玷污，他们为何不效法他维护合一？除了\Quote{我们偏要这样}之外，他们还有什么可辩护的呢？面对真理或正义的谴责，任何有罪或邪恶的人——例如好色者、醉汉、奸淫者、各种不洁者、窃贼、强盗、杀人者、掠夺者、作恶者、偶像崇拜者——当真理之声定他们的罪时，除了\Quote{我偏要做这事}和\Quote{我偏要这样喜欢}之外，他们还有什么回答呢？如果他们还有点基督徒的味道，还会进一步说：\BibleQuote{你是谁，竟敢判断别人的仆人？}\BibleRef{罗 14:4} 然而这些人尚存更多谦逊之心，当按照神法和人法，他们因放荡的生活和行为受到惩罚时，并不自称为殉道者；而\OriginalTerm{多纳徒派}{Donatists}却同时希望过着亵渎神的生活，享有无可指责的名声，不为恶行受罚，并在应受的惩罚中获取殉道者的荣耀。仿佛他们没有体验到更大的仁慈和忍耐，正如祂\BibleQuote{逐步执行审判，给他们悔改的机会}，\BibleRef{智 12:10} 且不断加重今生的鞭策；使他们思考自己所受的苦和受苦的原因，从而及时变得明智；并使那些为维护\OriginalTerm{多纳徒}{Donatus}的合一而接受了原属于\OriginalTerm{马克西米安}{Maximianus}派圣洗的人，更加乐意接受全世界的圣洗以维护基督的和平；使他们归于根源，与教会的合一和好，看到自己再无话可说，尽管仍有事要做；使他们为过去的行为，向长久忍耐的天主献上慈爱的祭献——他们曾以邪恶的罪过破坏了祂的合一，对祂的圣事造成了持久的伤害。因为\Quote{主是慈悲宽仁的，缓于发怒，富于慈爱和信实。} 愿他们今世拥抱祂的仁慈和忍耐，来世敬畏祂的信实。因为\BibleQuote{祂不愿罪人死去，而愿他离开恶道得以生活}，\BibleRef{则 23:11} 因为祂使审判转向针对加于祂的过犯。这就是我们的劝勉。\par

% structure: p0028 source=h2 target=subsection reason=relative-heading-rank; base=h2

\subsection{第十一章}

16. 因此，我们视他们为仇敌，因为我们宣讲真理，因为我们不敢缄默，因为我们害怕退缩，不尽力坚持我们的论点，因为我们服从宗徒所说的：\BibleQuote{宣讲圣言；无论顺境逆境，都要坚持；以各种教训，指摘、斥责、劝勉。}\BibleRef{弟后 4:2} 但正如福音所说：\BibleQuote{他们爱人间的称赞胜过天主的称赞。}\BibleRef{若 12:43} 他们害怕一时受责备，却不害怕永远受罚。他们自己也看到自己所犯的过错；他们看到自己无言以对，却以迷雾迷惑无经验者，而自己却活活被吞噬——也就是说，明知故犯地走向灭亡。他们看到人们惊骇、憎恶他们分裂成许多派系，尤其在非洲的首都和名城\Place{迦太基}；他们试图修补他们破衣烂衫的耻辱。他们以为可以消灭\OriginalTerm{马克西米安}{Maximianus}的追随者，便通过\OriginalTerm{奥普塔图斯·基尔多尼亚努斯}{Optatus the Gildonian}重重压迫他们；在最残酷的迫害中给他们施加了许多伤害。然后他们又接纳了一些人，以为所有人都能在同样的恐怖下转变；但他们对所接纳的人不愿做错事，即重新为那些在分裂中受他们施洗的人施洗，或者不如说由那些在外面施洗的同一个人在他们共融内重新施洗，从而立即对他们自己的不敬虔惯例开了一个例外。他们感到自己攻击全世界的圣洗是多么邪恶，因为他们已接受了\OriginalTerm{马克西米安}{Maximianus}追随者的圣洗。但他们害怕那些被他们自己重新施洗的人，唯恐自己向别人显示仁慈后，却得不到他们的仁慈；唯恐这些人在他们停止毁灭他人灵魂后，要他们为自己的灵魂交账。\par

% structure: p0030 source=h2 target=subsection reason=relative-heading-rank; base=h2

\subsection{第十二章}

17. 关于他们所接纳的马克西米安派追随者，他们能给出什么回答，他们无法猜测。如果他们說：\Quote{我们所接纳的那些人是无辜的}，那么显见的回答是：\Quote{那么你们曾定了无辜者的罪。}如果他们說：\Quote{我们出于无知而如此行}，那么你们就是妄加判断（正如你们曾对叛教者妄加判断），而你们宣称的\Quote{你们必须知道，他们曾被一个全体公会议的真实声音所定罪}是假的。因为真实的判断决不会定无辜者的罪。如果他们說：\Quote{我们并没有定他们的罪}，那么只需引用那次会议，引用主教和城市的名录即可。如果他们說：\Quote{那次会议本身不是我们的}，那么我们就引用总督省的文件，在那里他们不止一次引用同一次会议来证明将马克西米安派追随者排除出大殿是正当的，并用法官的喧嚣和他们盟友的力量来使他们困惑。如果他们說：后来被他们接纳的穆斯蒂的费利西安和阿萨瓦的普雷特克斯塔图斯并不属于马克西米安派，那么我们就引用记录，其中他们曾在法庭上要求将这些人从他们针对马克西米安派举行的会议中排除出去。如果他们說：\Quote{为了和平的缘故，我们接纳了他们}，我们的回答是：\Quote{那么你们为何不承认唯一真实而圆满的和平呢？是谁促使、是谁强迫你们接纳一个你们已定罪的裂教徒，以维护多纳图的和平，而违背基督的和平，未经听证就定全世界的罪？}真理从各方面困住他们。他们看见没有回答的余地，以为也没有别的事可做；他们不知该说什么。他们不被允许沉默。他们宁可悖逆地言语与真理争辩，也不愿借认罪而恢复和平。\par

% structure: p0032 source=h2 target=subsection reason=relative-heading-rank; base=h2

\subsection{第十三章}

18. 但谁不明白他们心中可能說些什么呢？他们说：\Quote{那么，对于那些我们已经重洗的人，我们该如何办呢？}与他们一同回到教会吧。将你们所伤的人带来，借和平之药医治；将你们所杀的人带来，借爱德之生命复活。兄弟般的合一在平息天主方面大有能力。我们的主說：\Quote{若你们中二人，在地上同心合意，无论为什么事祈祷，必蒙我在天之父的成全。} \BibleRef{玛 18:19} 如果二人同心，何况两个团体呢？让我们一同俯伏在主面前。你们与我们分享合一，我们与你们分享痛悔，让爱德遮盖许多罪过。\BibleRef{伯前 4:8} 请教真福的西彼廉自己吧。看他何等重视合一之结的恩赐，他并未使自己与意见不同者分离而避免与他们共融；尽管他认为在教会共融之外受洗者没有真正的洗礼，他仍愿意相信，仅仅通过进入教会，因着合一之结，他们就可以被接纳分享宽恕。因为他在写给犹巴安的信中这样解决了自己提出的问题：\Quote{但有人会说：\InnerQuote{那么，那些过去从异端来到教会，未经洗礼就被接纳的人，会怎样呢？}主能凭祂的仁慈赐予宽恕，不将那些因单纯而被接纳进入教会、且在教会中安眠的人，从祂教会的恩赐中分离。}\par

% structure: p0034 source=h2 target=subsection reason=relative-heading-rank; base=h2

\subsection{第十四章}

19. 但究竟是没有受洗更糟，还是受两次洗更糟，这很难断定。我确实看到哪一种更违反人的感受；然而当我求助于神圣的天平，在那里事物的分量不是由人的感觉，而是由天主的权威所决定时，我发现我们的主在两方面都有話说。祂对伯多禄說：\Quote{洗过澡的人，不需要再洗第二次。} 而对尼苛德摩說：\Quote{人除非由水和圣神而生，不能进天主的国。} \BibleRef{若 3:5} 至于天主更隐秘的旨意是什么，对我们这样有限的人来说或许难以得知；但就字面而言，任何人都能看到\Quote{不需要再洗}与\Quote{不能进天国}之间的差别。最后，教会本身也持守其传统，即没有洗礼，她不能接纳任何人到她的祭台前；但既然允许重洗者在补赎后得以被接纳，这显然证明他的洗礼被认为有效。因此，如果西彼廉认为那些他以为未受洗的人，因合一之结而获得部分赦免，那么主也有能力借单纯的合一与和平之结，与被重洗者修好，并借和平的同一补偿能力，减轻对那些重洗他们之人的不悦，赦免他们在错误中犯下的一切过失，只要他们献上爱德的祭献，这祭献遮盖许多罪过；这样，主不看那些因分离而受伤的人数，而看那因回头而脱离束缚的更大数目。因为在同一个和平之结中，西彼廉曾认为，那些他以为未经洗礼进入教会的人，因天主的仁慈，并未被剥夺教会的恩赐；我们也相信，借同一仁慈，被重洗者能获得祂的宽恕。\par

% structure: p0036 source=h2 target=subsection reason=relative-heading-rank; base=h2

\subsection{第十五章}

20. 既然天主教会，无论是在真福\Person{西彼廉}的时代，还是在他之前的更早时代，在她的怀中都包含了一些受过重行圣洗的人和一些未受圣洗的人，那么这两部分中，无论哪一部分，都只能借单纯的合一之力获得救恩。因为如果那些从异端者中过来的人没有受过圣洗，如同\Person{西彼廉}所断言，他们就没有被正当地接纳进入教会；然而他自己并没有对他们的获得天主的怜悯所赐的赦罪，因着教会的合一，感到失望。同样，如果他们已受过圣洗，就不应该对他们重行圣洗。那么，除了那以合一为乐的同一位爱德之外，还有什么能帮助另一部分人呢？以至于因人的软弱而隐藏的事，在圣事的考量上，可能不被天主的怜悯算为那些爱好和平之人的罪？那么，为什么你们一面害怕你们所重行圣洗过的人，一面又吝啬于让你们自己和他们都进入救恩呢？曾经有一段时间，关于圣洗的问题存在疑惑；持不同意见的人仍然保持合一。随着时间的推移，由于对真理的确实发现，那疑惑被消除了。那个尚未解决时没有吓倒\Person{西彼廉}使他与教会分离的问题，现在既然已经解决，就邀请你们再次回到羊栈内。来到和谐一致的天主教会，\Person{西彼廉}即使在疑惑困扰时也没有离弃它；或者，如果现在你们对\Person{西彼廉}的榜样不满，他与那些以异端的圣洗被接纳的人保持共融，公开声明我们应当\Quote{既不要判断任何人，也不要因某人意见与我们不同而剥夺他共融的权利}，你们这些可怜的人，要往哪里去呢？你们在做什么？你们必定要从你们自己那里逃离，因为你们已经前进到他曾经停留的位置之外了。但如果因他丰富的爱德、他对弟兄友爱的热爱以及和平的联系，他自己的罪和他人的罪都不能成为他的障碍，那么你们就回到我们这里来吧，在这里你们会发现在我们或你们前进的道路上，来自你们一派所虚构的障碍要少得多。\par

\SourceRecord{Translated by J.R. King and revised by Chester D. Hartranft.；Nicene and Post-Nicene Fathers, First Series；Vol. 4.；Edited by Philip Schaff.；Buffalo, NY: Christian Literature Publishing Co.,；1887.；<http://www.newadvent.org/fathers/14082.htm>.}

% structure: p0001 source=h1 target=section reason=page-title

\section{论圣洗，驳多纳徒派（第三卷）}

\Emph{奥斯定着手驳斥那些可能从\Person{西彼廉}写给\Person{犹拜安}的信中引申出来、用以粉饰异端者不能施行基督圣洗这一观点的论据。}\par

% structure: p0003 source=h2 target=subsection reason=relative-heading-rank; base=h2

\subsection{第一章}

我认为现在每个人都应清楚，真福\Person{西彼廉}的权威，为维护和平的纽带，避免违背那保存教会合一的最有益的爱德，可以被我们这一方用来支持，而非支持多纳徒派。因为，如果他们选择效法他的榜样，重洗天主教友，因为他认为异端者加入天主教会时应受洗，难道我们不应该更效法他的榜样吗？他树立了一个明确的规则：绝不可通过建立分离的共融而脱离大公共融，即分散于全世界的基督徒团体，即使接纳邪恶和亵圣的人也是如此，因为他甚至不愿意剥夺那些他视为曾无圣洗接纳亵圣者进入大公共融的人的共融权，他说：\Quote{不判断任何人，也不剥夺任何与我们意见相左者的共融权利。}\par

% structure: p0005 source=h2 target=subsection reason=relative-heading-rank; base=h2

\subsection{第二章}

尽管如此，我仍看到可能还需要我做什么，即，我应该回答那些貌似有理的论点，那些论点甚至在更早的时代，曾促使\Person{阿格里披努斯}，或\Person{西彼廉}本人，或\Place{非洲}与他们意见相同的人，或任何其他远隔重洋的远方的人，不是出于任何全教会性或甚至地区性公会议的权威，而仅仅通过书信往来，认为他们应当采纳一种既无教会古老习俗认可，又被大公世界极其一致的决议明确禁止的习俗，为的是使那种因这类讨论而开始潜入某些人心中的错误，能被更强大的真理和合一的全能治愈力量所治愈，而这力量是在安全一边的。因此他们可以看到我进行这番论述是多么安稳。如果我不能达到目的，不能表明如何反驳他们从\Person{西彼廉}的公会议和书信中提出的论点，即基督的圣洗不可由异端者之手施行，我仍将安全地留在教会内，而\Person{西彼廉}本人也正是在这教会的共融中，与那些与他意见相左的人一同留在这共融内。\par

3. 但如果他们说，那时天主教会存在，是因为有少数人，或者，如果他们愿意，甚至相当多人，否认任何在异端团体中施行的圣洗的有效性，并给所有从那里来的人重新施洗，那又怎么样？难道教会不存在于\Person{阿格里披努斯}之前吗？是他开始实行那种与以往一切习俗相悖的新体制的。或者，又在\Person{阿格里披努斯}之后，除非有回归原始习俗，否则\Person{西彼廉}何必召开另一个公会议？那时难道没有教会吗？因为这样一种习俗处处盛行：基督的圣洗应被视为无非是基督的圣洗，即使它被证明是在异端或分裂者团体中施行的。但如果那时教会存在，并未因连续性的断裂而消亡，反而屹立不倒，并在万国中增长，那么坚守这一习俗无疑是最安全的，这习俗那时在合一中一同包容好人和坏人。但如果那时教会不存在，因为亵圣的异端者无圣洗就被接纳，且这由普遍的习俗所盛行，那么\Person{多纳徒}是从哪里出现的？他从哪块土地冒出来？或从哪个海跃出？或从哪个天坠落？因此，正如我开始所说，我们在那教会的共融中是安全的，在该教会整个范围内，现今的习俗盛行，而该习俗在\Person{阿格里披努斯}之前也同样在整个范围内盛行，并在\Person{阿格里披努斯}与\Person{西彼廉}之间的时期也是如此，\Person{阿格里披努斯}、\Person{西彼廉}以及同意他们的人从未抛弃其合一，尽管他们与其余兄弟意见不同——他们所有人都与那些意见相左的人留在同一的合一共融中。但让多纳徒派自己考虑他们的真实立场是什么：如果他们既不能说他们的起源来自何处（如果教会已经因与无圣洗就被接纳到其怀抱中的异端者和分裂者共融的污点而被毁灭），也不能同意\Person{西彼廉}本人，因为他宣称他仍与那些接纳异端者和分裂者的人共融，也因而与那些被接纳者共融；而他们却因指控有人交出圣书而将自己与全世界的共融分离，这指控是他们针对在\Place{非洲}被他们诽谤的人提出的，但在海外受审时却未能定罪；尽管，即使他们指控的罪行属实，也远不如异端和分裂的罪那么严重；然而，这些罪并不能玷污\Person{西彼廉}，就那些从他那里来的人（按他的看法，他们无圣洗就被接纳进入大公共融）而言。还有，在他们声称效法\Person{西彼廉}的这一点上，他们对于承认\Person{马克西米安}的追随者的圣洗，连同那些人（尽管他们属于他们首先在自己的全教会性公会议上谴责、然后甚至通过世俗权力法庭起诉的那一党派，他们后来却在同一主教——即他们被谴责时所处的主教——的任期内将那些人接回共融）的圣洗，他们找不到任何答案。因此，如果恶人的共融在\Person{西彼廉}时代毁灭了教会，他们就无从获得自己的共融；如果教会未被毁灭，他们就没有分离的借口。此外，他们既没有效法\Person{西彼廉}的榜样——因为他们已打破合一的纽带，也没有遵守自己的公会议——因为他们已承认\Person{马克西米安}的追随者的圣洗。\par

% structure: p0008 source=h2 target=subsection reason=relative-heading-rank; base=h2

\subsection{第三章}

4. 因此，既然我们紧守西彼廉的榜样，现在继续来看西彼廉的主教会议。西彼廉说：\Quote{你们听到了，最亲爱的同僚，我们的同袍犹巴雅努斯主教写信给我，就关于异端者非法且亵渎的圣洗征询我这平庸的才智，以及我给他的答复——给予我们一再并经常作出的判决，即来到教会的异端者应当接受教会的圣洗而被圣化。犹巴雅努斯的另一封信也已向你们宣读，在这封信中，他基于真诚与虔诚的虔敬，回应我们的书信，不仅表示同意，而且表示感谢，承认他受到了教诲。}在这位蒙福的西彼廉的话语中，我们发现犹巴雅努斯曾向他征询，他如何回答他的问题，以及犹巴雅努斯如何感激地承认受到了教诲。那么，如果我们希望审视这封使犹巴雅努斯信服的书信，难道应被认为是不合理地固执吗？因为直到我们也信服（若有任何真理的论证可以做到这一点），西彼廉本人已藉由公教共融的权利为我们的安全建立了基础。\par

5. 因为他接着说：\Quote{现在剩下的是，我们各自就同一主题发表意见，不评判任何人，也不剥夺任何人与我们意见不同之人的共融权利。}因此，他在不失去共融权利的情况下，允许我不仅继续探讨真理，甚至可以持守与他不同的意见。他说：\Quote{因为我们中没有一人自封为主教的主教，或以暴虐的恐怖强迫同僚必须服从。}还有什么比这更仁慈？还有什么更谦卑？确实，这里没有任何权威阻止我们探讨何为真理。他说：\Quote{既然每位主教，在自由运用其自主权与权柄时，有权形成自己的判断，既不能受他人审判，也不能审判他人——}我想，这指的是那些尚未得到完全清晰解决的问题；因为他知道关于圣事的一个深刻问题当时正以各种辩论占据整个教会，他给每个人自由探讨的权利，以便通过研究使真理显明。他当然不会说假话，试图诱使较简单的同僚说话，以便当他们暴露出与他不同的意见时，他违反承诺提议将他们逐出教会。愿如此圣洁的灵魂远离如此可咒的诡诈；事实上，那些对这样一位人物持有这种看法、认为这有助于其赞美的人，只是表明这符合他们自己的本性。我绝不认为西彼廉，一位公教主教、公教殉道者，其伟大只使他在一切事上更加谦卑，以致在主面前蒙恩，\BibleRef{德 3:18}竟会在同僚的神圣会议中，尤其口说心中不存的话，特别是他又补充说：\Quote{但我们必须等候我们主耶稣基督的审判，唯有祂有权柄将我们安置在祂教会治理中，并判断我们在此中的行为。}当他如此提醒同僚这一庄严审判，期望从同僚那里听到真理时，他竟会先立下说谎的榜样吗？愿天主阻止这样的疯狂远离每位基督徒，更远离西彼廉！因此，我们拥有西彼廉最谦和、最真实的讲话所赐予的自由探讨权。\par

% structure: p0011 source=h2 target=subsection reason=relative-heading-rank; base=h2

\subsection{第四章}

6. 接下来，他的同僚逐一发表各自意见。但他们首先听取了写给犹巴雅努斯的信；正如前言所述，这封信被宣读了。因此，让我们也在自己中间宣读它，以便我们也能藉天主的帮助，从中发现当持守的看法。\Quote{怎么！}我想我听到有人说，\Quote{你要告诉我们西彼廉写给犹巴雅努斯的内容吗？}我读过这封信，我承认，若非那巨大的权威促使我更仔细地考虑这事，我肯定会转而接受他的观点——这权威源于那些被遍布世界、在如此多民族中（拉丁人、希腊人、蛮族，更不用说犹太民族本身）的教会所能产生的人——正是那生了西彼廉自己的同一教会——这些人我绝不能认为他们毫无道理地不愿持守此观点——不是因为在这如此棘手的问题上，一人或少数人的意见可能不接近真理，而是因为一个人不应轻率地、未经充分考虑和尽力研究，就因一人或甚至少数人的意见，反对如此众多同宗教同共融、且都拥有伟大才智和丰富学识之人的决定。因此，在更勤奋地探究中，西彼廉自己的信也向我提示了多少支持公教会现今所持观点的内容——即基督的圣洗应被承认和认可，不是根据施行者的功过，而是唯独根据祂的功过，关于祂曾记载：\Quote{他就是那施洗者}，\BibleRef{若 1:33}——这将在我们论述的过程中自然显明。因此，让我们设想西彼廉写给犹巴雅努斯的信已在我们的会议中被宣读，正如它在主教会议中宣读一样。我愿每个打算读我接下来所写内容的人都读这封信，免得他可能认为我遗漏了一些重要内容。因为若此时引用这封书信的所有话语，将耗费太多时间，且与阐明当前问题无关。\par

% structure: p0013 source=h2 target=subsection reason=relative-heading-rank; base=h2

\subsection{第五章}

7. 但若有人问起，在讨论此问题期间，我暂时持何立场，我回答：首先，\Person{圣西彼廉}的书信提示了我，在接下来开始讨论的问题的性质变得清晰之前，我应当持守什么。因为\Person{圣西彼廉}本人说：\Quote{但有人会说：\InnerQuote{那么，从前那些从异端来到教会、未经圣洗就被接纳的人，将会如何呢？}}他们是否真的未经圣洗，或者他们被接纳是因为接纳他们的人认为他们已经领受了圣洗，这是我们将来需要考虑的问题。无论如何，\Person{圣西彼廉}本人充分地表明了教会当时的普遍惯例，他说，在过去，那些从异端来到教会的人，是未经圣洗就被接纳的。\par

8. 因为在同一公会议中，\Person{西卡的卡斯图斯}说：\Quote{藐视真理却妄图追随惯例的人，要么是嫉妒或恶意对待真理已向其启示的弟兄，要么是对天主忘恩负义，因为祂的教会赖其启示而受教导。}真理是否已被启示，我们稍后将查考；无论如何，他承认教会的惯例是不同的。\par

% structure: p0016 source=h2 target=subsection reason=relative-heading-rank; base=h2

\subsection{第六章}

9. \Person{瓦加的利博苏斯}也说：\Quote{主在福音中说：\BibleQuote{我是真理。}} \BibleRef{若 14:6} 祂没有说：\Quote{我是惯例。}因此，当真理显明时，惯例必须让位于真理。显然，没有人能怀疑，当真理显明时，惯例必须让位于真理。但我们稍后将查看真理的显明问题。同时，他也表明惯例是在另一边。\par

% structure: p0018 source=h2 target=subsection reason=relative-heading-rank; base=h2

\subsection{第七章}

10. \Person{塔拉萨的佐西穆斯}也说：\Quote{当真理的启示已作出时，错误必须让位于真理；因为甚至连\Person{伯多禄}，起初行割礼，后来在\Person{保禄}宣布真理时也让步了。}他确实选择了说错误，而不是惯例；但他说\Quote{因为甚至连\Person{伯多禄}，起初行割礼，后来在\Person{保禄}宣布真理时也让步了}，这充分地表明，在圣洗问题上，也有一种与其观点相左的惯例。同时，他也提醒我们，\Person{圣西彼廉}可能持有与真理（如教会在他前后所持守的）相悖的关于圣洗的观点，这并非不可能，因为连\Person{伯多禄}也曾持有与宗徒\Person{保禄}教导我们的真理相悖的观点。\BibleRef{迦 2:11}\par

% structure: p0020 source=h2 target=subsection reason=relative-heading-rank; base=h2

\subsection{第八章}

11. \Person{布斯拉塞内的费利克斯}同样说：\Quote{在接纳异端者而不给予教会的圣洗时，谁也不可让惯例优先于理智和真理；因为理智和真理总是胜过惯例，排除惯例。}如果那真是理智和真理，那就再好不过；但我们稍后将查看这一点。同时，从此人的话中也可清楚看出，惯例是另一种方式。\par

% structure: p0022 source=h2 target=subsection reason=relative-heading-rank; base=h2

\subsection{第九章}

12. \Person{图卡的和诺拉图斯}同样说：\Quote{既然基督是真理，我们应当跟随真理而非惯例。}所有这些声明证明，我们并没有被排除在教会的共融之外，直到真理的本性（他们声称必须优先于我们的惯例）被清楚地表明为止。但如果真理已清楚表明，应当维持该惯例所规定的规条，那么显然，这一惯例并非徒然建立或确认，并且由于这些讨论，对如此伟大的圣事的最有益遵守（它其实从未在天主教会中被改变过），在得到公会议进一步确认后，得以更加谨慎地、以最严格的警惕被守护。\par

% structure: p0024 source=h2 target=subsection reason=relative-heading-rank; base=h2

\subsection{第十章}

13. 因此，\Person{圣西彼廉}致\Person{朱巴亚努斯}的信如下：\Quote{关于异端者的圣洗，他们置身于教会之外，并自认对既无权利也无能力之事拥有主权。对此，我们无法认为有效或合法，因为显然在他们之中这是非法的。}我们也不否认，一个人在异端者中或在教会外的任何分裂中受洗，就他分享异端者和分裂者的乖戾而言，他从中毫无所得；我们也不认为那些施洗者，尽管他们赋予真实的圣洗圣事，但在教会外聚集追随者并持有与教会相反的观点一事上，行事正确。但一件事是缺少圣事，另一件事是错误地拥有它并非法地篡夺它。因此，它们并不因为被非法使用（不仅被异端者，也被各种邪恶不虔之人使用）就不再是基督和教会的圣事。这些人确实应当被纠正和惩罚，但圣事应当被承认和敬畏。\par

14. \Person{圣西彼廉}确实说，关于此主题，举行了不止一次、而是两次或多次公会议；然而总是在非洲。因为有一次他提到，聚集了七十一位主教——对所有这些人的权威，我们毫不迟疑地，以对\Person{圣西彼廉}的一切应有敬意，更倾向于整个普世教会的权威，这权威有更多主教支持，\Person{圣西彼廉}本人也曾欣喜地自认是这教会不可分割的成员。\par

15. 水本身既非\Quote{亵渎而淫乱}，即使它是由亵渎和淫乱之人呼求天主之名而被祝圣的；因为无论是水这受造物本身，还是所呼求的名，都不是淫乱的。但基督的圣洗，藉着福音的话被祝圣，必然是圣的，无论其施行者何等污秽不洁；因为它固有的圣德不能被玷污，天主的德能在其圣事中常存，或为正确使用者的救恩，或为错误使用者的毁灭。难道你真要主张，太阳或烛光的光，即便照过污秽之地，本身也不沾染任何污秽，而基督的圣洗却可能被任何人的罪玷污，无论他是谁？因为如果我们思考圣事的可见质料本身，每个人都知道它们会朽坏。但如果我们思考它们所传递的，谁看不出它是不朽坏的，无论施行者的生活品行是得到赏报还是惩罚？\par

% structure: p0028 source=h2 target=subsection reason=relative-heading-rank; base=h2

\subsection{11.}

16. 但\Person{西彼廉}不被\Person{尤巴伊安}所写的话动摇，是正当的，因为他说：\Quote{\Person{诺瓦提安}的追随者重行圣洗那些从\Place{天主教会}来到他们那里的人。}因为，首先，异端者以乖谬的模仿精神所做的一切，并不表示天主教徒因此就要避免去做，因为异端者也做同样的事。而且，异端者和\Place{天主教会}各自应避免重行圣洗的理由是不同的。因为即使在天主教会中是适宜的，异端者也做是不对的；因为他们的论点是，在天主教徒中缺少他们自己还在教会范围内时所领受、离开时所带走的东西。而\Place{天主教会}不应再次施行在异端者中所给予的圣洗的理由是，以免它似乎断定属于基督独有的权能属于其成员，或宣布在异端者中缺少他们曾在教会范围内领受、且肯定不能因离开而失去的东西。因为\Person{西彼廉}本人与所有其他人共同确立了这一点：如果有人从异端回到教会，他们应被接纳回来，不是藉着圣洗，而是藉着补赎的纪律；由此可知，他们不能因退离而失去什么东西，因为在他们回来时并没有恢复给他们。也不应该片刻说，正如他们的异端是他们自己的，他们的错误是他们自己的，分裂的亵渎是他们自己的，同样圣洗也是他们自己的，而它实在是基督的。因此，当他们回来时，那些属于他们自己的错误得到纠正，而在那不属他们的东西上，应承认那一位的同在，祂是它的源头。\par

% structure: p0030 source=h2 target=subsection reason=relative-heading-rank; base=h2

\subsection{12.}

17. 但真福\Person{西彼廉}表明，他的决定并非新事或突发之事，因为此惯例早在\Person{阿格利皮努斯}时期就已开始。他说：\Quote{许多年，}他说，\Quote{许多时间已经过去，自从可敬的\Person{阿格利皮努斯}领导下，一大群主教决定了这一点。}因此，至少在\Person{阿格利皮努斯}时期，这事是新的。但我不明白\Person{西彼廉}说：\Quote{从那时直到今日，我们各省如此众多的异端者，归依了我们的教会，毫无犹豫或厌恶，反而以理性和意志的完全赞同，拥抱了生命之洗的恩宠和救恩的圣洗，}是什么意思，除非他确实说\InnerQuote{从那时直到今日，}因为从他们按照\Person{阿格利皮努斯}会议在教会领受圣洗时起，在任何受重行圣洗者身上都没有出现开除教籍的问题。然而，如果从\Person{阿格利皮努斯}时期到\Person{西彼廉}时期，给来自异端者施行圣洗的惯例一直有效，为什么\Person{西彼廉}还要为此召开新的会议呢？为什么他对同一个\Person{尤巴伊安}说，他并没有做什么新事或突发之事，而只是做了\Person{阿格利皮努斯}所确立的事？因为如果从\Person{阿格利皮努斯}到\Person{西彼廉}，这是教会持续的惯例，为什么\Person{尤巴伊安}会因新颖问题而困扰，以至于需要\Person{阿格利皮努斯}的权威来满足呢？最后，为什么他的这么多同僚极力主张理性和真理必须优先于习俗，而不是说那些想另做的人是与真理和习俗都相违背的呢？\par

% structure: p0032 source=h2 target=subsection reason=relative-heading-rank; base=h2

\subsection{13.}

18. 至于赦罪的问题，即是否通过异端者施行的圣洗而获得，我已在先前的一卷书中发表了我的意见；但此处我愿简要重述。如果借着圣洗的神圣性在那里赦免了罪，那么因固执地坚持异端或分裂，罪又会返回；因此这样的人必须返回天主教会内的平安，使他们不再作异端者和分裂者，并配得藉着在合一纽带中运作的爱德，清除那些返回的罪。但是，如果即使在异端者和分裂者中，那仍然是基督同一的圣洗，却因纷争的污秽和分歧的邪恶而不能施行赦罪，那么当他们来到教会的平安中时，同一圣洗便开始有效赦罪——[并非]说那已真正赦免的不应保留；也不是说异端的圣洗应被拒斥，如同属于另一种宗教，或与我们的圣洗不同，因而需重行圣洗；而是说同一圣洗，因教会外的纷争导致死亡，因教会内的平安带来得救。事实上，这正是同一馨香，宗徒说：\BibleQuote{我们在各处都是基督的馨香；}他又说：\BibleQuote{对得救的人，我们是致生的生味；对丧亡的人，我们是致死的死味。}\BibleRef{格后 2:15}。尽管他用这些话论及另一主题，我仍将其应用于此，使人明白，善的事物不仅能为正确使用它的人带来生命，也能为错误使用它的人带来死亡。\par

% structure: p0034 source=h2 target=subsection reason=relative-heading-rank; base=h2

\subsection{第十四章。}

19. 当我们考虑圣事的真确性与神圣性问题时，\Quote{领受圣事者相信什么、怀有何种信德}并不重要。这对于进入救恩至关重要，但就圣事问题而言则完全无关紧要。因为很可能一个人拥有真确的圣事而信德败坏，同样也可能他完整地持有信经的言辞，却对天主圣三、复活或任何其他点持有错误信仰。即便在天主教会内部，要持守完全与真理一致的信仰，即使关于天主本身，也不是小事，更不用说祂的任何受造物了。那么，是否应该主张：如果任何人在天主教会内受了圣洗，后来通过阅读、听讲或平静的辩论，因天主的启示发现自己以前信得不对，因此他需要重新受洗？但有什么属血肉的本性的人不因自己内心的虚妄幻想而迷失，按照自己血肉的感觉想象天主的本性，从而与天主的真观念相距如同虚妄与真理之远呢？最确实地，充满真理之光的宗徒说：\BibleQuote{属本性的人，不能领受天主圣神的事。}\BibleRef{格前 2:14}然而他这里指的是那些他表明已经受了洗的人。因为他对他们说：\BibleQuote{保禄曾为你们被钉在十字架上吗？或者你们曾因保禄的名受洗吗？}\BibleRef{格前 1:13}这些人因此领受了圣洗的圣事；然而，既然他们的智慧是属血肉的，他们对于天主除了按血肉的知觉能相信什么之外，又能相信什么呢？按此，\BibleQuote{属本性的人不能领受天主圣神的事。}他对这样的人说：\BibleQuote{我不能把你们当作属神的人说话，只能当作属血肉的人，当作基督内的婴孩。我给你们喝了奶，没有给肉吃；因为直到现在你们还不能吃，就是现在你们也不能，因为你们仍是属血肉的。}\BibleRef{格前 3:1}因为这样的人被各种教训之风摇动，他提到这种人说：\BibleQuote{使我们不再作小孩子，被各种教训之风摇动，飘来飘去。}\BibleRef{弗 4:14}那么，如果这些人进步到内在之人的属神年龄，并以其理解的完整性，认识到他们由于自己幻想的欺骗性而对天主所持的信仰与真理的要求相差多远，他们因此就必须重新受洗吗？因为按照这个原则，一位公教望教者有可能读到某位异端者的著作，缺乏辨别是非的知识，从而持有与公教信仰相悖的某种信念，但并不被信经的言辞所谴责，正如在相同言辞的外表下产生了无数异端错误。假设这位望教者以为自己正在研读某位伟大博学的公教学者的作品，并带着那种信念在天主教会受了洗，后来通过研究发现了他应当相信什么，从而拥抱公教信仰，摒弃先前的错误，那么他告明这一点后，应当重新受洗吗？或者假设在他自己学习和告明这一点之前，发现他持有这样的意见，并被教导应当摒弃什么和相信什么，而且弄清楚他受洗时确实持有这种错误信仰，那么他因此应当重新受洗吗？为什么我们坚持相反的意见？因为那在福音言辞中祝圣的圣事的圣德，完整地留在他的身上，正如他从施行者手中领受的一样，尽管他因深深地扎根于自己血肉思想的虚妄中，在受洗时持有不正确的信仰。因此，显然有可能，尽管信德有缺陷，圣洗的圣事仍可毫无缺陷地留在任何人身上；因此所有关于不同异端者的说法都是离题的。因为在每个人身上，纠正者发现什么错误就应当纠正什么。要治疗不健康的，要补充缺乏的，尤其是基督徒爱德的平安，没有它其余的都无益。然而，既然其余的存在，我们不应将其当作缺乏而施行，只应确保拥有它的人因和平的纽带和爱德的卓越而得益而非受损。\par

% structure: p0036 source=h2 target=subsection reason=relative-heading-rank; base=h2

\subsection{第十五章.}

20. 因此，如果\Person{玛尔西翁}用福音的话\BibleQuote{因父、及子、及圣神之名}\BibleRef{玛 28:19}祝圣了圣洗的圣事，那么圣事是完整的，尽管他用同样的话所表达的信德并不完整——因为他持有天主教真理所不教的见解——而是被虚妄的传说所玷污。因为在这些同样的话\Quote{因父、及子、及圣神之名}之下，不但是\Person{玛尔西翁}、\Person{瓦伦提努斯}、\Person{亚略}或\Person{欧诺米}，就是教会中属血肉的婴孩（宗徒曾对他们说：\InnerQuote{我从前不能把你们当作属神的人，只能当作属血肉的人}），若他们各自被问及对自己意见的准确解释，恐怕会显出与人数同样多的意见分歧，\Quote{因为属血气的人不接受天主圣神的事}。然而，能因此就说他们没有领受完整的圣事吗？或者，如果他们进步了，纠正了属血肉意见的虚妄，就必须重新寻求他们已经领受的吗？各人按自己信德的方式领受；然而，在信赖那同一宗徒所说\BibleQuote{你们若在什么事上存有别的思想，天主也要启示你们这件事}\BibleRef{斐 3:15}的天主慈悲引导下，他获得了多少呢？然而，异端者和分裂主义的陷阱对属血肉的人之所以极其有害，正是因为当他们的虚妄见解被肯定以反对天主教真理时，他们的进展被阻断了，并且他们分裂的悖逆被加强以反对天主教和平。然而，如果圣事是相同的，它们在各处都是完整的，即使被错误地理解并歪曲为不和的工具，正如同福音的经文，只要它们是相同的，在各处都是完整的，即使被引用来支持无尽错误的意见。至于\Person{耶肋米亚}所说的：\Quote{为什么迫害我的人占优势？我的创伤顽固，我怎能痊愈？它起初对我就像虚假的水，没有确定性}，——如果\Quote{水}这个词除了表示圣洗之外，从未在预言的形象和比喻语言中使用过，那么我们将难以发现\Person{耶肋米亚}这些话的意思；但事实上，既然在\WorkTitle{默示录}中\Quote{众水}\BibleRef{默 17:15}明确被用来表示\Quote{人民}，我就不明白为什么不能把\Quote{虚假的水，没有确定性}理解为\Quote{虚假的人民，我不能信任}。\par

% structure: p0038 source=h2 target=subsection reason=relative-heading-rank; base=h2

\subsection{第十六章。}

21. 但当人们说\Quote{圣神只在天主教会中藉按手礼赐予}时，我想我们的祖先是要我们理解宗徒所说的：\BibleQuote{因为天主的爱，藉着所赐予我们的圣神，已倾注在我们心中了。}\BibleRef{罗 5:5} 因为正是这爱，凡从天主教会共融中被切断的人所欠缺的；因缺乏这爱，\BibleQuote{即使他们能说人间的语言，并能说天使的语言；即使他们明白一切奥秘和一切知识；即使他们有先知之恩，并具备一切信德，甚至能移山；即使他们把自己所有的财产全施舍给穷人，并把身体交出被烧，对他们也毫无益处。}\BibleRef{格前 13:1} 但那些不关心教会合一的人，正缺乏天主的爱；因此，我们正确理解圣神可以说除天主教会外不在别处领受。因为圣神不仅仅在从前藉按手礼伴随暂时的、可感觉的奇迹而赐下，如同在初期为奠基的信德作证并扩展教会初期时那样。因为在今天，谁还期望那些按手领受圣神的人立刻说起方言呢？而是被理解为为了和平的纽带，无形且不可感知地，神圣的爱被吹入他们的心中，使他们能够说：\Quote{因为天主的爱，藉着所赐予我们的圣神，已倾注在我们心中了。}但圣神有许多运作，同一位宗徒在一段经文中足够长地列举后，总结说：\BibleQuote{这一切都是这同一圣神所推动，随他的心愿个别分配给人。}\BibleRef{格前 12:11} 既然圣事是一回事——连\Person{西满魔术士}也能领受\BibleRef{宗 8:13}——圣神的运作是另一回事——甚至常常在恶人身上发现，如\Person{撒乌耳}有先知之恩——而同一圣神的另一种运作只有善人才能拥有，如\BibleQuote{命令的终极就是出于纯洁的心、良好的良心和真诚的信德的爱}\BibleRef{弟前 1:5}——因此，无论异端者和分裂主义可能领受了什么，那\Quote{遮盖许多罪过}的爱正是天主教会合一与和平的特别恩赐；它也不在联合中的所有成员身上都找到，因为不是所有在联合之内的人都属于它，正如我们将在适当地方看到。无论如何，在联合之外，那爱不可能存在，没有这爱，其他一切甚至可被承认和称许的必备条件，都不能使人获益或脱离罪。但和好时按手在教会中并不像圣洗那样不可重复；因为它无非是一个为人的祈祷。\par

% structure: p0040 source=h2 target=subsection reason=relative-heading-rank; base=h2

\subsection{第十七章。}

22. \Quote{为了保存合一的形象，主赐给\Person{伯多禄}权能，凡他在世上释放的，就被释放，} \BibleRef{玛 16:19} 显然，那合一也被描述为一只无瑕疵的鸽子。 \BibleRef{歌 6:9} 那么，能说所有那些贪婪之人属于这同一只鸽子吗？\Person{西彼廉}本人曾在同一天主教会中如此悲痛地哀叹他们的存在。因为我相信，猛禽不能被称为鸽子，而是鹰。那么，他们如何给那些以诡诈欺骗掠夺庄园、以复利增加利润的人施洗，如果圣洗只由那不可分割、纯洁、完美的鸽子，即那只能理解存在于善人中的合一，来施行？是否可能，藉着教会内属神的圣人们的祈祷，如同鸽子频繁的哀鸣，一件伟大的圣事被施行，伴随着天主仁慈的隐秘安排，以至于那些受洗之人的罪也被赦免——不是由鸽子，而是由鹰——如果他们是在天主教会的合一的和睦中来到那圣事？但若是如此，为什么不能也同样发生：当每个人从异端或分裂来到天主教会的和平时，他的罪应藉着他们的祈祷而被赦免？但圣事的完整性在各地都被承认，尽管它在教会合一之外对不可撤销的罪赦无效。圣人们的祈祷，换言之，那只独一鸽子的叹息，也无法帮助那些处于异端或分裂中的人；正如它们不能帮助那在教会之内的人，如果他因邪恶的生活而自己保留罪债，尽管他受洗不是由这鹰，而是由鸽子本身的虔诚事工。\par

% structure: p0042 source=h2 target=subsection reason=relative-heading-rank; base=h2

\subsection{第十八章}

23. \BibleQuote{如同父派遣了我，} 我们的主说，\BibleQuote{我也派遣你们。他这样说了以后，就向他们嘘了一口气，并向他们说：你们领受圣神吧！你们赦免谁的罪，谁的罪就得到赦免；你们保留谁的罪，谁的罪就被保留。} \BibleRef{若 20:21} 因此，如果他们代表教会，而这话是对他们说的，如同对教会本身说的，那么，教会的和平赦免罪，与教会的疏离保留罪，这并非按照人的意愿，而是按照天主的旨意和属神的圣人们的祈祷，他们 \BibleQuote{判断一切，自己却不受任何人的判断} \BibleRef{格前 2:15}。因为磐石保留，磐石赦免；鸽子保留，鸽子赦免；合一保留，合一赦免。但这合一的和平只存在于善人中，存在于那些已经是属神的，或通过和谐服从而迈向属神事物的人中；它不存在于恶人中，无论他们在外面制造扰乱，或是在教会内被以哀叹容忍，施洗或受洗。但是，正如那些在教会内被以叹息容忍的人，虽然他们不属于同一鸽子的合一，也不属于那 \BibleQuote{荣耀的教会，没有瑕疵或皱纹，或任何类似的东西}，然而如果他们改正，并承认他们极不相称地接近了圣洗，他们不再重洗，而是开始属于那只鸽子，藉着它的叹息，那些在他们与她的和平疏离时被保留的罪得到赦免；同样，那些更公开地在教会之外的人，如果他们领受了相同的圣事，在改正后，来到教会的合一，不是因重复洗礼而从罪中获得自由，而是因同样的爱德法律和合一的纽带。因为如果 \Quote{只有那些被置于教会之上，并按照福音的法律和主的任命所设立的圣职者才能施洗}，那么这些以诡诈欺骗掠夺庄园、以复利增加利润的人，在任何意义上属于这一类吗？我不认为如此，因为那些由主的任命所设立的人，在宗徒给他们标准时说：\BibleQuote{不贪婪，不贪求不义之财} \BibleRef{铎 1:7}。然而，在 \Person{西彼廉} 的时代，这种人曾施行洗礼；他本人多次哀叹，他们是他同僚的主教，并以伟大的忍耐奖赏容忍了他们。但他们并没有赋予罪的赦免，这赦免是通过圣人们的祈祷，即鸽子的叹息而产生的，无论由谁施洗，如果得到赦免的人属于她的和平。因为主不会对强盗和高利贷者说：\BibleQuote{你们赦免谁的罪，谁的罪就得到赦免；你们保留谁的罪，谁的罪就被保留}。 \Quote{在教会之外，确实，什么都不能被束缚或释放，因为在外面没有人能束缚或释放；} 但他是被释放的，谁与鸽子缔造了和平；他是被束缚的，谁不与鸽子和平相处，无论他是公开地在外面，还是似乎在里面。\par

24. 但我们知道，\Person{达堂}、\Person{科辣黑}和\Person{阿彼兰}\EditorialNote{户 16}曾试图违背天主子民的合一，将献祭的权利据为己有，还有亚郎的儿子们在祭台上献了凡火\BibleRef{肋 10:1}，他们没有逃脱惩罚。我们也不说这罪行不会受罚，除非犯罪者自我改正，如果天主的忍耐引导他们悔改\BibleRef{罗 2:4}，给予他们改正的时间。\par

% structure: p0045 source=h2 target=subsection reason=relative-heading-rank; base=h2

\subsection{第十九章.}

25. 那些人确实说，不应重复施行圣洗，因为执事斐理伯给那些受洗者只覆手了，\BibleRef{宗 8:5}，他们所说的话完全不得要领；我们在寻求真理时决不可使用这样的论证。因此，我们更加远离\Quote{向异端者屈服}，如果我们否认他们所拥有的属于基督教会的一切是他们自己的财产，并且不因逃兵的罪行而拒绝承认我们统帅的旗帜；不仅如此，更因为\BibleQuote{主我们的天主是嫉妒的天主}，\BibleRef{申 4:24}，当我们看见他的任何事物在外人手中时，我们要拒绝让他认为那是他自己的。因为事实上，那嫉妒的天主亲自斥责那对他行淫的女子，作为误入歧途的子民的预表，并说她把她所拥有的东西给了她的情人们，又从他们那里收回那些不是他们的而是他的。在淫妇和通奸的情人手中，天主在愤怒中，作为嫉妒的天主，认出他的恩赐；难道我们说，以福音之言祝圣的圣洗属于异端者吗？难道我们愿意从他们的行为考虑，甚至把属于天主的也归给他们，好像他们有能力玷污它，或者好像他们能使属于天主的成为他们自己的，因为他们自己拒绝属于天主吗？\par

26. 那先知欧瑟亚所指出的淫妇是谁？她曾说：\Quote{我要跟随我的情人们，他们给我饼和水、羊毛和麻，以及一切适宜我的东西。}让我们承认，我们也可以这样理解那误入歧途的犹太子民；然而，假基督徒（所有异端者和分裂者都是如此）除了效法假以色列人之外，还会效法谁呢？因为也有真以色列人，正如主亲自为纳塔乃耳作证说：\BibleQuote{看，这确是一个以色列人，在他内毫无诡诈。}\BibleRef{若 1:47} 但谁是真基督徒，除非是那位主所说的：\BibleQuote{那拥有我命令并遵守的，他就是那爱我的人。}\BibleRef{若 14:21} 但遵守他的命令是什么意思，岂不是存留在爱中吗？因此他也说：\BibleQuote{我给你们一条新命令：你们要彼此相爱；}又说：\BibleQuote{如果你们彼此相爱，众人因此就会认出你们是我的门徒。}\BibleRef{若 13:34} 但谁能怀疑这话不仅是对那些在他肉身临在时用肉耳听他话的人说的，也是对那些通过福音学习他话语的人说的，当时他正坐在天上的宝座上？因为他来不是为废除法律，而是为成全。\BibleRef{玛 5:17} 但法律的成全就是爱。\BibleRef{罗 13:10} 而西彼廉在这方面非常丰富，以至于虽然他对圣洗持有不同见解，却仍未离弃教会的合一，并在主的葡萄树上是一根根深蒂固的枝条，结出果实，天上的农人用苦难的刀修剪它，使它结出更多的果实。\BibleRef{若 15:1} 但这兄弟之爱的仇敌，无论是公开在外面的，还是看似在里面的，都是假基督徒和假基督。因为他们一有机会就出去，正如经上所记：\Quote{一个想要与朋友分离的人，寻找机会。}但即使没有机会，当他们看起来在里面时，他们与那无形的爱之纽带是分离的。因此圣若望说：\BibleQuote{他们从我们中间出去了，但并不是属于我们的；因为如果他们属于我们，他们必定会继续和我们在一起。}\BibleRef{若一 2:19} 他没有说他们因出去就不再属于我们，而是说他们出去是因为不属于我们。宗徒保禄也谈到某些人，他们在真理上犯了错误，正在倾覆一些人的信德；他们的话像毒疮一样侵蚀。然而在说他们应当被避开时，他也暗示他们都在同一所大房子里，但作为羞辱的器皿——我想是因为他们还没有出去。或者如果他们已出去了，他怎么能说他们与尊贵的器皿在同一所大房子里，除非是因着圣事本身，即使在异端者分离的聚会中这些圣事也没有改变，他才把所有人都说成属于同一所大房子，尽管有不同的评价，有些是尊贵的，有些是羞辱的？因为他这样在致弟茂德书信中说：\BibleQuote{但要避开凡俗空谈；因为这些会使人更不敬虔。他们的话像毒疮一样侵蚀；其中有希默纳和斐肋托；他们在真理上犯了错误，说复活已经过去，倾覆了一些人的信德。然而天主的根基屹立不摇，有这印记：主认识属于他的人。并且：凡呼求基督名字的，应当远离不义。但在大房子里，不仅有金器和银器，也有木器和瓦器；有些是尊贵的，有些是羞辱的。所以，如果有人自洁脱离这些，他必成为尊贵的器皿，圣化，合乎主人的使用，预备行各样的善事。}\BibleRef{弟后 2:16}\par

27. 因此，那些邪恶、作恶、属血肉、属肉体、属魔鬼的人，以为他们从自己的诱惑者那里领受了唯独天主的恩赐，无论是圣事，还是关乎现世救恩的任何灵性运作。但这些人对天主没有爱德，反而忙于那些因骄傲而使他们误入迷途的人，并被比作先知所说的那个淫妇：\Quote{我要追随我的爱人，他们给我饼和水，给我羊毛和麻，给我油和一切合宜之物。}因为异端和分裂就是这样产生的，当那属肉体的百姓，不建立在天主的爱上，却说：\Quote{我要追随我的爱人}时，她便与这些人，或因信仰败坏，或因骄傲自大，可耻地犯了奸淫。但为了那些经历了困难、窘迫和诱惑者虚妄推理的障碍，随后感到恐惧的刺痛，回归平安之路，真诚寻求天主的人——为了他们，祂接着说：\Quote{因此，看哪，我必用荆棘堵塞你的路，筑起墙壁，使她找不着路径。她必追随她的爱人，却追不上；她必寻找他们，却寻不着；那时她必说：我要去，回到我前夫那里，因为那时比现在更好。}然后，为使他们不把他们从真理的教训中所持守的健全事物归功于他们的诱惑者——这些诱惑者用这些误导他们进入自己错误教义和分裂的虚假——为使他们不以为他们里面的健全属于那些人，祂立刻加上：\Quote{她不知道是我给了她五谷、酒和油，并增加了她的金钱；她却用它们为巴耳制造金器和银器。}因为上文她说：\Quote{我要追随我的爱人，他们给我饼}等等，完全不明白一切被她诱惑者所健全合法持守的，都是出于天主，而非出于人。就连他们自己也不会将这些事物据为己有，仿佛拥有权利一样，若不是他们反过来被一个迷失的百姓所迷惑，当人们信任他们、给予他们如此尊荣时，他们就能借此说这样的话，将这些东西据为己有，使他们的谬误被称为真理，他们的不义被视为义，凭借他们所持有的圣事和圣经——不是为了得救，而仅仅是外表。因此，那同一个淫妇通过厄则克耳的口受到责备：\Quote{你拿着我赐给你的金银珠宝，用它们制造了男人的偶像，并与它们行淫；又拿我的绣花衣服遮盖它们，将我的油和香陈列在它们面前。我赐给你的食物——细面、油和蜂蜜，我用以喂养你的——你也摆在它们面前作为馨香之祭。}（\BibleRef{则 16:17}）因为她将一切圣事和圣书的话语转向她自己偶像的肖像，她的血肉之心在其中喜爱打滚。然而，尽管那些偶像是虚假的，是魔鬼的道理，是伪善的谎言（\BibleRef{弟前 4:1}），但这些圣事和神圣话语却不因此失去它们应有的尊荣，以致被认为属于这样的人；因为主说：\Quote{我的金子、我的银子、我的绣花衣服、我的油、我的香、我的食物}等等。难道我们因为那些迷失的人以为这些东西属于他们的诱惑者，就不去承认它们真正属于谁吗？当祂亲自说：\Quote{她不知道是我给了她五谷、酒和油，并增加了她的金钱}？因为祂并没有说她因为是个淫妇就没有这些东西；而是说她拥有它们，而且这些东西不属于她自己或她的爱人，只属于天主。因此，尽管她行淫，但她用以装饰淫行的那些东西，无论她是被诱惑的还是反过来诱惑的，都不属于她，而属于天主。如果这些话是针对犹太民族说的——那时经师和法利塞人为要建立自己的传统而废除了天主的诫命，他们如同与一个离弃天主的百姓行淫——然而，那时百姓中的淫行，如同主通过谴责所揭露的，并未使那些奥迹属于他们；这些奥迹不是他们的，而是天主的，祂在向淫妇说话时宣称这一切都是祂的；因此主自己也打发那些被治愈的癞病人去遵行同样的奥迹，在司祭面前为自己献祭，因为那祭献对他们并未成为有效——主后来愿意在教会中为所有人纪念这祭献，因为祂亲自向所有人传报喜讯——既然如此，当我们发现新约的圣事存在于某些异端或分裂者中时，我们更当不把这些圣事归功于人，也不当谴责它们，仿佛我们不能认出它们一样。当承认真正丈夫的恩赐，即使在一个淫妇手中；并且用真理的话语改正那属于不洁女人的真正淫行，而不是挑剔那些完全属于怜悯之主的恩赐。\par

28. 根据这些和类似的考虑，我们的先辈，不仅在西彼廉和阿格里皮努斯时代之前，而且在此之后，都保持了一个极为健全的习惯：每当他们发现在任何异端或分裂中仍有任何神圣和合法的事物保持完整时，他们予以认可而非拒绝；但无论发现什么异质的东西，属于那虚假教义或分裂的特有之物，他们就在真理的光照下予以驳斥并医治。不过，在尤巴伊安努斯的信中仍需考虑的问题，鉴于本书的篇幅，我认为应当留待另作开端来处理。\par

\SourceRecord{Translated by J.R. King and revised by Chester D. Hartranft.；Nicene and Post-Nicene Fathers, First Series；Vol. 4.；Edited by Philip Schaff.；Buffalo, NY: Christian Literature Publishing Co.,；1887.；<http://www.newadvent.org/fathers/14083.htm>.}

% structure: p0001 source=h1 target=section reason=page-title

\section{论圣洗，驳多纳徒派（卷四）}

\Emph{在本卷中，他继续处理西彼廉致朱巴雅努斯同一书信中的后续内容。}\par

% structure: p0003 source=h2 target=subsection reason=relative-heading-rank; base=h2

\subsection{第一章}

1. 将教会比作乐园，向我们表明，人固然可以在教会之外领受她的圣洗，但没有任何人能在外边获得或保留永福的救恩。因为，正如圣经所言，从乐园源泉涌出的溪流甚至丰沛地流到它的疆界之外。它们的名字确有记载；它们流经哪些地区，以及它们位于乐园界限之外，是众所周知的；\BibleRef{创 2:8} 然而在美索不达米亚和埃及——这些河流所到的地区——并未发现乐园中所记载的那种生命之福。因此，虽然乐园的水可在其边界之外找到，但唯独乐园中才有幸福的恩赐。同样，教会的圣洗可能存在外部，但幸福生命的恩赐唯独在教会内才能找到，教会建立在磐石上，领受了束缚与释放的钥匙。\BibleRef{玛 16:18} \Quote{唯有她拥有新郎与主的一切权能为己特权；}凭借这种作为新娘的权能，她甚至能由婢女生育儿子。这些人若不心高气傲，将被召入继承的份额；但若心高气傲，将留在外面。\par

% structure: p0005 source=h2 target=subsection reason=relative-heading-rank; base=h2

\subsection{第二章}

2. 既然如此，因着\Quote{我们是为教会的尊荣与合一而战}，我们要小心，不可将我们承认在他们中间属于教会的任何东西归功于异端者；反之，我们要用论证教导他们，他们从合一所获得的一切，若不回到那合一之中，对他们的救恩便毫无效力。因为\Quote{教会的水对那些正确使用它的人来说，充满了信德、救恩和圣德。}然而，没有任何人能善用这水于教会之外。但那些邪用的人，无论在内在外，这水都用于惩罚，而无助于他们的赏报。因此，圣洗\Quote{不能被败坏或玷污}，即使由腐败者或奸淫者施予，也是如此，正如\Quote{教会自身是不朽的、纯洁的、贞洁的。}因此，贪婪者、盗贼、放高利贷者——许多这样的人，按西彼廉本人在他多处书信中所证明的，不仅存在于教会之外，而且实际上存在于教会之内——没有份于教会。然而，他们既受洗也施洗，内心毫无改变。\par

3. 因为，他在一篇关于向天主祈祷的书信中致司铎团也说了这样的话；在那书信中，他效法圣达尼尔的方式，将他子民的罪归到自己身上。因为在他所提到的许多其他罪恶中，他也说到他们\Quote{只在言语上而非在行为上弃绝世俗}；正如宗徒论及某些人所说：\Quote{他们承认认识天主，但在行为上却否认祂。}\BibleRef{铎 1:16} 因此，有福的西彼廉表明，这些人在教会自身之内，他们受洗时内心并未向善转变，因为他们只在言语上而非在行为上弃绝世俗，正如伯多禄宗徒所说：\Quote{这水所预表的圣洗，现在也拯救你们，并不是除掉肉体的污秽，而是向天主要求一良心的清白。}\BibleRef{伯前 3:21} 他们当然没有这良心的清白，正如论到他们所说\Quote{只在言语上而非在行为上弃绝世俗}；然而他尽己所能，通过斥责和劝诫，使他们最终走在基督的道路上，成为祂的朋友，而不是世俗的朋友。\par

% structure: p0008 source=h2 target=subsection reason=relative-heading-rank; base=h2

\subsection{第三章}

4. 倘若他们听从了他，开始正直地生活，不是作为假基督徒而是作为真基督徒，他会命令他们重新受洗吗？当然不会；相反，他们真正的悔改会为他们赢得这一点：即当他们尚未改变时，圣事有碍于他们的毁灭，但一旦他们改变，它便开始有助于他们的救恩。\par

5. 因为那些似乎在内却相反基督而生活（即违背祂诫命）的人，并不\Quote{忠于教会}；他们也不能被认为以任何方式属于那教会，就是祂以\OriginalTerm{圣洗的水}{the washing of water}所净化的，\BibleQuote{好使祂为自己呈现一个荣耀的教会，没有斑点、皱纹或任何这类东西。}\BibleRef{弗 5:26}但如果他们不在那教会内，即不属于其肢体，他们就不在那被称为\BibleQuote{我的\OriginalTerm{鸽子}{dove}，只有一只；她是她母亲的独一者}的教会中\BibleRef{歌 6:9}；因为她本身是没有斑点、没有皱纹的。否则，让他断言那些只在言语上而非行为上弃绝世界的人是这鸽子的肢体吧。同时，有一件事我们看见，我想为此才说：\BibleQuote{那看日子的，是为主看的。}\BibleRef{罗 14:6}因为天主判断每一天。因为，根据祂的先见，祂知道谁在创世以前就预定要肖似祂圣子的肖像，许多甚至在公开的外面、被称为异端者的人，比许多良好的天主教徒更好。因为我们看见他们今天怎样，明天怎样我们不知道。但在天主那里，未来已成为现在，他们已经是将来要成为的。但我们根据每个人目前的状况，询问他们今天是否应被算在那称为独一的鸽子、无斑点无皱纹的基督\OriginalTerm{新娘}{Bride}的教会肢体中——关于这教会，\Person{西彼廉}在他上面引用的书信中说，\Quote{他们没有遵守主的道路，也没有遵行为他们的救恩所赐予的诫命；他们没有履行他们主的旨意，热衷于自己的财产和利益，追随骄傲的驱使，屈服于嫉妒和纷争，不关心纯一和信德，只在言语上而非行为上弃绝世界，各人只取悦自己，而令众人不悦。}但如果鸽子不承认他们在她肢体中，并且主若对他们说——假设他们继续在同一邪僻中——\BibleQuote{我从来不认识你们：你们这些作恶的人，离开我去吧！}\BibleRef{玛 7:23}那么，他们看似在教会内，其实不在；\Quote{不，他们甚至反对教会。那么他们怎能以教会的圣洗施洗呢？}这圣洗对他们自己和对从他们领受的人都没有益处，除非他们以真正的悔改改变心灵，以致那圣事本身——当他们只在言语上而非行为上弃绝世界而领受时，对他们并无益处——当他们也开始在行为上弃绝世界时，就开始对他们有益了。同样，对于那些公开分离的教会的人也是如此；因为无论是这些人还是那些人，都尚未在鸽子的肢体中，但其中一些人或许将来会在某个时候。\par

% structure: p0011 source=h2 target=subsection reason=relative-heading-rank; base=h2

\subsection{第四章}

6. 因此，当我们拒绝在他们之后施洗时，我们并不\Quote{承认\OriginalTerm{异端者的洗礼}{baptism of heretics}}；但因我们认出那礼规甚至是出于基督的，即使在恶人中间，无论是公开与我们分离的，还是在我们内部隐秘地割裂的，我们以应有的尊重接受它，同时纠正那些在错误之处走入歧途的人。然而，当我似乎被逼问时，有人对我说：\Quote{那么，异端者能赋予罪过的赦免吗？}，反过来我也追问：\Quote{那么，违背天上命令的人、贪婪者、强盗、放高利贷者、嫉妒者、只在言语上而非行为上弃绝世界的人，能赋予这样的赦免吗？}如果你是指天主\OriginalTerm{圣事}{sacrament}的力量，那么两者都能；如果是指他自身的功德，两者都不能。因为那圣事，即使在恶人手中，也被知道是属于基督的；但两者都不是在那独一无二、无玷、圣洁、贞洁、无斑点无皱纹的鸽子的身体内。正如圣洗对那只在言语上而非行为上弃绝世界的人无益，同样，对在异端或分裂中受洗的人也无益；但当每个人改正自己的行为时，便开始从那以前无益而早已在他内的圣事获得益处。\par

7. \Quote{因此，那在异端中受洗的，并不是成为\OriginalTerm{天主的殿}{the temple of God}；但难道因此就认为他是不算受洗的吗？}因为那在教会内受洗的\OriginalTerm{贪吝}{covetousness}者，除非他离开贪吝，也不能成为天主的殿；因为那些成为天主的殿的人，必然继承天主的国。但宗徒在许多事中说：\BibleQuote{无论是贪婪者、勒索者，都不能继承天主的国。}\BibleRef{格前 6:10}因为在另一处，同一宗徒将贪婪比作\OriginalTerm{偶像崇拜}{idolatry}：\BibleQuote{贪婪者，即偶像崇拜者}\BibleRef{弗 5:5}同一\Person{西彼廉}在致\Person{安托尼亚努斯}的信中如此延伸此意，以至于他毫不犹豫地将贪吝之罪与那些在迫害时期书面声明要献香的人之罪相比。那么，那在异端中因圣三之名受洗的人，除非他放弃异端，也不能成为天主的殿，正如那在同样名下受洗的贪吝者，除非他放弃贪吝（即偶像崇拜），也不能成为天主的殿。因为同一宗徒也说：\BibleQuote{天主的殿与偶像有什么相合呢？}\BibleRef{格后 6:16}因此，当我们说他完全没有成为天主的殿时，就不应问我们\Quote{他成了哪位天主的殿}。然而，他并非因此没有受洗，他的污秽错误也不能使他所领受的、以福音言语祝圣的\OriginalTerm{圣事}{sacrament}不再是神圣的圣事；正如另一个人的贪吝（即偶像崇拜）和大污秽，即使他被另一个同样贪吝的人以同样的福音言语施洗，也不能阻止他所领受的是神圣的圣洗。\par

% structure: p0014 source=h2 target=subsection reason=relative-heading-rank; base=h2

\subsection{第五章}

8. 西彼廉接着说：\Quote{此外，那些被理性所胜的人，徒然以习惯对抗我们，仿佛习惯优于真理，或者属灵之事上，圣神所启示的更好之道不应被遵循。}这显然正确，因为理性与真理应优先于习惯。但当真理支持习惯时，没有比那更应坚决持守的了。然后他继续如下：\Quote{因一个人若仅是犯错，尚可宽恕，如宗徒保禄论及自己说：\BibleQuote{我先前是褻瀆者、迫害者、凌辱者，但因我无知而为之，我蒙受了怜悯；}\BibleRef{弟前 1:13} 但人在受默感和启示之后，仍故意并明知地坚持以前的错误，则其罪过因无知而被赦免的希望全无。因他被理性所胜时，仍凭借某种傲慢与固执。}这极其真实：明知故犯之罪远较无知之罪为重。因此，对于圣西彼廉，他不仅博学，且乐于受教——他本人充分理解这是宗徒所描述的主教值得称赞的一部分，\BibleRef{弟后 2:24} 故他说：\Quote{这也应在主教身上受到认可：不仅以知识施教，也以耐心学习。}——我毫不怀疑，如果他有机会与那些虔诚而博学的人士讨论这个问题（此问题在教会中长久且多有争论，我们正是因他们，那古老习惯后来被全体会议的权威所确认），他必毫不犹豫地展示，不仅在其确信把握的真理上是何等博学，而且在未能察觉之事上何等易于受教。然而，既然明知故犯比无知之罪严重得多，我愿有人告诉我：哪一样更糟？——入异端而不知其罪之大，还是明知贪念之重却不肯放弃？我甚至可如此设问：若一人无知而堕入异端，另一人明知却不肯离弃偶像崇拜——因宗徒亲自说：\BibleQuote{贪婪者即偶像崇拜者}；而西彼廉在致安多尼安书中也以同样方式理解同一段经文，他说：\Quote{新异端者莫以此自夸，称他们不与偶像崇拜者共融，而他们之中却有奸淫者和贪婪者，后者的罪过等于偶像崇拜；\BibleQuote{因为你们该知道并明白：凡淫乱者、污染者、贪婪者（即偶像崇拜者）在基督和天主的国度里都无分；}\BibleRef{弗 5:5} 又：\BibleQuote{所以，要治死你们在地上的肢体：淫乱、污秽、邪情、恶欲和贪婪，贪婪即偶像崇拜。}}我请问：谁罪更深？——无知而堕入异端者，还是明知却不肯放弃贪婪（即偶像崇拜）者？按那准则：明知之罪优先于无知之罪，则明知而贪婪者罪居首位。但既然实际罪之重大可能对异端产生与明知对贪婪相同的影响，我们姑且假设无知异端者在罪责上与明知贪婪者相等——虽然西彼廉本人从宗徒所引的证据似乎未证明这点。因我们对异端者所憎恶的，岂非他们的褻瀆？但当他欲显示无知之罪可获罪赦之易时，他以宗徒的例子为证：\BibleQuote{我先前是褻瀆者、迫害者、凌辱者，但我因无知而为之，蒙受了怜悯。}\BibleRef{弟前 1:13} 但若可能——如我先前所说——让两人的罪过被视为等量：无知者的褻瀆与明知者的偶像崇拜相等；且以同一判决审判两者——一人寻基督却落入似真的虚谎陈述，另一人明知而抗拒基督借祂宗徒所说：\BibleQuote{因凡淫乱者、污染者、贪婪者（即偶像崇拜者）在基督和天主的国度里都无分，}\BibleRef{弗 5:5}——然后我请问：为何在前者情形中，圣洗和福音之言被视同无效，而在后者情形中却被承认为有效，尽管两者同样被发觉与鸽子的肢体隔绝？是否因前者是公开在圈外争斗，故不得被接纳；后者是在圈内狡猾地赞同，故不得被驱逐？\par

% structure: p0016 source=h2 target=subsection reason=relative-heading-rank; base=h2

\subsection{第六章}

9. 但关于他所说的\Quote{谁也不可断言，凡他们从宗徒所领受的，他们就遵行；因为宗徒只传授了唯一的教会和唯一的圣洗，而且这圣洗只在这同一教会内施行：}——这句话并不足以使我敢于谴责在异端者中发现的基督的圣洗（正如当我发现他们中有福音时，虽然我厌恶他们的错误，却仍必须承认这福音本身一样），反而提醒我，即使在圣\Person{西彼廉}的时代，已有人将非洲会议所反对的那项习俗追溯到宗徒的权威，而关于这习俗，他自己在前面就曾说过：\Quote{那些在理性上被驳倒的人，徒然拿习俗的权威来反对我们。}我也找不到原因，为什么同一位\Person{西彼廉}发现这习俗（这习俗在他之后得到了普世大公会议的确认）在他之前已经如此牢固，以致他虽竭尽学识寻找一个值得遵循的权威来改变它，却只找到了不久前在\Place{非洲}召开的\Person{阿格里皮努斯}会议。而鉴于这会议对他而言并不足够，因为它违反普世习俗，他便抓住了那些理由，我们刚才仔细考虑过这些理由，并因习俗本身的古老以及随后大公会议的权威而确信它们似真而非真；然而，这些理由在他看来却是真的，因为他是在极大的晦暗中努力探索，并对罪过的赦免——即在基督的圣洗中是否可能不赐予，以及在异端者中是否可能赐予——感到疑惑。在此事上，如果真理的不完全启示赐予了\Person{西彼廉}，为要显明他不离开教会合一的爱德之伟大，那么，任何人都不应仅仅因为他从一次大公会议的权威中得到教导，看到了\Person{西彼廉}所没有看到的东西（因为当时教会尚未就此问题召开大公会议），就自以为可以超越\Person{西彼廉}那坚强的护盾、卓越的德行以及他身上的丰盛恩宠。正如没有任何人会如此疯狂，自以为超过宗徒\Person{伯多禄}的功绩，因为他受了宗徒\Person{保禄}书信的教导，并得到了教会本身习俗的确认，就不再迫使外邦人犹太化，而\Person{伯多禄}曾一度那样做。\BibleRef{迦 2:14}\par

10. 我们并没有\Quote{发现任何人，在异端者中受圣洗后，后来被宗徒以同一圣洗接纳并领受了共融}；但我们也未发现这一点：任何人从异端者的团体中来，曾在他们中受圣洗，而后被宗徒重新施圣洗。但这习俗（即使在当时，那些回顾以往时代的人也不能发现它是由后代人发明的）被合理地相信是从宗徒传下来的。还有许多其他同类的事情，要一一列举就太冗长了。因此，如果那些\Person{西彼廉}想要说服他们自己观点之真理的人，对他们说：\Quote{谁也不要说，凡我们从宗徒所领受的，我们就遵行}，那么，我们现在说这话就更有力了：\Quote{教会习俗一直所持守的、这论证未能证明为假的、以及大公会议所确认的，我们就遵行！}此外，我们还可以说，在对讨论双方的理性和圣经证据进行了仔细研究后，\Quote{真理所宣布的，我们就遵行。}\par

% structure: p0019 source=h2 target=subsection reason=relative-heading-rank; base=h2

\subsection{第七章}

11. 事实上，关于有人反对\Person{西彼廉}的论证，认为宗徒说：\BibleQuote{无论以何种方式，或出于假意，或出于诚意，让基督被传扬}；\Person{西彼廉}正确地揭示了他们的错误，指出这与异端者无关，因为宗徒是在谈论那些在教会内行动、怀有恶意嫉妒、寻求自身利益的人。的确，他们按照我们相信基督的真理宣讲基督，但并非以善良的福音宣讲者向鸽子的儿女宣讲基督的精神。\Person{西彼廉}说：\Quote{因为\Person{保禄}在他的书信中并非谈论异端者或他们的洗礼，以至于能证明他对此事有所规定。他是在谈论弟兄们，他们或行为无序，违反教会纪律，或在敬畏天主中遵守教会纪律。他宣称其中一些人坚定无畏地宣讲天主圣言，但另一些人出于嫉妒和纷争行事；有些人以和善的基督徒之爱围绕自己，但另一些人怀有恶意和纷争；然而他耐心忍受一切，为的是无论出于诚意或假意，\Person{保禄}所宣讲的基督之名能被最多的人知晓，并且当时尚属新颖不习惯的播种圣言工作，能因那些说话者的宣讲而更广泛传播。此外，在教会内以基督之名说话是一回事，在教会外、反对教会的人以基督之名施洗则是另一回事。}这些话似乎警告我们，必须区分教会外的恶人和教会内的恶人。而他说宗徒描述为不洁地、出于嫉妒地宣讲福音的人，他说他们确实是在教会内。然而，我认为可以这样说不算鲁莽：如果教会外无人能拥有任何属于基督的事物，那么教会内也无人能拥有任何属于魔鬼的事物。因为如果那关闭的园子能容纳魔鬼的荆棘，为何基督的泉源不能同样流出园外？但如果它不能容纳它们，那么即使在宗徒\Person{保禄}本人时代，教会内怎么会出现如此巨大的嫉妒与恶意纷争的邪恶呢？因为这些都是\Person{西彼廉}的话。难道嫉妒和恶意纷争是一种小恶吗？那么那些不和睦的人怎能算是合一呢？因为这并非我的声音，也非任何人的声音，而是主亲自的声音；而且在基督诞生时，声音并非来自人类，而是来自天使：\BibleQuote{光荣归于至高天主，并在世上，和平归于善人。}当基督诞生于大地时，天使的声音肯定不会如此宣告，除非天主希望我们理解：那些与基督的和平合一的人，就是在基督身体的合一中；而那些心怀善意的人，就是在基督的和平中。此外，正如善意表现在仁慈中，恶意也表现在怨恨中。\par

% structure: p0021 source=h2 target=subsection reason=relative-heading-rank; base=h2

\subsection{第八章}

12. 简言之，我们可以看到嫉妒本身是多么大的恶，它不可能不是恶意的。我们无需寻求其他见证。西彼廉本人足以为我们作证，主借他的口极其真实地发出了如此多的雷霆，并关于嫉妒和恶意发出了许多有益的训诫。因此，让我们阅读西彼廉关于嫉妒和恶意的书信，看看嫉妒那些比我们好的人是多么大的恶——他令人难忘地指出，这种恶的起源来自魔鬼本身。他说：\Quote{认为某事是好的而心生嫉妒，并嫉妒那些比你更好的人，最亲爱的弟兄们，在有些人看来似乎是一种轻微且微不足道的过错。}再过一会儿，当他探究这恶的来源和根源时，他说：\Quote{由此，魔鬼在世界之初首先自身灭亡，并引领他人走向毁灭。}在同一章稍后：\Quote{最亲爱的弟兄们，这是一种多么大的恶，竟使一位天使堕落！那崇高而显赫的尊荣竟被欺骗和推翻！那欺骗者竟被欺骗！从那时起，嫉妒在大地上潜行，当人因恶意而即将灭亡时，他服从了那毁灭之师；当那嫉妒者效法魔鬼时，正如经上所载：\BibleQuote{因魔鬼的嫉妒，死亡进入了世界，凡属于他的，就必找到它。}}\BibleRef{智 2:24} 西彼廉这些话语多么真实、多么有力，它们出自一封闻名世界的书信，我们不能不承认。西彼廉真正应当最有力地辩论和告诫关于嫉妒和恶意的事，他通过自己丰富的基督徒爱德证明自己的心远离这致命的恶；他小心翼翼地保持爱德，从而与他的同僚保持共融的合一，这些同僚对圣洗持有不同看法却没有恶意，而他本人与他们意见不同，并非出于任何恶意的争论，而是由于人性的软弱，在一个问题上犯了错误，天主会因他持守爱德而在适当的时候向他启示。因为他公开说：\Quote{我不论断任何人，也不剥夺任何人与我们不同意见者的共融权。因为我们中没有人将自己立为主教中的主教，或通过暴虐的恐怖强迫他的同僚服从。}在本信末尾，他说：\Quote{最亲爱的弟兄，根据我微薄的能力，我已将这些简要地写给你，不规定或预断任何人，以免阻止每位主教在自由运用自己判断时做他认为正确的事。就我们而言，我们不为了异端者与我们的同仁和同僚主教争辩，我们与他们保持着天主所命令的和谐和主的平安，尤其是因为宗徒说：\BibleQuote{若有人似乎好争辩，我们却没有这样的惯例，天主的各教会也没有。}\BibleRef{格前 11:16} 我们灵魂中的基督徒爱德、我们弟兄情谊的尊荣、信德的纽带、司祭职的和谐，这一切我们都以忍耐和温和持守。为此，我们也根据我们微薄的能力，在主允准和默感之下，现在写了一篇关于忍耐益处的论文，我们出于彼此的友爱已将其送给你。}\par

% structure: p0023 source=h2 target=subsection reason=relative-heading-rank; base=h2

\subsection{第九章}

13. 凭着这种基督徒爱德的忍耐，他不仅忍受了他善良的同僚们在一个模糊问题上所表现出的极其友善的意见分歧（他自己也受到容忍，直到时机成熟，当天主乐意时，在一个全体公会议上，通过真理的宣告，那始终最为健康的习俗得到进一步确认），他甚至容忍那些明显败坏的人，这些事他自己非常清楚，他们并非因问题的模糊性而持不同观点，而是像宗徒论到他们时所说：\Quote{你教导人不可偷窃，你自己却偷窃吗？}\BibleRef{罗 2:21} 因为西彼廉在他的书信中谈到他那个时代这样的主教，他的同僚，并与他们保持共融：\Quote{当他们在教会中拥有饥饿的弟兄时，他们却试图积聚大量金钱，通过欺诈手段占有产业，通过累积的高利贷倍增收益。}因为这里没有模糊的问题。圣经公开宣告：\BibleQuote{无论是贪婪的、勒索的，都不能承继天主的国；}\BibleRef{格前 6:10} 以及\BibleQuote{放债取利的人，}\BibleQuote{并且没有淫乱的、污秽的、贪婪的（就是拜偶像的），在基督和天主的国里承受产业。}\BibleRef{弗 5:5} 因此，他肯定不会在不知情的情况下指控如此贪婪的罪行，即人们不仅贪婪地囤积自己的财物，而且欺诈地侵占他人的财物，或者指控如此严重的偶像崇拜（他理解并证明其存在）；他也肯定不会对自己的同僚主教作假见证。然而，他以父爱和母爱般的胸怀容忍他们，以免在薅莠子之前将麦子也薅出来，\BibleRef{玛 13:29} 这无疑是效法了宗徒保禄，他以对教会的同样爱德容忍那些对他怀有恶意和嫉妒的人。\BibleRef{斐 1:15}\par

14. 然而，因为\BibleQuote{因魔鬼的嫉妒，死亡才进入了世界，凡属于他的人，都经验到它}\BibleRef{智 2:24}，不是因为他们由天主所造，而是因为他们自行误入歧途，正如\Person{西彼廉}自己也说的，魔鬼在成为魔鬼之前，是一位天使，而且是善的，那么，属于魔鬼的人怎能处于基督的合一之内呢？毫无疑问，正如主亲自所说的，\Quote{是仇敌做了这事}，他\Quote{在麦子中间撒了莠子}。因此，正如在羊栈内属于魔鬼的必须被定罪，同样，在羊栈外属于基督的也必须被承认。魔鬼在教会的合一之内拥有属于他的，难道天主在外就没有属于他的吗？这一点或许可以就个别人来说，就是魔鬼在圣天使中没有属他的人，同样，天主在教会共融之外也没有属他的人。但是，尽管魔鬼被允许在这仍然带着肉体可朽本性的教会中混杂莠子，即恶人，当教会远离天主流浪时，他得到这个允许正是由于教会本身的朝圣之旅，好使人们更热切地向往那天使所享有的安息之所，然而对于圣事却不能这样说。因为，教会内的莠子能够拥有并处理它们，虽不是为救恩，而是为注定被火焚烧的毁灭，同样，教会外的莠子也能如此，他们从内部的分离者那里领受了圣事；因为分离并没有使他们失去圣事。事实上，这一点从以下事实可以清楚地看出：当他们回来时，如果正是那些分离者中的任何人碰巧回来，圣洗不再重新授予。并且，不要有人说：那莠子有什么果实呢？如果是这样，就此而言，他们的情况无论在内还是在外都是一样的。因为在内莠子中发现谷粒，而在外莠子中没有谷粒，这当然是不可能的。但是当问题涉及圣事时，我们不考虑莠子是否结出果实，而是考虑他们是否分享天上的恩赐；因为无论在内还是在外，莠子都与麦子本身共享雨水，这雨水本身是天上的、甘甜的，即使在其影响下莠子荒芜地生长。同样，按照基督的福音，圣事是神圣而悦人的；也不应因其甘露所落之人的荒芜而视之为无物。\par

% structure: p0026 source=h2 target=subsection reason=relative-heading-rank; base=h2

\subsection{第10章}

15. 但有人可能会说，内在的莠子更容易转化为麦子。我承认是这样；但这与重复圣洗的问题有何关系？你肯定不会主张，如果一个从异端皈依的人，借其皈依的机会和机遇，在另一个在教会内较慢洗净其不义、从而改正和改变的人之前结出果实，那么前者因此不需要再领圣洗，而那个教会成员却需要再领圣洗，因为他比那个来自异端的人进步得更慢，因为他的悔改更为迟缓。因此，谁从自己特殊的乖僻转向信德、望德或爱德的直路时较晚或较慢，这与当前的问题无关。因为虽然羊栈内的恶人更容易变善，但有时会发生，某些外面的人会在皈依上超过某些里面的人；而当这些人仍然处于荒芜时，那些人恢复合一与共融，会以坚忍结出果实，有三十倍的、六十倍的、一百倍的。或者，如果只有那些顽固错误到底的人才被称为莠子，那么外面有许多麦穗，里面有许多莠子。\par

16. 但有人会争辩说，外面的恶人比里面的更坏。这确实是一个重要的问题，是否已经与教会分离的\Person{尼苛劳}\BibleRef{默 2:6}，以及仍然在教会内的\Person{西满}\BibleRef{宗 8:9}，谁更坏——一个是异端者，另一个是巫士。但如果仅仅因为分裂作为违反爱德的最明显标记而被视为更坏的恶，我承认是这样。然而许多人虽然失去了全部爱德之情，却因世俗利益而不分裂；并且他们寻求自己的事，而不是耶稣基督的事\BibleRef{斐 2:21}，他们不愿离开的，不是基督的合一，而是自己暂时的利益。因此，在赞美爱德时说，她\BibleQuote{不求自己的益处}\BibleRef{格前 13:5}。\par

17. 现在，问题在于：属于魔鬼党派的人怎能属于那没有瑕疵、皱纹或任何这类缺点的教会，关于她也曾说：\BibleQuote{我的鸽子是唯一的？}\BibleRef{歌 6:9} 但如果他们不能，显然她是在那些不属于她的人中间呻吟，有些人背信地在内设伏，有些人则在外对着她的门口吠叫。然而，这样的人们，即便在内，也领受圣洗、拥有圣洗，并将其本身神圣地传授；他们的邪恶，即使他们坚持到底，也丝毫不能玷污它。因此，同一位真福西彼廉教导我们，圣洗应被视为由福音的话语本身所祝圣，正如教会所领受的，不将施行者或领受者的任性与邪恶的任何考虑附加或混入其中；因为他亲自向我们指明了两项真理——既指出在教会内有些人不珍爱友善的基督徒爱德，反而施行嫉妒与不善的分裂，宗徒保禄曾论及这些人；也指出嫉妒者属于魔鬼的党派，正如他在关于嫉妒与恶意的书信中最公开地作证。因此，既然很可能属于魔鬼党派的人中，基督的圣事仍可能是神圣的——当然，不是为他们的救恩，而是为他们的定罪；不仅在他们领洗后被引入歧途时如此，甚至当他们在领受圣事时内心就是如此，只在言语上而非行动上弃绝世俗（如同同一位西彼廉所显示的那样）；并且即使后来他们被引上正道，他们曾在迷途中领受的圣事也不应重授；就我所能见，案情已清楚明了：在圣洗的问题上，我们必须考虑的，不是谁施予，而是他施予什么；不是谁领受，而是他领受什么；不是谁拥有，而是他拥有什么。因为如果属于魔鬼党派、因此绝不属于那唯一鸽子的人，竟能领受、拥有并施予那完全神圣、丝毫不被其任性所玷污的圣洗，正如西彼廉自己的书信所教导我们的，那么我们怎能将不属于异端者的东西归给他们呢？我们怎能说那真正属于基督的东西是他们的呢？我们岂不应当在他们身上认出我们君王的标记，并纠正他们背弃祂的行为吗？因此，如圣西彼廉所说：\Quote{在教会内的人因基督之名发言是一回事，在外且行事反对教会的人以祂之名施洗则是另一回事。}但许多在内的人因恶行和引诱软弱的灵魂效法他们的生活而行事反对教会；也有些在外的人因基督之名发言，并不被禁止行基督的工程，只是不得在内，因为为了医治他们的灵魂，我们设法抓住他们、与他们理论或劝告他们。因为那没有跟随基督与祂的门徒同行、却因基督之名驱魔的人也是在外，主曾吩咐不要阻止他；\BibleRef{路 9:49} 当然，在他不完善之处，他应按照主的话得以痊愈，主说：\BibleQuote{不与我相合的，就是反对我的；不同我收集的，就是分散的。}\BibleRef{玛 12:30} 因此，有些事是奉基督之名在外行作，并不反对教会；有些事是在内出于魔鬼的部分行作，反对教会。\par

% structure: p0030 source=h2 target=subsection reason=relative-heading-rank; base=h2

\subsection{第十一章}

18. 我们该怎样说那令人惊奇的事呢？细心观察的人可能发现，某些人虽不违反基督徒爱德，却可能教导无益之事，正如伯多禄曾想强迫外邦人遵守犹太习俗一样，\BibleRef{迦 2:14} 如同西彼廉本人会强迫异端者重领圣洗一般。因此，宗徒对这样一些根植于爱德、却在某些事上行走不正的好肢体说：\BibleQuote{你们若在什么事上存有异见，天主也必将这事启示给你们；}\BibleRef{斐 3:15} 而另有一些人，虽缺乏爱德，却可能教导一些有益的事？主论及他们说：\BibleQuote{经师和法利塞人坐在梅瑟的座位上：凡他们所吩咐你们遵守的，你们都要遵守并实行；但不要效法他们的行为，因为他们只说不做。}\BibleRef{玛 23:2} 因此，宗徒也论及那些嫉妒和恶意的人，他们仍借着基督宣讲救恩：\Quote{无论是出于虚假，还是出于真诚，总要让基督被传扬。}所以，无论在内外，人的任性都当被纠正，但神圣的圣事和言语不可归于人。因此，拒绝将明知不属于异端者的东西归于他们的人，并不是\Quote{异端者的庇护者}。我们不承认圣洗是他们的；我们承认那位的圣洗，关于祂曾说：\BibleQuote{那一位就是施洗者，}\BibleRef{若 1:33} 无论我们在何处找到它。但如果\Quote{背信和亵渎的人}继续行背信和亵渎，他无论在教会之外或之内都得不到\Quote{罪的赦免}；或者，如果因圣事的能力他暂时得到赦免，那么同样的能力在内外都运作，正如基督之名的能力过去甚至在教会之外也能驱逐魔鬼。\par

% structure: p0032 source=h2 target=subsection reason=relative-heading-rank; base=h2

\subsection{第十二章}

19. 但是他辩称，\\Quote{我们发现宗徒们在所有书信中都咒骂和憎恶异端者的渎神邪恶，以致说：\\InnerQuote{他们的言语如同毒疮蔓延。}} 那么怎样？难道\\Person{保禄}不是也显示，那些说\\Quote{我们吃吃喝喝吧，因为明天就要死了}的人，以不良的交往败坏了良好的风俗，并立即补充说\\Quote{不良的交往败坏良好的风俗}；然而他暗示这些人在教会内，当他说\\Quote{你们中有人说没有死人复活}呢？但他何时没有表达对贪婪之人的憎恶？或者，有什么能比同一位宗徒所宣称的，贪婪应被称为偶像崇拜更严厉的话呢？\\BibleRef{弗 5:5} \\Person{西彼廉}对这句话的理解也无不同，在需要时将其插入自己的书信中；虽然他承认，在他当时的教会里，不仅有普通的贪婪之人，更有强盗和放高利贷者，而且这些人并非出现在平信徒中，而是出现在主教中。然而，我本愿理解宗徒所说的\\Quote{他们的言语如同毒疮蔓延}是指教会之外的人，但\\Person{西彼廉}本人不允许我这样理解。因为，在他致\\Person{安多尼安}的书信\\WorkTitle{致安多尼安书}中，他指出，任何人不应仅仅因为教会中混杂了恶人，就在义人与不义之人最终分离之前脱离教会的合一；当他表明自己何等圣洁，配得上他所赢得的辉煌殉道时，他说：\\Quote{这是何等狂妄的傲慢，何等的谦卑与温良的遗忘，竟有人敢相信或自认能做主都没有赐予宗徒们的事——即认为他能分清莠子和麦子，或以为主赐给了他簸箕和扬净禾场的能力，试图将糠秕从麦粒中分离出来！而宗徒说：\\InnerQuote{但在大户人家，不但有金器银器，也有木器瓦器}，\\BibleRef{弟后 2:20} 他就似乎应该选择金器银器，而鄙视、丢弃和谴责木器瓦器；其实木器只应在主的日子被天火的烈焰焚烧，瓦器则应由那已获得\\InnerQuote{铁杖}的那一位打碎。} 因此，藉此论证，\\Person{西彼廉}针对那些以逃避恶人交往为借口而脱离教会合一的人表明，宗徒所说的\\OriginalTerm{大屋}{great house}，其中不但有金器银器，也有木器瓦器，他所理解的不是别的，正是教会，其中既有善人也有恶人，直到最后的日子它会被簸箕净化为打谷场。如果这样，在教会自身，即大屋本身，就有不体面的器皿，他们的言语如同毒疮蔓延。因为宗徒论到他们说：\\Quote{他们的言语如同毒疮蔓延；其中有\\Person{希默纳约}和\\Person{斐理托}，他们偏离了真理，说复活已过，因而颠覆了一些人的信德。然而，天主的稳固基础屹立不动，且带着这印记：\\InnerQuote{主认识属于祂的人。}并且\\InnerQuote{凡称呼基督之名的人，应远离不义。}但在大户人家，不但有金器银器，也有木器瓦器。} \\BibleRef{弟后 2:17} 因此，如果那些言语如同毒疮蔓延的人，好比大屋中不体面的器皿，而\\Person{西彼廉}以那\\Quote{大屋}理解为教会本身的合一，那么他们的毒疮当然不可能污染基督的圣洗。因此，无论在教会之外还是之内，属于魔鬼阵营的人，都不能在自己身上或任何人身上玷污属于基督的圣事。因此，并不是\\Quote{那如同毒疮蔓延入听者耳中的言语能给予罪过的赦免}；但当圣洗是按福音的话语施行时，无论施行者或领受者的理解何等悖谬，圣事本身由于它是属于那位的圣事，而自为圣洁。如果有人从迷途者手中领受圣洗，却没有接受施礼者的谬误，而只接受奥秘的神圣，并籍着信德、望德和爱德紧密联结于教会的合一，他就获得罪过的赦免——不是凭借那如同毒疮侵蚀的言语，而是凭借来自天上泉源的福音圣事。但如果领受者本人是迷途的，那么所施与的圣事对迷途者的救恩毫无益处；然而另一方面，他所领受的仍在他内保持神圣，若他日后归于正途，也不再为他重行。\par

% structure: p0034 source=h2 target=subsection reason=relative-heading-rank; base=h2

\subsection{第十三章}

20. 因此，\\Quote{正义与不义之间没有交情}，不仅在教会之外，也在教会之内；因为\\Quote{主认识属于祂的人}，并且\\Quote{凡称呼基督之名的人，应远离不义}。此外，\\Quote{光明与黑暗之间也没有相通}，\\BibleRef{格后 6:14} 不仅在教会之外，也在教会之内；因为\\Quote{恼恨自己弟兄的，仍在黑暗中}。\\BibleRef{若一 2:9} 那些人确实恼恨\\Person{保禄}，他们出于嫉妒和纷争宣讲基督，以为这样能增加他锁链的苦楚；\\BibleRef{斐 1:15} 然而同一位\\Person{西彼廉}认为这些人仍在教会之内。因此，既然\\Quote{黑暗不能光照，不义不能成义}，正如\\Person{西彼廉}再次所说，我问：那些人怎能在教会自身之内施行圣洗呢？我问：那大屋中所包容的不体面器皿，怎能在大屋本身之内施行神圣之事来圣化人们呢？除非是因为圣事的神圣性甚至不能被不洁者所污染，无论是他们施行时，还是那些内心和生活没有改善的人领受时，都一样。关于这些身处教会之内的人，\\Person{西彼廉}自己说：\\Quote{他们只在口头上弃绝世界，而非在行动上。}\par

21. 因此，在教会内也有\Quote{心被假基督之灵所占据的天主仇敌}；然而他们\Quote{处理属灵和神圣的事}，只要他们保持原样，这些事便不能为他们的救恩带来益处；他们也不能用自己的不洁玷污它们。因此，关于他所说的：\Quote{他们在教会的救恩恩宠中没有分得任何部分，那些分散并攻击基督教会的人，被基督亲自称为仇敌，被祂的宗徒称为假基督者}，这一点必须考虑到，这类人既有在教会内的，也有在教会外的。但那将教会内的人与鸽子的完美和合一分离，不仅在某些人身上为天主所知，甚至在某些人身上也为他们的同伴所知；因为，通过观察他们公开堕落的生活和根深蒂固的邪恶，并将其与天主诫命的规则相比较，他们明白，对于大批的莠子和糠秕，如今有些在教会内，有些在教会外，但注定要在末日最明显地分开，那时主将说：\BibleQuote{你们这些作恶的人，离开我吧！}\BibleRef{玛 7:23} 和 \BibleQuote{你们该往永火里去，那是为魔鬼和他的使者预备的。}\BibleRef{玛 25:41}\par

% structure: p0037 source=h2 target=subsection reason=relative-heading-rank; base=h2

\subsection{第十四章\newline{}}

22. 但我们不可对任何人的悔改绝望，无论他在教会内或在教会外，只要\BibleQuote{天主的恩慈引导他悔改}\BibleRef{罗 2:4}，且\Quote{用杖责罚他们的过犯，用鞭责罚他们的罪孽}。因为这样\Quote{祂必不将祂的慈爱全然收回}，如果他们自己有时\Quote{爱自己的灵魂，悦乐天主}。但正如那\BibleQuote{忍耐到底的必然得救}\BibleRef{玛 24:13}，同样，那恶人，无论在内或在外，若坚持作恶到底，必不得救。我们也不是说\Quote{任何地方、以任何方式受洗的人都获得圣洗的恩宠}，如果圣洗的恩宠是指通过圣事庆典所赐予的实际救恩；但许多人即使在教会内也未能获得这救恩，虽然他们显然拥有圣事本身，而圣事本身是圣的。主在福音中警戒我们，不要与恶顾问交往，\BibleRef{谷 13:21}他们以基督的名号行走；但这类人既有在教会内的，也有在教会外的，事实上，除非他们先在教会内怀有恶意，否则不会从教会出去。我们知道宗徒论到大户人家中的器皿时说：\BibleQuote{人若自洁，脱离这些事，就必成为贵重的器皿，成圣，合乎主用，预备行各样的善事。}\BibleRef{弟后 2:21}但他更上面一点说明了各人应如何自洁，说：\BibleQuote{凡称呼主名的人，务要远离不义。}\BibleRef{弟后 2:19}免得在末日，与糠秕——无论是已从禾场吹走的，或是最后将要分离的——一同听到命令：\BibleQuote{你们这些作恶的人，离开我吧！}\BibleRef{玛 7:23}由此确实显明，正如\Person{西彼廉}所说：\Quote{我们不可即刻接纳并认可凡以基督之名所宣称的一切，而只可接纳那在基督真理中所行的事。}但以下行为并非在基督真理中所行：人们\Quote{以欺诈的借口强占田产，用累积的高利贷增加收益}，或他们\Quote{只在言语上弃绝世俗；}然而，这一切在教会内发生，\Person{西彼廉}本人便提供了充分的见证。\par

% structure: p0039 source=h2 target=subsection reason=relative-heading-rank; base=h2

\subsection{第十五章\newline{}}

23. 接着，他长篇大论地追究这一点，即\Quote{那些亵渎基督的圣父的人不能在基督内领受圣洗}，因为显然他们因错误而亵渎（因为前来领受基督圣洗的人不会公开亵渎基督的圣父，而是因持有与关于基督的圣父的真理教导相反的观点而被引向亵渎），我们已经充分证明，以福音之言祝圣的圣洗，不受任何人（无论是施行者还是领受者）错误的影响，无论此人持有与天主圣父、圣子或圣神的神圣启示相反的见解。因为即使在教会内，也有许多属血肉和属本性的人领受圣洗，正如宗徒明确说的：\BibleQuote{属血气的人不接受天主圣神的事。}\BibleRef{格前 2:14} 而他们领受圣洗后，他又说他们\BibleQuote{仍是属血肉的。}\BibleRef{格前 3:3} 但根据血肉的理智，屈服于肉欲的灵魂只能对天主持有血肉的智慧。因此，许多人在领受圣洗后，特别是在婴儿或幼年时期领受圣洗的人，随着理智日益清晰明亮，\BibleQuote{内心的人却是一天比一天更新}\BibleRef{格后 4:16}，便唾弃并痛悔以往被虚妄想象所愚弄时对天主持有的见解，承认自己的错误。然而，他们并不因此被认为没有领受圣洗，或领受了与其错误相符的圣洗；相反，在他们身上，圣事的完整性受到尊重，而思想的迷惑得以纠正，即使这迷惑因长期确认而根深蒂固，或曾在许多争论中被坚持。因此，即使是在明显外边的异端者，如果他在那里领受了按福音规定的圣洗，他当然不是领受了与他所盲目的错误相符的圣洗。所以，当他返回智慧之路时，他认识到应当放弃他所持的错误，但同时也绝不能放弃他已领受的善；也不应因为他的错误应受谴责，就因此熄灭他内在的基督的圣洗。因为从那些在教会内领受圣洗却持有关于天主的错误观点的人的例子中，已经足够清楚地表明，圣事的真理应与错误相信者的错误区分开来，尽管两者可能同在一人身上。因此，当任何根植于错误（即使是在教会外）的人，已领受了真正的圣事圣洗，当他被恢复至教会的合一中时，真正的圣洗不能取代真正的圣洗，如同真正的信仰取代虚假的信仰，因为一个事物不能取代自身，既不能取代，也不能让位。因此，异端者加入天主教会，正是为了纠正他们自身所有的恶，而不是重复他们从天主那里已有的善。\par

% structure: p0041 source=h2 target=subsection reason=relative-heading-rank; base=h2

\subsection{第十六章。}

24. 有人说：如果两个人同样根植于错误和邪恶，在没有改变生活和内心的情况下领受圣洗，一个在教会外，一个在教会内，难道没有区别吗？我承认有区别。因为在教会外领受圣洗的人更恶劣，除了他的其他罪之外——但不是因为他的圣洗，而是因为他在外边；因为分裂的恶本身绝非微不足道或无关紧要。然而，这种区别只在于：如果在教会内领受圣洗的人渴望在教会内，不是为任何属世或暂时的利益，而是因为他更喜欢普世教会的合一而非分裂的纷争；否则，他也应被视为那些在外边的人。因此，我们这样设定两种情况。假设一个人，为了论证，对基督持有与\Person{Photinus}相同的看法，并在他的异端中，在天主教会的共融之外领受圣洗；另一个人持有相同的看法，但在天主教会内领受圣洗，相信他的看法正是天主教信仰。我认为他尚未成为异端者，除非当天主教信仰的教义向他阐明时，他选择抵抗，并偏向自己已有的见解；直到这种情况发生之前，显然在外边领受圣洗的人更恶劣。因此，在一种情况下，只有错误的见解需要纠正；在另一种情况下，还有分裂的罪也需要纠正；但在这两种情况下，圣事的真理都不应重复。但是，如果有人与第一个人持相同见解，并且知道这样的见解只在脱离教会的异端中才被教导或学习，却为了一些暂时的利益而渴望在天主教合一中领受圣洗，或者已经领受圣洗后，因着上述利益而不愿脱离教会，那么他不应仅被视为分裂者，而且他的过犯更严重，因为他在异端错误和分裂合一之外，还加上了伪善的欺诈。因此，每个人的堕落，越是危险和缺乏正直，就越应当以更大的热忱和力量加以纠正；然而，如果他身上有任何善，尤其不是出于自己，而是来自天主，我们不应因他的堕落而轻视这善，或像对待堕落那样谴责它，或将它归于堕落，而应归功于天主的丰盛美善，祂甚至向那卖淫、追随情人的灵魂，仍赐予她的面包、她的酒、她的油，以及其它食物或装饰，这些既不是来自她自己，也不是来自她的情人，而是来自祂，祂出于怜悯甚至愿意警告她应当归向谁。\BibleRef{欧 2:5}\par

% structure: p0043 source=h2 target=subsection reason=relative-heading-rank; base=h2

\subsection{第十七章。}

25. 西彼廉说：\Quote{圣洗的能力能比宣认或殉道更大或更好吗？即人在人前承认基督，并以自己的血受洗？}然而，他接着说：\Quote{即使为了承认基督而在教会外被处死，这圣洗也不能使异端者获益。}这是千真万确的；因为，在教会外被处死，证明他没有爱德，关于爱德，宗徒说：\BibleQuote{我若舍身被焚，却没有爱德，于我毫无益处。}\BibleRef{格前13:3}但如果殉道因没有爱德而毫无益处，那么那些如保禄所说、西彼廉进一步阐述的、在教会内生活却没有爱德、心怀嫉妒和恶意的人，也同样不能获益；然而他们却能领受和传递真正的圣洗。他说：\Quote{救恩不在教会之外。}谁说是呢？因此，人所拥有的一切属于教会的东西，在教会外对救恩毫无益处。但一件事是没有，另一件事是拥有却毫无用处。没有的人必须受洗以拥有；但拥有却毫无用处的人必须被纠正，使他所拥有的能有益于他。异端者的圣洗之水也不是\Quote{淫乱的}，因为天主所造的受造物本身不是邪恶的，而且在任何走上歧途的人口中，福音的话语也不能被指摘；但过错在于他们身上有淫乱的精神，即使这精神从合法的配偶那里领受了圣事的装饰。因此，圣洗可以\Quote{为我们和异端者所共有}，正如福音可以为我们共有，无论我们的信德与他们的错误之间有何差异——无论他们对圣父、圣子或圣神的看法与真理不同；还是从合一中被剪除，不与基督一同聚集，反而分散，\BibleRef{玛12:30}——看到圣洗的圣事可以为我们（如果我们是主的麦子）与教会内贪婪的人、强盗、醉汉以及类似的有害之人所共有，正如经上所说：\BibleQuote{他们不能继承天主的国}\BibleRef{格前6:10}，然而将他们与天主的国分开的恶习并不为我们所共有。\par

% structure: p0045 source=h2 target=subsection reason=relative-heading-rank; base=h2

\subsection{第十八章。}

26. 当然，宗徒说\BibleQuote{行这样事的人不能继承天主的国}，并不仅仅指异端。但值得花片刻看看他所列举的事。他说：\BibleQuote{肉体的行为是明显的，即：淫乱、不洁、放荡、崇拜偶像、行邪术、仇恨、纷争、嫉妒、忿怒、争吵、分党、异端、妒忌、凶杀、醉酒、宴乐，以及类似的事。关于这些，我先前已告诉过你们，现在再次告诉你们：行这样事的人不能继承天主的国。}\BibleRef{迦5:19}因此，假设有人贞洁、节制、不贪婪、不崇拜偶像、好客、周济穷人、不与人为敌、不争竞、忍耐、安静、不嫉妒、不忌恨、清醒、节俭，但他是异端者；当然，所有人都清楚，仅因这一个过犯——他是异端者——他就不能继承天主的国。再假设另一个人是行淫者、不洁者、放荡者、贪婪者，甚至更公开地崇拜偶像、研习邪术、喜爱争吵和纷争、嫉妒、暴躁、好结党、忌恨、醉酒、宴乐，但他是天主教徒；难道仅因他是天主教徒这一功绩，他就能继承天主的国吗？尽管他的行为属于宗徒所断言的那一类：\BibleQuote{关于这些，我先前已告诉过你们，现在再次告诉你们：行这样事的人不能继承天主的国？}如果这样说，我们就自欺了。因为天主的言语不会欺骗我们，它既不沉默，也不宽容，更不因任何奉承而骗人。事实上，它在别处也这样说：\BibleQuote{因为你们知道，凡是行淫乱的、不洁的、贪婪的（即是崇拜偶像的），在基督和天主的国里，都得不到产业。不要让人用空言欺骗你们。}\BibleRef{弗5:5}因此，我们没有理由抱怨天主的言语。它确实说了，而且公开自由地说了：活在邪恶生活中的人，在天主的国里无份。\par

% structure: p0047 source=h2 target=subsection reason=relative-heading-rank; base=h2

\subsection{第十九章。}

27. 因此，我们不要奉承这位被所有这些恶习包围的天主教徒，也不要仅因他是天主教基督徒就冒险向他许诺圣经并未许诺他的豁免；如果他有上述任何过失，我们也不应许诺他分享那永恒的天国。因为，在写给格林多人的书信中，宗徒列举了几项罪恶，每一项都隐含地表明不能继承天主的国。他说：\BibleQuote{不要受骗：无论是淫乱者、拜偶像者、奸淫者、作娈童者、同性恋者、偷窃者、贪婪者、醉酒者、辱骂者、勒索者，都不能继承天主的国。}\BibleRef{格前 6:9} 他并没有说，同时拥有所有这些恶习的人不能继承天主的国；而是说这些或那些都不能；因此，当每个都被提及时，你应当明白，其中任何一个都不能继承天主的国。因此，正如异端者不能拥有天主的国，贪婪的人也不能继承天主的国。而且我们确实不能怀疑，那些不能继承天主的国的人所遭受的惩罚本身，会按照其过犯的差异而有所不同，有些惩罚会比另一些更严厉；以致在永火本身中，根据罪过的不同重量，惩罚中会有不同的折磨。因为主并非无故说：\BibleQuote{在审判的日子，索多玛地方所受的惩罚，也要比你们容易忍受。}\BibleRef{玛 11:24} 但是，就未能继承天主的国而言，无论你选择这些恶习中哪一个较轻的，都同样确定，如同你选择多个，或者你选择某个你认为更可憎的；因为那些被审判者安置在祂右边的人将继承天主的国，而对于那些不配被安置在右边的人，除了被安置在左边之外，没有别的结局，他们像山羊一样从牧羊人口中所听到的唯一宣告是：\BibleQuote{可咒骂的，离开我，到那给魔鬼和他的使者预备的永火里去！}\BibleRef{玛 25:41} 尽管在那火中，正如我先前所说，惩罚可能会根据罪过的不同而有不同的分配。\par

% structure: p0049 source=h2 target=subsection reason=relative-heading-rank; base=h2

\subsection{第二十章。}

28. 但关于我们是否应当将一个品行最败坏的天主教徒优先于一个除异端外生命无可指责的异端者，我不敢草率判断。但如果有人说，因为他是异端者，所以他不可能只此一项而没有其他恶习随之而来——因为他是属血肉和属本性的，因此也必定是嫉妒的、易怒的、猜忌的、与真理本身为敌且完全背离真理的——那么让他公正地理解，在他被认为已选择某一不太严重的其他过失中，没有一项能在任何人身上单独存在，因为他反过来也是属血肉和属本性的；就以醉酒为例，人们现在不仅不感到恐怖，反而有些高兴地谈论它，它能独自存在于发现它的任何人身上吗？因为哪个醉汉不是好争论的、易怒的、嫉妒的、与所有健全的劝告不合、且严厉敌视责备他的人呢？此外，他很不容易避免成为淫乱者和奸淫者，尽管他可能不是异端者；正如异端者可能不是醉汉、奸淫者、淫乱者、纵欲者、贪财者、行巫术者，而且不可能同时拥有所有这些。事实上，没有一项恶习会随之带来所有其他恶习。因此，假设两个人——一个是具有所有这些恶习的天主教徒，另一个是免受异端者所能避免的所有恶习的异端者——虽然他们不是都对抗信仰，但每个人都违背信仰而生活，每个人都受虚幻希望的欺骗，每个人都远离属神之爱德，因此每个人都与独一鸽子的身体断绝联系；那么为什么我们在一个人中承认基督的圣事，而在另一个人中不承认，好像它属于这个人或那个人，而实际上它在两者中是相同的，只属于天主，即使在最坏的人中也是好的呢？如果拥有它的人中，一个比另一个更坏，这并不意味着他们所拥有的圣事在一个中比在另一个中更坏，因为即使在两个不好的天主教徒中，如果一个比另一个更坏，他并不拥有更坏的圣洗，如果一个是好的，另一个是坏的，圣洗在坏的中不是坏的，在好的中也不是好的；而是在两者中都是好的。就像太阳光甚至灯的光，在不好的人眼前无疑不比在较好的人眼前黯淡；但它在两者中是相同的，尽管它根据他们能力的差异而使他们受益或伤害不同。\par

% structure: p0051 source=h2 target=subsection reason=relative-heading-rank; base=h2

\subsection{第二十一章。}

29. 关于西彼廉所受的反对意见，就是那些在殉道中被捕、为基督之名而被杀的人，即使没有领洗也获得了冠冕，我不太明白这与本题有何相干，除非他们实在主张，异端者可以带着洗礼更容易地进入基督的王国，而慕道者没有洗礼也能进入，因为主亲自说过：\BibleQuote{人除非由水和圣神而生，不能进天主的国。}\BibleRef{若 3:5} 在此事上，我毫不迟疑地将公教慕道者——就是那怀着对天主的炽热爱情者——置于已受洗的异端者之前；我们也并不因此对后者已领受的圣洗圣事有所不敬，尽管前者尚未领受；我们也不认为慕道者的圣事优于圣洗圣事，因为我们承认有些慕道者比某些已受洗的人更好、更忠实。例如百夫长科尔乃略，在受洗前已比已受洗的西满更好。因为科尔乃略在受洗前就充满了圣神\BibleRef{宗 10:44}；西满却在受洗后被邪灵所充满。然而，科尔乃略若在受了圣神后仍拒绝受洗，就必被定以轻视如此神圣圣事之罪。但他受洗时，所领受的圣事绝不低于西满；但不同的人的不同功过，在同一圣事的同等神圣性下显明出来——如此真实的是，领受者的善功或恶行并不增加或减少圣洗的神圣性。但正如一位善良的慕道者因缺少圣洗而不能进入天国，同样一个已受洗的恶人因缺少真正的皈依也不能进去。因为那位说过\BibleQuote{人除非由水和圣神而生，不能进天主的国}的主，也亲自说过：\BibleQuote{除非你们的义德超过经师和法利塞人的义德，你们决进不了天国。}\BibleRef{玛 5:20} 为使慕道者的义德不自以为安稳，经上记载：\BibleQuote{人除非由水和圣神重生，不能进天主的国。} 再者，为使已受洗者的不义不因领受了圣洗而自恃安稳，经上记载：\BibleQuote{除非你们的义德超过经师和法利塞人的义德，你们决进不了天国。} 两者缺少其一都不够；两者齐全，方能使那继承产业的嗣子臻于完美。因此，正如我们不应轻视一个人的义德——这义德在他加入教会之前就已开始存在，如同科尔乃略的义德在他进入基督徒团体之前就已存在——这义德并非毫无价值，否则天使就不会对他说：\BibleQuote{你的祈祷和你的施舍已升到天主面前，获得纪念}\BibleRef{宗 10:4}；但这也还不足以使他获得天国，否则就不会告诉他派人去找伯多禄了。同样，我们也不应轻视圣洗圣事，即使它是从教会之外领受的。但除非那已完满地领受了圣洗的人，同时也纠正自己的邪恶，被纳入教会，否则这圣洗对救恩毫无益处。因此，让我们纠正异端者的错误，好能认出在他们身上哪些不是出于他们自己，而是属于基督的。\par

% structure: p0053 source=h2 target=subsection reason=relative-heading-rank; base=h2

\subsection{第二十二章}

30. \OriginalTerm{圣洗}{baptism}的地位有时可以由\OriginalTerm{殉道}{martyrdom}来补充，这是真福西彼廉从那个强盗的例证中提出的一个绝非微不足道的论据；那人虽然没有受洗，但主却对他说：\BibleQuote{今天你就要同我一起在乐园里。}\BibleRef{路 23:43}。我一再思考这一点，发现不仅为基督而殉道可以补足圣洗的欠缺，而且\OriginalTerm{信德}{faith}和心灵的归化也可以，若是因时间不足而无法举行圣洗圣事的奥迹。因为那个强盗并非为基督的名被钉，而是为他自身行为的报应；他也不是因信而受苦，而是在痛苦中相信了。因此，在那强盗身上显示了\BibleQuote{心里相信，可成义；口里宣认，可获救恩}的伟大能力，正如\OriginalTerm{宗徒}{apostle}所说的（\BibleRef{罗 10:10}），即使没有可见的圣洗圣事也是如此。但是，这种欠缺只是无形地得到补足，只有当施行圣洗不是被对宗教的轻视所阻止，而是被当时的必要性所阻止时才会发生。因为在\Person{科尔乃略}及其朋友的情形中，比那强盗更甚，似乎他们再用水受洗是多余的，因为在他们身上，\OriginalTerm{圣神}{Holy Spirit}的恩赐（按照圣经的见证，其他人只在领洗后才领受）已经通过那时代一切确切无疑的标记（他们说起各种语言）显明出来。然而，他们还是受了洗，而这一行动我们有宗徒的权威作为依据。我们所有人都应当远避这样的想法：若在领洗之前，一个人因虔诚心灵的运作而内在有所进展，达到了属神的领悟，就因而轻视那由施行者之手施予身体、但却是天主自己用来在精神上使人奉献于祂的\OriginalTerm{圣事}{sacrament}。我也不认为施洗的职务被交给\Person{若翰}、以致被称为若翰的洗礼，是出于其他原因，而只是因为主自己（设立了这洗礼）没有轻看接受祂仆人的洗礼（\BibleRef{玛 3:6}），以祝圣谦卑的道路，并通过这样的行动极为清楚地表明祂自己的圣洗（祂后来要用这圣洗施洗）应受到何等高度的重视。因为祂像一位卓越的永恒救恩的医生一样，预见到有些人会存有过分的骄傲——这些人对真理的理解和品德的端正已经取得如此进步，以至于在生活和知识上毫不犹豫地将自己置于许多受洗者之上——认为他们无需受洗，因为他们觉得自己已经达到了许多受洗者仍在努力提升的心灵状态。\par

% structure: p0055 source=h2 target=subsection reason=relative-heading-rank; base=h2

\subsection{第二十三章。}

31. 但是，\OriginalTerm{圣事}{sacrament}的圣化（那强盗没有领受，并非出于他意愿的欠缺，而是因为不可避免的省略）究竟有何确切的价值，以及它实质性的施用在人身上有何效果，是不容易说清的。尽管如此，如果它没有极大的价值，\OriginalTerm{主}{Lord}就不会接受仆人的\OriginalTerm{圣洗}{baptism}。但是，既然我们必须就其本身来看待它，而不涉及其旨在成就的领受者的\OriginalTerm{救恩}{salvation}问题，那么它清楚地表明：无论在恶人身上，还是在那些言语上弃绝世界而行动上不然的人身上，圣事本身是完整的，但他们若不改善生活，便不能获得救恩。然而，如同在那强盗身上，圣事的实质性施行是必然欠缺的，但他的救恩却是完整的，因为通过他的虔诚，救恩在精神上临在了；同样，当圣事本身临在时，救恩也是完整的，如果那强盗所拥有的东西是必然欠缺的话。这是普世\OriginalTerm{教会}{Church}的坚定传统，关于婴儿的圣洗尤其如此：这些婴儿当然还不能\BibleQuote{心里相信以成义，口里宣认以获得救恩}，像那强盗所能做的；相反，当他们接受奥迹时，甚至以哭号和呻吟反对那些神秘的言语，然而，没有一个基督徒会说他们受洗是徒劳的。\par

% structure: p0057 source=h2 target=subsection reason=relative-heading-rank; base=h2

\subsection{第二十四章。}

32. 若有人在这件事上寻求神圣权威，虽然整个教会所持守的，并非经大公会议制定，而是作为不变的习惯，理应被视为由权威所传授，但我们仍可从割损礼的类比，对婴孩领受圣洗圣事之价值作出正确的推断。割损礼是天主在早期子民中所接受的，而亚巴郎在领受割损礼之前已经成义，正如科尔乃略在受圣洗之前已蒙赐圣神的恩惠。然而，宗徒论及亚巴郎本人说：他领受了割损的标记，作为因信德而成义之印证，因为他心中早已相信，以致\BibleQuote{这就算为他的正义}。那么，为何天主命令他，所有男婴在第八天都要受割损礼（\BibleRef{创 17:9}），尽管他们尚不能用心中相信，使之算为正义，因为圣事本身具有极大的效用？这在梅瑟之子的事上，藉天使的传报得以显明：当他的母亲带着尚未受割损的婴孩时，藉着明显的眼前危险，要求他必须受割损（\BibleRef{出 4:24}），这事一完成，死亡的危险便消除了。因此，正如在亚巴郎身上，信德的成义在先，割损礼后来加添作为信德的印证；同样，在科尔乃略身上，属神的圣化在先，藉圣神的恩赐，而重生的圣事后来加添在圣洗之水中。又如依撒格出生后第八天受割损，这信德正义的印证首先赐下，后来当他成长并效法父亲的信德时，正义本身随之而来，而这印证在他婴孩时已预先赐下；同样，在受圣洗的婴孩身上，重生的圣事首先赐下，如果他们持守基督徒的虔诚，日后心灵的皈依也会随之而来，其神秘的标记已先行于外在的身体。又如盗贼的情况，全能者的慈爱弥补了圣洗圣事之所缺，因其缺失并非出于骄傲或轻蔑，而是因为缺乏机会；同样，对于受圣洗后去世的婴孩，我们必须相信全能者的同一恩宠弥补了缺失，使他们不是因意志的悖逆，而是因年龄不足，既不能心中相信以致成义，也不能口里宣认以致得救。因此，当别人代他们宣誓，使圣事的举行得以在他们身上完成时，这无疑为他们的奉献于天主有益，因为他们不能自己回答。但如果别人代一个能自己回答的人回答，那就没有同样的益处。与此规则相符，我们在福音中看到，每个人读到时都感到自然的那句话：\BibleQuote{他已是成年人，让他为自己说话吧。}（\BibleRef{若 9:21}）\par

% structure: p0059 source=h2 target=subsection reason=relative-heading-rank; base=h2

\subsection{第二十五章}

33. 所有这些考虑都证明：圣洗圣事是一回事，心灵的皈依是另一回事；但人的得救需两者共同完成。我们也不应认为，若两者中缺一，另一必定也缺失；因为圣事可能存在于婴孩身上，而无心灵的皈依；而心灵皈依在无圣事的情况下也是可能的，如盗贼的情形，天主在两种情况下都填补了非自愿的缺失。但如果这两者中有任何一样是故意缺少的，那么这人就要为缺失负责。圣洗可能存在而心灵的皈依却缺失；但就这种皈依而言，它确实可能在未领受圣洗的情形下找到，但绝不会在轻视圣洗的情形下找到。当人轻视天主的圣事时，绝不能说是心灵的转向天主。因此，我们正确地谴责、弃绝、憎恶并痛恨异端者心灵的悖逆；但这并不意味着他们没有福音的圣事，因为他们没有使圣事生效的东西。所以，当他们来到真信仰，并藉补赎寻求赦罪时，我们不是在奉承或欺骗他们，而是用天上的训练教导他们，使他们适于天国，纠正并改造他们心中的错误和悖逆，目的是绝不破坏他们身上健全的部分，也不因人的过犯，而宣布任何来自天主的恩赐是无效或有缺陷的。\par

% structure: p0061 source=h2 target=subsection reason=relative-heading-rank; base=h2

\subsection{第二十六章}

34. 在致犹巴雅奴斯的书信中，还有几点尚待注意；但这些点将同时引发对教会过去习惯以及若翰洗礼的讨论——对那些对此明显之事不够留意的人来说，这极易引起不小的疑惑，因为宗徒曾命令那些受过若翰洗礼的人再次受圣洗（\BibleRef{宗 19:3}）。这些问题不宜草率处理，最好留待另一本书，以免本卷篇幅过大。\par

\SourceRecord{Translated by J.R. King and revised by Chester D. Hartranft.；Nicene and Post-Nicene Fathers, First Series；Vol. 4.；Edited by Philip Schaff.；Buffalo, NY: Christian Literature Publishing Co.,；1887.；<http://www.newadvent.org/fathers/14084.htm>.}

% structure: p0001 source=h1 target=section reason=page-title

\section{\WorkTitle{论圣洗，驳多纳徒派（卷五）}}

\Emph{他考查了西彼廉写给犹巴伊安书函的最后部分，连同他写给奎恩图斯的信函、非洲主教会议写给努米底亚主教的信，以及西彼廉写给庞培的信函。}\par

% structure: p0003 source=h2 target=subsection reason=relative-heading-rank; base=h2

\subsection{第一章}

1. 我们有真福西彼廉的见证：目前天主教会仍保留着这一习惯，即从异端或分裂者方面来的人，若已按福音经文领受了圣洗，就不再重新受洗。因为他本人曾向自己提出这个问题，而这问题来自那些或寻求真理、或为真理争辩的弟兄之口。在他试图证明异端者应重新受洗的论证过程中——我们在前几卷已充分考察了这些论证，就当前目的而言——他说：\Quote{但有人会说：\InnerQuote{那些在以往从异端来到教会，未受圣洗就被接纳的人，以后会怎样呢？}}这个问题包含了多纳徒派整个论点的颠覆，我们与他们正是为此争辩。因为如果那些从异端来而被如此接纳的人并未真正领受圣洗，他们的罪仍留在身上，那么，当这些人被接纳进入共融时，无论是由西彼廉之前的前辈，还是由西彼廉本人，我们必须承认发生了两种情况之一：要么教会因与这样的人共融的污染而当场灭亡，要么任何留在合一中的人不会因他人昭彰的罪而受害。但他们不能说当时的教会因那些——如西彼廉所说——未受圣洗就被接纳的人所带来的污染而灭亡，否则他们就无法维护自己起源的有效性，因为执政官年表证明，从西彼廉殉道到焚烧圣书（他们借此制造分裂，散布诽谤的烟雾）之间相隔四十多年。因此，他们只能承认，基督的合一是不会被任何这样的共融——甚至与已知的罪犯共融——所污染的。承认这一点后，他们将无法找到任何理由来证明他们必须与普世教会分离，而我们知道普世教会同样由宗徒建立。因为其他教会不可能因任何种类的罪人掺杂而灭亡，而他们，即使留在与他们的合一中也不会灭亡，却因与弟兄分离、打破和平的纽带而自取灭亡。因为分裂的亵圣在他们身上最为明显，倘若他们没有充分的分离理由。显然他们没有充分的分离理由，因为即使已知罪犯的存在，也不会污染仍在合一中的好人。至于好人留在合一中不会受已知罪犯的污染，我们有西彼廉的教导，他说：\Quote{以往从异端来到教会的人，未受圣洗就被接纳了。}然而，如果他们亵圣的邪恶（仍在其身上，因未被圣洗涤除）不能污染并毁灭教会的圣洁，那么教会也不会因恶人的任何感染而灭亡。因此，如果他们承认西彼廉说的是真话，那么根据他的见证，他们就被定为分裂之罪；如果他们坚持他说的不是真话，那么他们就不要在洗礼问题上使用他的见证。\par

% structure: p0005 source=h2 target=subsection reason=relative-heading-rank; base=h2

\subsection{第二章}

2. 不过，现在我们既然开始与西彼廉这位和平之士辩论，就让我们继续下去。当他提出一个他知道是弟兄们所质疑的异议：\Quote{那些在以往从异端来到教会，未受圣洗就被接纳的人，以后会怎样呢？}他答道：\Quote{主能凭祂的仁慈宽恕，并使他们不致与祂教会的恩赐分离，这些人诚心诚意地被接纳进祂的教会，并在教会内安眠。}他很好地假设爱德能遮盖许多罪过。但如果他们确实领受了圣洗，而认为他们应重新受洗的人没有正确认识到这一点，这一错误被合一的爱德所遮盖，只要这爱德不包含不和谐和魔鬼的精神，而仅包含人性的软弱，直到如宗徒所说：\BibleQuote{若他们另有想法，主必指示他们。}\BibleRef{斐 3:15} 但那些因亵圣的分裂而撕裂合一的人有祸了，他们要么重洗（如果圣洗在我们和他们都存在），要么根本不施圣洗（如果圣洗只存在于天主教会）。因此，无论他们重洗还是未施圣洗，他们都不在和平的纽带中；因此，让他们对这两处创伤中的任何一处加以治疗吧。但如果我们将人未受圣洗就接纳入教会，我们就属于那些如西彼廉所假设的那样，因保持合一而可获得宽恕的人。但如果（我认为根据前几卷所说已清楚）基督徒的圣洗即使在异端的谬误中也能保持其完整性，那么即使当时有人重洗，但只要不离开合一的纽带，他们仍能因那种和平之爱而获得宽恕，正如西彼廉所见证的，那些甚至未受圣洗就被接纳的人也能获得恩赐而不致分离。此外，如果异端者和分裂者中没有基督的圣洗，那么其他人的罪又如何能伤害那些固守在合一中的人呢？既然人们即使未受圣洗来到教会，他们自己的罪也能被赦免！因为如果按照西彼廉，合一的纽带如此有效，那么那些不愿与合一分离的人怎么能被别人的罪伤害呢？即使未受圣洗、从异端希望来到教会的人，也能因此逃脱自己应得的灭亡？\par

% structure: p0007 source=h2 target=subsection reason=relative-heading-rank; base=h2

\subsection{第三章}

3. 然而，在\Person{西彼廉}附加的话中，他说：\Quote{人既曾一度犯错，并非必须永远犯错；明智而敬畏天主的人，更宜欣然且毫不犹豫地追随那清楚地展现在眼前的真理，而非顽固而固执地为异端者争辩，反对自己的弟兄和同司铎}，他道出了最完美的真理；抗拒明显真理的人，反对的是他自己而非他的邻人。但是，就我所能判断的，从我已引述的诸多论据来看，基督的圣洗即使由异端者施行或在他们当中领受，也不会因他们的邪恶而无效，这一点是非常清楚和确定的。然而，即使这尚未确定，至少任何考虑过上述论点的人，即使从敌对的立场出发，也不会断言问题已被相反地决定了。因此，我们并非与明显的真理作对，而是——依我之见——在为清楚无误的真理辩护，或者，至少如那些认为问题尚未解决的人可能持有的观点那样，我们是在寻求真理。因此，即使真理与我们想的有所不同，我们仍以同样的诚实之心接纳那些由异端者施洗的人，如同\Person{西彼廉}认为那些因紧守教会合一而能获得宽恕的人所受的接纳一样。但如果基督的圣洗，如上面许多论据所表明的，无论生活或信德上有任何缺陷——无论是那些看似在教会内却不属于独一鸽子肢体的人，还是那些与教会分离以致公开在外的人——都能保持其完整性，那么，早期那些寻求重行圣洗的人，因他们紧守合一的爱情，应得同样的宽恕，正如\Person{西彼廉}认为那些他相信未经圣洗便被接纳的人因同样种类的爱情所应得的一样。因此，那些毫无理由（因为\Person{西彼廉}自己表明，在教会的合一中，坏人不能伤害好人）就将自己与这合一中所显示的爱德割裂的人，丧失了所有宽恕的余地；他们即使没有给在天主教会受过洗的人重新施洗，也已因分裂的罪过招致毁灭；那么，他们该受何等痛苦的惩罚！他们要么试图将那\Person{西彼廉}声称自己没有的东西给予拥有它的天主教徒，要么——正如事实所表明的——指控天主教会没有连他们自己也拥有的东西。\par

% structure: p0009 source=h2 target=subsection reason=relative-heading-rank; base=h2

\subsection{第四章}

4. 但是，既然如先前所说，我们已开始与\Person{西彼廉}的书信进行辩论，我认为，即使他本人若在场，也不应指责我\Quote{顽固而固执地为异端者争辩，反对自己的弟兄和同司铎}，当他得知那促使我们相信即使异端者在其恶毒的谬误中顽固执拗，基督的圣洗在其本身仍是至圣至敬的有力理由时。鉴于他本人（其证言对我们极有分量）见证，过去通常接纳异端者而不予重洗，我愿任何被\Person{西彼廉}的论据说服而认为异端者应当重新受洗的人，也考虑以下事实：那些因对方论据有力而尚未被说服应如此行的人，处于与过去那些人同样的地位——他们诚实地接纳在异端中受洗的人，仅纠正其个人的错误，并因合一的联系而与这些人一同获得救恩。任何被教会过去的习俗、随后全体会议的权威、圣经中如此多有力的证据、\Person{西彼廉}本人的许多证言以及真理的清晰推理所引导，而理解基督的圣洗——以福音的言辞祝圣——不能被世上任何人的错误所玷污的人，应当明白，那些当时持不同看法但仍保存爱德的人，能因同样的合一联系而得救。在此，他也应当明白那些分散在世界各地的教会团体中的人，原本不可能被任何莠子或任何糠秕所玷污，只要他们自己愿意成为结实的麦粒；他们却毫无离婚理由地自同一合一的联系中分离出来；那么，无论这两种观点哪一种是真实的——是\Person{西彼廉}当时所持的，还是普世天主教会（\Person{西彼廉}并未离弃它）一致所持的——这两种情形下，他们既已公开置身于分裂的亵渎之中，绝不可能得救；他们所有的一切圣事恩典以及唯一合法新郎的无偿恩赐，在他们继续如此时，只能使他们灵魂混乱而非得救。\par

% structure: p0011 source=h2 target=subsection reason=relative-heading-rank; base=h2

\subsection{第五章}

5. 因此，即使异端者真正渴望纠正其错误并来到教会，仅仅因为他们相信除非在教会中领受否则就没有圣洗，我们也不应屈从他们要求重行圣洗的愿望；相反，应当教导他们：一方面，圣洗本身虽然完善，但若他们不愿受纠正，则绝无益于他们的执拗；另一方面，圣洗的完善也不会因他们拒绝纠正的执拗而受损；再者，当他们接受纠正时，圣洗在他们身上并未增添额外的完善；而是当他们离开不义时，那先前在他们内导致毁灭的，现在开始有益于救恩。因为学习了这一点，他们必将认识到在公教合一中得救的必要，并停止将本是基督的据为己有；他们将不会把真理的圣事（尽管存在于他们之中）与他们个人的错误相混淆。\par

6. 此外我们还可补充一点：人藉某种隐秘的上天启示，对任何第二次领受已在任何地方领受过的圣洗者，都本能地畏避；甚至异端者自己在辩论开始时，也困惑地擦着额头，他们几乎全体平信徒——即使那些在他们的团体中已年老、并对天主教会心怀顽固敌意的人——也承认，他们制度中的这一点令他们不悦；许多人为获得某种世俗利益或避免某种损失而想投奔他们，便暗中竭力争取，希望能作为特殊和个人的特权获准不被重领圣洗；有些人受他们其他虚幻谬误和诬蔑天主教会的虚假指控的影响而轻信，却因这一考量——不愿与他们为伍以免被迫重领圣洗——而被召回合一。多纳徒派惧怕这种已彻底占据人心的情绪，便同意承认在他们所谴责的马克西米安追随者中所施行的圣洗，从而截断自己的舌头并封住自己的嘴，而不是重新为他们所接纳的、曾与\Person{斐理西安}、\Person{普莱泰克斯塔图斯}及其他曾被他们谴责后又返回的人一同归来的\Place{姆斯提}、\Place{阿苏雷}及其他地区的众多民众重领圣洗。\par

% structure: p0014 source=h2 target=subsection reason=relative-heading-rank; base=h2

\subsection{第六章}

7. 因为当这种事偶尔发生在个别人身上，且时隔久远、地隔千里时，其严重性不为人同等地察觉；但假如所有曾随时间由上述马克西米安追随者施洗的人，或在死亡危险逼迫下、或在他们的逾越节庆典中领洗的人，突然被聚集起来，并被告知必须重领圣洗，因为他们在分裂的亵圣中已领受的圣洗是无效的，那么他们或许会说——那是顽固坚持谬误所迫使的——以任何虚假的伪装的一致性来掩藏他们硬心肠的严酷和冰冷，以逃避真理的温暖。但事实上，因为马克西米安派不能忍受这事，并且那些将要强制执行的人也无法忍受在如此众多的人身上必须做的事，尤其是那些人在普里米安派中重施圣洗时，正是当初在马克西米安派中为他们施洗的同一批人——因此他们的圣洗被接纳了，多纳徒派的骄傲被挫败。而他们肯定不会选择采取这种做法，除非他们认为，重领圣洗所引起的反感比放弃辩护所失去的声誉对他们的立场伤害更大。我说这话，并非认为若真理迫使那些从异端者来的人重领圣洗时，我们应受人情顾虑的约束；而是因为圣\Person{西彼廉}说：\Quote{如果异端者必须在天主教会重领圣洗，他们可能更被逼向皈依的必要性}，为此我愿记录下这种几乎深植于每人内心的强烈反感——我相信这是天主亲自感动的，使教会能借此反感的本能，对抗软弱者无法驱散的种种可能论点。\par

% structure: p0016 source=h2 target=subsection reason=relative-heading-rank; base=h2

\subsection{第七章}

8. 确实，当我注视\Person{西彼廉}的原话时，我受到提醒要说一些对解决此问题非常必要的话。他说：\Quote{因为他们若看到，我们用正式决议和投票已确定并规定，他们在异端中所领受的圣洗被视为合法且正当，他们就会以为他们也合法且正当地拥有教会及其一切恩赐。}他没有说\Quote{他们以为拥有}，而是\Quote{合法且正当地拥有教会的恩赐}。但我们说，我们不能承认他们合法且正当地拥有圣洗。我们不能否认他们拥有它，因为我们在福音的话语中认出主的圣事。因此他们有合法的圣洗，但他们不合法地拥有它。因为无论谁，既在天主教会的合一中，又活出相称的生活，他就既有合法的圣洗，又合法地拥有它；但无论谁，或在天主教会内部（像混在麦子中的糠秕），或在外部（像被风吹走的糠秕），他确有合法的圣洗，但不合法地拥有它。因为他如何用它，就如何拥有。但谁不合法地使用它——即违犯法律使用它——他就是不合法地拥有它。每个领洗后却仍过着放荡生活的人，无论教会内外，都是如此。\par

% structure: p0018 source=h2 target=subsection reason=relative-heading-rank; base=h2

\subsection{第八章}

9. 因此，正如宗徒论法律所说：\BibleQuote{法律原是好的，只要人用得合法}\BibleRef{弟前 1:8}，我们也可以恰当地论圣洗说：圣洗是好的，只要人用得合法。正如那些不合法地使用法律的人，并不能因此使法律本身不再是好的，也不能使之无效；同样，任何人不合法地使用圣洗，或因生活在异端中，或因过着最恶劣的生活，都不能使圣洗本身不是好的，或完全无效。因此，当他归向天主教合一，或归向与此伟大圣事相称的生活方式时，他开始拥有的不是另一种合法的圣洗，而是那同一个圣洗，只是合法地拥有。不可挽回之罪的赦免并不随着圣洗而必然发生，除非一个人不仅拥有合法的圣洗，而且合法地拥有它；然而，这并不导致：若一个人不合法地拥有它，以致他的罪要么未被赦免，要么即或赦免了却又回到他身上，因此圣洗圣事在受洗者身上就是坏的或无效的。因为正如犹达斯，主给了他一块饼，他却在自己的心里给魔鬼留了地方，不是因为他接受了坏的东西，而是因为他接受得不好（\BibleRef{若 13:27}），同样，每个人在接受主的圣事时，并不因为他自己是坏的而使那圣事成为坏的，也不因为他没有为得救而接受就什么也没有接受。因为那仍然是主的身体和主的血，即使是在那些宗徒所说的人身上：\BibleQuote{那不相称地吃和喝的人，是吃喝自己的定罪}\BibleRef{格前 11:29}。所以，异端者当在天主教会中寻求的不是他们所拥有的，而是他们所没有的——即命令的终极，没有它，虽然可以拥有许多神圣的事物，但无法获益。\BibleQuote{但命令的终极是爱德，出自纯洁的心、善良的良心和真诚的信德}\BibleRef{弟前 1:5}。因此，让他们赶快走向天主教会的合一与真理，不是为获得圣洗的圣事（如果他们已在异端中沐浴过），而是为健康地拥有它。\par

% structure: p0020 source=h2 target=subsection reason=relative-heading-rank; base=h2

\subsection{第九章}

10. 现在我们必须看看关于若翰的洗礼说了什么。因为\Quote{我们在\WorkTitle{宗徒大事录}中读到，那些已经受过若翰洗礼的人，又被保禄施洗了}，只是因为若翰的洗礼不是基督的洗礼，而是基督允许若翰施行的洗礼，因此特别称为若翰的洗礼；正如同一位若翰所说：\BibleQuote{人不能领受什么，除非由天上赐给他}\BibleRef{若 3:27}。并且，免得他可能显得是从天主父领受此恩，却不是从子领受，他接着论到基督自己说：\BibleQuote{从祂的满盈中，我们都领受了}\BibleRef{若 1:16}。但借着一项特殊安排的恩宠，若翰领受了这恩，这恩不会持久，只够为预备主的道路，因为他需要作前驱。正如我主即将以极度的谦卑踏上这条路，并带领那些谦卑跟随祂的人达到成全，就如祂洗了仆人的脚（\BibleRef{若 13:4}），同样祂也愿意接受仆人的洗礼（\BibleRef{玛 3:13}）。因为正如祂亲自俯身去洗那些祂所引导之人的脚，祂也把自己交托给若翰的恩赐，而这恩赐原是祂自己所赐予的，为叫众人明白，他们若藐视各人应由主领受的洗礼，将显出何等亵渎的傲慢，因为主亲自接受了祂自己所赐予仆人的，使那仆人能以他自己的名义施予；并且，当若翰——妇女所生者中没有比之更大的（\BibleRef{玛 11:11}）——给基督作了这样的见证，承认自己连解祂的鞋带也不配（\BibleRef{若 1:27}）时，基督既可以因接受他的洗礼而被视为人中最谦卑的，又可以因废除了若翰洗礼的地位而被相信为至高的天主，祂既是谦卑的导师，又是高举的施予者。\par

11. 因为我们在圣经中读到，没有一位先知，没有任何人，曾获赐以悔改的水施洗，使罪过得以赦免，如同赐给若翰的那样；以致藉此奇妙的恩宠，使人民的心归向他，好为他们预备道路，迎接那位他宣称比自己无限伟大者。但主耶稣基督藉着这样的圣洗来洁净祂的教会，使领受者无需再受其他洗礼；而若翰所施行的首次洗涤，则使领受者仍需领受主的圣洗——并非要重复第一次的洗礼，而是使那些已领受若翰洗礼的人，能进一步领受基督的圣洗，因为若翰正是为祂预备道路。因为如果基督的谦卑不值得我们注意，那么若翰的洗礼也就没有必要；反之，如果终结在于若翰，那么在他洗礼之后就不再需要基督的圣洗。但因为\Quote{基督是法律的终向，使一切信仰的人获得正义}（\BibleRef{罗 10:4}），所以若翰指示人应往何处去，并在到达祂时安息于祂内。因此，若翰既显示了主的崇高（他将自己远置祂之下），也显示了祂的谦卑（他如同最低微的人一样给祂施洗）。但如果若翰只为基督一人施洗，那么人们可能会认为他的洗礼（基督独享的洗礼）比基督徒所领受的基督圣洗更优越；反之，如果所有人都应先受若翰的洗礼，然后再受基督的圣洗，那么基督的圣洗就会显得不完整、不完美，不足以带来救恩。因此，主领受了若翰的洗礼，为要使人类骄傲的颈项屈服于祂那救恩的圣洗；但祂并非独自领受，以免祂的圣洗显得不如那只有祂才配领受的洗礼；祂也没有让若翰的洗礼持续更久，以免祂所施行的唯一圣洗似乎需要另一洗礼作为预备。\par

% structure: p0023 source=h2 target=subsection reason=relative-heading-rank; base=h2

\subsection{第十章}

12. 因此，我问：如果若翰的洗礼能赦免罪过，那么基督的圣洗还能给那些领受了若翰洗礼后、又蒙宗徒保禄以基督的圣洗施洗的人，增添什么呢？但如果若翰的洗礼不能赦免罪过，那么居普良时代那些他自己所说\Quote{以诡诈的欺诈霸占田产，以累积的利息增加收益}的人，难道比若翰更好吗？然而藉着他们的施洗，罪过的赦免却得以施行？或者是因为他们仍在教会的合一之内？那么若翰呢？他不也在那合一之内吗？他是新郎的朋友，预备主的道路，甚至亲自为主施洗。谁会疯狂到断言这一点呢？因此，虽然我相信若翰以悔改的水施洗，为赦免罪过，但那些受他洗礼的人所领受的是罪过赦免的期望，实际赦免发生在主的圣洗中——正如我们在希望中期待末日的复活，就像宗徒所说：\Quote{祂使我们一同复活，一同坐在天上基督耶稣内}（\BibleRef{弗 2:6}），又说：\Quote{我们得救是在于希望}（\BibleRef{罗 8:24}）；又如若翰自己说：\Quote{我以水洗你们，为叫你们悔改，以赦免你们的罪过}（\BibleRef{玛 3:11}），但一看见主就说：\Quote{请看天主的羔羊，除免世罪者}（\BibleRef{若 1:29}）——尽管如此，我不愿与任何人激烈争论，即使有人坚持若翰的洗礼也赦免了罪过，而保禄命令那些人再受圣洗，是藉基督的圣洗赋予他们某种更丰厚的圣化（\BibleRef{宗 19:3}）。\par

% structure: p0025 source=h2 target=subsection reason=relative-heading-rank; base=h2

\subsection{第十一章}

13. 因为我们必须关注与当前问题特别相关的一点（无论若翰的洗礼性质如何，既然他显然属于基督的合一），即：为何人在领受了圣若翰的洗礼之后，应当再受圣洗，而在领受了贪婪主教的洗礼之后，却不该再受圣洗？因为没有人否认，在主的田地里，若翰如同麦子，结出百倍的果实（如果那是最高收成的话）；也无人怀疑，贪婪（即偶像崇拜）在主庄稼中被视为糠秕。那么，为何一个人领受了来自麦子的洗礼后要再受圣洗，而领受了来自糠秕的洗礼后却不重洗？如果是因为保禄比若翰更好，所以他在若翰之后施洗，那么为何居普良不比他那些贪婪的同僚更好（他远胜于他们）而也在他们之后施洗？如果是因为那些同僚与他处于合一之中，所以他不给他们施洗，那么保禄也不该在若翰之后施洗，因为他们也处于同一合一之中。难道欺诈者和勒索者属于那一只鸽子的肢体，而那位藉鸽子形状的圣神显现而领受了主耶稣基督全权者，却不属于它吗？诚然，他最深切地属于它；但那些或因某次丑闻、或因末日扬场而必须与之分离的人，绝不属它。然而，人们在若翰之后重洗，在那些人之后却不重洗。那么原因何在？除非保禄命令他们领受的圣洗，与若翰所施的洗礼并不相同。因此，在教会同一合一中，即使由贪婪的职员施行，基督的圣洗也不可重复；但那些领受了若翰洗礼的人，即使是出自若翰本人之手，也应当后来再受基督的圣洗。\par

% structure: p0027 source=h2 target=subsection reason=relative-heading-rank; base=h2

\subsection{第十二章}

14. 因此，我也可以借用真福西彼廉的话，引导听者思考一件真正奇妙的事，假如我说：\Quote{那在先知中算为最大的若翰——他在母胎中就已充满天主的恩宠；他怀着厄里亚的灵和能力；他不是敌对者，而是主的先驱和宣报者；他不仅用言语预言了我们的主，还把他指给人看；他亲自为基督付洗，而基督是众人受洗的根源——}他竟不配如此施洗，以致凡受他洗的人，在他以后还需再受洗？难道有人会认为，一个人在教会中受洗后，若曾被贪吝者、欺诈者、勒索者、放高利贷者施洗，就不该再受洗吗？对于这个令人反感的问题，答案不就明摆着吗：你为什么认为这不合适？仿佛若翰受了羞辱，或贪吝者反得了尊荣？但若翰所说的\BibleQuote{那一位就是祂，祂要以圣神施洗}\BibleRef{若 1:33}，这圣洗本不该重复。因为无论藉谁的手施行，这圣洗都属于那被说成\Quote{那一位就是祂，祂要施洗}的主。但若翰自己的圣洗也未被重复，尽管宗徒保禄曾命令那些受若翰洗的人要因基督受洗。因为他们从新郎的朋友那里没有领受的，就理应从新郎本人那里领受，而那位朋友曾说：\BibleQuote{那一位就是祂，祂要以圣神施洗}。\par

% structure: p0029 source=h2 target=subsection reason=relative-heading-rank; base=h2

\subsection{第十三章}

15. 因为主耶稣若愿意，本可把他施洗的权柄赐给一位或几位他早已视为朋友的首要仆人，就是他对他们说\BibleQuote{我不再称你们为仆人，而是朋友}\BibleRef{若 15:15}的那些人；这样，正如亚郎因发芽的棍杖被显明为司祭\BibleRef{户 17:8}，同样，在教会中，当更多更大的奇迹施行时，那些圣德更卓越的仆人和他奥秘的分施者，或许会藉某种记号被显明为唯一有权施洗的人。但若是这样，尽管施洗的权柄是由主赐予他们，却必然被称为他们自己的圣洗，如同若翰的圣洗一样。因此，保禄感谢天主，因为他没有给那些人中的任何人施洗，这些人仿佛忘记了他们因谁的名受洗，竟要按不同人的名字分裂成派系\BibleRef{格前 1:12}。因为当圣洗在卑贱者手中与在宗徒手中同样有效时，它就被认出既不是这个人的圣洗，也不是那个人的圣洗，而是基督的圣洗；正如若翰藉着鸽子的显现，在主自己身上所领受的教训。因为他在别的方面说\BibleQuote{我也不认识祂}\BibleRef{若 1:32}，我无法清楚明白。因为如果他根本不认识祂，就不会在祂前来受洗时对祂说：\BibleQuote{我该受祢的洗}\BibleRef{玛 3:14}。那么，他说\BibleQuote{我看见圣神仿佛鸽子从天降下，停在祂身上。我不认识祂，但那派遣我来以水施洗的对我说：你看见圣神降下并停在谁身上，谁就是以圣神施洗的那位}\BibleRef{若 1:32}，这又是什么意思？鸽子显然是在祂受洗后降下的。但祂前来受洗时，若翰已说：\BibleQuote{我该受祢的洗}。所以他早已认识祂。那么他说\BibleQuote{我不认识祂}以及后面的话，既然这事发生在祂受洗之后，除非他在某些方面认识祂，在其他方面又不认识祂？他确实认识祂是天主之子、新郎，众人将从他的满盈中领受；但是，关于他从祂的满盈中领受了施洗的权柄，以致这圣洗被称为若翰的圣洗，他并不知道主是否也要如此把权柄赐给其他人，或者主要使自己的圣洗成为这样：无论由谁施行——无论是结百倍、六十倍或三十倍果实的人，无论是麦子还是糠秕——这圣洗都被知道唯独属于祂；而这一点，他是藉着像鸽子般降下并停在祂身上的圣神而学到的。\par

% structure: p0031 source=h2 target=subsection reason=relative-heading-rank; base=h2

\subsection{第十四章}

16. 因此，我们发现宗徒们使用这样的说法：\BibleQuote{我的夸耀}\BibleRef{格前 9:15}，虽然这夸耀无疑是在主内；\BibleQuote{我的职务}\BibleRef{罗 11:13}，\BibleQuote{我的知识}\BibleRef{弗 3:4}，\BibleQuote{我的福音}\BibleRef{弟后 2:8}，尽管这些都明明白白是由主赐予和恩赐的；但他们中没有一个人曾说过一次\Quote{我的圣洗}。因为他们的夸耀并不相等，他们也不是以同等的权柄服务，他们也不都拥有同等的知识，在宣讲福音时，一人比另一人更有能力，因此一人可以说在得救的道理上比另一人更有学问；但没有人能说比另一人更受洗或较少受洗，无论他是被一个更大或更不值得的施行者所洗。所以，当\BibleQuote{肉体的行为是明显的，即奸淫、污秽、淫荡、崇拜偶像、行邪术、仇恨、纷争、嫉妒、恼怒、结党、分裂、异端、嫉妒、醉酒、荒宴及类似的事}\BibleRef{迦 5:19}时，如果说\Quote{人在若翰之后受洗，却不在异端者之后受洗}很奇怪，那么为什么说\Quote{人在若翰之后受洗，却不在嫉妒者之后受洗}不也同样奇怪？因为西彼廉自己在论嫉妒与恶意的书信中作证，贪吝者属于魔鬼的党羽，西彼廉自己也从宗徒保禄的话中显明，正如我们上面已指出的，在宗徒时代，在基督之名的宣讲者中，已有人在基督的教会里怀着嫉妒？\par

% structure: p0033 source=h2 target=subsection reason=relative-heading-rank; base=h2

\subsection{第十五章}

17. 因此，若翰的圣洗不是基督的圣洗，我想这已经足够清楚地表明了；所以不能由此推论说，在异端者之后圣洗应当重复，因为曾在若翰之后重复过：因为若翰不是异端者，他能有一个圣洗，这圣洗虽由基督授予，却尚不是基督自己的圣洗，因为他拥有基督的爱；而一个异端者可以同时拥有基督的圣洗和魔鬼的乖戾，正如教会内的另一个人可以同时拥有基督的圣洗和魔鬼的嫉妒。\par

18. 但有人会坚持说，在异端者之后的圣洗更加必要，因为若翰不是异端者，而在他的圣洗之后却重复了圣洗。按此原则，一个人可以说，在醉汉之后我们更必须重行圣洗，因为若翰是清醒的，而在他之后圣洗却被重复了。我们对这样的人将无话可答，只能指出：那些受若翰洗礼的人领受了基督的圣洗，因为他们还没有得到它；但在那些拥有基督圣洗的人身上，他们方面的任何不义都不能使基督的圣洗在他们内失效。\par

19. 因此，说\Quote{异端者首先施洗，从而获得圣洗的权利}是不正确的；因为并非他用他自己的圣洗施洗，并且尽管他没有施洗的权利，他所给予的却是基督的，而领受它的人也是基督的。因为许多事物被不正当的给予，却并不因此被认为不存在或未曾给予。因为那些仅在言语上而非在行为上弃绝世俗的人，并未合法地领受圣洗，然而他确实领受了它。因为西彼廉记载，在他那个时代的教会中就有这样的人，而我们自己也经历并哀叹这一事实。\par

20. 但奇怪的是，在何种意义上可以说\Quote{圣洗和教会绝不可能彼此分离和脱离}。因为如果圣洗不可分离地存留于受洗者内，那么他如何能被分离于教会，而圣洗却不能呢？但清楚的是，圣洗确实不可分离地存留于受洗者内；因为无论受洗者跌入何种罪恶的深度，或陷入何种可怕的罪恶漩涡，甚至堕落到背信的毁灭中，他仍不会失去他的圣洗。因此，如果他藉悔改回头，圣洗不再重授，因为认为他不可能已失去它。但谁能怀疑一个受过洗的人可以与教会分离呢？因为所有异端都是由此而产生的，它们藉着使用基督教术语来欺骗人。\par

% structure: p0038 source=h2 target=subsection reason=relative-heading-rank; base=h2

\subsection{第十六章}

因此，既然显然圣洗在受洗者与教会分离时仍存留于他内，那么他内的圣洗必然与他一同分离。因此，并非所有持守圣洗的人都持守教会，正如并非所有持守教会的人都持守永生。或者，如果我们说只有那些遵守天主诫命的人才持守教会，那么我们立刻承认，有许多人持守圣洗，却不持守教会。\par

21. 因此，异端者并非\Quote{先占有了圣洗}，因为他已从教会领受了它。而且，尽管他分裂出去，圣洗也不能被他所失去——我们断言他不再持守教会，却允许他持守圣洗。也没有人\Quote{放弃他的长子权，并将它交给异端者}，因为他说他带走了他不能合法给予、但给予时却合于法律的东西；或者他不再合法地拥有它，尽管它在他手中是合于法律的。但长子权只在于圣善的品行和良好的生活，所有那些组成那位没有瑕疵或皱纹的新娘，或那只在众多乌鸦的邪恶中呻吟的鸽子的肢体都属于此——除非，厄撒乌因贪恋一碗红豆汤而失去了他的长子权，\BibleRef{创 25:29}而我们却认为它被欺骗者、强盗、高利贷者、嫉妒者、醉汉等类人持守着，对于他们的存在，西彼廉在他的书信中哀叹在他那个时代的教会里。因此，要么持守教会和持守神圣事物的长子权不是同一回事，要么，如果每一个持守教会的人也持守长子权，那么所有那些我们中间每个人都似乎承认在教会内能施洗的恶人并不持守教会，因为除非是完全不懂神圣事物的人，否则没有人会说他们在神圣事物中持守长子权。\par

% structure: p0041 source=h2 target=subsection reason=relative-heading-rank; base=h2

\subsection{第十七章}

22. 但是，在考虑并处理了所有这些要点之后，我们现在终于来到了西彼廉在书信末尾那段和平的言语，我读它千遍也不厌倦——其中所散发出的弟兄之爱的愉悦如此之大，所洋溢的爱德之甘饴如此之丰。他说：\Quote{这些事，最亲爱的弟兄，我们按我们微薄的能力简短地写给你们，不设定规范，也不预先判断任何人，以免每位主教不能按自己自由意志的运用做他认为正确的事。就我们而言，我们不与我们的同僚和主教弟兄们在异端问题上争论，我们与他们在主内保持和谐与和平；尤其是宗徒也说过：\BibleQuote{若有人似乎好争辩，我们没有这样的习惯，天主的各教会也没有。}\BibleRef{格前 11:16}我们耐心而温和地持守着精神的愛德、弟兄情谊的尊荣、信德的纽带、司祭职的和谐。为此，我们也凭着我们微薄的能力，在天主的允许和默感下，写了这篇论\WorkTitle{忍耐的美益}的论文，鉴于我们彼此的爱，我们把它送给你们。}\par

23. 在这些话语中有许多事值得考虑，基督的爱德的光辉在这人身上闪耀出来，他\Quote{喜爱主殿宇的美丽，及祂住所帐幕的居所}。首先，他没有隐瞒自己的感受；其次，他如此温柔和平地表达出来，因为他与那些持不同意见的人保持了教会的和平，因为他明白和平的纽带中蕴含着多么大的健康，他如此热爱和平，并清醒地维护它，看见并感觉到即使见解不同的人也可以在救恩的爱德中持有各自的意见。因为他不会说他能与恶人维持神圣的和谐或主的和平；因为善人可以对恶人保持和平，但不能与他们联合在恶人所没有的和平中。最后，他不规定任何人，也不预判任何人，以免每位主教不能按照自己自由意志认为正确的去行事，他为我们留下了——无论我们是什么人——一个与他和平讨论这些事的地方。因为他不仅在他的书信中临在，而且通过那在他身上如此非凡且永不熄灭的爱德而临在。因此，藉着他的祈祷的帮助，我渴望依附并与他结合，如果我的罪过的不配不阻止我，我将尝试通过他的书信来学习主是如何藉着他以极大的和平与安慰管理祂的教会；并且，通过他的讲论的感动，穿上谦卑的心肠，如果与普世教会一起，我持有比他更正确的道理，我不会把我的心灵置于他之上，即使在那一点上他虽然持有不同观点，却没有与普世教会分离。因为在那时，当那个问题因缺乏充分讨论而尚未决定时，尽管他的观点与许多同僚不同，他却保持了如此大的节制，以至于他不愿因任何分裂的污点而损害天主教会的神圣共融，一种更大的卓越力量在他身上显现出来，胜过如果他没有那种美德，而在每一点上都持有不仅正确而且与他们一致的观点。我也不会按照他的愿望行事，如果我假装将他的才华、流利的言辞和丰富的学识置于所有民族的圣教会之上，他通过他的精神合一肯定临在于其中，尤其是他现在处于如此圆满的真理之光中，能够完全确知他在这里以完全和平的精神所寻求的。因为从那个丰盛的富足中，他嘲笑我们这里看似雄辩的一切，仿佛那是婴孩的初次尝试；在那里他看见他在这里是按照何种虔诚的规则行事，即在教会中没有什么比合一更可贵。在那里，他以难以言喻的喜悦看见主以何等的预先和至慈的眷顾，为了医治我们的肿胀，\BibleQuote{拣选了世上愚拙的，叫有智慧的羞愧}\BibleRef{格前 1:27}，并在祂教会肢体的安排中，将一切置于如此有益的方式，使人不会说他们因自己的才华或学识而被选来帮助福音，而他们却不认识这些的来源，从而被致命的骄傲所膨胀。哦，西彼廉是多么喜乐！他以何等更完美的平静看见，这对人类健康是多么有益，即使在基督徒和虔诚的演说家的著作中也应发现应受谴责之处，而在渔夫的著作中却找不到任何这类东西！因此，我完全确信这圣洁灵魂的喜乐，绝不敢认为或说我的著作没有任何错误，也不在反对他的那个意见时——即他认为从异端者来的人应当以不同于往日（正如他亲自见证）或如今（如合理的习惯，由整个基督教世界的全体会议所确认）的方式被接纳——将我的观点，而是将神圣天主教会（他如此热爱并仍热爱的教会，他在其中以宽容结出如此丰硕的果实，他本身虽不是教会的全部，却留在教会的全部中；他从未离开她的根，虽然已经从根上结出果实，却被天上的园丁修剪，\BibleQuote{使他结出更多的果实}\BibleRef{若 15:2}）的观点置于他之上；为了教会的和平与安全，以免麦子连同莠子一起被拔除，他以真理的自由责备，并以爱德的恩宠忍受了那么多来自与他同在合一中的人们的恶行。\par

% structure: p0044 source=h2 target=subsection reason=relative-heading-rank; base=h2

\subsection{第十八章}

24. 为此，\Person{Cyprian}本人再度极其充分地提醒我们：许多死于过犯和罪恶中的人，虽不属于基督的身体，也不属于那只纯洁无邪的鸽子的肢体（因此，若唯独她施洗，那些人肯定不能施洗），但在外表上却似乎既在教会内受洗，又在教会内施洗。在他们中间，无论他们多么死寂，他们的圣洗却仍然活着，因为圣洗并非死物，死亡不能再统治它。因此，既然教会内有死人，而且他们并非隐藏的，否则\Person{Cyprian}就不会如此多地谈论他们，这些人要么根本不属于那只活的鸽子，要么至少尚未属于她；既然教会外也有死人，他们更明显地根本不属于她，或尚未属于她；既然确实\Quote{一个自己不活的人不能叫另一个人活}，——那么显然，那些在教会内由这样的人施洗的人，如果怀着真心的悔改领受圣事，便由那圣洗所归属的主所复活。但若他们只在言语上而非行动上弃绝世界，正如\Person{Cyprian}所宣称的某些在教会内的人那样，那么显然，除非他们悔改，他们自己不会复活，但他们即使不悔改，也拥有真实的圣洗。由此同样清楚的是，教会外的死人，虽然\Quote{他们自己不活，也不能叫别人活}，却拥有活的圣洗，这圣洗一旦他们悔改归向和平，便能有益于他们获得生命。\par

% structure: p0046 source=h2 target=subsection reason=relative-heading-rank; base=h2

\subsection{第十九章}

25. 因此，对于那些以基督的同一圣洗（即他们在教会外领受的圣洗）接纳来自异端的人，并说\Quote{他们遵循了古代惯例}——正如现今教会也如此接纳他们——而言，以下反对意见是徒劳的：\Quote{古人那里只有异端和分裂的初步开端，因此那些涉及其中的人是脱离教会者，并且最初是在教会内受洗的，所以他们回归并做补赎时不必重新受洗。}因为每当某个异端出现并脱离大公教会的共融时，它不仅可能在次日，甚至就在当日，其追随者就可能为一些投奔他们的人施洗。因此，如果这是古老的惯例，即他们应如此被接纳进入教会（即使是在辩论中持相反意见的人也无法否认），那么任何一个仔细思考此事的人都毫无疑问：那些在异端中受洗的人也是同样被接纳的。\par

26. 但我看不出以下情形有何道理可言：某人因在寻求基督内的救恩时陷入异端的错误并在他们的团体中受洗，就否认他是\Quote{迷途的羊}；而另一个人在教会内仅在言语上而非行动上弃绝世界，并以这种虚伪之心领受圣洗，却已被视为在教会内的羊。或者，如果这样的人除非真心归向天主也不能成为羊，那么当他成为羊时并未受洗——如果他已受洗但尚未成为羊——同样，从异端前来要成为羊的人，若已受过同一圣洗，则那时也不应受洗，尽管他尚未成为羊。因此，既然教会内所有恶人——贪婪者、嫉妒者、酗酒者、以及违背基督纪律生活的人——都可正当地被称作说谎者、在黑暗中、死人、以及敌基督者，那么他们难道因此就不施洗，理由是\Quote{真理与谎言、光明与黑暗、死亡与不朽、基督与敌基督之间不可能有任何共通之处}吗？\par

27. 因此，他并非依据\Quote{纯粹的惯例}，而是依据\Quote{真理本身的理性}断言：天主的圣事不能因任何人的错误而变为错误，因为它被宣布甚至存在于那些错误的人身上。确实，宗徒若望最清楚地说：\Quote{凡恼恨自己弟兄的，就是在黑暗中，直到如今。}\BibleRef{若一 2:9}又说：\Quote{凡恼恨自己弟兄的，就是杀人犯。}\BibleRef{若一 3:15}那么，为什么他们给教会内那些\Person{Cyprian}本人宣称怀着恶毒嫉妒的人施洗呢？\par

% structure: p0050 source=h2 target=subsection reason=relative-heading-rank; base=h2

\subsection{第二十章}

一个杀人犯怎能洁净和祝圣水呢？黑暗怎能祝福油呢？但如果天主临在于祂的圣事中，无论圣事由谁施行，祂都确认祂的话语，那么天主的圣事处处有效，而恶人虽无益于圣事，却处处乖谬。\par

28. 但这是一种什么样的论据呢？即\Quote{一个异端者必须被视为没有圣洗，因为他没有教会。}并且必须承认，\Quote{当他受洗时，他被问及关于教会的事。}仿佛在教会内以言语而非行为弃绝世界的人，在洗礼中不是同样被问及关于教会的事。因此，正如他虚假的答复并不阻止他所领受的成为圣洗，同样，另一个人关于圣教会所作的虚假答复也不阻止他所领受的成为圣洗；并且正如前者，如果他后来真诚地履行了他所虚假承诺的事，并不会领受第二次洗礼，而只是修正的生活，同样，就后者而言，如果他后来来到教会——他曾就所提的问题给出虚假回答，认为自己拥有教会，而实际并未拥有——那么他所不曾拥有的教会本身被赐予他，但他所领受的并不重复。但我无法解释这是为什么：天主可以藉着一个杀人者口中所出的话语\Quote{祝圣油}，却\Quote{不能在异端者所立的祭台上祝圣它，}除非是因为那位在教会内不被人心虚假皈依所阻碍的天主，却被外面虚假竖立的一些木头所阻碍，而不屑于临在于祂的圣事中，尽管人的任何虚假都不能阻碍祂。因此，如果福音所说的\BibleQuote{天主不听罪人} \BibleRef{若 9:31} 这句话的意思是罪人不能举行圣事，那么祂如何听一个杀人者祈祷，无论在洗礼的水上、油上、圣体上，还是在他按手的人的头上呢？这一切事仍然被执行，并且有效，甚至由杀人者执行，也就是由那些恨自己弟兄的人执行，甚至在教会内部也是如此。既然\Quote{没有人能给予他自己所没有的，}一个杀人者如何赐予圣神呢？然而这样的人甚至在教会内也施洗。因此，是天主赐予圣神，即使当这样的人施洗时。\par

% structure: p0053 source=h2 target=subsection reason=relative-heading-rank; base=h2

\subsection{第二十一章}

29. 但至于他所说的，\Quote{来到教会的人应当受洗并更新，好使他在内藉着圣洁者而成圣，}如果他在内部也遇到不圣洁的人，那他该怎么办？或者，杀人者可能是圣洁的吗？如果他在教会内受洗的理由是\Quote{他也应脱去这一件事，即他作为一个寻求归向天主的人，因错误的欺骗而落在一个不洁的人身上，}那么他后来在哪里脱去这件事，才能使他有机会在教会内部寻求一位天主的人时，因错误的欺骗而落在一个杀人者身上呢？如果\Quote{一个人身上不能同时存在无效和有效的事物，}那么为什么在杀人者身上，圣事可以是圣洁的，而他的心却不圣洁呢？如果\Quote{凡不能赐予圣神者就不能施洗，}那么为什么杀人者在教会内施洗呢？或者，杀人者如何拥有圣神？因为凡拥有圣神的人都充满光明，但\Quote{恨自己弟兄的人仍在黑暗中。} \BibleRef{若一 2:9} 如果因为\Quote{只有一个圣洗，一个圣神，}所以没有那一个圣神的人就不能拥有那一个圣洗，那么为什么教会内的无辜者和杀人者拥有同一个圣洗，却没有同一个圣神呢？同样，异端者和公教徒可以拥有同一个圣洗，却没有同一个教会，正如在天主教会内，无辜者和杀人者可以拥有同一个圣洗，尽管他们没有同一个圣神；因为如同只有一个圣洗，同样也只有一个圣神和一个教会。因此结果是，在每个人身上，我们必须承认他已经拥有的，并且给予他所没有的。如果\Quote{凡被天主称为祂的仇敌和敌人的人所做的事，在天主面前不能得确认和批准，}那么为什么由杀人者所施行的洗礼得确认呢？我们难道不把杀人者称为主的仇敌和敌人吗？但\Quote{恨自己弟兄的人就是杀人者。}那么，那些恨保禄——耶稣基督的仆人——并因此恨耶稣本人的人，他们是如何施洗的呢？因为主自己曾对扫罗说：\BibleQuote{你为什么迫害我？} \BibleRef{宗 9:4} 那时他正在迫害祂的仆人，并且最后主将亲自说：\BibleQuote{你们没有给我这些最小兄弟中的一个做，就是没有给我做。} \BibleRef{玛 25:45} 因此，凡从我们中间出去的，不属于我们；但并非所有与我们同在的人都属于我们；如同人打谷时，一切从禾场上飞走的都显明不是麦子，但一切留在那里的也未必是麦子。因此若望也说：\BibleQuote{他们从我们中间出去了，但他们不属于我们；如果他们属于我们，就必会留在我们当中。} \BibleRef{若一 2:19} 因此，天主甚至通过恶人的手赐予恩宠的圣事，但恩宠本身只由祂自己或通过祂的圣人赐予。因此，祂或由自己，或通过那鸽子的肢体——祂对他们说：\BibleQuote{你们赦免谁的罪，就给谁赦免；你们保留谁的罪，就给谁保留。} \BibleRef{若 20:23} ——赐予罪过的赦免。但既然无人能怀疑圣洗——即罪过得赦免的圣事——甚至被杀人者所拥有，而这些人仍在黑暗中，因为他们心中没有除去对弟兄的仇恨，因此，要么如果他们受洗时没有伴随内心的改变，就没有罪过的赦免赐予他们；要么如果罪过得赦免，它们立刻又回到他们身上。我们得知圣洗本身是圣洁的，因为它是属于天主的；无论它是由这类人所施或由他们所领受，都不能因他们的任何乖谬而受玷污，无论是在教会内还是教会外。\par

% structure: p0055 source=h2 target=subsection reason=relative-heading-rank; base=h2

\subsection{第二十二章}

30. 因此，我们同意\Person{西彼廉}的看法：\Quote{异端人不能赦罪}；但我们坚持，他们能施行圣洗——实际上在他们身上，无论他们施予或领受，这圣洗只对他们的毁灭有益，因为他们滥用了天主如此伟大的恩赐；正如心怀恶意和嫉妒的人（\Person{西彼廉}自己也承认他们是在教会内）不能赦罪，而我们都承认他们能施行圣洗。因为如果论到那些得罪我们的人，曾说过：\BibleQuote{如果你们不宽恕别人的过犯，你们的父也必不宽恕你们的过犯，}\BibleRef{玛 6:15}，那么，那些恨那爱他们的弟兄，并在这仇恨中受洗的人，他们的罪更不可能被赦免；然而当他们归正时，圣洗并不重新赋予他们，而是在他们真正的皈依中，赐予他们当时不应得的赦免。所以，甚至\Person{西彼廉}写给\Person{昆图斯}的信，以及他与同僚\Person{利贝拉里斯}、\Person{卡尔多尼乌斯}、\Person{尤尼乌斯}等人共同写给\Person{撒图尔尼努斯}、\Person{玛克西穆斯}等人的信，在适当考虑之后，都绝不适合被用来反对整个天主教会的共识，而他们曾以身为教会成员为荣，既没有将自己与教会切断，也没有允许那些持相反意见的人被切断，直到最后，按主的旨意，多年后在一次全体主教会议上，显明了更完美的道路，这并不是建立任何新奇的事物，而是确认了古老的传承。\par

% structure: p0057 source=h2 target=subsection reason=relative-heading-rank; base=h2

\subsection{第二十三章。}

31. \Person{西彼廉}也就此事写信给\Person{庞培}，并在那封信中清楚表明，\Person{斯德望}（我们得知他当时是罗马教会的主教）不仅不赞同他在这几点上的看法，反而写了并教导相反的观点。但\Person{斯德望}当然没有\Quote{与异端人共融}，仅仅是因为他不敢攻击基督的圣洗，他知道这圣洗在他们内心的错谬中仍然保持完整。因为如果那些对天主持有错误看法的人都没有圣洗，那么在我看来，已经充分证明，即使在教会内也可能发生这种情况。诚然，\Quote{宗徒们没有对此事给予训令}；但是与\Person{西彼廉}意见相反的习俗，可以推测源自宗徒传统，正如有许多为整个教会所遵守的事，因此相当公平地被认为是由宗徒们所命令的，虽然它们并未出现在他们的著作中。\par

32. 但有人会辩称，经上论到异端人说：\BibleQuote{他们是自定罪的。}\BibleRef{铎 3:11} 那么怎么样？那些被说\BibleQuote{因为你论断人，正是定自己的罪？}\BibleRef{罗 2:1} 的人，不也是自定罪的么？但宗徒对他们说：\BibleQuote{你教导人不可偷盗，自己却偷盗么？}\BibleRef{罗 2:21} 等等。而这些人正是如此，他们是主教，与\Person{西彼廉}本人一同在天主教会合一中，却以诡诈的欺骗掠夺产业，同时一直向百姓宣讲宗徒的话，他说：\BibleQuote{勒索的不能承受天主的国。}\BibleRef{格前 6:10}\par

33. 因此，我只需简要地回顾一下基于同样原则的其他意见，这些意见见于前述写给\Person{庞培}的信中。凭什么圣经的权威表明：\Quote{如果一个已经从异端团体领受过基督圣洗的人，不再重新受洗，这是违反天主的诫命}？但很清楚的是，许多假基督徒，虽然他们没有与圣人们同处于爱德的联系中（没有这联系，他们可能拥有的任何圣洁事物都对他们无益），却与圣人们共同拥有圣洗，正如已经充分证明的那样。他说：\Quote{教会、圣神和圣洗彼此不可分离}，因此他希望\Quote{那些与教会和圣神分离的人，也应被理解为与圣洗分离}。但如果是这样，那么当一个人在天主教会领受圣洗后，只要他自己留在教会内，圣洗就留在他内——但事实并非如此。因为当他回归时，圣洗并不重新赋予他，正是因为当他离开时并没有失去它。然而，正如不肖儿子不像爱子那样拥有圣神，却仍有圣洗；同样，异端人不拥有教会像天主教徒那样，却仍有圣洗。因为\BibleQuote{训诫的圣神必逃避欺诈，}\BibleRef{智 1:5} 而圣洗却不会逃避它。所以，正如圣洗可以继续存在于圣神离开的人身上，圣洗也可以继续存在于教会不在的地方。但是，如果\Quote{覆手}不\Quote{给从异端来的人应用}，他就会被视为完全无咎；但为了爱德的合一（这是圣神最大的恩赐，没有它，一个人可能拥有的任何其他圣物都对他的得救无益），当异端人被带到认识真理时，就给他们覆手。\par

% structure: p0061 source=h2 target=subsection reason=relative-heading-rank; base=h2

\subsection{第二十四章。}

34. 我记得我已充分讨论了\Quote{天主的宫殿}这句话，以及\BibleQuote{你们凡受洗归于基督的，就是穿上了基督}\BibleRef{迦 3:27} 应如何理解。因为贪婪的人不是天主的宫殿，正如经上记载：\BibleQuote{天主的宫殿与偶像怎能相合？}\BibleRef{格后 6:16} 而西彼连引证保禄的话，指出贪婪就是崇拜偶像。但人穿上基督，有时仅限于领受圣事，有时更进一步至于获得生活的圣德。前者为善人与恶人所共有；后者则为善人与虔敬者所独有。因此，如果\Quote{洗礼不能没有圣神}，那么异端者也有圣神——但是为毁灭，而非为救恩，正如撒乌耳的情况\BibleRef{撒上 19:23}。因为藉着圣神，因基督之名驱魔，就连教会之外的人也能做到，这曾引起门徒向他们的主提出疑问\BibleRef{谷 9:38}。正如贪婪的人有圣神，但却不是天主的宫殿。因为\BibleQuote{天主的宫殿与偶像怎能相合？}因此，如果贪婪的人没有天主的圣神，却拥有洗礼，那么洗礼可能没有天主的圣神而存在。\par

35. 因此，如果异端\Quote{不能藉基督为天主生子，因为它不是基督的新娘}，那么教会内那群作恶的人同样也不能，因为他们也不是基督的新娘；因为基督的新娘被描述为没有瑕疵或皱纹。因此，要么不是所有受洗的人都是天主的子女，要么连那非新娘的也能为天主生子。但正如有人问：\Quote{在异端中领受基督洗礼的人是否属灵重生？}同样可以问：在天主教会中领受基督洗礼却没有以真诚的心归向天主的人，是否属灵重生？对这样的人，我们不能说他没有领受洗礼。\par

% structure: p0064 source=h2 target=subsection reason=relative-heading-rank; base=h2

\subsection{第二十五章}

36. 我不愿再处理西彼连对斯德望激动地倾泻而出的言论，况且这也完全没必要。因为那些论点正是已经充分讨论过的；最好略过那些涉及有害分裂危险的内容。斯德望认为，我们甚至应当远离那些企图破坏在接纳异端者方面的古老习俗的人；而西彼连则因问题本身的困难，且怀有基督徒爱德的圣洁心肠，认为我们应当与那些与我们意见分歧的人保持合一。因此，尽管他在愤怒中并非没有激动，但那确实是兄弟般的激动，然而基督的平安在他们心中得胜，以至在那样一场争论中，他们之间没有产生任何分裂的恶果。但并未发现\Quote{由此生出更多的异端和分裂}，因为其中属于基督的部分被认可，属于他们自己的部分被谴责；相反，那些坚持重洗礼律法的人反而分裂成更小的派别。\par

% structure: p0066 source=h2 target=subsection reason=relative-heading-rank; base=h2

\subsection{第二十六章}

37. 至于他所说\Quote{主教应当\InnerQuote{可教}}，并补充说：\Quote{可教的人是指温良、柔和、善于学习的人；因为主教不仅应教导，也应学习，因为每日藉学习更好的事物而成长进步的人，才是更好的教师。}——在这些话中，这位怀有虔诚爱德的圣人无疑指出，我们应当毫不犹豫地按这样的意义来读他的书信，即如果教会后来藉更深入、更长期的讨论确认了所发现的内容，我们不应感到为难；因为正如博学的西彼连可以教导许多事物，同样，可教的西彼连仍有某些事物需要学习。但他给我们的劝告：\Quote{我们应当回到泉源，即宗徒传统，并由此将真理的河道引至我们的时代}——这是极好的，应毫不犹豫地遵循。因此，正如他自己所记载的，宗徒们传下来的就是\Quote{同一的天主、同一的基督、同一的望德、同一的信德、同一的教会、同一的洗礼}。因此，既然我们发现，即使在宗徒时代，也有人没有同一的望德，却拥有同一的洗礼，那么真理本身就由泉源传至我们，使我们清楚看到：有可能虽然只有一个教会，正如只有一个望德和一个洗礼，但那些没有同一教会的人却可能有同一洗礼；正如在早期时代，那些没有同一望德的人却可能有同一洗礼。因为那些说：\BibleQuote{我们吃喝吧，反正明天我们就要死了}\BibleRef{格前 15:32}、断言没有死人复活的人，又如何与圣人和义人有同一的望德呢？然而他们正是同一位宗徒对其说：\BibleQuote{保禄为你们钉死在十字架上吗？或者你们受洗归于保禄名下吗？}\BibleRef{格前 1:13} 的同一批人。因为他极其明确地写信给他们说：\BibleQuote{你们中怎有人说没有死人复活呢？}\BibleRef{格前 15:12}\par

% structure: p0068 source=h2 target=subsection reason=relative-heading-rank; base=h2

\subsection{第二十七章}

38. 既然教会如此在 \WorkTitle{雅歌} 中被描述：\BibleQuote{我的妹妹，我的新娘，是关闭的花园，是封锁的泉源，是封印的活水井；你的栽种是石榴果园，结有佳美的果子。}\BibleRef{歌 4:12} 我不敢将此理解为别的，只理解为圣洁和正义的人——而不是贪婪者、欺诈者、抢劫者、高利贷者、醉酒者、嫉妒者；关于这些人，我们既从 \Person{Cyprian} 的书信中充分得知（如同我常指出的），自己也训导说，他们与义人同有圣洗，但肯定与义人无分于基督的爱德。因为我希望有人告诉我，\Person{Cyprian} 曾证明他们口头上而非行动上弃绝世界，却仍在教会内，他们是如何 \InnerQuote{潜入那关闭的花园和封印的泉源} 的？因为若他们自身都在那里，且自身是基督的新娘，那么她怎能如所描述的 \BibleQuote{没有瑕疵或皱纹，}\BibleRef{弗 5:27} 而洁白的鸽子竟被如此一部分肢体玷污呢？这些难道就是她如同百合花在其间的荆棘，正如同一 \WorkTitle{雅歌} 中所说的吗？\BibleRef{歌 2:2} 因此，百合花延伸多远，\InnerQuote{关闭的花园和封印的泉源} 也就延伸多远，即遍及所有内心受割损的犹太人\BibleRef{罗 2:29}（因为 \BibleQuote{君王的女儿内在全是光荣}），在他们身上有在世界奠基之前预定的固定圣人数目。但那众多的荆棘，无论是暗中还是公开分离的，从外部无数地压近它。经上说：\BibleQuote{我若宣告他们并谈论他们，他们比可数的更多。}因此，义人的数目，就是 \BibleQuote{按祂旨意蒙召的人，}\BibleRef{罗 8:28} 关于他们经上说：\BibleQuote{主认识属于祂的人，}\BibleRef{弟后 2:19} 这数目本身就是 \InnerQuote{关闭的花园、封印的泉源、活水的井、结满美果的石榴园。}这数目中，有些人随从圣神生活，走上爱德的卓越道路；当他们 \BibleQuote{以温和的精神纠正一个被过犯所胜的人，要留心自己，免得自己也受诱惑。}\BibleRef{迦 6:1} 而当他们自己也被胜过时，爱德的情感只是稍微受挫，并未熄灭；再次兴起并重新点燃，恢复原先的进程。因为他们知道说：\BibleQuote{我的灵魂因困苦而融化：求祢照祢的话语坚固我。}但当 \BibleQuote{他们在任何事上存有不同心意，天主也会将此启示给他们，}\BibleRef{斐 3:15} 只要他们坚守在爱德的烈焰中，不破坏和平的纽带。但有些尚属血肉、充满肉身欲望的人，致力于追求进步；为能适宜于天上的食粮，他们以神圣奥迹的乳汁被喂养，他们在敬畏天主中避免任何即便在世人眼中也显然腐坏的事物，并最警醒地努力，使自己越来越不贪恋世俗和暂时的事物。他们最恒常地遵守被勤奋寻求出的信仰准则；若在任何事上偏离，他们便在公教权威下接受迅速纠正，尽管，如 \Person{Cyprian} 所说，由于肉身的欲望，他们被幻想的各种冲突所摇动。还有一些人起初生活邪恶，甚至陷于异端或外邦人的迷信，但即使那时 \BibleQuote{主认识属于祂的人。}因为，在天主那不可言喻的先知中，许多看似在外的人其实在内，许多看似在内的人其实在外。因此，所有那些，我可以说，内在且隐藏在内的人，组成了那个 \InnerQuote{关闭的花园}，\InnerQuote{封印的泉源、活水的井、结满美果的石榴园。}这些人的神圣赐予，部分是他们独有的，如现世中永不衰败的爱德和来世的永生；部分是与邪恶和乖戾之人共有的，即所有构成神圣奥迹的其他事物。\par

% structure: p0070 source=h2 target=subsection reason=relative-heading-rank; base=h2

\subsection{第二十八章。}

39. 因此，我们现在得以更轻易、更简单地思考\Person{诺厄}所建造并掌舵的那只方舟。因为\Person{伯多禄}说，在\Person{诺厄}的方舟内，\BibleQuote{藉著水获救的并不多，只有八个人。这水所预表的圣洗，如今也拯救你们，并不是除掉肉体的污秽，而是向天主恳求一颗纯洁的良心。}\BibleRef{伯前 3:20}因此，如果那些只是在言语上而非行动上弃绝世界的人，在人看来是在大公教会的合一中受洗的，那么，在他们内没有良心的应答，他们又如何属于这方舟的奥迹呢？或者，那些滥用圣洗，虽然似乎是在内，却坚持过一种邪恶放荡的生活直至末日的人，他们如何能藉水得救呢？或者，\Person{西彼廉}本人记载，有些人过去曾带着他们在异端中所受的圣洗被单纯地接纳进入教会，他们又怎能不藉水得救呢？因为方舟的同一合一拯救了他们，在方舟内，除了藉水，没有人得到拯救。因为\Person{西彼廉}本人说：\Quote{主能因祂的仁慈赐予宽恕，并不将那些因单纯被接纳进入教会而在其怀抱中安眠的人，与祂教会的恩赐割裂开来。}若不藉水，如何在方舟内？若不在方舟内，如何在教会内？但若在教会内，必然在方舟内；若在方舟内，必然藉水。因此，有可能那些在外受洗的人，因天主的预知，被认为实际上是在内受洗的，因为在内，水开始为他们有益于救恩；也不能说他们除了藉水以外，还有其他方式在方舟内得救。同样，有些似乎在内受洗的人，因同样的天主的预知，被认为更真实地是在外受洗的，因为他们滥用圣洗，被水所灭，而这样的事从未发生在任何不在方舟外的人身上。诚然，显然的是，当我们论及与教会相关的\InnerQuote{在内}和\InnerQuote{在外}时，我们必须考虑心灵的处境，而非身体的处境，因为所有在心里在内的人，都在方舟的合一中藉着同一的水得救，而所有在心里在外的人，无论他们身体是否也在外，都作为合一的敌人而灭亡。因此，正如拯救方舟内的人并毁灭方舟外的人的不是别的，乃是同一的水，同样，拯救好公教徒并毁灭坏公教徒或异端者的，不是不同的圣洗，而是同一圣洗。但是，至圣\Person{西彼廉}对天主教会持何看法，以及异端者如何被他的权威彻底击溃；尽管我已说了很多，我仍决定，若蒙天主允许，在我首先就他的会议说出我认为应当说的话后，再以稍为更充分和更明确的方式单独阐述；这一切，在天主的旨意下，我将在下一卷书尝试。\par

\SourceRecord{Translated by J.R. King and revised by Chester D. Hartranft.；Nicene and Post-Nicene Fathers, First Series；Vol. 4.；Edited by Philip Schaff.；Buffalo, NY: Christian Literature Publishing Co.,；1887.；<http://www.newadvent.org/fathers/14085.htm>.}

% structure: p0001 source=h1 target=section reason=page-title

\section{论圣洗，驳多纳徒派（第六卷）}

\Emph{此处考虑的是在迦太基大公会议，由西彼廉主持并在其权威下召开，以决断关于异端者的圣洗的问题。}\par

% structure: p0003 source=h2 target=subsection reason=relative-heading-rank; base=h2

\subsection{第一章}

1. 或许已经足够，那些理由已经如此反复地陈述、考量并以各种方式讨论过，再加上从圣经得来的证明，以及\Person{西彼廉}本人如此多段落的共同见证，即使是心迟钝的人也应当理解（我相信他们确实理解）：\OriginalTerm{基督的圣洗}{baptism of Christ}不会因人的任何邪僻——无论施洗或受洗者——而归于无效。当我们发现，在那时，当所讨论的问题以一种违反古代惯例的方式解决之后，经过未损害救恩的爱德与合一的讨论，在一些卓越的人物（他们是基督的主教）看来，其中\Person{真福西彼廉}尤为突出，认为基督的圣洗不可能存在于异端者或分裂者中。这仅仅是因为他们没有区分圣事本身与圣事的效果或应用；并且因为圣事的效果和应用在异端者中并未显现在赦免其罪过和纠正其心志上，所以圣事本身也被认为在他们那里是欠缺的。但是如果我们转目注视教会内大量的糠秕，因为这些邪僻之人、即使在合一中过着放纵生活的人，似乎也没有能力赦免或保留罪过，因为不是对恶人说的，而是对善人说的：\BibleQuote{你们赦免谁的罪，就赦免他们；你们保留谁的罪，就保留。}\BibleRef{若 20:23}然而，这些人既有、也施行、也领受圣洗圣事，这对于散布全世界的\OriginalTerm{天主教会}{Catholic Church}牧者们来说是足够清楚的，通过他们，最初的惯例后来由全体大公会议的权威所确认。所以，即使迷途在外、从外来的假掠夺者那里接受了主印记的羊，若它寻求基督徒合一的救恩，它就从错误中得到净化，从俘囚中获得释放，伤口得医治，而那主的印记在它身上被承认，而非被拒绝；因为这印记本身常被狼印在狼身上，它们看起来确实在羊栈内，但却因它们行为的果实（它们坚持到底）而被证明不属于那在众多中合成一体的羊；因为按天主的预知，有多少羊游荡在外，就有多少狼阴险地潜伏在内，在这群中主认识那些属于祂的，它们只听从牧者的声音，即使牧者通过像法利塞人那样的人的声音来召唤，关于这些人曾说过：\BibleQuote{凡他们吩咐你们的，你们要遵守并实行。}\BibleRef{玛 23:3}\par

2. 因为属神的人，持守\BibleQuote{诫命的终向，即来自纯洁的心、良好的良心和真诚信德的爱德}，\BibleRef{弟前 1:5}，从一个尚且\BibleQuote{可朽坏的重压着灵魂的身体}中，\BibleRef{智 9:15}，有些事就看得不那么清楚，并且容易在有些事情上持不同的见解，这些事天主若他持守同一个爱德，会在祂适当的时候向他启示\BibleRef{斐 3:15}。同样，在一个属血肉且邪僻的人身上，也可能发现某些良善而有益的东西，其来源不在于人本身，而在于别的源头。因为正如在结果的枝条上会有需要修剪的部分，使它结更多的果实；同样，也常有葡萄挂在贫瘠、干枯或被捆绑的藤上。因此，正如喜爱结果枝条上需要修剪的部分是愚昧的，而明智的人不拒绝任何地方悬挂的甜葡萄；同样，如果有人因重行圣洗而自绝于合一，仅仅因为\Person{西彼廉}认为应当给从异端者来的人再次施洗，那么这样的人就偏离了那位伟人值得称赞的部分，而效法了需要纠正的部分，甚至达不到他所效法的东西本身。因为\Person{西彼廉}，出于对天主的热情，深切憎恶所有自绝于合一的人，认为他们因此也与圣洗本身分离；而这些人，认为自己自绝于基督的合一最多是轻微的过错，甚至坚持认为基督的圣洗不在合一之内，而是随他们一同出去了。因此，他们与\Person{西彼廉}的丰硕果实相差甚远，甚至比不上他身上需要修剪的部分。\par

% structure: p0006 source=h2 target=subsection reason=relative-heading-rank; base=h2

\subsection{第二章}

3. 再者，若有人没有爱德，行走在最邪恶生活的放纵道路上，似乎在内而实则在外，并且即使对异端者也不寻求重领圣洗，这丝毫无助于他的不结果实，因为他并非因自己的果实而结果实，却背负着他人的果实。但可能有人爱德之根茂盛，且在\Person{西彼廉}所持异议的事上持守极正，然而\Person{西彼廉}可能比他结出更多果实，而他比\Person{西彼廉}需要更多洁净。因此，我们不仅不将坏的公教徒与蒙福的\Person{西彼廉}相比，就连好的公教徒，我们也轻率地不将其与那位我们虔诚慈母教会列为少数出类拔萃、卓越恩宠的圣人相提并论，尽管这些公教徒可能在异端者中也承认基督的圣洗，而\Person{西彼廉}却持不同看法；如此，借着\Person{西彼廉}的实例（他虽在一件事上不太明白，却仍在教会合一中最稳固），可更清楚地向异端者显示：撕毁和平的锁链是何等亵渎的罪行。因为盲眼的法利塞人虽有时命令做正确的事，也不可与宗徒\Person{伯多禄}相比，尽管\Person{伯多禄}有时命令做不正确的事。不仅他们的干枯不能与他的青翠相比，甚至别人的果实也不可视为与他的丰饶相等。因为如今没有人强迫外邦人犹太化，然而如今在教会中，无论一个人在善行上进步多大，也不能与\Person{伯多禄}的宗徒职相比。因此，我怀着应有的敬意，并尽我所能，向和平的主教和光荣的殉道者\Person{西彼廉}致以相称的荣耀，但我仍敢说，他对分裂者和异端者洗礼的看法，与后来并非凭我个人而是凭整个教会的决定（由全体大公会议的权威所确认和加强）所显明的真理相反；正如我向首位宗徒和殉道者中最高贵的\Person{伯多禄}致以应有的敬意，但我仍敢说，他强迫外邦人犹太化做得不对；这也是我所说的，并非出于我自己的教导，而是根据整个教会所保存并持守的\Person{保禄}宗徒健全的教义。\BibleRef{迦 2:14}\par

4. 因此，在讨论\Person{西彼廉}的意见时，虽然我自己远不如\Person{西彼廉}的功德，但我断言：好人和坏人都能拥有、能赋予、能领受圣洗圣事——好人为了他们的健康和益处，坏人为了他们的毁灭和丧亡——而圣事本身在两者中具有同等的完美；并且，对于圣事在所有人中的同等完美而言，拥有圣事的坏人越坏并无影响，正如拥有圣事的好人越好也无影响。同样，赋予圣事的人越坏并无影响，正如赋予圣事的人越好也无影响；领受圣事的人越坏也无影响，正如领受圣事的人越好也无影响。因为圣事凭借其本身的卓越，在不平等地正义的人中，以及在不平等地不义的人中，都是同样神圣的。\par

% structure: p0009 source=h2 target=subsection reason=relative-heading-rank; base=h2

\subsection{第三章}

5. 但我认为，我们从圣经的正典以及\Person{西彼廉}本人的书信中已充分表明：坏人（绝未转变成较好的心志）能够拥有、赋予并领受圣洗，而这些人显然不属于天主的圣教会，尽管他们似乎在内，因为他们贪婪、强盗、放高利贷、嫉妒、怀恶意等；而教会是一只鸽子\BibleRef{歌 6:8}，端庄贞洁，无玷无皱的新娘，是关闭的花园、封闭的泉源、结佳美果实的石榴园\BibleRef{歌 4:12}，以及所有归于她的类似特性；这一切只能理解为存在于良善、圣洁、正义的人中——即不仅遵循天主恩赐的运作（这些恩赐好人和坏人相同），而且遵循内在的爱德之链，这爱德显于那些拥有圣神的人，主对他们说：\Quote{你们赦免谁的罪，就给谁赦免；你们保留谁的罪，就给谁保留。}\BibleRef{若 20:23}\par

% structure: p0011 source=h2 target=subsection reason=relative-heading-rank; base=h2

\subsection{第四章}

6. 因此，显然此处没有显示任何充分的理由，说明为何拥有圣洗的坏人不能也赋予圣洗；正如他拥有圣洗是为了丧亡，他也能赋予圣洗是为了丧亡——这不是因为所赋予之事的性质，也不是因为赋予者的性质，而是因为领受者的性质。因为当一个坏人将圣洗赋予一个好人（即处于合一之链中、真心悔改的人）时，赋予者的邪恶并不使所赋予的善圣事与领受圣事的好教会成员分离。当他真心悔改归向天主时，他的罪由那些因他真心悔改而与之联合的人所赦免。因为那同一位圣神赦免它们，这圣神赐给所有在爱中彼此相系的圣人，无论他们在肉身中是否彼此认识。同样，当一个人的罪被保留时，它们确实被那些他因生活不同和败坏的心背离而与之分离的人所保留，无论他们在肉身中是否认识他。\par

% structure: p0013 source=h2 target=subsection reason=relative-heading-rank; base=h2

\subsection{第五章}

7. 因此，一切恶人在精神上已与善人分离；但如果他们因明显的分歧也在身体上分离，他们就变得更坏了。但是，如前所述，对圣洗圣事的神圣性而言，拥有它的人多么恶劣，或施予它的人多么恶劣，都无关紧要：然而，那分离的人可以施予它，正如那分离的人可以拥有它；但正如他拥有它而趋于灭亡，他施予它也同样趋于灭亡。但是，那接受施予的人，如果他自己不是在分离中领受，则可以领受它而有益于灵魂的健康；正如许多人发生过的那样：怀着大公的精神，心底并未与平安的合一疏离，在死亡迫近的压力下，仓促地转向某个异端者，从他那里领受了基督的圣洗，却没有分享他的邪僻；因此，无论他们死去或重获生命，他们绝不愿与那些他们心底从未归附的人保持共融。但是，如果领受者自己是在分离中领受圣洗，那么他领受它便更趋于灭亡，与他自己未能妥善领受的善的大小成正比；而这在分离中对他越趋于灭亡，就如它在合一中越有益于一个人的得救一样。因此，如果他革除自己的邪僻，离开自己的分离，回归大公的平安，他的罪便因平安的纽带和那同一圣洗而得赦免——过去他的罪曾因分离的亵渎而被存留，因为圣洗在义人与不义之人中永远是神圣的，它不因任何人的义而增加，也不因任何人的不义而减少。\par

8. 既然如此，许多与他同为主教的人都同意\Person{西彼廉}的那个意见，并各自发表了支持同一立场的看法，这对如此明确的真理有什么影响呢？只是他的爱德对基督合一的热忱变得越发显著。因为如果他只是唯一持有那意见的人，没有人同意他，那他留在合一中可能会被认为是出于对分裂之罪的退缩，因为他在自己的错误中找不到同伴；但当这么多人同意他时，他通过与其他与他想法不同的人保持合一，表明他保存了普世大公性最神圣的纽带，不是出于对孤立的恐惧，而是出于对和平的爱。因此，现在考虑那公会议中其他主教的各个意见，似乎确实是多余的；但是，既然那些心思迟钝的人认为，如果对任何论述中的任何段落，本来可以当场给出的回答却没有在那里给出，而是在别处给出，那么他们就会觉得根本没有得到回答，那么与其让他们因理解甚少而留有抱怨说辩论没有得到公平进行的余地，不如让他们通过大量阅读而变得锐利。\par

% structure: p0016 source=h2 target=subsection reason=relative-heading-rank; base=h2

\subsection{第六章}

9. 首先，让我们记录\Person{西彼廉}本人提出的供审议的案件，他以此开始公会议的程序，并表明一种和平的精神，充满基督徒爱德的丰硕成果。他说：\Quote{最亲爱的同僚们，你们已经读过了我们的同为主教者\Person{犹巴安}写信给我，就异端者非法而亵渎的圣洗一事询问我卑微的能力，以及我所写给他的回信，向他表达了我一再、反复、常常表达的意见，即异端者来到教会，应当受洗，并以教会的圣洗而得圣化。\Person{犹巴安}的另外一封信也被念给你们听了，在这封信中，他根据其真诚而虔诚的热忱，在答复我们的信函时，不仅表示同意，而且感激地承认获得了教导。剩下的，就是我们应当各自就同一主题发表意见，不判断任何人，也不剥夺任何人的共融权利，如果他持有不同意见。因为我们中没有人将自己立为主教的主教，或以暴君的恐怖强迫同僚服从的必要，因为每位主教，在自由运用其自主权和权力时，都有自由判断的权利，他既不能受他人判断，也不能判断他人。但我们都在等待我主耶稣基督的审判，只有祂有权在治理祂的教会中提升我们，并审判我们的行为。}\par

% structure: p0018 source=h2 target=subsection reason=relative-heading-rank; base=h2

\subsection{第七章}

10. 我想，在前几卷中，我已尽己所能，为天主教会的同心合意与商议辩护——在这些单位中，他们如同虔敬的肢体得以存留——回应了\Person{西彼廉}写给\Person{犹巴安}的信，以及他写给\Person{昆图斯}的信，还有他偕同某些同僚写给另一些同僚的信，以及他写给\Person{庞培}的信。因此，现在似乎也适合考虑其他各位分别持何种见解，且以他所不会剥夺我们的自由来考量，正如他所说：\Quote{不判断任何人，也不因见解不同而将任何人从共融的权利中移除。}而他这样说，并非旨在将同僚们隐藏的心思从秘密角落中挖掘出来，而是因为他真正热爱和平与合一，这从同类其他段落中极易看出，例如他写给\Person{犹巴安}个人的信中说：\Quote{最亲爱的弟兄，这些是我们根据自己微薄的能力，简要答复你的，并不藉我们的著作或判断阻止任何主教按自己认为正确的行事，拥有自由运用自己判断的权利。}并且，为免有人因在这自由判断中持不同见解而被逐出弟兄的团体，他接着说道：\Quote{就我们而言，我们并不为异端者反对我们的同僚和主教同仁，我们与他们共同持守虔敬的合一和主的平安。}稍后他又说：\Quote{精神的爱德、对弟兄之情的尊重、信德的纽带、司铎职的和谐，都藉忍耐和温良由我们持守。}同样，在他写给\Person{玛克努斯}的信中，当被问及洒水或浸水圣洗的有效性是否有区别时，他说：\Quote{在此事上，我过于谦虚和谨慎，不愿以我的判断阻止任何人按他认为正确的去想，按他认为正确的去做。}这些言论清楚表明，这些主题在他们中间被处理时，尚未被承认为毫无疑义的定论，而是作为尚未启示的事，以极大的谨慎在探究。因此，我们在此维持关于所有圣洗同一性的立场，这必须被承认为普世教会的普遍惯例，且由大公会议的决定所确认；并且，我们从\Person{西彼廉}的话语中也获得更大信心，他允许我即使持与他不同的见解也不丧失共融的权利，因为合一被赋予更大的重要性和赞美，如真福\Person{西彼廉}及其同僚（他与他们召开那次会议）与持不同见解者所维持的那样，从而以主耶稣基督之名，动摇并推翻异端者和分裂者的恶意毁谤——主藉其宗徒说：\BibleQuote{凡事谦虚、温柔、忍耐，用爱心互相宽容，以和平彼此联络，竭力保守圣神所赐合而为一的心。}（\BibleRef{弗 4:2}）又藉同一宗徒的口说：\BibleQuote{我们若在什么事上存不同的心意，天主也必启示你们这事。}（\BibleRef{斐 3:15}）——我们，我说，提出真福主教们的见解以供考虑和讨论，但不破坏与他们之间的合一与平安纽带，在持守这纽带时，我们尽可能藉主自己效法他们。\par

% structure: p0020 source=h2 target=subsection reason=relative-heading-rank; base=h2

\subsection{第八章}

11. \Place{比尔塔}的\Person{塞西略}说：\Quote{我知道在唯一的教会中只有一种圣洗，在教会以外没有。唯一的一种将存在于那里，那里有真正的望德和确定的信德。因为经上记载：\BibleQuote{一个信德，一个望德，一个圣洗。}（\BibleRef{弗 4:4}）不在异端者中，那里没有望德，且有虚假的信德；那里一切事都藉谎言而行；那里一个附魔的人驱魔；圣事的问题由那口出癌症之言的人提出；无信者赐予信德；有罪者赦免罪过，且假基督以基督之名施圣洗；天主所诅咒者祝福；死人应许生命；不和者赐予平安；亵渎者呼求天主；不敬者执掌司铎职；亵圣者设立祭台。凡此种种还有此一恶上加恶，即魔鬼的仆人竟敢举行圣体圣事。若非如此，让那些支持他们的人证明这一切关于异端者都是假的。请看教会被迫同意的这类事，她被迫没有圣洗或罪过的赦免而与之共融。弟兄们，我们应当逃避并避免这些，使自己与如此大罪分离，并持守唯独赐予教会的唯一圣洗。}\par

12. 对此我回答：凡在教会内宣称认识天主，但行为上否认祂的人，诸如贪婪、嫉妒者，以及因仇恨弟兄而被宣告为杀人者——不是凭我的见证，而是基于圣宗徒\Person{若望}的见证，\BibleRef{若一 3:15}——所有这些人都没有望德，因为他们良心有亏；没有信德，因为他们不履行对天主的誓愿；是撒谎者，因为他们作虚假的宣认；是附魔者，因为他们让魔鬼及其使者进入心中；他们的话语腐蚀人心，因为恶交败坏了善行；他们是不信者，因为他们嘲笑天主对这类人的警告；是可咒骂者，因为他们生活邪恶；是敌基督者，因为他们生活与基督对立；是天主所诅咒的，因为圣经到处咒诅这样的人；是死人，因为他们没有义的生命；是不和善者，因为他们与天主的诫命相悖；是亵渎者，因为他们堕落的行为辱没了基督徒的名号；是亵圣者，因为他们被灵性上排除在天主的内殿之外；是渎圣者，因为他们邪恶的生活玷污了自身内天主的殿宇；是魔鬼的奴仆，因为他们事奉欺诈和贪婪，即偶像崇拜。这一类人有些，甚至很多，在教会内也是如此，这由\Person{保禄}宗徒和\Person{西彼廉}主教见证。那么，他们为何施圣洗？为何有些\Quote{在言语而非行为上弃绝世俗}的人，没有脱离这样的生活就受了圣洗，而回头后却不再受圣洗？至于他满怀义愤所说的话：\Quote{请看教会被迫同意何等事情，竟被强制在没有圣洗或罪赦的情况下与它共融}，若不是有其他主教在其他地方强迫人做这些事，他绝不会说出这样的话。由此也显明，那时那些没有离开远古习惯的人持有更正确的观点，这习惯后来由大公会议的共识所确认。但他又说：\Quote{弟兄们，我们应当逃避并避免这事，远离如此大罪。}这是什么意思？因为若他意指自己不做也不赞同此事，那是另一回事；但若他意指谴责并与持相反意见者分离，则他背离了西彼廉先前的话：\Quote{不判断任何人，也不剥夺任何人与我们共融的权利，即使他与我们意见不同。}\par

% structure: p0023 source=h2 target=subsection reason=relative-heading-rank; base=h2

\subsection{第九章.}

\Place{米吉珥帕}的长老\Person{费利克斯}说：\Quote{我认为每个从异端来的人都应受圣洗。因为若有人以为在那里已受了圣洗，那是徒劳的，因为在教会之外没有圣洗，只有教会内那唯一的真圣洗；因为只有一位主、一个信德、一个教会，其中驻有唯一的圣洗、圣洁及其他一切。凡在外面举行的事，都没有能力赋予救恩。}\par

对于\Place{米吉珥帕}的\Person{费利克斯}所说的话，我们回答如下：如果唯一的真圣洗只存在于教会内，那么它必定不会存在于那些脱离合一的人身上。但它确实存在于他们身上，因为他们回归时并没有领受，只因他们离开时并未失去它。至于他所说的\Quote{凡在外面举行的事，都没有能力赋予救恩}，我同意他，并认为这十分正确；因为圣洗不在那里是一回事，它没有能力赋予救恩是另一回事。因为当人们进入天主教会和平时，他们在加入前就已拥有但未能获益的事物，便立刻开始有益于他们。\par

% structure: p0026 source=h2 target=subsection reason=relative-heading-rank; base=h2

\subsection{第十章.}

对于\Place{阿德鲁梅图姆}的\Person{波利卡普}的声明：\Quote{那些声称异端者的圣洗有效的人，使我们的圣洗无效}，我们回答：如果异端者的圣洗是由异端者施行的，那么教会内由贪婪者和杀人者施行的也是他们的圣洗。但如果这不是他们的圣洗，那么另一种也不是异端者的圣洗；所以无论由谁施行，它都是基督的圣洗。\par

% structure: p0028 source=h2 target=subsection reason=relative-heading-rank; base=h2

\subsection{第十一章.}

\Place{塔穆加迪斯}的\Person{诺瓦图斯}说：\Quote{虽然全部圣经都为赋予救恩的圣洗作证，但我们仍应表达我们的信念：异端者和分裂者带着似乎已受圣洗的样貌来到教会，他们应在永不枯竭的泉源中受圣洗；因此，根据圣经的见证和那些最神圣的人、我们的同僚的谕令，凡归向教会的分裂者和异端者都应受圣洗；此外，所有似乎已领受圣职的人，都应作为普通平信徒接纳。}\par

\Place{塔穆加迪斯}的\Person{诺瓦图斯}陈述了他所做的，但他没有提出任何证据来表明他应该如此行。因为他提到了圣经的见证和同僚的谕令，但他没有从中引用任何我们可以考量的内容。\par

% structure: p0031 source=h2 target=subsection reason=relative-heading-rank; base=h2

\subsection{第十二章.}

18. 图布纳的\Person{奈美西亚努斯}说：\Quote{在圣经中处处宣布，异端者和分裂者所施行的圣洗不是真的，因为他们的首领本身就是假基督和假先知，正如主藉所罗门的口所说的：\InnerQuote{凡信赖谎言的，就是饲风；他追随飞鸟。因为他离开了自己葡萄园的道路，迷失了自己田地的路径。他行走在无路干旱之地，注定干渴之地；他手中收集无果的杂草。}又说：\InnerQuote{你要远离陌生的水，不要喝陌生的泉源之水，好使你活得长久，并增加年岁。}而在福音中，我们的主耶稣基督亲自开口说：\BibleQuote{人除非由水和圣神而生，不能进天主的国。}\BibleRef{若 3:5}这是从起初\BibleQuote{在水面上运行}的圣神。\BibleRef{创 1:2}因为圣神不能无水而行，水也不能无圣神而行。因此，有些人对自己作了错误的解释，声称他们藉覆手领受圣神，并得以进入教会，然而显然他们应当在天主教会内藉这两件圣事重生。因为到那时他们才能成为天主的子女，正如宗徒所说：\BibleQuote{尽力以和平的联系，保持圣神的合一：只有一个身体和一个圣神，正如你们蒙召，同有一个希望一样；只有一个主，一个信德，一个洗礼；只有一个天主。}\BibleRef{弗 4:3}凡此种种，天主教会都加以断言。他又在福音中说：\BibleQuote{那由肉生的属于肉，那由圣神生的属于神；因为圣神就是天主，并且由天主而生。}因此，异端者和分裂者所做的一切都是属肉的，正如宗徒所说：\BibleQuote{本性私欲的作为是显而易见的：即淫乱、不洁、放荡、崇拜偶像、施行邪法、仇恨、竞争、嫉妒、忿怒、争吵、不睦、异端，以及与这些相类似的事。我以前劝戒过你们，如今再劝戒你们：做这种事的人，决不能承受天主的国。}\BibleRef{迦 5:19}宗徒与一切恶人一同谴责那些制造分裂的人，即分裂者和异端者。因此，除非他们领受那唯一的、只在天主教会内才有的救恩圣洗，他们不能得救，反要在主的审判中与属肉的人一同受罚。}\par

19. \Place{图布纳}的\Person{内梅西阿努斯}引用了许多圣经章节来证明他的论点；但实际上，他所说的许多话恰好支持了我们所承担阐述和维护的天主教会的观点。诚然，除非我们必须认为，那信赖暂时事物之望的人并不\Quote{信赖虚假的事}（像一切贪婪之人和强盗，以及那些\Quote{在言语上弃绝世界，但在行为上却没有}的人），关于这些人，\Person{西彼廉}仍作证说，他们不仅在教会内施洗，甚至也在教会内受洗。因为他们自己也\Quote{追随飞鸟}，因为他们得不到他们所渴望的。但不仅是异端者，而且每一个过邪恶生活的人都\Quote{遗弃了自己葡萄园的道路，偏离了自己田地的路径。他行走在无路和干旱之处，一片注定干渴的土地；他手中收集无果的杂草；}因为一切正义都是多产的，一切不义都是不毛的。再者，那些\Quote{从奇异的水泉喝奇异之水}的人，不仅在异端者中间找到，也在所有不按天主的教导生活、而按魔鬼的教导生活的人中间找到。因为如果他在谈论圣洗，他就不会说：\Quote{不要喝奇异的水泉}，而是说，不要在奇异的水泉中洗浴。再者，我完全看不出他从我们主的话\BibleQuote{人除非由水和圣神而生，不能进入天主的国}\BibleRef{若 3:5}那里得到什么帮助来证明他的论点。因为说凡要进入天国的人必须先由水和圣神重生（因为除非人由水和圣神而生，他不能进入天国——这是主的话，是真实的）是一回事；说凡由水和圣神而生的人都必进入天国，则是另一回事，那无疑是假的。因为\Person{西满术士}也由水和圣神而生，\BibleRef{宗 8:13}然而他并没有进入天国；异端者也可能如此。或者，如果只有那些以真正的悔改而改变的人才是由圣神所生，那么所有\Quote{在言语上弃绝世界而在行为上却没有}的人，肯定不是由圣神所生，而只是由水所生，然而按照\Person{西彼廉}的见证，他们仍在教会内。因为我们不得不承认两件事之一：要么那些诡诈地弃绝世界的人是由圣神所生，尽管这导致他们的毁灭而非得救，因此异端者也可以这样生；要么，如果所记载的\BibleQuote{圣神的训诲会逃避诡诈}\BibleRef{智 1:5}足以证明，那些诡诈地弃绝世界的人不是由圣神所生，那么人可以受水洗而不由圣神所生，而\Person{内梅西阿努斯}说圣神不能离水工作、水也不能离圣神工作，就是徒劳的。事实上，我们已经多次表明，人如何可能共享同一圣洗却并不共享同一教会，正如在教会自身这个身体内，那些因正义而成圣的人和那些因贪婪而被染污的人，可能没有同一位圣神，却共享同一圣洗。因为经上说\Quote{一个身体}，即教会，正如它说\Quote{一个圣神}和\Quote{一个圣洗}一样。他所引用的其他论据反而有利于我们的立场。因为他从福音中引证了一个证据，就是这句话：\BibleQuote{由肉生的属于肉，由圣神生的属于神；因为圣神是天主，并由天主而生}\BibleRef{若 3:6}；并且他提出论据说，因此任何异端者或分裂者所做的一切都是属血肉的，正如宗徒所说：\BibleQuote{血肉的工作是明显的，即是：奸淫、不洁}等等；他列举了宗徒在那里列出的清单，其中他算上了异端，因为\BibleQuote{行这样事的人，不得继承天主的国}\BibleRef{迦 5:19}。然后他继续说：\Quote{因此宗徒与所有恶人一同谴责那些引起分裂的人，即分裂者和异端者。}在这方面他做得很好，因为当他列举血肉的工作（其中也包括异端）时，他发现并宣布宗徒同样谴责它们全体。因此，让他质问圣\Person{西彼廉}本人，从他那里了解在教会内有多少人按照血肉的邪恶工作生活（宗徒将这些工作与异端一同谴责），而这些人既施洗又受洗。那么，为什么唯独说异端者不能拥有圣洗，而他们的同受定罪者却拥有呢？\par

% structure: p0034 source=h2 target=subsection reason=relative-heading-rank; base=h2

\subsection{第十三章}

20. \Place{兰贝塞}的\Person{雅努阿里乌斯}说：\Quote{遵循圣经的权威，我宣布所有异端者都应受洗，然后被接纳进入圣教会。}\par

21. 我们回答他说：遵循圣经的权威，全世界的普世大公会议曾规定，基督的圣洗即使是在异端者中发现，也不应否认。但如果他从圣经中提出了任何证据，我们本该表明它们要么不反对我们，甚至支持我们，正如我们随后对下一位所做的那样。\par

% structure: p0037 source=h2 target=subsection reason=relative-heading-rank; base=h2

\subsection{第十四章}

22. \Person{Lucius of Castra Galbæ}说：\Quote{主曾在祂的福音中说：\BibleQuote{你们是地上的盐；盐若失了味，那从它而来的盐就毫无用处，只好被丢在外面，被人践踏。}而且再次，在祂复活后，当派遣祂的宗徒时，祂命令他们说：\BibleQuote{天上地下的一切权柄都交给了我，所以你们要去使万民成为门徒，因父、及子、及圣神之名给他们付洗，}\BibleRef{玛 28:18}——既然很明显，异端者，即基督的敌人，没有完整地宣认这圣事，而且分裂者不能用属神的智慧来推理，因为他们自己，在失了味后从唯一的教会中退出，变得与教会相反，那么，就应当施行经上所写的：\InnerQuote{敌对法律者的房屋必须被洁净。}因此，那些被与基督为敌的人付洗而被玷污的人，应当先被洁净，然后才领洗。}\par

23. \Person{Lucius of Castra Galbæ}从福音中引用了主的这句话作为证据：\Quote{\BibleQuote{你们是地上的盐；盐若失了味，那从它而来的盐就毫无用处，只好被丢在外面，被人践踏；}}就好像我们坚持认为被丢在外面的人无论对自己还是对别人的救恩都没有任何益处。但是那些虽然看起来在里面，却属于这一类的人，不仅在精神上在外面，最终在身体上也会被分开。因为所有这样的人都是无用的。但这并不意味着在他们里面的圣洗圣事是虚无的。因为即使在被丢在外面的人中，如果他们恢复理智并回来，那离开他们的救恩就会回来；但洗礼不会回来，因为它从未离开。当主说：\Quote{\BibleQuote{所以你们要去使万民成为门徒，因父、及子、及圣神之名给他们付洗，}}祂只允许好人付洗，因为祂没有对坏人说：\Quote{\BibleQuote{你们赦免谁的罪，谁就得赦免；你们保留谁的罪，谁就被保留。}}\BibleRef{若 20:23}那么，那些不能赦免罪恶的恶人如何在教会内付洗呢？还有，那些心未改变、罪恶仍在身上的恶人，正如若望所说：\Quote{\BibleQuote{凡恨自己弟兄的，就是在黑暗中，直到现在，}}\BibleRef{若一 2:9}他们又怎能付洗呢？但如果这些人的罪在他们是借由爱与善人正义之人紧密联合时被赦免，而罪恶在教会内是通过善人正义之人被赦免的，尽管他们是由恶人付洗的，那么那些从外面来并通过和平的内在纽带与基督身体的同一架构联合的人的罪也同样被赦免。然而，基督的洗礼在两者中都应被承认，并且在任何人中都不应被视为无效，无论在他们悔改之前（那时对他们毫无益处），还是在他们悔改之后（这样对他们才有益处），正如他所说：\Quote{因为他们自己，在失了味后从唯一的教会中退出，变得与教会相反，那么，就应当施行经上所写的：\InnerQuote{\BibleQuote{敌对法律者的房屋必须被洁净。}}因此，}他接着说，\Quote{那些被与基督为敌的人付洗而被玷污的人，应当先被洁净，然后才领洗。}那么，难道盗贼和杀人犯不反对那说\Quote{\BibleQuote{不可杀人；不可偷盗}}的法律吗？\Quote{因此他们必须被洁净。}谁会否认呢？然而，不仅是那些在教会内被这样的人付洗的人，而且那些自己就是这样的人、没有真心改变而受洗的人，当他们改变时，也不需再受洗。纯洗礼圣事的力量是如此之大，以至于尽管我们承认一个已受洗却继续过恶行的人需要被洁净，但我们禁止他再次受洗。\par

% structure: p0040 source=h2 target=subsection reason=relative-heading-rank; base=h2

\subsection{第十五章}

24. \Person{Crescens of Cirta}说：\Quote{我们最亲爱的\Person{西彼廉}写给\Person{Jubaianus}以及\Person{斯德望}的书信，在这么多最神圣的司祭弟兄的聚会上被宣读，这些书信包含了如此大量来自赐给我们天主的圣经的神圣见证，因此我们有充分的理由同意它们，我们全都借由天主的恩宠联合一致，我作出判断：所有异端者或分裂者愿意来到天主教会者，除非他们先被驱魔并领洗，否则不得进入；有明显例外，即那些原本在天主教会中受洗的人，这些人通过覆手而和好，并被接纳进入教会的告解圣事。}\par

25. 这里我们再次被提醒要探究为何他说：\Quote{当然，那些最初在天主教会内领受圣洗的人除外。}这是因为他们没有失去先前所领受的吗？那么，为什么他们不能将在教会之外所能拥有的事物也在教会之外传授呢？是因为在教会之外传授是非法的吗？但在教会之外拥有也是非法的，然而他们却拥有；所以，在教会之外给予是非法的，但却给予了。但是，给予从异端归来、曾在教会内领受圣洗的人，与给予来到教会、曾在教会外领受圣洗的人是一样的——即，使他能在教会内合法地拥有先前在教会外非法拥有的事物。但或许有人会问，在真福西彼廉致斯德望的书信中对此点有何论述，这封书信在这次公断中被提及，尽管不在议会开幕词中——我想是因为认为无必要。因为克雷森斯声称那封信已在集会中被宣读，我毫不怀疑，如果我没记错，这是照惯例做的，以使已聚集的主教们同时能从那封信所包含的主题中获取一些信息。因为那封信确实与当前主题无关；我更惊讶于克雷森斯竟认为应当提及它，而不是它被开幕词忽略。但如果有人认为我避而不谈那封信中任何与当前论点相关的重要内容，就让他去读一读，看看我所说的是否属实；或者，如果他发现并非如此，就让他指控我说谎。因为那封信丝毫没有涉及在异端者或分裂者中施行的圣洗，而这正是我们当前争论的主题。\par

% structure: p0043 source=h2 target=subsection reason=relative-heading-rank; base=h2

\subsection{第十六章}

26. 塞格米人尼科美德说：\Quote{我的判断是：来到教会的异端者应当领受圣洗，因为他们在罪人之中不能获得罪的赦免。}\par

27. 对此的答复是：整个天主教会的主张是，已经接受了基督的圣洗、尽管是在异端中的异端者，在来到教会时不应再次领受圣洗。因为如果在罪人中没有罪的赦免，那么教会内的罪人也不能赦罪；然而，那些由他们所施圣洗的人并不再次领受圣洗。\par

% structure: p0046 source=h2 target=subsection reason=relative-heading-rank; base=h2

\subsection{第十七章}

28. 吉尔巴人蒙努鲁说：\Quote{弟兄们，我们母亲、天主教会之真理，在我们中间已经延续并仍在延续，尤其在圣洗的三重本质上，正如我们的主说：\InnerQuote{你们要去，使万民成为门徒，因父、及子、及圣神之名给他们授洗。}\BibleRef{玛 28:19} 因此，}他接着说，\Quote{既然我们清楚知道异端者既没有父，也没有子，也没有圣神，他们在来到我们的母亲教会时，应当真正地重生并领受圣洗，以便他们曾有的癌症、谴责之忿怒以及毁灭性的错误力量，能藉着神圣而天上的洗濯被圣化。}\par

29. 对此我们答复：所有藉着以福音之言祝圣的圣洗而领受圣洗的人，仅在圣事中拥有圣父、圣子和圣神；但在内心和生活中，那些在教会内过着放荡而受诅咒生活的人，也不拥有祂。\par

% structure: p0049 source=h2 target=subsection reason=relative-heading-rank; base=h2

\subsection{第十八章}

30. 塞迪亚人塞昆迪努斯说：\Quote{既然我们的主基督说：\InnerQuote{不与我在一起的，就是反对我。}\BibleRef{玛 12:30} 而宗徒若望声明那些从教会出去的是假基督，\BibleRef{若一 2:18} 所以毫无疑问，基督的敌人、被称为假基督的人，不能施行赐予救恩之圣洗的恩宠；因此我的判断是：从异端的网罗逃到教会的人，应当由我们——因祂的屈尊俯就被称为天主的朋友——来为他们施圣洗。}\par

31. 对此的答复是：所有反对基督的人，就是那些在末日对他说\Quote{主啊，我们不是因你的名字行了许多奇迹吗？}以及所有那里记载的其他话的人，祂将回答：\Quote{我从来不认识你们，你们这些作恶的人，离开我吧！}\BibleRef{玛 7:22}——所有这些种类的糠秕注定被投入火中，如果他们坚持作恶到底，无论其中一部分在簸扬之前飞到外面，还是似乎留在里面。因此，如果那些来到教会的异端者要再次领受圣洗，以便由天主的朋友施圣洗，那么那些贪婪者、强盗、杀人犯，是天主的朋友吗？还是必须由他们所施圣洗的人重新领受圣洗？\par

% structure: p0052 source=h2 target=subsection reason=relative-heading-rank; base=h2

\subsection{第十九章}

32. 巴盖人费利克斯说：\Quote{正如盲人领盲人，两人都掉进坑里，\BibleRef{玛 15:14} 同样，当异端者为异端者施圣洗时，两人都一同陷入死亡。}\par

33. 这是真的，但他随后所加的并不真确。他说：\Quote{因此，异端者必须领受圣洗并得以活过来，免得我们这些活着的人与死人共融。}那些说\Quote{我们吃喝吧，因为明天要死了}的人岂不是死的吗？\BibleRef{格前 15:32} 因为他们不信死人的复活。那么，那些被他们恶毒的交往所败坏、并跟随他们的人，不也同样与他们一起掉进坑里吗？然而在他们中间，有些人已领受圣洗，宗徒正给他们写信；因此，如果他们悔改，也不会再次领受圣洗。同一宗徒不也说：\Quote{随从肉性的思念，乃是死亡}吗？\BibleRef{罗 8:6} 而贪婪者、欺诈者、强盗——西彼廉自己也在他们中间叹息——确实是随从肉性的。那么，死人能伤害那活在合一中的人吗？或者，谁能说，因为这样的人拥有或施予基督的圣洗，这圣洗就因他们的不义而被玷污了呢？\par

% structure: p0055 source=h2 target=subsection reason=relative-heading-rank; base=h2

\subsection{第二十章}

34. 米勒乌人波利安努斯说：\Quote{异端者在圣教会中领受圣洗是正确的。}\par

35. 诚然，没有什么比这更简短的了。但我认为这也同样简短：在基督的教会中，基督的圣洗不应被轻视，这是正确的。\par

% structure: p0058 source=h2 target=subsection reason=relative-heading-rank; base=h2

\subsection{第二十一章}

36. 希波勒基斯人特奥根斯说：\Quote{根据我们所领受的天上恩宠之圣事，我们相信在圣教会中只有一个圣洗。}\par

37. 这也可能是我自己的判断。因为它是如此平衡，以至于其中不包含任何与真理相悖的内容。因为我们也相信在圣教会内只有一个洗礼。假如他确实说：\Quote{我们相信只在圣教会内的那一个洗礼}，那么就必须给他和其余人相同的答复。但既然他这样表达：\Quote{我们相信在圣教会内的那一个唯一洗礼}，以至于肯定了洗礼存在于圣教会内，却没有否认它也可能在其他地方，无论他的本意如何，都无需辩驳这些话。因为如果我被分别询问这几个问题：首先，是否只有一个洗礼，我会回答是的。其次，是否这洗礼在圣教会内，我会回答是的。第三，是否我相信这洗礼，我会回答是的；因此我会回答我相信在圣教会内的那一个洗礼。但如果被问是否这洗礼只在圣教会内，而不在异端者和分裂者中，我会回答，与整个教会一致，我相信相反的观点。但由于他在自己的判断中没有加入这一点，我认为如果我为了辩驳而添加我没有看到的话，那只是轻率之举。因为如果他说：\Quote{\Place{幼发拉底河}的水只有一个，它在\Place{乐园}里}，没有人能反驳他所说的真理。但如果他被问那水是否只在\Place{乐园}里而不在别处，并说确实如此，那他就是在说假话。因为除了\Place{乐园}，它还存在于它从源头流出的那些地方。但谁敢轻率地说他可能会断言假话，而完全可能他是在断言真理呢？因此，这个判断的话无需反驳，因为它们绝不与真理相悖。\par

% structure: p0061 source=h2 target=subsection reason=relative-heading-rank; base=h2

\subsection{第二十二章}

38. \Place{Badiæ}的\Person{Dativus}说：\Quote{就我们能力所及，我们拒绝与异端者共融，除非他已在教会内受洗，并获得罪过的赦免。}\par

39. 对此的回答是：如果你之所以希望他受洗是因为他没有获得罪过的赦免，那么假设你在教会内发现一个已经受洗的人，却对自己的弟兄怀有仇恨，既然主不会说谎，祂说：\BibleQuote{如果你们不宽恕别人的过犯，你们的父也不会宽恕你们的过犯}\BibleRef{玛 6:15}，你会命令这样一个人在改正后重新受洗吗？肯定不会；同样你也不应命令异端者。我们不可忽略他为何没有简略地说\Quote{我们不与异端者共融}，而是加了\Quote{就我们能力所及}。因为他看到有更多的人同意这观点，但他和他的朋友们不能脱离这些人的共融，以免损害合一，所以他加了\Quote{就我们能力所及}——明确表明他并不情愿与那些他认为没有洗礼的人共融，但为了和平与合一，一切都应忍受；正如那些认为\Person{Dativus}及其一派错了的人所做的那样，他们持有后来通过更充分的真理宣告和古代习俗所教导的观点，并得到了后来大公会议更强有力的确认；然而反过来，他们怀着焦虑的热忱彼此容忍，虽然持有不同意见却没有违反基督徒的爱德，努力\BibleQuote{以和平的联系保持圣神的合一}\BibleRef{弗 4:3}，直到天主启示他们中的一人，即使他有不同想法，也认识这错误。\BibleRef{斐 3:15} 我希望那些因这个同一大公会议的权威而攻击合一的人注意这点，这会议正宣告合一应受到多大的珍爱。\par

% structure: p0064 source=h2 target=subsection reason=relative-heading-rank; base=h2

\subsection{第二十三章}

40. \Place{Abbir Germaniciana}的\Person{Successus}说：\Quote{异端者要么什么都不能做，要么什么都能做。如果他们能施洗，他们也能赐予圣神；但如果他们不能赐予圣神，因为他们不拥有圣神，那么他们也不能属灵地施洗。因此我们作出判断：异端者应该受洗。}\par

41. 对此我们几乎可以逐字回答：杀人者要么什么都不能做，要么什么都能做。如果他们能施洗，他们也能赐予圣神；但如果他们不能赐予圣神，因为他们不拥有圣神，那么他们也不能属灵地施洗。因此我们作出判断：由杀人者施洗的人，或者已经受洗但未悔改的杀人者本人，当他们改正后，应该受洗。但这并不真实。因为\BibleQuote{凡恨自己弟兄的，就是杀人者}\BibleRef{若一 3:15}；而\Person{Cyprian}知道在教会内就有这样的人，他们确实施洗。因此，用这类关于异端者的话是徒劳的。\par

% structure: p0067 source=h2 target=subsection reason=relative-heading-rank; base=h2

\subsection{第二十四章}

42. \Place{Thuccabori}的\Person{Fortunatus}说：\Quote{耶稣基督我们的主天主，圣父及创造者天主的圣子，将祂的教会建立在磐石上，而非异端上，并将施洗的权力赐给了主教，而非异端者。因此，那些在教会之外，与基督对立，分散祂的羊群的人，不能在外面施洗。}\par

43. 他加上\Quote{外面}一词，是为了避免像对\Person{苏克塞苏斯}那样以同样简短的话回答他。否则，他也可以逐字回答：\Quote{耶稣基督、我们的主和天主、圣父和造物主的天主之子，将祂的教会建立在磐石上，而非不义之上，并将施洗的权力赐予主教，而非不义之人。}因此，那些不属于他们建造其上的磐石、听了天主的话而实行的（\BibleRef{玛 7:24}），却相反基督而生活，听了话而不实行，因此建造在沙土上，这样以败坏的品行作为榜样而分散祂的羊群的人，不能施洗。这话岂非貌似真实？然而它是虚假的。因为不义之人的确施洗，\Person{西彼廉}认为那些强盗（也是不义之人）与自己共融。但为此，多纳徒派说，他加了\Quote{外面}。那么，为什么他们不能在外面施洗呢？是否因为他们在外面，所以本身就更坏？但就圣洗的有效性而言，施行者有多坏并无区别。因为坏与更坏之间的差别，不如好与坏之间的差别大；然而，当坏人施洗时，他给出与好人相同的圣事。因此，当更坏的人施洗时，他也给出与较不坏者相同的圣事。或者，是否因为不是基于人的功德，而是基于圣洗圣事本身，它不能在外面被授予？如果这样，它也不能在外面被拥有，那么一个人每次离开教会又回来，就必须重新受洗。\par

44. 再者，如果我们更仔细地探究\Quote{外面}的含义，尤其他自己提到教会建立其上的磐石，那么，那些在磐石上的人不就是在教会内吗？而那些不在磐石上的人，也不在教会内吗？因此，现在让我们看看，那些听了基督的话而不实行的人，是否将房屋建造在磐石上。主亲自宣告相反的情况，说：\BibleQuote{凡听了我的话而实行的，我要把他比作一个聪明人，他把房屋建造在磐石上；}稍后又说：\BibleQuote{凡听了我的话而不实行的，好比一个愚昧人，他把房屋建造在沙土上。}因此，如果教会是在磐石上，那么那些在沙土上的人，因为他们是在磐石之外，必然是在教会之外。所以，让我们回想，\Person{西彼廉}提到多少在内部的人却是建造在沙土上的，也就是那些听了基督的话而不实行的人。因此，因为他们是在沙土上，就被证明是在磐石之外，也就是在教会之外；然而即使他们处于这种状态，要么尚未、要么永远不会改好，他们不仅施洗和受洗，而且他们所拥有的圣洗在他们身上仍然有效，尽管他们注定要受罚。\par

45. 在此也不能说：然而谁能实行主在福音讲道中所写的一切话呢？在该讲道的结尾，祂说，听了这些话而实行的，建造在磐石上；听了而不实行的，建造在沙土上。因为，即使某些人没有完全实现所有的话，但在这同一讲道中，祂指定了补救方法，说：\BibleQuote{你们宽恕，也必蒙宽恕。}\BibleRef{路 6:37}在详细记录了同一讲道中的主祷文之后，祂说：\BibleQuote{我告诉你们：如果你们宽恕别人的过犯，你们的天父也必宽恕你们；但如果你们不宽恕别人的过犯，你们的父也不会宽恕你们的过犯。}\BibleRef{玛 6:14}因此\Person{伯多禄}也说：\BibleQuote{因为爱德遮盖许多罪过。}\BibleRef{伯前 4:8}他们肯定没有这爱德，因此他们建造在沙土上。关于这些人，\Person{西彼廉}说，他们在教会内，甚至在宗徒时代，以不合基督爱德的恶意仇恨来往；所以他们看似在内，实则在外，因为他们不在那代表教会的磐石上。\par

% structure: p0072 source=h2 target=subsection reason=relative-heading-rank; base=h2

\subsection{第二十五章}

46. \Person{塞达图斯}（来自\Place{图布尔波}）说：\Quote{正如水，在教会内因司铎的祈祷而被祝圣，洗去罪过；同样，当它被异端者的话感染，如被癌症污染时，就增多罪过。因此，必须尽一切促进和平的努力，使任何被异端错误感染玷污的人，不要拒绝接受唯一真实的圣洗；凡不受此圣洗的，不能继承天国。}\par

47. 对此我们回答：如果因司铎祈祷时缺乏技巧而念出某些错误的言辞，以致水未得圣化，那么在教会本身内，不仅许多恶人，连许多善弟兄也未能使水圣化。因为许多祈祷文在向更有学问的人诵念时，每日都被改正，其中发现许多违反公教信仰的内容。那么，假设有人是在这些祈祷文已诵念于水上之后受洗的，他们会被命令重新受洗吗？为何不呢？因为通常祈祷中的缺失远为献祷者的意向所弥补；而那些没有它们圣洗便不能成圣的固定福音言辞具有如此效力，以致凭借它们的力量，祈祷中所念出的任何违反信仰规则的错误内容都归于无效，就如因基督之名驱逐魔鬼一样。因为显然，如果一个异端者念出错误的祈祷，他并没有爱的良好意向来弥补这种缺乏技巧，因此他就像天主教会本身内部任何怀有嫉妒或恶意的兄弟一样——西彼廉证实这类人存在于教会之内。或者，有人可能献上某种祈祷（这并非罕见），其中他竟违反信仰规则发言，因为许多人鲁莽地使用并非仅由不善讲话者、甚至由异端者所写的祈祷文，并在无知中因单纯而以为它们良好而使用它们；然而其中的错误并不败坏正确的部分，反而使之归于无效，正如在一位怀抱良好希望且信德被认可的人身上（但他仍是一个人），如果他在某方面有不同见解，他所持有的正确部分并不因此受损，直到天主也向他启示他所有不同见解之处。\BibleQuote{\BibleRef{斐 3:15}}但假设那人本身邪恶乖僻，那么，他若献上正直的祈祷（任何部分都不违反公教信仰），并不因为祈祷正确，那人本身也正确；而如果他向某些人献上错误的祈祷，天主亲自临在，以支持祂福音的言辞（没有这些言辞，基督的圣洗不能成圣），祂亲自圣化祂的圣事，使在接受者身上——无论在他受洗之前、受洗之时，或在他将来真正归向天主的时刻——那同一圣事有益于救恩，而若他不归正，它就会有力于他的毁灭。但谁不知道，如果没有构成外在有形标记的福音言辞，就没有基督的圣洗呢？然而，你更容易找到根本不施洗的异端者，而难以找到不念那些言辞就施洗的人。因此，我们不说每一个洗礼（因为在许多偶像崇拜的可憎仪式中，人们被说成是受洗），而是说基督的圣洗，即任何在福音言辞中成圣的洗礼，处处相同，且不能因任何人的乖僻而受玷污。\par

48. 我们不可轻忽地放过在这判决中他插入了一处，说：\Quote{因此，我们必须以一切有助于和平的努力，竭力使没有一个人被感染……}等等。因为他顾及到真福西彼廉在开场白中的话：\Quote{不判断任何人，也不剥夺任何人的共融权利，若他持不同见解。}请看，在教会的善子们内，合一与和平的爱具有何等力量：他们宁愿对那些他们称之为亵圣者不洁者（他们认为这些人未得圣洗圣事就被接纳了）表示容忍，如果他们不能如他们所认为的正确那样纠正他们，也不愿因他们的缘故而打破那神圣的纽带，以免因莠子连麦子也一起拔出来\BibleQuote{\BibleRef{玛 13:29}}——在他们力所能及的范围内，如同在撒罗满那最高贵的判决中，允许婴孩的身体宁可由假母亲喂养，也不被劈成两半\BibleQuote{\BibleRef{列上 3:26}}。然而，这既是那些对圣洗圣事持更正确观点者的意见，也是那些天主因他们极大的爱而打算向他们启示任何他们所有不同见解之点的人的意见。\par

% structure: p0076 source=h2 target=subsection reason=relative-heading-rank; base=h2

\subsection{第26章}

49. \Person{Privatianus} of \Place{Sufetula}说：\Quote{若有人说异端者拥有施洗的能力，他应首先说谁创立了异端。因为若异端出自天主，它可能享有天主的恩宠；但若不出自天主，它怎能拥有或赋予任何人天主的恩宠呢？}\par

50. 此人可以逐字回答如下：若有人说心怀恶意和嫉妒的人拥有施洗的能力，他应首先说谁是恶意和嫉妒的创始人。因为若恶意和嫉妒出自天主，它们可能享有天主的恩宠；但若不出自天主，它们怎能拥有或赋予任何人天主的恩宠呢？然而，正如这些话同样明显虚假，那些被说出以反驳它们的话也是如此。因为心怀恶意和嫉妒的人施洗，就连西彼廉自己也承认，因为他证明他们也在教会之内。因此，连异端者也能施洗，因为洗礼是基督的圣事；但嫉妒和异端是魔鬼的作为。然而，一个人虽拥有它们，并不因此使他所拥有的基督圣事也被算在魔鬼的作为数量中。\par

% structure: p0079 source=h2 target=subsection reason=relative-heading-rank; base=h2

\subsection{第27章}

51. \Person{Privatus} of \Place{Sufes}说：\Quote{对于认可异端者洗礼的人，除了说他与异端者共融外，还能说什么呢？}\par

52. 对此我们回答：我们在异端者身上认可的并非异端者的洗礼，如同我们在贪婪者、背信者、欺诈者、强盗或嫉妒者身上认可的也不是他们的洗礼；因为所有这些都是不义之人，但基督是正义的，祂的圣事存在于他们身上，他们并不在本质上破坏它。否则，另一个人可以说：对于认可不义者洗礼的人，除了说他与不义者共融外，还能说什么呢？如果这个反对意见被用来反对天主教会本身，它也会得到答复，正如我答复上述问题那样。\par

% structure: p0082 source=h2 target=subsection reason=relative-heading-rank; base=h2

\subsection{第28章}

53. \Place{拉雷斯}的\Person{霍滕修斯}说：\Quote{有多少圣洗，让那些支持或偏向异端者去决定吧。我们宣认教会只有一个圣洗，我们只在教会内知道它。或者，那些基督亲自宣称是祂仇敌的人，又怎能因基督的名给人施洗呢？}\par

54. 给这个人以类似言辞的回答，我们说：让那些支持或偏向不义者的人看着办吧：我们将那唯一的圣洗——我们知道它只属于教会，无论它在哪里——引领回教会。或者，那些基督亲自宣称是祂仇敌的人，又怎能因基督的名给人施洗呢？因为祂对所有不义者说：\BibleQuote{我从来不认识你们，你们这些作恶的人，离开我吧！}（\BibleRef{玛 7:23}）然而，当他们施洗时，并不是他们自己在施洗，而是祂——\Person{若翰}论及祂时说：\BibleQuote{就是那位施洗的。}（\BibleRef{若 1:33}）\par

% structure: p0085 source=h2 target=subsection reason=relative-heading-rank; base=h2

\subsection{第二十九章}

55. \Place{马科马德斯}的\Person{卡修斯}说：\Quote{既然不能有两个圣洗，那将圣洗赐予异端者的人，就从自己身上夺走了圣洗。因此，我声明我的判断：当异端者——这些令我们流泪的对象，这些腐朽之物——开始来到教会时，应当给他们付洗，好藉神圣而属天的洗礼洗净，并得生命之光照耀，然后被接纳进入教会——如今他们不再是敌人，而是和平的人；不再是外人，而是主内信仰的家人；不再是私生子，而是天主的子女；不再参与错谬，而是分享救恩——但那些从教会中迁出的信者转而投奔异端者除外，这些人应藉覆手而恢复。}\par

56. 另一个人可能会说：既然不能有两个圣洗，那将圣洗赐予不义者的人，就从自己身上夺走了圣洗。但连我们的对手也会与我们一同抵制此人，当他说我们将圣洗赐予不义者时，这圣洗并非属于不义者，如同他们的不义，而是属于基督——义是从祂而来，祂的圣事即使在不义者之间，也不是不义的。因此，他们本应一同对我们所说的关于不义者的话，让他们对异端者也这样说吧。所以他更应该这样说：因此我声明我的判断：当异端者——这些令我们流泪的对象，这些腐朽之物——开始来到教会时，如果他们已拥有基督的圣洗，就不应再给他们付洗，而应纠正他们的错谬。因为我们可以同样对不义者（异端者是其中一部分）说：因此我声明我的判断：不义者——这些令我们流泪的对象，这些腐朽之物——如果他们已经受过洗，当他们开始来到教会时，就不应再受第二次洗；就是说，来到那磐石——其外都是那些听了基督的话却不实行的人；但既然已藉神圣而属天的洗礼洗净，如今又得真理之光的照耀，他们应被接纳进入教会，不再是敌人，而是和平的人，因为不义者没有和平；不再是外人，而是主内信仰的家人，因为有不义者对之所说的话：\BibleQuote{你怎么竟变成了我这棵异样葡萄树的坏枝子？}（\BibleRef{耶 2:21}）不再是私生子，而是天主的子女，因为不义者是魔鬼的子女；不再参与错谬，而是分享救恩，因为不义不能救人。我所指的教会，就是那磐石，那鸽子，那关闭的园子和封住的泉源，它只在小麦中被认出，不在糠秕中——无论这些糠秕是被风远远吹散，还是直到最后簸扬时似乎仍与麦子混在一起。因此，卡修斯所加的这句：\Quote{但那些从教会中迁出的信者转而投奔异端者除外}，是白费的。因为如果他们自己因分裂而失去了圣洗，就让他们也恢复圣洗；但如果他们没有失去它，那么他们给予的圣洗就应得到承认。\par

% structure: p0088 source=h2 target=subsection reason=relative-heading-rank; base=h2

\subsection{第三十章}

57. \Place{凯撒村}的另一位\Person{雅努阿里乌斯}说：\Quote{如果错谬不服从真理，那么真理更不会同意错谬；因此我们坚守我们所主持的教会，以便她只将自己的圣洗据为己有，我们给那些教会未曾付洗的人付洗。}\par

58. 我们回答：教会所施洗的人，就是那磐石所施洗的人，在磐石之外的都是那些听了基督的话却不实行的人。因此，凡被这样的人施洗的，都应重新受洗。但若不这样做，那么，正如我们在这些人身上承认基督的圣洗，同样也应在异端者身上承认，尽管我们谴责或纠正他们的不义和错谬。\par

% structure: p0091 source=h2 target=subsection reason=relative-heading-rank; base=h2

\subsection{第三十一章}

59. \Place{卡尔皮斯}的另一位\Person{塞昆迪努斯}说：\Quote{异端者是基督徒还是不是？若是基督徒，为何不在天主的教会里？若不是基督徒，就让他们成为基督徒。否则，主的话语中祂所说的那句话又指何意？祂说：\InnerQuote{不跟我在一起的，就是反对我的；不跟我收集的，就是分散的。}（\BibleRef{玛 12:30}）由此可知，圣神不能仅藉覆手而降临在异邦子女和假基督的苗裔身上，因为显然异端者没有圣洗。}\par

60. 对此我们回答：不义的人是不是基督徒？如果是基督徒，为何他们不在那盘石上，教会就建立在其上？因为他们听了基督的话，却不实行。如果不是基督徒，就让他们成为基督徒。否则，主的话语中，祂说\BibleQuote{不与我相合的，就是反对我的；不同我收集的，就是分散的？}又有什么意义呢？他们因伪仿效主而引羊群走向生命的败坏，是在分散主的羊。由此清楚可见，对于外邦之子（所有不义之人如此称呼），以及对于假基督的苗裔（所有反对基督的人都是），单靠按手，圣神并不能降临，除非加上真心的悔改；因为显然，不义之人，只要他们还是不义的，固然可以拥有圣洗，却不能拥有圣洗所象征的救恩。让我们看看，在那篇圣咏中，是否以如下的话描述异端者，论到外邦之子：\BibleQuote{上主，求你救我脱离外邦之子之手，他们的口谈论虚妄，他们的右手是欺骗的右手；他们的儿子好比茁壮的树苗，他们的女儿好似雕像般修饰的宫殿；他们的仓库充满，供应各种存储；他们的羊群繁衍，在街上多产；他们的肥牛强壮；没有篱笆的倒塌，没有通路的开启，街市上没有怨言。人们认为这样的民族有福；但真正有福的是以上主为天主的人民。}因此，如果那些将幸福寄于暂时事物和地上丰裕、轻看主诫命的是外邦之子，那么让我们看看，西彼廉所谈论的是否正是这些人，他将他们也变成自己，以表明他是在谈论那些与他共享圣事的人。他说：\Quote{不遵守主的道路，也不遵守赐给我们为救恩的天上诫命。我们的主遵行了天父的旨意，而我们却不遵行主的旨意，热衷于我们的祖业或盈利，追求骄傲，等等。}但如果这些人既能拥有圣洗也能传递圣洗，为何要否认圣洗可以存在于外邦之子中，而西彼廉却劝告他们，通过遵守藉独生子传给他们的天上诫命，应配得成为祂的兄弟和天主的儿子呢？\par

% structure: p0094 source=h2 target=subsection reason=relative-heading-rank; base=h2

\subsection{第三十二章}

61. \Person{塔布拉卡的维克托里库斯}说：\Quote{如果异端者可以施洗，并赐予罪过的赦免，我们为何要破坏他们的信誉，称他们为异端者呢？}\par

62. 如果另一个人说：如果不义的人可以施洗，并赐予罪过的赦免，我们为何要破坏他们的信誉，称他们为不义的人呢？我们对于不义之人所应给出的回答，也同样可以给予对方关于异端者——那就是，首先，他们所施行的圣洗不是他们自己的；其次，凡拥有基督的圣洗的人，也未必确定得到自己罪过的赦免，如果他只有外在的标记，而没有真心悔改的转变，以至于那赐予赦免的人自己应当得到罪过的赦免。\par

% structure: p0097 source=h2 target=subsection reason=relative-heading-rank; base=h2

\subsection{第三十三章}

63. 另一位\Person{乌提纳的费利克斯}说：\Quote{最神圣的司祭职的弟兄们，没有人能怀疑，人的僭妄没有如此大的能力，如同我们主耶稣基督可敬可畏的威严。因此，记住这危险，我们不仅自己应当遵守这一点，还应当通过我们的全体同意来确认，即所有来到慈母教会怀抱的异端者都应受洗，以便那因长期败坏而被玷污的异端思想，在藉水洗的圣化得到洁净后，得以革新。}\par

64. 或许这位将他对异端者施洗的论据建立在清除长期败坏之上的人，会放过那些一头扎进某种异端、只停留短暂时间、随后得到纠正并转至天主教会的。此外，他本人没有看到，可以说所有来到那盘石（意指教会）的不义之人都应受洗，以便那不义的心思，它在盘石之外把房子建在沙土上，听了基督的话却不实行，可藉水洗的圣化得到洁净而革新；然而，如果他们已受过洗，这便不能发生，即使能证明他们受洗时的情形正是如此，即他们\Quote{在言语上弃绝了世界，却在行为上没有。}\par

% structure: p0100 source=h2 target=subsection reason=relative-heading-rank; base=h2

\subsection{第三十四章}

65. \Person{布鲁格的奎埃图斯}说：\Quote{我们因信德生活的人，应当以信德的服从，遵守从前预先告诉我们为教训的事。因为撒罗满记载：\BibleQuote{被死人洗过的人，他的洗涤有何益处？}他这话显然是指那些被异端者所洗的人，以及那些施洗的人。因为如果他们在异端中受洗的人能通过罪过的赦免得永生，他们为何要来到教会？但如果从死人那里得不到救恩，因而他们承认先前的错误，带着补赎回归真理，他们就应当被那惟一的、在天主教会内的赋予生命的圣洗所圣化。}\par

66. 我们已在先前充分说明了被死者受洗是什么意思，但对同一段圣经的更仔细考虑并无妨碍。但我想要问的是，为什么他们认为只有异端者才是死的，而宗徒\Person{保禄}却普遍地论及罪说：\BibleQuote{罪的工价乃是死。}（\BibleRef{罗 6:23}）又说：\BibleQuote{属肉体的心思乃是死。}（\BibleRef{罗 8:6}）当他说\BibleQuote{一个享乐的寡妇是死的}（\BibleRef{弟前 5:6}）时，那些\Quote{在言语上而非在行为上弃绝世俗}的人怎么不是死的呢？因此，由他们授洗的人所受的圣洗有什么益处呢？除非他自己也是同样的人，他确实有圣洗，但于救恩无益。但如果授洗者是那样的人，而受洗者以不虚伪的心归向主，那么他不是由那个死者受洗，而是由那活着的、经上说\BibleQuote{这就是那施洗者}（\BibleRef{若 1:33}）的一位受洗。至于他关于异端者所说的，如果他们在其中受洗的人借由罪的赦免获得永生，为什么他们要来到教会呢？我们回答：他们来是为了这个原因：虽然他们接受了基督的圣洗，直至举行圣事的地步，但他们不能借着合一的爱德获得永生；正如那些嫉妒和恶意的人不能获得永生一样，即使他们只怀有对所伤害之人的仇恨，他们的罪也不会被赦免；因为真理说：\BibleQuote{如果你们不宽恕别人的过犯，你们的父也必不宽恕你们的过犯}（\BibleRef{玛 6:15}），更何况当他们仇恨那些以善报恶的人呢？然而这些人，虽然\Quote{在言语上而非在行为上弃绝世俗}，如果他们后来被纠正，不会重新受洗，而是由那唯一的活圣洗成圣。这确实在天主教会内，但不仅限于此，正如它不仅在那些建基于磐石、\BibleQuote{并由那只鸽子所组成的圣人们内一样}（\BibleRef{歌 6:9}）。\par

% structure: p0103 source=h2 target=subsection reason=relative-heading-rank; base=h2

\subsection{第三十五章}

67. \Person{卡斯图斯} of \Place{西卡}说：\Quote{凡不顾真理而傲慢地遵行习惯的人，要么是对那些被启示真理的弟兄怀有嫉妒和恶意，要么是对那以默感教导祂教会的主忘恩负义。}\par

68. 如果此人证明那些与他意见不同、并且持有后来在基督徒大公会议批准下被普世接受的观点的人，是遵行习惯而轻视真理，那么我们有理由惧怕这些话；但既然看到这习惯既源于真理又被真理所证实，我们在这场审判中无所畏惧。然而，如果他们对弟兄怀有嫉妒或恶意，或对主忘恩负义，请看看他们是愿意与什么样的人共融；看看他们对待与自己意见不同的人，正如\Person{西彼廉}起初所命令的，不将他们从共融的权利中移除；看看在保持合一中他们不被什么样的人玷污；看看和好的纽带何等值得珍爱；看看那些基于主教会议（他们的前辈）而控告我们的人持有什么观点，他们并不效法这些前辈的榜样，而当案件的权利被考虑时，他们的榜样反倒定了他们的罪。如果按照这个判决所见证的，习惯是接受来到教会的异端者，并保留他们已经有的圣洗，那么要么这样做是正确的，要么恶人在合一中不会玷污善人。如果这样做是正确的，为什么他们因为世界如此接受而控告世界？但如果恶人在合一中不会玷污善人，他们又如何为自己犯下的亵渎分裂之罪辩护呢？\par

% structure: p0106 source=h2 target=subsection reason=relative-heading-rank; base=h2

\subsection{第三十六章}

69. \Person{欧克拉提乌斯} of \Place{特尼}说：\Quote{我们的天主和主耶稣基督，亲口教导宗徒们，完全奠定了我们的信仰、圣洗的恩宠和教会律法的规则，说：\InnerQuote{你们要去，教训万民，因父、及子、及圣神之名给他们授洗。}（\BibleRef{玛 28:19}）因此，异端者的虚假不义的圣洗，我们当以庄严的见证予以拒绝和反驳，因为从他们的口中发出的不是生命，而是毒物，不是天上的恩宠，而是对天主圣三的亵渎。所以显然，来到教会的异端者应当以完全而大公的圣洗受洗，好使他们从他们狂妄的亵渎中被洁净，并由圣神的恩宠得以改造。}\par

70. 显然，如果圣洗不是因父、及子、及圣神之名祝圣的，那么它应被视为异端者的，并由我们以庄严的见证拒绝为不义的；但如果我们在这名号中辨别出它来，我们最好将福音的言语与异端者的错误区分开来，认可其中健全的部分，纠正有缺陷的部分。\par

% structure: p0109 source=h2 target=subsection reason=relative-heading-rank; base=h2

\subsection{第三十七章}

71. \Person{利博苏斯} of \Place{瓦加}说：\Quote{主在福音中说：\InnerQuote{我就是真理；}（\BibleRef{若 14:6}）祂没有说，我是习惯。因此，当真理显明时，让习惯屈服于真理；这样，即使过去在教会中有人没有给异端者授洗，现在可以开始给他们授洗。}\par

72. 在此他完全没有试图说明他所声称的习俗应当屈从的真理究竟是什么。但更重要的是，他承认习俗确实存在，这反而帮助我们对抗那些自合一中分离的人，胜过他认为习俗应屈从于一个他并未指明的真理。因为习俗的性质是如此：若它容许亵圣之人未经圣洗的净化就接近基督的祭台，且不玷污那留在合一中的任何善人，那么所有自这同一合一中割裂出来的人，既然在合一中不可能被任何恶人的沾染所玷污，就毫无理由地自分离，并犯了明显是分裂的亵圣之罪。但若所有人都因那习俗的污染而丧亡，那么他们又从哪个洞穴里钻出来，既缺乏原始的真理，又带着诽谤的一切诡计呢？然而，若那习俗（即如此接纳异端者）是正确的，那么他们应放弃疯狂，承认错误，来到天主教会——并非为了重新以圣洗圣事沐浴，而是为了从割裂的创伤中得到痊愈。\par

% structure: p0112 source=h2 target=subsection reason=relative-heading-rank; base=h2

\subsection{第三十八章}

73. \Person{Lucius of Thebaste}说：\Quote{我宣告我的判断：异端者、亵圣者及不义之人，他们用各种言辞窃取圣经神圣可敬的话语，应被视作受诅咒的，因此要驱魔并受圣洗。}\par

74. 我也认为他们应被视作受诅咒的，但并非因此就要驱魔并受圣洗；因为我诅咒的是他们自己的谎言，但我尊敬基督的圣事。\par

% structure: p0115 source=h2 target=subsection reason=relative-heading-rank; base=h2

\subsection{第三十九章}

75. \Person{Eugenius of Ammedera}说：\Quote{我也宣告同样的判断：异端者应该受圣洗。}\par

76. 我们回答他：但这并非教会所宣告的判断；如今天主已在一个全体会议中启示了你们那时尚持异议之点（\BibleRef{斐 3:15}），但因着救恩的爱德在你们内，你们仍留在合一中。\par

% structure: p0118 source=h2 target=subsection reason=relative-heading-rank; base=h2

\subsection{第四十章}

77. 另一位\Person{Felix of Ammacura}说：\Quote{我也依从圣经的权威宣告我的判断：异端者应受圣洗，连同那些声称在分裂者中已受圣洗的人。因为，按基督的警告，我们的泉源是关闭的，只属于我们自己（\BibleRef{歌 4:12}），让我们教会所有的敌人都明白，它不能属于他人；我们羊群的牧者也决不会将救恩之水赐给两个不同的民族。因此显然，异端者和分裂者都不能领受任何天上的事物，他们竟敢从有罪和教会之外的接受。当施予者没有立足之地时，接受者也决然无法获益。}\par

78. 我们回答他：圣经从未命令要在异端者中受圣洗的异端者重新受圣洗，反而在许多地方指出：所有不在磐石上、不属于鸽子肢体的人都是教会之外的，然而他们仍施圣洗和受圣洗，并拥有救恩的圣事却没有救恩。但我们的泉源如何像乐园的泉源——正如它同样流到乐园的边界之外——上面已充分论述；至于\Quote{我们羊群的牧者决不会将救恩之水赐给两个不同的民族}，即赐给祂自己的和赐给外人的，我完全同意。但这难道意味着，因为水不是为救恩，它就不是同一的水吗？因为洪水的水对在方舟内的人是救恩，但对在外的人却带来死亡，然而水是同一的。还有许多外人，就是说心怀嫉妒的人——\Person{Cyprian}从圣经中宣告并证明他们属于魔鬼的一党——他们好像在方舟之内，然而他们若非在方舟之外，就不会被水所灭。因为这样的人被圣洗所杀，如同基督的馨香对那些使徒所说的成了死亡的气味（\BibleRef{格后 2:15}）。那么，为什么异端者或分裂者不能领受任何天上的事物呢？正如荆棘和莠子，像那些在方舟外的人确实接受了从天而来的洪水之雨，但为毁灭而非为救恩。因此，我不费心去反驳他最后所说的：\Quote{当施予者没有立足之地时，接受者也决然无法获益。}因为我们也说，当接受者在异端中、与异端者合意时，它对他们无益；所以他们来到天主教会的和平与合一，并非为了接受圣洗，而是为了使他们所领受的能开始有益于他们。\par

% structure: p0121 source=h2 target=subsection reason=relative-heading-rank; base=h2

\subsection{第四十一章}

79. 另一位\Person{Januarius of Muzuli}说：\Quote{我感到惊奇，既然众人都承认只有一个圣洗，为何众人却不明白同一圣洗的合一。因为教会与异端是两件不同的事。若异端者有圣洗，我们就没有；但若我们有，异端者就不能有。然而无疑只有教会拥有基督的圣洗，因为只有教会拥有基督的恩宠与真理。}\par

80. 另一个人可以同样不真实地说：\Quote{我感到惊奇，既然众人都宣认只有一个圣洗，为何众人却不明白圣洗的合一。因为义与不义是两件不同的事。若不义者有圣洗，义人就没有；但若义人有，不义者就不能有。然而无疑只有义人拥有基督的圣洗，因为只有义人拥有基督的恩宠与真理。}这无疑是假的——正如他们自己所承认的。因为那些心怀嫉妒的人——他们属于魔鬼的一党，虽在教会内（如\Person{Cyprian}告诉我们的），并为使徒保禄所深知——他们有圣洗，却不属于那只安然栖身于磐石之上的鸽子的肢体。\par

% structure: p0124 source=h2 target=subsection reason=relative-heading-rank; base=h2

\subsection{第四十二章}

81. \Person{Adelphius of Thasbalte}说：\Quote{他们用虚假和恶意之词指责真理，说我们重新施圣洗，这毫无理由，因为教会并非给异端者重新施圣洗，而是给他们施圣洗。}\par

82. 诚然，它并未给他们重新付洗，因为它只给那些先前未受洗者授洗；而这一较早的惯例只是在后来的会议中，通过更审慎地完善真理才得以确认。\par

% structure: p0127 source=h2 target=subsection reason=relative-heading-rank; base=h2

\subsection{第四十三章}

83. \\Person{小勒普蒂斯的德默特琉斯}说：\\Quote{我们坚持一个圣洗，因为我们只将属于天主教会本身的归于她。但那些说异端者真实合法地施洗的人，正是他们自己制造了不止两个，而是许多圣洗；因为既然异端为数众多，圣洗也将按它们的数目来计算。}\par

84. 我们回答他说：倘若如此，那么有多少本性的行为，就会算作有多少圣洗，宗徒说\\BibleQuote{行这样事的人决不能承受天主的国}（\\BibleRef{迦 5:21}），在这些行为中也算有异端；在教会内，这些行为中有许多像在糠秕中被容忍，然而为它们全体只有一个圣洗，这圣洗不被任何不义的行为所败坏。\par

% structure: p0130 source=h2 target=subsection reason=relative-heading-rank; base=h2

\subsection{第四十四章}

85. \\Person{提巴里的味增爵}说：\\Quote{我们知道异端者比外邦人更坏。如果他们悔改，愿意来到天主面前，他们必定有真理的准则，就是主以祂神圣的命令托付给宗徒们的，说：\\InnerQuote{你们往普天下去，因我的名按手，驱逐魔鬼。}（\\BibleRef{谷 16:15}）又在另一处说：\\InnerQuote{你们要去使万民成为门徒，因父、子及圣神之名给他们授洗。}（\\BibleRef{玛 28:19}）因此，首先借着驱魔时的按手，其次借着圣洗中的重生，他们才能到达基督的恩许；但依我的判断，不应以其他方式这样做。}\par

86. 我不知道他根据什么规则说异端者比外邦人更坏，因为主说：\\BibleQuote{若他不听从教会，你就将他看作外邦人或税吏。}（\\BibleRef{玛 18:17}）难道异端者比这样的人还坏吗？我不否认。然而，我不承认因为那人本身比外邦人，即比外邦人和异教徒更坏，因此凡圣事中属于基督的任何东西，都与他的恶习和品性混在一起，并因这样的混合而败坏。因为即使那些离开教会、不是成为追随者而是成为异端创立者的人，在脱离之前已经受过洗，他们继续拥有圣洗，尽管按照上述规则，他们比外邦人更坏；因为如果他们在改正后回来，他们不会重新领受圣洗，而如果他们已经失去了它，他们当然会这样做。因此，一个人可能比外邦人更坏，但基督的圣事不仅仍在他内，而且丝毫不逊于它在圣洁和义人内的状态。因为尽管就他的能力而言，他没有保存这圣事，反而在内心和意志上施以暴力，但就圣事本身的性质而言，它在鄙视和拒绝它的人身上仍然完整无损。难道索多玛人不就是外邦人，即异教徒吗？那么犹太人更坏，主对他们说：\\BibleQuote{在审判的日子，索多玛地所受的惩罚要比你容易承受。}（\\BibleRef{玛 11:24}）先知也对他们说：\\BibleQuote{你使索多玛成为了义人。}（\\BibleRef{则 16:51}）也就是说，与你相比，索多玛是义人。难道我们要因此主张，存在于犹太人中的神圣圣事也具有了犹太人自身的性质——那些圣事是主自己也接受的，祂还打发被治好的癞病人去履行它们（\\BibleRef{路 17:14}），当\\Person{匝加利亚}施行这些圣事时，天使站在他旁边，宣告他在圣殿中献祭时祈祷已蒙垂听？同样的圣事，既存在于那个时代的好人身上，也存在于那些比外邦人更坏的坏人身上，因为他们在邪恶上被列为比索多玛人更甚，但那些圣事在两者中都是完美而神圣的。\par

87. 因为即使外邦人在他们的教义中可能有一些圣善和正确的东西，我们的圣人并没有谴责它，无论外邦人本人因他们的迷信、偶像崇拜、骄傲以及其他败坏而多么可憎，并且若不受矫正，将受到来自天上的审判惩罚。因为当宗徒\\Person{保禄}在雅典人面前谈论天主时，他引证了他所说的话，证明他们中有些人曾说过类似的话（\\BibleRef{宗 17:28}），如果他们来到基督面前，这些话肯定不会被谴责，反而会得到承认。\\Person{圣西彼廉}也引用了类似的证据来反驳同样的外邦人；因为，在谈到术士时，他说：\\Quote{然而，他们中的首领\\Person{霍斯塔内斯}断言，真天主的形状无法被看见，而且真天使站立在祂座旁。\\Person{柏拉图}也同样同意这一点，并在坚持一位天主存在的同时，称其他的为天使或魔鬼。\\Person{赫耳墨斯·特里斯墨吉斯忒斯}也谈论一位天主，并承认祂是不可领悟的，超出我们的估计能力。}因此，如果他们要获得在基督内的救恩，肯定不会对他们说：你们所有的这个是不好的或虚假的；但显然应当对他们说：尽管这在你们内是完善和真实的，但除非你们来到基督的恩宠前，它毫无益处。因此，如果即使在十分外邦人身上也能找到并正确认可任何圣善的事，尽管源于基督的救恩尚未赐予他们，我们就不应当，即使异端者比他们更坏，被激发起纠正他们自身中属于他们自己的坏事的愿望，而不愿意承认他们内属于基督的善事。但我们将从一个新的前言开始，考虑这会议的其余判决。\par

\SourceRecord{Translated by J.R. King and revised by Chester D. Hartranft.；Nicene and Post-Nicene Fathers, First Series；Vol. 4.；Edited by Philip Schaff.；Buffalo, NY: Christian Literature Publishing Co.,；1887.；<http://www.newadvent.org/fathers/14086.htm>.}

% structure: p0001 source=h1 target=section reason=page-title

\section{《论圣洗圣事，驳斥多纳徒派（卷七）》}

\Emph{其中审查了迦太基会议的其他判决。}\par

% structure: p0003 source=h2 target=subsection reason=relative-heading-rank; base=h2

\subsection{第一章}

1. 但愿读者不嫌我们烦扰，因为我们常从不同角度讨论同一问题。因为，虽然圣而公教会遍及万邦，凭借原始习惯和全体会议的权威，已巩固自身，以抵御那些为圣洗圣事问题笼罩些许黑暗的论点——即异端者和分裂者中的圣洗圣事是否可能与教会中的相同——但是，在教会合一内，一些绝不可轻视的人，尤其是\Person{西彼廉}，曾持有不同意见，而那些人企图利用他的权威来反对我们，我们远不及他的爱德。因此，我们被迫利用机会，审查并思考我们在他会议和书信中发现的关于此问题的一切，以便，可以说，相当详细地处理这同一问题，并说明整个公教会如何更真实地作出了决定：即已经从其原先所属团体领受过圣洗的异端者或分裂者，应当带着此圣洗被接纳进入公教会的共融，只要他们的错误得到纠正，并在信德中扎根和立定；如此，就圣洗圣事而言，不是添加所欠缺的，而是使他们已有的变成有益的。确实，如今可朽坏的身体不再压迫灵魂，地上的居所不再压迫那沉思许多事物的心神（\BibleRef{智 9:15}），圣\Person{西彼廉}更清楚地看见那真理，因他的爱德使他配得达到那真理。因此，愿他以祈祷帮助我们，当我们在这必死肉身的云翳中劳苦时，若主恩准，我们可以尽己所能效法他身上的善。但如果他在某点上想得不正确，并说服某些兄弟和同僚接受他的观点——如今他因所爱之主启示而清楚看见的事——那么，我们虽远不及他的功德，却按自己软弱所能，追随他本人曾是的公教会权威——他是该教会卓越而最尊贵的成员——当竭尽全力反对异端者和分裂者。因为他们既脱离了\Person{西彼廉}所持守的合一，失去了他藉以结果实丰盈的爱德，远离了他所站立的谦卑，就会被\Person{西彼廉}更多地否认和谴责，因为他知道他们怀着背叛的意图查阅他的著作，却不愿效法他为维护和平所做的一切——就像那些自称为纳匝肋基督徒的人，仿效犹太人割去肉体的包皮，本是生来就在那错误中的异端者——\Person{伯多禄}曾因偏离真理被\Person{保禄}责斥（\BibleRef{迦 2:11}）——而他们至今仍顽固不化。因此，正如他们留在自己的悖逆中，从教会身体被剪除，而\Person{伯多禄}藉殉道的光荣在宗徒首席职位上得了冠冕；同样，这些人虽因悖逆而孤独，但\Person{西彼廉}因丰盈的爱德，藉苦难的荣光被接纳入圣人行列，他们必须承认自己是被逐于合一之外的人，并为了维护自己的诽谤，竟然将一个合一之公民树立为反对合一之家本身的对头。让我们因此继续以同样方式审查那次会议的其他判决。\par

% structure: p0005 source=h2 target=subsection reason=relative-heading-rank; base=h2

\subsection{第二章}

2. \Place{马克塔里斯}的\Person{马库斯}说：\Quote{异端者既是真理的仇敌和反对者，却将托付并赐予他人的东西据为己有，这并不奇怪。令人惊异的是，我们中的一些人，背叛真理，却支持异端者，反对基督徒；因此我们决定：应给异端者施洗。}\par

3. 我们回答他说：确实，更令人惊奇、更值得大为称颂的是，西彼廉和他的同僚对合一如此热爱，以致他们继续与那些他们视为真理背叛者的人保持合一，毫不担心被他们玷污。因为当马尔谷说\Quote{令人惊奇的是，我们中有些真理的背叛者，竟然支持异端者并反对基督徒}时，他自然应当加上\Quote{因此我们断定，不应与他们共融。}他没有这样说；他实际说的是\Quote{因此我们断定，异端者应受洗。}这是遵循了平和的西彼廉先前所嘱咐的话：\Quote{不判断任何人，也不因任何人有不同意见而剥夺其共融的权利。}因此，当多纳徒派诽谤我们，称我们为\Emph{\OriginalTerm{交付者}{traditors}}时，我很想知道，假如发现任何犹太人或外邦人，在阅读那次公议的记录后，按照他们自己的规则，称我们和他们都为真理的背叛者，我们该如何共同辩护，以驳斥并洗刷如此严重的指控。他们称那些人为\Emph{交付者}，而这些人，他们过去从未能定罪，现在也不能证明其有罪，反而他们自己更显出该受同样的指控。但这与我们何干？对于那些他们自己所表明的、毋庸置疑的背叛者，我们又该说什么？因为，如果我们，无论多么不实，被称为\Emph{交付者}，是因为，如他们所说，我们与\Emph{交付者}参与了同一共融，那么我们都与这些\Emph{交付者}有份，因为在真福西彼廉的时代，多纳徒派尚未从合一中分裂出去。交付圣书（他们由此开始被称为\Emph{交付者}）发生在西彼廉殉道后约四十多年。因此，如果我们因源自\Emph{交付者}而成为\Emph{交付者}，如他们所信或所声称，那么我们双方都同样源自那些其他的背叛者。因为不能说他们没有与这些背叛者共融，因为他们称这些人为自己的一方。在他们最喜欢引用的公议话语中，它宣告：\Quote{我们中有些人，}它说，\Quote{真理的背叛者，支持异端者。}对此又加上西彼廉的见证，清楚地表明他与他们保持共融，当他说：\Quote{不判断任何人，也不因任何人有不同意见而剥夺其共融的权利。}因为那些有不同意见的人，正是马尔谷称为真理背叛者的人，因为他们支持异端者，如他所坚持，通过不接受洗礼就接纳他们进入教会。而且，这种接纳方式本是惯例，西彼廉本人多处证实，该公议中几位主教也证实。由此可见，如果异端者没有洗礼，那时的基督教会就充满了背叛者，他们通过这种方式接纳异端者而支持他们。因此，我主张，我们应共同辩护，应对他们无法否认的背叛指控，而在此中，我们的特殊案件将针对交付圣书的指控进行辩论，这一指控他们无法证明。但让我们假设他们定了我们的罪；我们如何共同回答那些向我们双方提出我们祖先普遍背叛的人，我们就如何回答那些向我们提出我们祖先交付圣书的人。因为正如我们因祖先交付圣书而死了，这导致他们与我们分裂，同样，我们和他们双方都因祖先的背叛而死了，我们和他们都源自这些祖先。但既然他们说他们活着，他们认为那背叛丝毫不影响他们，因此交付圣书也不影响我们。还应注意，根据他们，背叛是无可争辩的；而根据我们，无论是前面的背叛指控（因为我们说异端者也可以有基督的洗礼），还是后面的交付圣书指控（因为在这件事上他们自己被打败了），都没有真理。因此，他们没有理由以可恶的分裂之罪分离自己，因为，如果我们的祖先没有犯交付圣书之罪，如我们所说，那么根本没有任何指控影响我们；但如果他们犯了罪，如这些人所说，那么这罪丝毫不影响我们，正如那些其他的背叛者的罪不影响我们或他们一样。因此，既然没有任何因祖先的不义而牵连我们的指控，那么针对他们因自身分裂而来的指控就明显成立了。\par

% structure: p0008 source=h2 target=subsection reason=relative-heading-rank; base=h2

\subsection{第三章}

4. \Place{西奇利巴}的\Person{萨提乌斯}说：\Quote{如果异端者在自己受洗时获得了罪的赦免，那么他们来教会就没有理由了。因为既然人因罪在审判日受罚，那么异端者若已获得罪的赦免，在基督的审判中就无所畏惧了。}\par

5. 这也可能是我们自己的判断；但让它的作者当心，这话是在什么精神下说的。因为它的措辞如此重要，以至于我在同样精神下同意并赞同它时，不会感到任何良心的不安，因为我也相信异端者确实可以拥有基督的洗礼，但不能获得他们罪过的赦免。但他没有说：如果异端者施洗或受洗，而是说：\Quote{如果异端者}，他说，\Quote{在他们自己的洗礼中获得罪过的赦免，那么他们来到教会就没有理由。}因为如果我们把那些西彼廉所认识的、在教会内\Quote{仅以言语而非以行为弃绝世界}的人置于异端者的位置，我们也可以以完全真实的同样话语表达这一判断。如果那些似乎只是悔改的人在他们自己的洗礼中获得了罪过的赦免，那么后来引导他们走向真正的悔改就没有理由。因为既然人在审判之日因罪受罚，\Quote{那些仅以言语而非以行为弃绝世界}的人，如果获得了罪过的赦免，在基督的审判中就无所畏惧。但这一推理只有在添加了某些上下文时才得以完善。但他们在基督的审判中应当畏惧，并应立即在其心的真理中悔改；当他们这样做时，肯定不需要再次受洗。因此，他们有可能接受洗礼，要么没有获得罪过的赦免，要么立即被已获赦免的罪担所重压；异端者的情况也是如此。\par

% structure: p0011 source=h2 target=subsection reason=relative-heading-rank; base=h2

\subsection{第四章}

6. 高尔的维克多说：\Quote{既然罪过只在教会的洗礼中得以赦免，承认异端者未经洗礼而共融的人，犯有违背理性的两个错误：一方面，他不洁净异端者，另一方面，他玷污了基督徒。}\par

7. 对此我们回答说：教会的洗礼甚至存在于异端者中间，尽管他们自己不在教会内；正如天堂的水在埃及地也能找到，尽管那地本身不在天堂。因此，我们不承认异端者未经洗礼而共融；既然他们带着纠正的执拗而来，我们接纳的不是他们的罪，而是基督的圣事。并且关于罪过的赦免，我们在此再次说出与以上完全相同的观点。当然，关于他在判决结束时所说的话，他宣称\Quote{他犯了两个违背理性的错误，因为他一方面不洁净异端者，另一方面又玷污了基督徒}，西彼廉本人首先与同意他的同僚们最热切地驳斥这一点。因为他并不认为当他为了和平的纽带而决定与这样的人共融是正当的，并说出\Quote{不判断任何人，也不因任何人持有不同意见而将其从共融的权利中移除}时，他自己被玷污了。或者，如果异端者因未经洗礼而被接纳共融玷污了教会，那么整个教会都因这里多次记载的惯例而被玷污了。正如那些人因我们的祖先而称我们为\OriginalTerm{叛教者}{traditor}，尽管他们在指控时无法证明祖先有任何此类行为，同样，如果每个人均与其共融者的品性有份，那么所有人都成了异端者。如果每个断言这一点的人都是疯狂的，那么维克多所说的\Quote{承认异端者未经洗礼而共融的人，不仅不洁净异端者，而且也玷污了基督徒}必然是假的。或者，如果这是真的，那么那些人并未未经洗礼而被接纳，而是那些曾按这些当事人承认存在的惯例如此被接纳的人拥有基督的洗礼，尽管它是在异端者中施与和领受的；基于同样理由，现在他们同样被正当地以同样方式接纳。\par

% structure: p0014 source=h2 target=subsection reason=relative-heading-rank; base=h2

\subsection{第五章}

8. 乌提卡的奥勒留说：\Quote{既然宗徒说我们不应与别人的罪有分\BibleRef{弟前 5:22}，那么与异端者共融而未接受教会洗礼的人，除了使自己与别人的罪有分外，还做了什么呢？因此我宣告我的判决：异端者应受洗，好使他们获得罪过的赦免，然后才允许他们共融。}\par

9. 回答是：因此，西彼廉和所有那些主教与别人的罪有分，因为他们仍与这样的人共融，而他们并未将任何持有不同意见的人从共融的权利中移除。那么，教会在哪里呢？那么，暂时不谈异端者——因为这一判决的话也适用于其他罪人，就如西彼廉悲哀地看到在教会内与他同在的、他一面驳斥一面容忍的那些人——教会在哪里呢？按照这些话，它必被认为从那一刻起就因他们罪的感染而消亡了。但如果，正如最坚定地确立的真理，教会既已存留且仍存留，那么宗徒所禁止的与别人的罪有分必须仅仅被理解为同意它们。但让异端者再次受洗，好使他们获得罪过的赦免，如果执拗和嫉妒的人再次受洗，这些人既然\Quote{仅以言语而非以行为弃绝世界}，固然能够接受洗礼，却没有获得罪过的赦免，正如主说：\Quote{如果你们不宽恕别人的过犯，你们的父也必不宽恕你们的过犯。}\BibleRef{玛 6:15}\par

% structure: p0017 source=h2 target=subsection reason=relative-heading-rank; base=h2

\subsection{第六章}

10. 日耳曼尼基阿纳的杨布斯说：\Quote{那些赞同异端者洗礼的人，就是反对我们的洗礼，以致否认那些在教会外——我姑且不说被洗净，而是被玷污——的人应在教会内受洗。}\par

11. 答覆他：我们中没有人赞同异端者的圣洗，但都赞同基督的圣洗，即使它存在于身为教会外之糠秕的异端者中，如同它可以存在于教会内身为糠秕的其他不义之人中。因为如果那些在教会外受洗的人不是被洗净而是被玷污，那么那些在教会所建于其上的磐石之外受洗的人肯定也不是被洗净而是被玷污。但所有听见基督的话而不实行的人，都在这磐石之外。或者，如果他们在圣洗中确实被洗净了，却仍停留在不义的玷污中，因为他们不愿改过迁善，那么异端者也是如此。\par

% structure: p0020 source=h2 target=subsection reason=relative-heading-rank; base=h2

\subsection{第七章}

12. 卢库马的卢西安说：\Quote{经上记载：\BibleQuote{天主看见光好，就将光与黑暗分开。}\BibleRef{创 1:4}如果光与黑暗可以相合，那么我和异端者之间就可以有共通之处。因此，我发表我的判断：异端者应该受洗。}\par

13. 答覆他：如果光与黑暗可以相合，那么正义的人与不义的人之间就可以有共通之处。因此，让他对那些西彼廉在教会内所驳斥的不义之人发表他的判断，认为他们应该重新受洗；或者，谁能说那些\Quote{只在言语上而非在行为上弃绝世俗}的人不是不义的呢？\par

% structure: p0023 source=h2 target=subsection reason=relative-heading-rank; base=h2

\subsection{第八章}

14. 卢佩尔恰纳的佩拉吉安说：\Quote{经上记载：\InnerQuote{要么天主是神，要么巴耳是神。}同样，要么教会是教会，要么异端是教会。此外，如果异端不是教会，那么教会的圣洗如何能存在于异端者中呢？}\par

15. 我们可以如下答覆他：要么乐园是乐园，要么埃及是乐园。此外，如果埃及不是乐园，乐园的水如何能在埃及呢？但有人会对我们说，水从乐园流出而延伸至那里。同样地，圣洗也延伸至异端者。我们也可以说：要么磐石是教会，要么沙是教会。此外，既然沙不是教会，圣洗如何能存在于那些听见基督的话而不实行、在沙上建造的人中呢？\BibleRef{玛 7:24} 然而，圣洗确实存在于他们中；同样，也存在于异端者中。\par

% structure: p0026 source=h2 target=subsection reason=relative-heading-rank; base=h2

\subsection{第九章}

16. 米迪拉的雅德尔说：\Quote{我们知道，在天主教会中只有一个圣洗，因此我们不应接纳异端者，除非他在我们身体内受洗，免得他认为自己在天主教会外受洗了。}\par

17. 我们答覆他说，如果这句话是针对那些在磐石之外的不义之人说的，那么它无疑就是错误的。所以对于异端者也是如此。\par

% structure: p0029 source=h2 target=subsection reason=relative-heading-rank; base=h2

\subsection{第十章}

18. 同样，马拉扎纳的另一位费利克斯说：\Quote{只有一个信德，一个圣洗\BibleRef{弗 4:5}，但这是属于天主教会，因为只有它被授予施洗的权柄。}\par

19. 假如另一个人这样说：只有一个信德，一个圣洗，但只属于正义的人，因为只有他们被授予施洗的权柄？正如这些话可以被驳斥，同样也可以驳斥费利克斯的判断。那些在圣洗中心志未改、\Quote{只在言语上而非在行为上弃绝世俗}的不义之人，难道还属于教会的肢体吗？让他们思考，这样的教会是否是真正的磐石、真正的鸽子、无瑕无疵的新娘本身。\par

% structure: p0032 source=h2 target=subsection reason=relative-heading-rank; base=h2

\subsection{第十一章}

20. 博巴的保禄说：\Quote{我个人并不动摇，即使有些人未能持守教会的信德与真理，因为宗徒说：\BibleQuote{纵然有些人不相信，他们的不信能使天主的信德失效吗？断乎不能！不如说，天主是真实的，人人都是撒谎者。}\BibleRef{罗 3:3}但如果天主是真实的，天主的真理如何在异端者的团体中呢？那里没有天主。}\par

21. 我们答覆他：贪婪的人中有什么？然而圣洗存在于他们中；同样也存在于异端者中。因为天主所在的人，是天主的宫殿。\BibleQuote{天主的宫殿与偶像怎能有什么一致呢？}\BibleRef{格后 6:16} 此外，保禄认为，西彼廉也同意，贪婪就是偶像崇拜；而西彼廉本人又与他的同僚联合，这些人虽是强盗却仍然施洗，且以极大的忍耐报酬相待。\par

% structure: p0035 source=h2 target=subsection reason=relative-heading-rank; base=h2

\subsection{第十二章}

22. 狄奥尼西纳的蓬波尼说：\Quote{显然，异端者不能施洗并给予罪赦，因为没有任何权柄赐予他们，使他们在世上能束缚或释放任何事。}\par

23. 答覆是：这权柄也没有赐给谋杀者，即那些恨自己兄弟的人。因为这样的人并没有得到\Quote{你们赦免谁的罪，就给谁赦免；你们保留谁的罪，就给谁保留。}\BibleRef{若 20:23} 但他们却施洗，保禄在圣洗的共融中容忍他们，西彼廉也承认他们。\par

% structure: p0038 source=h2 target=subsection reason=relative-heading-rank; base=h2

\subsection{第十三章}

24. 蒂尼萨的韦南提说：\Quote{如果一个丈夫出外远行，将他妻子的监护权托付给一位朋友，那么他必定会以最大的谨慎守护所托付的妻子，使她的贞洁与圣洁不被任何人玷污。我们的主、天主基督，当祂往父那里去时，将祂的新娘托付给我们照顾；我们是保持她纯洁无瑕，还是将她的纯洁与贞操出卖给奸淫者和败坏者呢？因为谁使基督的圣洗与异端者共有，谁就是将基督的新娘出卖给奸淫者。}\par

25. 我们答覆：那些受洗时只用嘴唇而不用心归向主的人，他们不是怀着奸淫的心思吗？他们自己不是爱世俗、只在言语上而非在行为上弃绝世俗的人吗？他们不是通过邪恶的交往败坏良好的品行，说：\Quote{我们吃喝吧，因为明天我们就要死了}吗？难道宗徒的话不是警戒这样的人，他说：\BibleQuote{但我怕你们的心意，或许被蛇的诡诈腐化，失去那在基督内的纯朴。}\BibleRef{格后 11:3} 因此，当西彼廉认定基督的圣洗与这样的人共有时，他是将基督的新娘出卖给奸淫者，还是即使在一个奸淫者身上也认出新郎的项链呢？\par

% structure: p0041 source=h2 target=subsection reason=relative-heading-rank; base=h2

\subsection{第十四章}

26. \Quote{我们领受了同一个圣洗圣事，我们也施与这同一个；但若有人说，权柄也赐给了异端者，使他们可以施洗，那他就制造了两个圣洗圣事。}\par

27. 我们回答他说：那么，主张不义的人也能施洗的，岂不也制造了两个圣洗圣事么？因为义人和不义人本身是彼此对立的，然而义人所施的圣洗——如保禄或息普黎安所施的——与那些仇恨保禄的不义之人通常所施的圣洗并不相悖；息普黎安认为那些人不属异端者，而是恶劣的天主教徒。虽然息普黎安身上所见的节制与其同僚身上的贪婪本身是彼此对立的，然而息普黎安惯常施与的圣洗与反对他的同僚所施的圣洗并不相悖，反而是同一而相同的，因为在这两种情形中，施洗的正是那一位，经上论祂说：\BibleQuote{那一位就是施洗的。\BibleRef{若 1:33}}\par

% structure: p0044 source=h2 target=subsection reason=relative-heading-rank; base=h2

\subsection{第十五章}

28. \Quote{如果异端者可以施洗，他们就行非法之事而被宽恕和辩护；我也不明白为何基督教他们为祂的敌人，或宗徒称他们为假基督。}\par

29. 我们回答他说：我们说异端者无权施洗，其意义与我们说诈骗者无权施洗相同。因为天主不仅对异端者，也对罪人说：\Quote{你怎敢传述我的诫命，你的口怎敢提及我的盟约？}祂确实对同一人说：\Quote{你看见盗贼，便与他同伙。}何况那些不与盗贼同伙，却惯用诡诈欺骗侵夺田产的人，情况岂不更糟？然而息普黎安并未与他们同伙，虽然他在天主教会的麦田中容忍他们，免得连麦子一同拔除。与此同时，他们自己所施的圣洗正是同一个圣洗，因为它不属于他们，而属于基督。正如他们一样——虽然基督的圣洗在他们内被承认——却未因行非法之事而被宽恕和辩护，而基督也正确地称那些坚持此类行为、最终将听到判决的人为祂的敌人：\BibleQuote{离开我，你们这些作恶的人！\BibleRef{玛 7:23}}因此他们也被称为假基督，因为他们违反基督的话而生活，与基督对立。同样，异端者也是如此。\par

% structure: p0047 source=h2 target=subsection reason=relative-heading-rank; base=h2

\subsection{第十六章}

30. \Quote{外邦人虽崇拜偶像，却也承认并明认至高神、父及创造者。马尔基盎亵渎祂，有些人竟不羞于认可马尔基盎的圣洗。那些不给天主的敌人施洗、又与未受洗的他们保持共融的司铎，如何维护或辩护天主的司铎职？}\par

31. 回答如下：当使用这类言词时，实已越过一切节制；他们也未考虑到，连他们自己\Quote{既不论断任何人，也不因有人持异议而将其排除在共融权利之外}，如此与这样的人保持共融。但撒杜尔尼诺斯在他这次判决中使用了一个论点，可为他提供劝诫的素材（若他肯注意的话），即应在每个人身上纠正错误，认可正确；因为他说：\Quote{外邦人虽崇拜偶像，却也承认并明认至高神、父及创造者。}那么，若有这样一个外邦人来到天主面前，他难道愿意在他身上纠正并改变他承认并明认天主父及创造者这一点么？我想不会。但他会修正他身上的偶像崇拜——那是他内部的邪恶——并将他未拥有的基督圣事赐给他；凡在他身上发现的任何乖谬，他都予以纠正；凡缺乏的，他都予以补充。同样，对待马尔基盎派的异端者，他会承认圣洗的完整性，纠正他的乖谬，教导他天主教真理。\par

% structure: p0050 source=h2 target=subsection reason=relative-heading-rank; base=h2

\subsection{第十七章}

32. \Quote{既然罪只在教会的圣洗中被赦免，那么，不给异端者施洗的人，就是与罪人保持共融。}\par

33. 那么，与做这事的人保持共融的，难道不是与罪人保持共融么？但他们所有的人不正是\Quote{不论断任何人，也不因有人持异议而将其排除在共融权利之外}么？那么教会何在？那些忍耐、容忍莠子以免连麦子一同拔除的人，难道这些事对他们并非障碍？我愿那些因与整个世界分离而犯了分裂亵渎之罪的人说明，他们口中虽有息普黎安的判决，心中却为何没有息普黎安的忍耐？但对于这位玛尔克路斯，我们已在上文论及圣洗与罪的赦免时回答过了，解释了一个人内可以有圣洗，虽在他内并无罪的赦免。\par

% structure: p0053 source=h2 target=subsection reason=relative-heading-rank; base=h2

\subsection{第十八章}

34. \Quote{如果教会不给异端者施洗，因为据说他已受洗，那么异端就更大了。}\par

35. 回答是：依同样原则可以说：因此，如果教会不给贪婪者施洗，因为据说他已受洗，那么贪婪就更大了。但这是假的，因此另一个也是假的。\par

% structure: p0056 source=h2 target=subsection reason=relative-heading-rank; base=h2

\subsection{第十九章}

36. \Quote{我承认一个教会，以及与之相属的一个圣洗。若有人说，圣洗的恩宠存在于异端者中，他就必须首先指明并证明教会与他们同在。}\par

37. 我们回答他说：如果你说圣洗的恩宠等同于圣洗，那么它就存在于异端者中；但若圣洗是恩宠的圣事或外在标记，而恩宠本身是罪的消除，那么圣洗的恩宠不存在于异端者中。然而，如此则有一个圣洗和一个教会，正如有一个信德一样。因此，正如善人与恶人虽没有同一望德，却能有同一个圣洗，同样，那些没有同一共同教会的人，也能有同一共同的圣洗。\par

% structure: p0059 source=h2 target=subsection reason=relative-heading-rank; base=h2

\subsection{第二十章}

38. 塔拉萨的佐齐穆说：\Quote{当真理被启示出来之后，错误必须向真理让步；因为伯多禄也曾如此，他先前惯常行割损，但当保禄宣示真理时，他便让步了。}\par

39. 我们回答说：这也可被视为我们判断的表达，而这正是就圣洗问题所发生的事。因为当真理被更清楚地启示之后，错误便向真理让步了，那最有益的惯例藉着全体会议的权威得以进一步确认。然而，他们不断记得，连宗徒之长伯多禄也可能一度有与真理要求不合的见解，这很好；我们相信，在西彼廉身上也是如此，我们对他毫无不敬，且以全部爱德待他，因为他不该被爱得超过伯多禄。\par

% structure: p0062 source=h2 target=subsection reason=relative-heading-rank; base=h2

\subsection{第二十一章}

40. 特莱普特的尤利安说：\Quote{经上记载：\BibleQuote{人不能领受什么，除非是从天上赐予他的。}\BibleRef{若 3:27}如果异端是从天上来的，那么它就能施行圣洗。}\par

41. 让他也听听别人说的：如果贪财是从天上来的，那么它就能施行圣洗。然而，贪财之人确实施行它；所以异端者也可以。\par

% structure: p0065 source=h2 target=subsection reason=relative-heading-rank; base=h2

\subsection{第二十二章}

42. 提米达雷吉亚的福斯托说：\Quote{让那些偏袒异端的人不要自夸。谁为异端的缘故干涉教会的圣洗，谁就把他们变成基督徒，而把我们变成异端。}\par

43. 我们回答他说：假如有人说，一个人领受圣洗时没有获得罪的赦免，因为他心中对弟兄怀有仇恨，但当他从心中消除那仇恨时，却不需要再受洗；那么这人难道是为杀人者干涉教会的圣洗吗？还是他使他们成为义人，而把我们变成杀人者？因此，让他在异端的事上也明白同样道理。\par

% structure: p0068 source=h2 target=subsection reason=relative-heading-rank; base=h2

\subsection{第二十三章}

44. 富尔尼的格米尼说：\Quote{我们某些同僚或许将异端者看得比自己高，但他们不能将他们看得比我们高；因此我们持守我们曾经决定的：我们应该为那些从异端来到我们这里的人施洗。}\par

45. 这个人也最公开地承认，他某些同僚的见解与他自己的相反；因此，合一之爱一再得到确认，因为他们之间并未因分裂而分离，直到天主向其中一人启示他们另有心思的任何事。\BibleRef{斐 3:15}但我们回答他说，他的同僚并没有把异端者看得比自己高，而是如同在贪财者、欺诈者、强盗、杀人者身上承认基督的圣洗一样，在异端者身上也承认它。\par

% structure: p0071 source=h2 target=subsection reason=relative-heading-rank; base=h2

\subsection{第二十四章}

46. 诺瓦的罗加提安说：\Quote{基督建立了教会，魔鬼建立了异端；撒旦的会众怎能拥有基督的圣洗呢？}\par

47. 我们回答他说：难道因为基督建立了善意之人，魔鬼建立了嫉妒之人，所以魔鬼的党派——被证明存在于嫉妒者之中——就不能拥有基督的圣洗吗？\par

% structure: p0074 source=h2 target=subsection reason=relative-heading-rank; base=h2

\subsection{第二十五章}

48. 布拉的特拉皮说：\Quote{若有人将基督的圣洗放弃并出卖给异端者，除了说他是基督净配的犹达斯之外，还能说他是什么呢？}\par

49. 这里对所有分裂者是何等大的谴责！他们以邪恶的亵渎使自己与分散在全世界的基督产业分离，如果西彼廉曾与像叛徒犹大那样的人共融，却未被他们玷污；或者如果他曾被玷污，那么所有人都变得像犹大一样；或者如果他们不是，那么前人的恶行并不属于后人，即使他们是同一共融的后裔。那么，他们为什么向我们控诉\Emph{\OriginalTerm{背教者}{traditores}}，却没有证明他们的指控，却不向自己控诉犹大，西彼廉及其同僚曾与他共融？看这会议，这些人常以此夸耀！我们确实说，那认可基督的洗礼甚至在异端者中的人，并没有向异端者出卖基督的洗礼；正如那认可基督的洗礼甚至在谋杀者中的人，并没有向谋杀者出卖基督的洗礼。但既然他们声称从这会议的法令中为我们规定应当持守的见解，就让他们先自己同意这些法令。看，其中所有说异端者虽在异端中受洗却不需再受洗的人，都被比作叛徒犹大。然而西彼廉愿意与他们共融，当他说：\Quote{我不判断任何人，也不剥夺任何持有相反意见者的共融权利。}但从前在教会内曾有这样的人，从他的话中显明：\Quote{但有人会说：\InnerQuote{那么，对那些从前未受洗就被接纳进入教会的人，该怎么办呢？}}这正是教会的习惯，这会议的成员们一再见证。因此，如果有人这样做，\Quote{除了是基督净配的犹大之外，不能称为别的}，正如泰拉皮乌斯的判决所言；但根据福音的教导，犹大是叛徒；那么，所有那些当时发出这些判决的人，以及在发出判决之前因当时教会保留的习惯而成为叛徒的人，都与叛徒共融。因此，所有人——即我们和那些作为那合一的后裔的他们自己——都是叛徒。但我们以两种方式为自己辩护：第一，正如西彼廉在开场白中所言，在不损害合一权利的前提下，我们不同意这会议的判决；第二，我们认为恶人在天主教会合一中绝不影响善人，直到最后糠秕从麦子中分离。但我们的对手，既躲在这会议的法令之下，又主张善人会因与恶人共融而像被感染一样灭亡，他们无法避免承认：早期的基督徒，作为他们的祖先，是叛徒，因为他们被自己的会议定罪；并且前人的行为确实影响他们，因为他们向我们控诉我们祖先的行为。\par

% structure: p0077 source=h2 target=subsection reason=relative-heading-rank; base=h2

\subsection{第二十六章}

50. 另一位门布雷萨的路济乌斯说：\Quote{经上记载：\BibleQuote{天主不听罪人。} \BibleRef{若 9:31} 一个罪人怎能被垂听而受洗呢？}\par

51. 我们回答：贪吝的人、强盗、放高利贷者和谋杀者，他们怎么被垂听？他们岂不都是罪人？然而西彼廉虽在天主教会中指责他们，却容忍他们。\par

% structure: p0080 source=h2 target=subsection reason=relative-heading-rank; base=h2

\subsection{第二十七章}

52. 另一位布斯拉切尼的费利克斯说：\Quote{在接纳异端者未经洗礼进入教会的事上，不要让任何人以习惯取代理性与真理；因为理性与真理总是排除习惯。}\par

53. 我们回答他：你没有显示真理；你承认那习惯的存在。因此，即使真理仍被隐藏，我们坚持那后来由全体会议所确认的习惯，也是正确的，我们相信真理已经显明。\par

% structure: p0083 source=h2 target=subsection reason=relative-heading-rank; base=h2

\subsection{第二十八章}

54. 另一位阿比提尼的撒杜尔尼努斯说：\Quote{如果假基督能将基督的恩宠赐予任何人，那么被称为假基督的异端者也能施洗了。}\par

55. 假如有人说：如果谋杀者能赐予基督的恩宠，那么那些恨自己弟兄、被称为谋杀者的人也能施洗了。他这样说似乎有道理，但他们仍能施洗；同样，异端者也能施洗。\par

% structure: p0086 source=h2 target=subsection reason=relative-heading-rank; base=h2

\subsection{第二十九章}

56. 阿吉雅的昆图斯说：\Quote{拥有某物的人才能给予；但众所周知，异端者一无所有，他们能给予什么呢？}\par

57. 我们回答他：如果一个人拥有某物才能给予，那么显然异端者能给予洗礼：因为他们脱离教会时，仍然拥有在教会内领受的洗礼圣事；当他们回归时，并不再领受，因为他们离开教会时并没有失去它。\par

% structure: p0089 source=h2 target=subsection reason=relative-heading-rank; base=h2

\subsection{第三十章}

58. 另一位马尔切利亚纳的尤利安说：\Quote{若有人能事奉两个主，天主和\OriginalTerm{玛门}{mammon}，\BibleRef{玛 6:24}那么洗礼也能事奉两个，基督徒和异端者。}\par

59. 诚然，如果洗礼能事奉节制者和贪吝者、清醒者和醉酒者、仁爱者和谋杀者，为何不能事奉基督徒和异端者呢？——其实它并不真正事奉他们，而是给他们施行，由他们施行，为正确使用的人得救恩，为错误使用的人得审判。\par

% structure: p0092 source=h2 target=subsection reason=relative-heading-rank; base=h2

\subsection{第三十一章}

60. 霍雷亚·切利亚的特纳克斯说：\Quote{只有一个洗礼，但属于教会；哪里没有教会，哪里也不能有洗礼。}\par

61. 我们回答他：那么，当没有磐石而只有沙土时，洗礼怎能存在呢？因为教会建立在磐石上，而非沙土上。\par

% structure: p0095 source=h2 target=subsection reason=relative-heading-rank; base=h2

\subsection{第三十二章}

62. 另一位阿苏拉斯的维克多说：\Quote{经上记载：\InnerQuote{只有一位天主，一位基督，一个教会，一个洗礼。}那么，在没有天主、基督或教会的地方，怎能有人施洗呢？}\par

63. 怎能有人在那沙土上施洗呢？那里既没有教会（因为教会建立在磐石上），也没有天主和基督（因为那里没有天主和基督的殿宇）。\par

% structure: p0098 source=h2 target=subsection reason=relative-heading-rank; base=h2

\subsection{第三十三章}

64. 卡普斯的多纳图卢斯说：\Quote{我也一向持此见解：凡在教会之外一无所获的异端者，当他们回头归向教会时，应受圣洗。}\par

65. 对此的答复是：他们在教会之外确实一无所获，但那是指救恩方面一无所获，而不是指圣事方面一无所获。因为救恩是善人所特有的，而圣事则为善人与恶人所共有。\par

% structure: p0101 source=h2 target=subsection reason=relative-heading-rank; base=h2

\subsection{第34章}

66. 鲁西卡德的维鲁卢斯说：\Quote{一个异端者不能给予他所没有的东西；分裂者更是如此，因为他已经失去了他所拥有的。}\par

67. 我们已经表明，他们仍然拥有它，因为他们分离时并未失去它。因为当他们回来时，他们并没有重新领受它。所以，如果他们曾因被认为没有而无法给予，那么现在应当明白，他们可以给予，因为已经明白他们也拥有它。\par

% structure: p0104 source=h2 target=subsection reason=relative-heading-rank; base=h2

\subsection{第35章}

68. 库库利的普登提亚努斯说：\Quote{我最近被祝圣为主教，这使我在等待并听取长辈们的决定。因为显然异端没有、也不可能有任何东西；因此，如果有人从他们那里来，公义的决定是应当给他们施圣洗。}\par

69. 因此，正如我们已经回答过那些先我们而去的人——该人正在等待他们的判断——同样，应当明白我们已回答了他本人。\par

% structure: p0107 source=h2 target=subsection reason=relative-heading-rank; base=h2

\subsection{第36章}

70. 希波迪亚里图斯的伯多禄说：\Quote{既然在天主教会中只有一个圣洗，显然人不能在教会之外受圣洗；因此我断定，那些在异端或分裂中受过圣洗的人，在来到教会时应重新受圣洗。}\par

71. 在天主教会中只有一个圣洗，其意义是：当有人离开她时，它并不在离开者中变成两个，而是保持同一个。因此，在回来的人身上所承认的，也应当在那些从分离者那里领受的人身上承认，因为他们分离到异端时并未失去它。\par

% structure: p0110 source=h2 target=subsection reason=relative-heading-rank; base=h2

\subsection{第37章}

72. 另一位奥萨法的路齐说：\Quote{按照我的心灵和圣神的感动，既然只有一位天主，我们主耶稣基督的父，一位基督，一个希望，一位圣神，一个教会，那么也只应有一个圣洗。因此我说，凡在异端者中所开始或施行的事都应废止；并且，从异端者中出来的人应在教会中受圣洗。}\par

73. 因此，让那些听了天主的话而不实行的人所施行的圣洗被宣告无效吧，当他们开始从不义转向正义，即从沙土转向磐石时。但如果这没有发生，因为他们里面属基督的部分并未因他们的不义而受损，那么也应在异端者身上理解这一点。因为在不义者中，只要他们还在沙土上，就没有与在磐石上者相同的希望；然而两者之中却有同一个圣洗，尽管如同只有一个希望一样，也只有一个圣洗。\par

% structure: p0113 source=h2 target=subsection reason=relative-heading-rank; base=h2

\subsection{第38章}

74. 古尔吉特斯的费利克斯说：\Quote{我断定，按照圣经的诫命，那些在教会之外由异端者非法受圣洗的人，若愿投靠教会，应在合法施行圣洗的地方获得圣洗的恩宠。}\par

75. 我们的答复是：让他们确实开始合法地、为得救而拥有他们先前不合法地、为灭亡而拥有的东西；因为同一个人在同一种圣洗下，当他以真实的心转向天主时就称为义，正如他在领受圣洗时\Quote{只在言语上弃绝世界，而非在行为上}就受谴责。\par

% structure: p0116 source=h2 target=subsection reason=relative-heading-rank; base=h2

\subsection{第39章}

76. 拉马斯巴的普西鲁斯说：\Quote{我相信，除非在天主教会内，圣洗不能带来救恩。凡在天主教会之外的，都只是虚伪。}\par

77. 这确实是真的：\Quote{除非在天主教会内，圣洗不能带来救恩。}因为圣洗本身也可以存在于天主教会之外；但在那里它不能带来救恩，因为在那里它不产生救恩；正如基督的馨香在灭亡的人中确实不是带来救恩，\BibleRef{格后 2:15}尽管其过失不在于它本身，而在于他们。但\Quote{凡在天主教会之外的，都只是虚伪，}只是就它不是公教而言。但在天主教会之外也可能有公教的事物，如同基督的名字可以存在于基督的会众之外，那个没有跟随门徒的人正是因这名驱魔。\BibleRef{谷 9:38}因为在天主教会之内也可能有虚伪，正如那些\Quote{只在言语上而非在行为上弃绝世界}的人毫无疑义就是如此，然而这虚伪并非公教的。因此，正如在天主教会中有某些不是公教的事物，同样，在天主教会之外也可能有某些公教的事物。\par

% structure: p0119 source=h2 target=subsection reason=relative-heading-rank; base=h2

\subsection{第40章}

78. 加扎乌法拉的萨尔维亚努斯说：\Quote{众所周知，异端者一无所有；因此他们来到我们这里，好领受他们先前没有的东西。}\par

79. 我们的答复是：按此理论，创立异端的人本身就不是异端者，因为他们从教会分离出来，而他们之前当然拥有在那里领受的东西。但如果认为那些使其余人成为异端者的人不是异端者，这是荒谬的；因此，异端者可能拥有某种东西，却因滥用而导向毁灭。\par

% structure: p0122 source=h2 target=subsection reason=relative-heading-rank; base=h2

\subsection{第41章}

80. 图卡的奥诺拉图斯说：\Quote{既然基督是真理，我们应当跟随真理胜过习惯；这样，我们用教会的洗礼圣化来到我们这里的异端者，只因为他们在外一无所获。}\par

81. 这个人也是习惯的见证，在这方面他给了我们最大的帮助，尽管他可能说了其他反对我们的话。但异端者来到我们这里，并非因为他们在外面一无所获，而是为了使他们所领受的能开始对他们有益：因为这在外面绝不可能。\par

% structure: p0125 source=h2 target=subsection reason=relative-heading-rank; base=h2

\subsection{第42章}

82. 奥克塔夫的\Person{维克托}说：\Quote{正如你们自己也知道，我作主教不久，因此我等候了长辈们的意见。故此我表达我的意见：凡从异端而来的人，应当无疑地受洗。}\par

83. 因此，对于他所等候的那些人所作的回答，也可视为对他的回答。\par

% structure: p0128 source=h2 target=subsection reason=relative-heading-rank; base=h2

\subsection{第四十三章}

84. 马斯库拉的\Person{克拉鲁斯}说：\Quote{我们的主耶稣基督的话是清楚的，当祂派遣祂的宗徒，并将从祂的圣父所领受的权能赐给他们独有；我们是他们的继承者，以同一权能治理主的教会，并给相信信德的人施洗。因此，异端者既在教会以外，没有权能，也没有基督的教会，不能以祂的圣洗给任何人施洗。}\par

85. 那么，心存恶意的凶杀者竟是宗徒的继承者吗？那他们为何施洗？是因为他们不在教会之外吗？但他们却在磐石之外，主将钥匙交给了这磐石，并说祂要在其上建立祂的教会。\BibleRef{玛 16:18}\par

% structure: p0131 source=h2 target=subsection reason=relative-heading-rank; base=h2

\subsection{第四十四章}

86. 坦贝的\Person{塞昆迪亚努斯}说：\Quote{我们不应因过分热心而欺骗异端者；他们既未在吾主耶稣基督的教会中受洗，因此也未获得罪过的赦免，免得在审判之日到来时，他们将我们没有给他们施洗、没有获得天主恩宠的宽恕归咎于我们。为此，既然只有一个教会、一个圣洗，当他们归向我们时，让他们连同教会也接受教会的圣洗。}\par

87. 不，当他们被移入磐石，加入鸽子的团体时，让他们获得罪过的赦免；这赦免他们在磐石之外、鸽子之外无法获得，无论他们是像异端者一样公开在外，还是像堕落的公教徒一样貌似在内；然而显而易见的是，这些人既拥有又施行不带有罪过赦免的圣洗，因为即使从他们自己手中，那些并未变得更好、只用嘴唇尊敬天主、心却远离祂的人接受了圣洗。\BibleRef{依 29:13} 然而，诚然只有一个圣洗，正如只有一个鸽子；尽管那些不在鸽子唯一共融中的人，却可能共有圣洗。\par

% structure: p0134 source=h2 target=subsection reason=relative-heading-rank; base=h2

\subsection{第四十五章}

88. 另有楚拉比的\Person{奥勒留}说：\Quote{宗徒若望在他的书信中立下以下诫命：\InnerQuote{若有人到你们那里，不传这个道理，不要接他到家里，也不要向他问安；因为向他问安的，就在他的恶行上有份。}这样的人怎能不加考虑就被接入天主的殿呢？他们被禁止接入我们的私宅，我们怎能与他们分享基督的圣洗而不共融呢？如果我们只向他们问安，就在他们的恶行上有份了。}\par

89. 关于若望的这段证言，无需再辩论，因为它与我们目前讨论的圣洗问题毫无关系。他说：\BibleQuote{若有人到你们那里，不传基督的道理。}但异端者离开他们错误的道理，归向基督的道理，以便被纳入教会，开始属于那只鸽子的肢体，而他们之前已经拥有它的圣事。因此，他们先前所缺乏的属于它的东西赐给了他们，即出自纯洁的心、良好的良心和真诚的信德的平安与爱德。\BibleRef{弟前 1:5} 但他们先前已经拥有的属于鸽子的东西，被承认并毫无贬损地接受；正如在淫妇身上，即使她追随她的情人，天主仍承认祂的恩赐；因为当她改正了淫行，回头归向贞洁时，那些恩赐不被归咎于她，而是她本人被纠正。\EditorialNote{欧瑟亚第二章} 但是，正如\Person{西彼廉}在他与这样的人保持共融时，如果这段若望的证言被用来指责他，他可以为自己辩护；同样，让那些被指责的人自己辩护。因为对于眼前的问题，正如我先前所说，它毫无关系。因为若望说我们不可向持异端道理的人问安；但宗徒保禄更强烈地说：\BibleQuote{若有称为弟兄的人，是贪吝者、或醉酒者、或什么的，这样的人，连与他一同吃饭也不可。}\BibleRef{格前 5:11} 然而\Person{西彼廉}曾容许他的同僚，那些放高利贷者、背信者、欺诈者和强盗，不只与他私人的餐桌共餐，而是与天主的祭台共融。至于这种作法如何可以辩护，已在其他书中充分论及。\par

% structure: p0137 source=h2 target=subsection reason=relative-heading-rank; base=h2

\subsection{第四十六章}

90. 格梅利的\Person{利特乌斯}说：\Quote{\BibleQuote{若瞎子领瞎子，两人都要掉进坑里。}\BibleRef{玛 15:14} 既然显然异端者自己是瞎子，不能给任何人光明，因此他们的圣洗无效。}\par

91. 我们并非说他们的圣洗对救恩有效，只要他们仍是异端者；正如它对我们所说的那些凶杀者毫无价值，只要他们仍恨自己的弟兄：因为他们自己也身处黑暗，若有人跟随他们，便一同掉进坑里；但这并不等于他们或者没有圣洗，或者不能施洗。\par

% structure: p0140 source=h2 target=subsection reason=relative-heading-rank; base=h2

\subsection{第四十七章}

92. 奥厄的\Person{纳塔利斯}说：\Quote{不仅我自己在场，还有萨布拉塔的\Person{庞培}和大雷普提斯的\Person{迪奥加}，他们委托我代表他们的意见；身体虽不在，精神却在；他们与我们的同僚作出同样的判断：除非异端者用教会的圣洗受洗，否则不能与我们共融。}\par

93. 我想，他是指属于鸽子团体的共融；因为在共享圣事上，他们无疑与他们共融，不判断任何人，也不因某人持不同意见而剥夺他的共融权利。但无论他说这话的所指为何，都没有太大必要反驳这些话。因为异端者除非已经以教会的圣洗受了洗，否则肯定不会被接纳进入共融。但显然，教会的圣洗甚至存在于异端者中，如果它以福音之言被祝圣；正如福音本身属于教会，与他们的乖戾毫无关系，而是确实保存了自身的圣洁。\par

% structure: p0143 source=h2 target=subsection reason=relative-heading-rank; base=h2

\subsection{第四十八章}

94. \Person{犹尼乌斯}·\Place{纳波利}的\Person{犹尼乌斯}说：\Quote{我并没有离弃我们曾经宣告的判断，即当异端者来到教会时，我们应该为他们施洗。}\par

95. 既然此人没有从圣经提出任何论据或证明，他不必让我们久留。\par

% structure: p0146 source=h2 target=subsection reason=relative-heading-rank; base=h2

\subsection{第四十九章}

96. \Person{西彼廉}·\Place{迦太基}的\Person{西彼廉}说：\Quote{我的意见在写给我们的同僚\Person{犹巴安}的信中已最为充分地阐述，即异端者按照福音和宗徒们的见证被称为基督的敌人和假基督，当他们来到教会时，应以教会唯一的圣洗为他们施洗，使他们从敌人变为朋友，从假基督变为基督徒。}\par

97. 这里还有什么必要进一步辩论呢？因为我们已极其仔细地处理了他所提及的那封致\Person{犹巴安}的信。至于他在这里所说的，我们不要忘记，这话也可以用在所有不义的人身上，正如他自己所见证的，他们存在于天主教会中，而他们拥有和施行洗礼的能力并不受我们任何人质疑。因为那些从魔鬼的阵营转向基督，建立在磐石上，与鸽子联合，安置在封闭的花园和封缄的泉源中的人，来到教会；在那园中，那些违背基督诫命生活的人，无论他们似乎在哪里，都找不到一个。因为在他写给\Person{玛格诺}的信中，讨论这个问题时，他本人已足够充分且毫不含糊地告诫我们，应理解教会由何种团体组成。因为他在谈到某个人时说：\Quote{让他成为外人和亵渎者，成为主之和平与统一的敌人，不住在天主的家中，也就是说，不住在基督的教会中，因为只有在其中居住的人才是同心合意的人。}因此，那些想凭借\Person{西彼廉}的权威对我们施加吩咐的人，请暂时留意我们这里所说的话。因为如果只有同心合意的人才住在基督的教会中，那么毫无疑问，那些出于嫉妒和纷争，没有爱德而宣讲基督的人，并不住在基督的教会中，无论他们表面上似乎是在教会内；保禄宗徒所提到的不是异端者和分裂者，而是与他同在教会内交往的假弟兄，他们肯定不该施洗，因为他们不住在教会中——他自己也说，只有同心合意的人才住在其中；除非有人远离真理到如此地步，竟说那些嫉妒、恶毒、纷争而无爱德的人是同心合意的；然而他们曾施洗；他们所显示的可憎乖戾丝毫没有破坏或减损由他们处理并施行的基督圣事。\par

% structure: p0149 source=h2 target=subsection reason=relative-heading-rank; base=h2

\subsection{第五十章}

98. 确实值得考虑上述致\Person{玛格诺}的信中他所述的全部段落，内容如下：\Quote{不住在天主的家中——也就是说，在基督的教会中——因为只有在其中居住的人才是同心合意的人，正如圣神在圣咏中论到\BibleQuote{使人在家中同心合意的天主}时所说的。最后，主的祭献本身也宣告，基督徒彼此之间藉着坚定而不可分离的相爱而联合。因为主称那由许多谷粒联合压缩而成的饼为祂的身体，\BibleRef{若 6:51}，祂是在象征祂所承载的、结合成一个身体的子民；当祂称那由许多枝条和串串葡萄压榨而汇聚成一的酒为祂的血时，\BibleRef{玛 26:26}，祂也是在象征由众多合而为一的群体所联结成的一群羊。}这些话表明真福\Person{西彼廉}既理解又爱慕天主之家的荣耀，他断言这家庭由同心合意的人组成，并以先知们的见证和圣事的含义证明之；而在这家庭中，肯定找不到那些嫉妒、恶毒而无爱德的人，然而他们曾施洗。由此可见，基督的圣事可以存在于那些不在基督教会中的人身上，并由他们施行，而\Person{西彼廉}自己见证，只有同心合意的人才住在教会中。确实不能说，只要他们未被发现就可以施洗，因为保禄宗徒并未未能发现那些人，他在书信中明确见证他们的职事，说他也因他们宣讲基督而欢喜。因为论到他们，他说：\BibleQuote{无论是出于假意还是真诚，基督被传扬了，为此我欢喜，而且还要欢喜。} \BibleRef{斐 1:18}\par

% structure: p0151 source=h2 target=subsection reason=relative-heading-rank; base=h2

\subsection{第五十一章}

99. 既思及这一切，我以为并非轻率地说：在天主的家中，有些人并非那本身即为天主之家的存在。这房屋被说成是建在磐石上的\BibleRef{玛 16:18}，被称为独一的鸽子\BibleRef{歌 6:9}，被称为毫无瑕疵或皱纹的美丽新娘、关闭的花园、封住的泉源、活水之井、结满佳果的石榴园\BibleRef{歌 4:12}；这房屋也领受了钥匙和束缚与释放的权柄\BibleRef{玛 16:19}。若有人轻忽这房屋的逮捕与纠正，主说：\Quote{让他对你们如同外邦人和税吏一样。}\BibleRef{玛 18:17} 关于这房屋，经上写道：\Quote{主，我喜爱你所居住的殿宇，以及你荣耀所居之处}；\Quote{祂使人在一家之中同心合意}；\Quote{我对那向我说\InnerQuote{我们进入主的殿宇}的人，就欢喜}；\Quote{住在你殿宇中的，真是有福；他们仍要赞美你}；还有无数类似的经文。这房屋也被称为麦子，以忍耐结出果实，有三十倍、六十倍、一百倍的。这房屋也在金器银器\BibleRef{弟后 2:20}、宝石和不朽坏的木器中。对这房屋说：\Quote{彼此以爱德宽容，竭力保持圣神所赐的合一，在和平的联系中}\BibleRef{弗 4:2}；\Quote{因为天主的殿是圣的，这殿就是你们。}\BibleRef{格前 3:17} 因为这房屋是由那些善良、忠信、散布全世界的天主圣仆组成的，藉圣神的合一相连，无论他们彼此是否相识。但我们认为，另有一些人是在这房屋之中，却并不属于这房屋的实质，也不属于那多结果实与和平的正义之团体，而仅仅是糠秕存在于麦子中间；因为当他们仍在房屋内时，我们不能否认，正如宗徒所说：\Quote{在大户人家，不但有金器银器，也有木器瓦器；有作贵重的，有作卑贱的。}\BibleRef{弟后 2:20} 在这无数的人群中，不仅有那在教会内使圣徒之心忧伤的群众（与如此庞大的数目相比，圣徒是很少的），而且还有那些裂开网眼而存在的异端与分裂，现在可以说他们是出于房屋而非在房屋内，关于他们经上说：\Quote{他们从我们中间出去，却不是属于我们的。}\BibleRef{若一 2:19} 因为他们如今在身体上也与我们分离，比那些在教会内属血肉与世俗生活、只在精神上与我们分离的人，分离得更为彻底。\par

% structure: p0153 source=h2 target=subsection reason=relative-heading-rank; base=h2

\subsection{第五十二章}

100. 对于所有这些不同的类别，关于那些首先在天主家中以至于他们本身即为天主之家的人，无论他们已是属灵的人，或仍是婴孩仅以奶为养料，却以热诚之心向着属灵之事进步，无人怀疑他们不仅自己拥有有益的圣洗，并将其传给效法他们而有益的人；但当他们传给虚伪之人时——圣神远离这些人——尽管他们自己尽其所能地施予圣洗使之有益，但后者领受却徒劳，因为他们不效法施予者。然而，那些在大户人家中如同卑贱器皿的人，他们拥有圣洗对自己无益，传给效法他们的人也无益；但那些心性联合于天主的圣家而非其仆役的人，领受圣洗却有益。至于那些分离得更彻底以至于出于房屋而非在房屋内的人，他们拥有圣洗对自己毫无益处；此外，从他们那里领受圣洗的人也无益，除非他们因迫切的需要而领受，且在领受时心中不离开合一的联系；然而他们仍然拥有圣洗，尽管拥有无益；而且即使领受者领受时无益，也从他们那里领受了圣洗；但为使圣洗变得有用，他们必须离开自己的异端或分裂，依附于天主的家。这不仅异端者和分裂者应当如此，而且那些因圣事共融而在房屋内、却因品性乖戾而在房屋外的人也当如此。因为这样，圣洗才开始甚至对他们自己有益，而此前它是无益的。\par

% structure: p0155 source=h2 target=subsection reason=relative-heading-rank; base=h2

\subsection{第五十三章}

101. 通常也会提出一个问题：从一位自己未领受圣洗的人那里领受的圣洗，是否有效——如果此人出于某种好奇心的驱使，偶然学会了应当如何施行圣洗；以及，领受者以何种心态领受，是嘲弄还是真诚，是否有所不同：如果是嘲弄，那么，当嘲弄是出于欺诈时，是在教会内，还是在他们以为是教会的地方；或者，当嘲弄是出于玩笑时，如在戏剧中；并且，哪一种是更可咒骂的：在教会内欺诈地领受，还是在异端或分裂中无欺诈地、即全心全意地领受？或者，在异端中欺诈地领受，与在戏剧中真诚地领受——如果有人在他的表演中被突然的宗教情感所触动——何者更糟？然而，如果我们即使将这样一个与在天主教会本身之内欺诈领受的人相比较，我怀疑是否有人会犹豫应该优先选择哪一个；因为我实在看不出，那位真诚施洗者的意向对欺诈领受者有何益处。但让我们考虑另一种情况：某人也是欺诈地施洗，当施洗者和领受者都在天主教会本身的合一中欺诈地行动时，这种圣洗是否更应该被承认，还是戏剧中施予的圣洗——如果有人发现有人在戏剧中因突然的宗教冲动而真诚地领受了？或者，是否真实情况是：就人本身而言，在戏剧中信服的领受者与在教会中嘲弄的领受者之间有很大的差别，但就圣洗的真实性而言，并无差别？因为，如果在天主教会本身之内，某些人是真诚地还是欺诈地举行圣洗，只要两者仍举行同一件事，在圣洗的真实性上没有差别，那么我看不出为什么在教会之外就该有差别，既然领受者并非被他的欺诈所掩盖，而是被他的宗教冲动所改变。或者，那些真诚者——在他们当中举行圣洗——对于确认圣洗有更大的能力，超过那些欺诈者——由他们和在他们当中举行圣洗——所施加的使圣洗无效的能力？然而，如果欺诈后来被揭露，没有人会要求重复圣洗，而是欺诈要么被绝罚所惩罚，要么被补赎所纠正。\par

102. 但稳妥的做法是，我们不要以任何判断的轻率来提出一种观点，这种观点既没有在天主教会的任何地区性会议中被提出，也没有在全体会议中被确认；而是要满怀着不容置疑的声音的自信，来坚持那已由普世教会的共识所确认的，在我们的主天主和救主耶稣基督的指引下。然而，如果有人催促我——假设我正适当坐在一次会议中，会上提出了这类问题——要我在没有参考其他人之前表达的观点的情况下声明我自己的意见，我宁愿遵循他人的判断，如果我受到同样的感受影响，正如导致我先前所说的话那样，我会毫不犹豫地说，所有人在任何地方、从任何种类的人那里领受的圣洗，只要是以福音的言辞祝圣，并且是领受者没有欺诈、带着某种程度的信德而领受的，他们都拥有圣洗；尽管如果他们没有爱德，无法被接枝到天主教会中，那么这圣洗对于他们灵魂的救恩毫无益处。因为宗徒说：\BibleQuote{我若有信德，甚至能移山，但我若没有爱德，我什么也不是。}\BibleRef{格前13:2}正如从我们前辈已确立的法令，我毫不迟疑地说，所有那些虽然欺诈地领受，却是在教会内领受圣洗的人，或者是在那些领受者所属的团体认为是教会的地方领受圣洗的人——关于他们曾说过：\BibleQuote{他们从我们中间出去了。}\BibleRef{若一2:19}——都拥有圣洗。但当时没有那些如此相信的人的团体，并且领受者本人也不持这样的信仰，而整个事情是作为一场闹剧、喜剧或玩笑进行的——如果我问，这样施予的圣洗是否应被认可，我会声明我的意见，即我们应当通过某种启示祈求天主审判的宣告，同心合意地祈祷并以恳切哀求的呻吟寻求，同时谦卑地顺服在后于我发言者的决定，以防他们提出某些已经知道和确定的事情。因此，我现在的意见更应被视为不影响更勤勉的研究或比我更高的权威的发表！\par

% structure: p0158 source=h2 target=subsection reason=relative-heading-rank; base=h2

\subsection{第五十四章}

103. 但现在，我以为正是时候结束这些关于圣洗圣事的书卷。在这些书卷中，我们的主天主通过和平的主教西彼廉以及与他意见一致的兄弟们的话语，向我们显示了应当对至公合一怀有多大的爱；以致即使他们在某些事上心思不同，直到天主也向他们启示这事（\BibleRef{斐 3:15}），他们也宁可容忍那些与自己意见不同的人，而不是通过邪恶的分裂与自己分离；由此，多纳徒派的口全然被封住，即使我们不说玛克西米安派的人。因为如果恶人在合一中玷污善人，那么即使是西彼廉自己，也已经找不到他可以加入的教会了。但如果恶人不在合一中传染善人，那么亵圣的多纳徒派就没有摆在自己面前的分离理由。但如果圣洗被那些行肉体之事（关于此事，经上说：\BibleQuote{行这样事的人不得继承天主的国}（\BibleRef{迦 5:19}））的众多其他人所拥有并转移，那么它也当为异端者所拥有并转移，因为异端者也被列于这些事之中；因为他们若留下，本可以转移它，而并未因他们的分离而失去它。但这类人将圣洗施行于他们的同伴，是同样无效且无益的，就像其他与他们相似的人一样，因为他们将不得继承天主的国。并且，正如当那些其他人被导入正途时，并非圣洗开始临在（此前是缺席的），而是它开始有益于他们（它早已在他们内）；也是如此，异端者亦然。因此，西彼廉和与他意见一致者不能限制天主教会，她不愿损伤他们。但他们在某些事上心思不同，我们并不惧怕，因为我们和他们一样，都分享对伯多禄的敬礼；然而他们未离开合一，我们为此喜乐，因为我们像他们一样，是建立在磐石上的。\par

\SourceRecord{Translated by J.R. King and revised by Chester D. Hartranft.；Nicene and Post-Nicene Fathers, First Series；Vol. 4.；Edited by Philip Schaff.；Buffalo, NY: Christian Literature Publishing Co.,；1887.；<http://www.newadvent.org/fathers/14087.htm>.}
